\documentclass[11pt,italian]{book}
%\usepackage[a4paper,centering,paperwidth=17.9truecm, paperheight=25.1truecm]{geometry}
\usepackage[text={5in,205mm},paperwidth=17.9truecm, paperheight=25.1truecm,centering]{geometry}
\usepackage[italian]{babel}
\topmargin -7mm
\usepackage[a4,
  % set width and height to a4 width and height + 6mm
  % width=18truecm, height=24truecm,
  % use any combination of these options to add different cut markings
  %  cam, axes, frame, cross,
  cross,
  % set the type of TeX renderer you use
  pdftex,
  % center the contents
  center,
  noinfo
]{crop}
\usepackage{amsfonts}
\usepackage{mathtools}
\DeclareMathOperator{\atantwo}{atan2}
\usepackage{fmta}
%\usepackage{psfig}
\usepackage{epsfig}
\usepackage{wrapfig}
\usepackage{mathabx}
\usepackage[multiple]{footmisc}
\usepackage{graphicx}
\usepackage{caption}
\usepackage{xcolor}
\usepackage{makeidx}
\usepackage{bm}
\usepackage[makeroom]{cancel}
\makeindex
\usepackage[totoc]{indexgarz}
%\usepackage[nottoc]{tocbibind}
\begin{document}
\renewcommand\indexname{Indice Analitico}

\logoa{
\begin{center}
\includegraphics[width=0.16\textwidth]{logoMUP.pdf}
\end{center}
}
\logob{
%\begin{center}
\noindent MUP Editore \`e il marchio editoriale di Monteparmaservizi Srl,\\
\noindent impresa strumentale della\\ \vskip 0mm
\noindent \includegraphics[width=0.3\textwidth]{logofond.pdf}
\vskip -8.5mm \hfill www.mupeditore.it
\null
\vskip 5mm
\null
%\end{center}
}

\autore{Rinaldo Garziera}
%\affiliazione{Dipartimento di Ingegneria e Architettura \\
%             Universit\`a di Parma }
%\titolo{Meccanica delle Macchine}
%\titolooo{ \em Cinematica}

%\pubby{MUP}
%\pubin{{\sf Parma}}
%\date{2022}
\marchireg{\noindent Stampato con un Contributo del Dipartimento di Ingegneria e Architettura dell'Universit\`a di Parma.
\newline
\newline
\noindent Il testo e le figure sono stati composti usando i 
programmi \LaTeX, \TeX, Vectornator, GNU C, Ghostscript, Gnuplot e Octave.
 }
\Copyright{\noindent {\copyright \hskip 1 mm 2023 \hskip 3mm Monte Universit\`a Parma Editore} \\ 
\vskip 2mm \noindent {\em Nessuna parte di questo libro pu\`o
essere riprodotta o trasmessa in qualsiasi forma o
con qualsiasi mezzo elettronico, meccanico o altro senza
l'autorizzazione scritta dei proprietari dei diritti e
dell'editore.}}
\isbn{978--88--7847--634--9}
\ristampa{}
%\finito{{\sf Finito di stampare a Parma. 2022.}}

\occhiello{Meccanica delle Macchine}
\occhielloo{Cinematica}
\frontespizio
\thispagestyle{empty}
\vfill
\begin{center}
\includegraphics[width=0.16\textwidth]{logoMUP.pdf}
\end{center}
\null
\newpage
\thispagestyle{empty}
\null
\newpage
\pagenumbering{gobble}
\pagestyle{headings}
\tableofcontents
\newpage
\null
\pagestyle{plain}
\newpage
\pagestyle{empty}
\hphantom{1 pt}
\vspace{1 in}
\begin{flushright}
Questo lavoro \`e dedicato al ricordo del professor Franco Giordana.
Senza i suoi consigli e il suo affettuoso  incoraggiamento
il presente volume non avrebbe avuto luce di stampa.
Un particolare ringraziamento va al professor Alessandro Tasora, per suoi preziosi suggerimenti.
\end{flushright}
\newpage
\null
\newpage
\pagestyle{headings}
\pagenumbering{arabic}

\def\thesection {\arabic{section}}
\introduction*{Note Introduttive}
\addcontentsline{toc}{chapter}{\protect\numberline{}{Note Introduttive}}
\markboth{Note Introduttive}{pagenumber}
%\endinput

\noindent Al termine {\em Meccanica} non corrisponde una sola disciplina. Il dominio
abbracciato da questa materia \`e assai vasto e una classificazione, forse leggermente
arbitraria e imprecisa, sicuramente incompleta,
dispone questa parte della {\em Fisica} in quattro branche, a loro volta molto estese:
la {\em Meccanica Classica}, la {\em Meccanica Statistica},
la {\em Meccanica Relativistica} e la {\em Meccanica Quantistica}. Questi quattro rami,
i nomi dei quali si apprendono, unitamente ai loro principali risultati,
nel corso della formazione
secondaria, sono tra loro intrecciati: \`e comune trovare, nelle librerie
scientifiche, trattati di {\em Meccanica Quantistica Relativistica}
oppure di {\em Meccanica Statistica  Relativistica}.

\noindent Prendendo le mosse da \cite{finzi},
se si volesse rispondere alla domanda ``che cos'\`e la {\em Meccanica}?''
 si potrebbe dire:
essa si propone lo scopo di descrivere i movimenti che avvengono nello spazio e nel
tempo in relazione alle cause che provocano tali moti; 
definizione che ognuno vede quanto sia ampia. Infatti, tutte
le interazioni fisiche, avendo luogo
nello spazio e nel tempo, implicano {\em movimento}. Prendendo quindi tale definizione
per  buona, la {\em Meccanica} dilagherebbe in quasi tutti i campi della {\em Fisica}. 
D'altra parte, i fenomeni naturali appartengono a tutte le branche della {\em Fisica} e,
a rigore, sta a chi si occupa di tali problemi decidere di trascurare
certuni o cert'altri aspetti. 
Quando, nello studio del movimento, sono trascurabili o vengono trascurati
gli scambi e le variazioni termiche, le relazioni elettromagnetiche,
chimiche, nucleari e non si palesano velocit\`a prossime a quelle della luce
ricadiamo, in generale, 
nell'ambito della Meccanica Classica\index{meccanica classica}.

\noindent Limitandoci quindi alla branca della Meccanica Classica avremo una copiosissima
biblioteca di opere importanti e ispiratrici, se non altro come fonti del
modello espositivo {\em standard} dell'introduzione alla {\em Cinematica},
modello che, per ragioni che cercher\`o di esporre
pi\`u avanti, non seguiremo.
\noindent Riporto soltanto due trattati di Meccanica Classica, 
ben conscio per\`o di proporre
una  porzione minuscola della letteratura effettivamente disponibile in questo campo. 
Sia \cite{goldstein} sia \cite{levicivita} entrano in argomento definendo e descrivendo il
moto di un punto. Come anticipato, non sar\`a per\`o questo il nostro {\em incipit}. Uno 
dei motivi che mi hanno persuaso a non cominciare lo studio della Meccanica
dalla cinematica del punto riguarda il fatto
che il concetto di velocit\`a angolare compare in questo studio in
modo forzato, tramite
altri oggetti, in quanto sui punti non si possono misurare angoli. Ma il motivo 
principale, che da solo potrebbe bastare a farci prendere un'altra strada, \`e che,
pur riconoscendo allo studio della cinematica del punto un notevole valore
formativo, tale studio ci sembra fuorviante dal solco maggiormente
 aderente alle applicazioni
tecniche della Meccanica, che \`e la traccia che invece
vorremmo seguire in questo spazioso campo.

\noindent Il nostro percorso, infatti, si snoder\`a tra le {\em macchine} e i loro componenti: esse saranno studiate
utilizzando i metodi di indagine, i risultati, e quanto possibile dell'impianto
concettuale e del complesso di idee che formano la Meccanica Classica,
scienza con vocazione prevalentemente matematica. Evitando
di entrare in un secondo girone di definizioni che rispondono a ``cosa sono le macchine?'',
desiderosi anzi di toglierci quanto prima da tali gineprai,
baster\`a qui dire che non esistono macchine o parti di esse che siano punti isolati.
Sicuramente la nozione di {\em punto} \`e fondamentale, intuitiva e presa volentieri
in prestito, ringraziando, dalla geometria. Ma il punto, nella nostra trattazione,
apparterr\`a sempre a un corpo o, pi\`u precisamente, a una entit\`a
che definiremo come {\em corpo rigido}
\index{corpo rigido} e che sar\`a il protagonista dei nostri moti. Partiremo quindi
dallo spostamento dei corpi, come in \cite{whittaker}. I risultati e le conclusioni
ai quali giungeremo saranno comunque in sintonia
con tutte le altre trattazioni della Meccanica.

\section{Cenni Storici}
Riprendiamo nel presente paragrafo quanto riporta \cite{finzi} nella sua introduzione, opera
insigne, questa, che l'autore ha avuto la fortuna di dover studiare per
preparare l'esame di {\em Meccanica Razionale}.
Pare che, prima dell'intervento leonardesco, pochi autori si fossero espressi,
e soltanto in considerazioni filosofiche, sulle cause del moto e sulle condizioni di
equilibrio, seguendo i resti della Filosofia Greca conservata e giunta in Europa
con l'eredit\`a della ``Scuola Alessandrina''. Tra i diversi
frammenti, a noi pervenuti da tale ``patrimonio'', segnaliamo
un'opera preziosa per chi volesse esplorare le sorgenti del
pensiero scientifico: \cite{aristotele}.
{\em Leonardo da Vinci} apr\`i il primo sentiero, per lo studio della Meccanica,
pressoch\'e scevro 
dai pregiudizi della Filosofia Antica e dalle opinioni teologiche, tanto autorevoli
e tanto in voga durante i secoli precedenti.
Il periodo moderno della Meccanica ha per\`o inizio con le nuove concezioni 
circa il moto dei pianeti elaborate da {\em Copernico} e {\em Keplero}; 
e tali idee
portarono un radicale sconvolgimento nel campo dell'astronomia.
\`E appena posteriore a questa notevolissima ``riforma astronomica'' la fondazione,
da parte di {\em Galileo}, 
del cosiddetto {\em metodo di indagine sperimentale}, che
negli ultimi tre secoli \`e stato spesso indicato come l'unica strada
sicura per il progresso della Scienza.
{\em Galileo} formula per primo la nozione precisa di
accelerazione e intuisce che questa grandezza dovr\`a comparire tra i protagonisti sul
palcoscenico della Meccanica. Nel frattempo {\em Cartesio} pone le basi
esatte per lo studio quantitativo della geometria e, siccome il
movimento dei corpi ha luogo
nelle dimensioni dello spazio, forgia alcuni strumenti
che renderanno il cammino della Meccanica pi\`u rapido,
 {\em in primis} la
{\em Geometria Analitica}.
Su queste solide fondamenta si sviluppano alcuni rami della {\em Statica},  
per opera di {\em Varignon} e {\em Stevin},
e tali studi -- senza dimenticare quelli pionieristici e ancora fondamentali
di {\em Archimede} -- costituiscono ancora oggi un notevole bagaglio di sapere tecnico.
\begin{wrapfigure}{r}{0.5\textwidth}
	  \begin{center}
	  \includegraphics[width=0.42\textwidth]{part0/FIG/mechanique.pdf}
	  \end{center}

	  {\em Il frontespizio del libro di Lagrange.}
\end{wrapfigure}
\noindent La definitiva sistemazione della legge d'inerzia (gi\`a intravista da
{\em Galileo}) da parte di {\em Newton} e il ``suo''
{\em Calcolo Infinitesimale}\footnote{
{\em Philosophi\ae Naturalis Principia Mathematica}, 1687, opera non visionata dall'autore.
}
dotano il secolo XVIII di strumenti per affrontare un ventaglio amplissimo
di problemi di meccanica. Nel corso di quel secolo, i
{\em Bernoulli} ({\em Giacomo}, {\em Giovanni} e {\em Daniele})
prima, {\em Eulero}  poco dopo, raffinando ulteriormente 
i concetti dell'Analisi Matematica,
posero le basi per quello che oggi si chiama
{\em Calcolo delle Variazioni}. Essi impostarono 
correttamente e risolsero svariati
problemi paradigmatici di statica e di dinamica. Ormai il {\em Principio dei Lavori
Virtuali} aveva ottenuto una formulazione precisa e {\em Monsieur de la Grange}, noto
poi come {\em Lagrange}, scriveva il suo bel trattato {\em M\'echanique Analitique}
(il cui frontespizio appare qui a lato) nel
quale, facendosene un curioso vanto, non inser\`i alcuna figura.
{\em Lagrange} portava cos\`i a compimento e sintetizzava  quel dovizioso
novero di metodi matematici disponibili  
alla fine del secolo XVIII.
Al termine di tale secolo, la Meccanica si trova ormai ripartita
in diversi rami dovendo fornire risposte a specifiche
necessit\`a dei vari ambiti applicativi, vertiginosamente
crescenti con l'avvento della Rivoluzione Industriale:
problemi di
{\em elasticit\`a}, di {\em idrodinamica}, di {\em suono} e {\em vibrazioni}, di
{\em meccanica delle macchine}. 

\noindent \`E proprio durante il secolo XIX che la Meccanica subisce, a livello teorico, ulteriori
e ``definitive'' sistematizzazioni dovute a {\em Laplace}, {\em Gauss}, {\em Cauchy}
e {\em Hamilton}, nonch\'e alcune importanti rivisitazioni critiche per opera di {\em 
Poincar\'e}, {\em Hertz} e soprattutto {\em Mach}. A stimolare la revisione dei
fondamenti
di questa scienza sono prevalentemente il concetto di {\em forza} e di
{\em spazio} e
{\em tempo assoluti}. {\em Mach}, nel suo prezioso libro \cite{mach}, 
prepara il terreno per la Meccanica Relativistica di {\em Einstein}, che dovr\`a 
entrare poi in quasi tutte le branche della Fisica. Consiglio agli
appassionati di meccanica, di scienza e di tecnica una lettura di \cite{mach}, se
non altro
per rendersi conto di quanto un pensiero, molto indipendente e rigoroso, possa mettere
in discussione qualsiasi teoria che, per et\`a, per diffusione, per abbondanza di
risultati, tenda al dogmatismo.

\noindent Oltre a questo, come si \`e gi\`a accennato, il secolo XIX porta alla
ribalta le esigenze della nascente Rivoluzione Industriale. La Meccanica scopre
allora la propria vocazione {\em applicativa}. Alcuni dei frutti delle ricerche svolte da
{\em Poisson}, {\em Cauchy}, {\em Carnot}, {\em Coriolis} sono squisitamente
applicativi e appartengono alla {\em Meccanica Applicata}. 
Questa, a sua volta, si \`e occupata e si occupa di un 
campo molto esteso di problematiche. Si va dalla {\em Statica} dei corpi all'{\em
Elasticit\`a}, studi che ormai spettano alla {\em Scienza delle Costruzioni},
si procede con la {\em Meccanica dei Fluidi}, di estremo interesse
per gli ingegneri civili quando il fluido \`e l'acqua e per gli ingegneri aeronautici quando
il fluido \`e l'aria. Si spazia poi dall'{\em Acustica} alle {\em Vibrazioni}, dalla
{\em Biomeccanica} alla moderna {\em Balistica}.

\noindent Giungiamo  quindi alle {\em Macchine}, con i loro organi e i loro azionamenti, 
i loro movimenti speciali, le loro 
criticit\`a e la loro straordinaria rilevanza economica.
Eccoci all'imbocco di una nuova restrizione del campo di questa
scienza; \`e l'unico strato che riguarder\`a da vicino la materia del
nostro corso. 
Anzi, nel dominio stesso della meccanica delle macchine, saremo portati o costretti
a fare delle scelte che, per forza di cose, limiteranno ulteriormente 
lo scopo di queste note.

\section{Il Nostro Approccio alla Meccanica delle  Macchine}
Questo lavoro si rivolge, in particolare, agli studenti dei corsi di
Ingegneria che prevedono l'insegnamento della Meccanica Applicata alle Macchine.
Alcuni di tali corsi non sono preceduti
da insegnamenti che abbraccino argomenti di meccanica,
fatta eccezione per l'insegnamento
della fisica del primo anno. Si rende cos\`i necessario anteporre alle
nozioni di {\em Cinematica Applicata}, contenute nella {\em Seconda Parte} e
approfondite nella {\em Terza Parte}
di questo volume, un'introduzione alla {\em Cinematica Teorica}.
La {\em Prima Parte}, infatti, richiama
i concetti di {\em corpo rigido} e di {\em moto relativo},
proponendo anche alcune nozioni circa la composizione
dei movimenti. Come gi\`a accennato, partiremo dal concetto di 
{\em velocit\`a}  e {\em accelerazione} di punti appartenenti a corpi rigidi,
dal concetto di {\em velocit\`a angolare} e di {\em moto roto-traslatorio}. 
Il {\em teorema di Coriolis}, indispensabile per la comprensione
dei {\em moti relativi} dei vari organi
delle macchine, in particolare dei {\em sistemi articolati},
ci mostrer\`a quindi la via corretta per l'interpretazione delle quantit\`a
cinematiche provenienti da differenti osservatori. 


\noindent Nella seconda parte di questo lavoro,
entrando nel campo pi\`u proprio alla meccanica
applicata e in particolare alla cinematica applicata,
ci \`e parso opportuno
presentare una panoramica sulle {\em ruote dentate}. Esse sono qui
studiate esclusivamente dal punto di vista cinematico, prendendo comunque
in considerazione tutti i loro aspetti geometrici, anche quelli connessi
con la loro
fabbricazione. L'esposizione dei concetti legati alle ruote dentate
si snoda quindi attraverso un ``percorso'' che,
in maniera induttiva, mette in luce
l'ineluttabilit\`a della effettiva forma dei denti,
che di fatto risultano profilati a evolvente di cerchio.
Sempre in questo capitolo, vengono analizzati alcuni
aspetti geometrico-funzionali rilevanti per gli ingranaggi,
{\em in primis} le loro
correzioni mediante lo spostamento del profilo.

\noindent Il capitolo che riguarda i {\em sistemi articolati} fornisce,
in ristretto, alcune indicazioni
circa la loro analisi e le possibilit\`a che essi offrono per ottenere
leggi di moto e traiettorie particolari.

\noindent Pi\`u esteso \`e il capitolo che tratta i {\em manovellismi}.
Sia per la loro importanza nel mondo industriale, sia per la loro tradizionale
presenza nei corsi di Meccanica Applicata, i manovellismi hanno meritato
in queste note uno spazio che comprende diversi metodi per la loro
analisi. Anche i manovellismi non ordinari vengono posti,
in questo capitolo,  al centro dell'attenzione, giungendo a dare conto
di qualche loro applicazione notevole.


\noindent I meccanismi a {\em camma} sono senza dubbio tra quelli maggiormente
usati nell'ambito delle macchine automatiche, soprattutto quando
sono implicati movimenti particolari, da eseguirsi con la massima precisione.
Nel capitolo dedicato alla loro profilatura,
lo studio delle leggi di moto e le tecniche per la loro
sintesi precedono pertanto
l'analisi degli aspetti cinematici
(quindi geometrici) di questi sofisticati eccentrici,
 in modo tale da fornire al lettore una discreta
base per la loro progettazione.

\noindent La terza parte di queste note \`e destinata ad alcuni approfondimenti e di essa
fanno parte tre ulteriori capitoli.
Il primo
fornisce qualche nozione circa il movimento nello spazio tridimensionale.
Il secondo capitolo mostra alcune applicazioni
delle ruote dentate con profilo cicloidale dei denti,
la pi\`u interessante delle quali riguarda
gli ingranaggi per orologi.
L'ultimo capitolo aggiunge qualche ulteriore notizia
circa la sintesi degli azionamenti a camma.
Qui si accenna alla profilatura
delle camme speciali {\em offset}, e vengono esemplificate la progettazione
di una camma tastata da un cedente a bilanciere
e quella di una camma tastata da un piattello.
\vskip .1cm

\noindent 
Il presente volume, come annuncia il suo titolo, \`e limitato
allo studio 
della {\em Meccanica Applicata alle Macchine} in un perimetro
che riguarda solamente la {\em Cinematica}. Nelle 
intenzioni dell'autore risiede per\`o la volont\`a di esporre,
in ulteriori lavori,
altre parti fondamentali di questa materia. 
Il cuore dello studio della Meccanica delle Macchine \`e rappresentato
 dall'analisi 
del loro comportamento {\em dinamico} e di tutte le problematiche legate al loro
funzionamento: dovrebbe allora essere questo il contenuto di un altro
volume, completamente
dedicato alla {\em Dinamica delle Macchine} e allo studio dei loro
principali componenti.
``Nell'accezione moderna una macchina \index{macchina} \`e qualcosa in
grado di produrre beni o servizi elaborando energia e pu\`o essere elementarmente
schematizzata come una scatola (o blocco, o {\em black box}) in cui entrano uno o
pi\`u flussi di energia (eventualmente di natura diversa) ed escono uno
o pi\`u flussi energetici (in generale di natura diversa)'', \cite{giordana}.
In questo blocco generale si riconoscono spesso importanti sotto-sistemi:
il motore, gli organi di
trasmissione e le parti che, utilizzando il lavoro meccanico fornito loro 
dai primi due, producono effetti di varia utilit\`a.
La descrizione della macchina, vista come somma di
tre sotto-sistemi, \`e ovviamente incompleta
e grossolana, e lascia sicuramente notevoli eccezioni fuori dal suo perimetro.
Tuttavia tale definizione comprende una vasta generalit\`a di casi ed \`e
sufficiente per prendere
le mosse nello studio di questa disciplina. La {\em Meccanica Applicata alle Macchine}
si propone proprio l'analisi generale di queste ultime trovando, ove possibile,
un comune denominatore tra gli approcci alle loro problematiche.
Questa disciplina ha, di fatto, anche il compito di studiare
il funzionamento di particolari sotto-sistemi presenti nella stragrande
maggioranza 
delle macchine, come {\em sistemi articolati}, {\em manovellismi},
{\em eccentrici}, {\em freni}, {\em frizioni}, {\em ruote dentate},
{\em pulegge}, {\em cinghie}, {\em catene}, {\em volani}, {\em supporti}...
Tutti questi dispositivi dovranno trovare adeguato spazio
nel volume che riguarder\`a la Dinamica.
\vskip 2mm

\noindent Un Terzo Volume dovrebbe invece
riguardare  lo studio delle  {\em Vibrazioni Meccaniche}.
Tradizionalmente, per motivi di tempo,
le lezioni che toccano questo argomento
si fermano all'analisi di sistemi discreti con pi\`u gradi di libert\`a.
Con un pizzico di ottimismo crediamo che il volume sulle vibrazioni
possa per\`o comprendere anche argomenti maggiormente specialistici,
giungendo all'analisi delle 
{\em Vibrazioni
dei Continui}, allo studio delle vibrazioni nel dominio della frequenza,
e ad alcuni aspetti legati alla loro misura.
\vskip 2mm

\noindent L'esperienza, accumulata in pi\`u di trent'anni di insegnamento,
mi ha vieppi\`u confermato nel credere in 
una distribuzione di interessi molto
variegata per questa materia.
\noindent Le parti scritte sotto forma di note, o in caratteri ridotti,
unitamente alla parte terza del volume,
integrano
la sintetica e spesso mutilata presentazione di alcuni argomenti
e potranno essere lette dagli studenti pi\`u volenterosi
a complemento delle nozioni
impartite frontalmente.
\vskip .2cm
\noindent Le figure che compaiono nella seconda e nella terza parte del volume,
eccezion fatta per una decina di esse, sono state prodotte da 
codici numerici, sviluppati integralmente dall'autore per tale proposito.
Non avendo in mente fini didattici o di ricerca, lo stile dei programmi (se
cos\`i si pu\`o chiamare)
\`e molto rudimentale, con pochi commenti, i quali contengono
un'infinit\`a di refusi. Ciononostante, dato che tali codici
danno origine a figure che appaiono plausibili, ci sembrava un peccato
non mettere l'intero {\em software} a disposizione di qualche studente
curioso. Il tutto si trova all'indirizzo {\tt gitlab.com/public}
dove si deve cercare il nome dell'autore, il quale sar\`a grato a
chiunque lasci qualsiasi osservazione in merito.
\vskip 2mm
\noindent Abbiamo ritenuto opportuno
fornire una bibliografia di dimensioni ragionevoli,  citata
nel testo laddove necessario. La scelta, che potrebbe apparire arbitraria,
delle opere citate aderisce al principio secondo il quale chi scrive dovrebbe
avere confidenza e dimestichezza
con ci\`o che propone come fonti, o come esempi,
 insomma con le opere citate. Talvolta uniremo alla citazione
brevi commenti circa l'opera indicata e il suo contenuto, 
con la speranza di suscitare la curiosit\`a attiva di almeno uno studente
ogni due anni.
\newpage
\thispagestyle{empty}
\null

\def\thesection {\thechapter.\arabic{section}}


\endinput


\part{Corpi Rigidi e loro Moti Relativi}
\chapter{Il Corpo Rigido: Spostamenti e Velocit\`a}

\section{Introduzione}

\begin{wrapfigure}{r}{0.5\textwidth}
     \begin{center}
     \includegraphics[width=0.42\textwidth]{part1/cinematica/FIG/f11.jpeg}
     \end{center}
\vskip -7.3mm
	\caption{\em Biella di motore d'automobile.}\index{biella}
     \label{fig:f11}
\end{wrapfigure}

Le macchine sono costituite da un certo numero di parti, molte delle
quali 
sono corpi rigidi. {\em Bielle},  \index{biella} {\em manovelle}, \index{manovella}
{\em aste}, {\em bilancieri}, \index{bilanciere}{\em ruote dentate},
\index{ruote!dentate} {\em telai}, \index{telaio} sono tutti {\em corpi rigidi}. 
Per definizione, un corpo \`e rigido \index{corpo rigido} se risulta essere 
invariante nel tempo la distanza tra due suoi punti scelti ad arbitrio. Siano $P_1$ e $P_2$
punti appartenenti al nostro corpo: perch\'e tale corpo sia rigido deve essere
\begin{equation}
	d(P_1, P_2, t_1) = d(P_1, P_2,t_2)\,,
	\label{e11}
\end{equation}
\noindent dove $d(.,.,.)$ \`e la funzione distanza tra due punti al tempo $t$.
Come vedremo tra poco, la
definizione di corpo rigido appena introdotta
risulter\`a essere di grande utilit\`a
nello studio della cinematica, in quanto i corpi rigidi possono dare luogo
esclusivamente a {\em spostamenti rigidi}\footnote{
	Non \`e vero il contrario: anche corpi deformabili come gas e liquidi possono muoversi rigidamente. 
La condizione \ref{e11} fa s\`i che un tetraedro, formato prendendo
come vertici quattro punti qualsiasi del corpo, rimanga invariato
nella forma
anche dopo un eventuale spostamento.
Non \`e difficile immaginare una
bolla d'aria all'interno di un liquido che scorra in
una condotta trasparente (in modo da poter essere osservata), la quale si muova
senza cambiare forma.
}.
\noindent La definizione enunciata poc'anzi ci offre, peraltro, l'opportunit\`a di riflettere sui modelli che mettiamo
in campo per utilit\`a nello studio dei fenomeni naturali: in questo caso
si intuisce che nessun corpo ``reale'' potr\`a
rispettare rigorosamente la \ref{e11},
dato che non esistono corpi indeformabili.
Ma il modello, cio\`e la rappresentazione
ideale e parziale, che si deduce dalla \ref{e11} indica anche la via che si desidera
percorrere nell'indagine scientifica: il corpo rigido \`e utile per
lo studio della cinematica
degli organi delle macchine, dove risulta comodo e appropriato (ma questo si deve verificare 
a posteriori) trascurare le elasticit\`a.
Tale modello non sarebbe affatto opportuno 
qualora il campo di ricerca
fosse quello della resistenza dei materiali, dove gli oggetti considerati sono
posti sotto carico
e subiscono delle variazioni (di norma piccole) della loro forma:
in questo caso il modello del {\em corpo rigido} non
solo sarebbe fuori luogo, ma risulterebbe un 
ostacolo insormontabile nel considerare le deformazioni che subiscono tali corpi.
La biella di un motore endotermico, figura \ref{fig:f11}, sar\`a considerata come un
{\em corpo rigido} nello studio della cinematica dei manovellismi, ma 
nessuno negher\`a che tale biella si deformer\`a sotto gli effetti
del carico laterale dovuto alle sue accelerazioni,
come si studia di fatto
nella disciplina che ha nome {\em Costruzione di Macchine}\footnote{
Questa nota vorrebbe
invitare lo studente a una riflessione circa la natura, 
l'utilit\`a e i limiti dei modelli che si impiegano ai fini di studiare
e di prevedere il comportamento, anche oltre l'immediato
presente, degli oggetti e dei fenomeni
che popolano il mondo circostante. La {\em Teoria dei Modelli}\index{teoria
dei modelli}
\`e una scienza a s\'e stante, dunque questa digressione
non ha alcuna pretesa
di completezza n\'e di profondit\`a.

\noindent Modellare il comportamento di un sistema reale
significa tentare un parallelo tra la realt\`a, complessa e 
ricchissima di dettagli, e altre strutture (i modelli appunto), talvolta
matematiche, altre
volte informatiche, oppure anche esse fisiche che, 
essendo intrinsecamente pi\`u sintetiche e meno doviziose di particolari,
probabilmente inessenziali,
risultano pi\`u comode da indagare e spesso feconde di risultati.
Tali risultati potranno essere considerati validi per il sistema reale soltanto
se confrontabili positivamente
con l'esperienza, nel perimetro
di tutti gli aspetti che nel modello sono stati considerati.
Ad esempio, sembrerebbe strano poter ricavare da un {\em modello urbanistico
fisico},
cio\`e da quelle tavole di legno compensato su cui gli architetti costruiscono
le loro citt\`a in scala ridotta, ricavare, dicevo, informazioni riguardanti
la risposta sismica dei fabbricati; sarebbe impossibile persino scuotendo
energicamente la tavola.

\noindent I modelli matematici e numerici sono ancora pi\`u austeri: fecondi di risultati, semplici
e sicuri, ma intrinsecamente refrattari a essere infarciti di particolari.
Trascurare il colore di un oggetto reale che si modella con il {\em metodo
degli elementi finiti} \`e un comportamento che tutti i tecnici 
approvano, e con ragione. Ma la legittimit\`a di questo comportamento, cio\`e
di trascurare ``arbitrariamente'' nei modelli alcune caratteristiche della
realt\`a, necessita comunque del conforto sperimentale o con modelli di altra
natura.

\noindent I modelli insomma spogliano la realt\`a di molti dettagli e caratteristiche
ritenuti inutili ai fini di prevedere ``ci\`o che accadr\`a'': 
l'impatto urbanistico di un nuovo quartiere o la frequenza delle
oscillazioni del pendolo di una certa lunghezza.
Quanto pi\`u
la realt\`a ``alleggerita'' che si ottiene da un modello
\`e semplice, priva di 
informazioni inessenziali, tanto pi\`u sar\`a
facile interpretarne i risultati,
sempre da verificare, limitatamente all'ambito per il quale il modello \`e stato pensato.
}.

\section{Spostamenti e Velocit\`a}

\noindent {\em Posizione} \index{posizione} e {\em velocit\`a} \index{velocit\`a}
appaiono a chiunque come concetti ampiamente
intuitivi. Del resto siamo abituati a conoscere la velocit\`a
dei nostri spostamenti in automobile semplicemente
leggendo il {\em tachimetro} \index{tachimetro} del cruscotto.
Per quel che riguarda poi la nostra posizione, avvertiamo la sensazione 
che tutto ci\`o che ci circonda funga da {\em sistema di riferimento}
naturale. Sappiamo, in generale, renderci conto senza sforzo della nostra
posizione: ci troviamo a Milano, nella Piazza del Duomo, giusto all'ingresso...
Posizione e velocit\`a sono peraltro a disposizione di tutti gli
utenti di dispositivi dotati di GPS \index{GPS}({\em Global Positioning System}). 

\noindent Qualche difficolt\`a emerge qualora, ad esempio, si volesse conoscere la
velocit\`a del piede di un calciatore nell'istante in cui calcia il pallone.
Tale difficolt\`a nasce dall'intuizione che la velocit\`a del piede del
calciatore
difficilmente coincider\`a con la velocit\`a del calciatore stesso, cio\`e
con quella che potrebbe
essere dedotta da un orologio con GPS indossato dall'atleta.
Allo stesso modo nessuno troverebbe corretto
identificare la velocit\`a delle dita di un pianista durante un concerto
con la velocit\`a del pianista stesso, il quale, di solito, sta seduto su di
uno sgabello.
Eccoci venuti al punto: i corpi, gli oggetti, sono i protagonisti
del movimento ma i concetti 
di {\em spostamento}\index{spostamento} e di {\em velocit\`a}\index{velocit\`a}
si mostrano estremamente adatti a essere applicati a punti precisi di tali oggetti
in movimento. Locuzioni come ``in questo istante
il modulo della velocit\`a della punta dell'indice del
pianista \`e di $0.5 m/s$'' assumono immediatamente carattere 
di indicazione molto accurata, ammesso che la punta dell'indice
sia un riferimento tanto preciso da renderci soddisfatti.
\noindent \`E subito chiaro che,
seguendo la strada sulla quale ci siamo incamminati, non ci accontenteremo
di punta delle dita o di altri riferimenti approssimativi, perci\`o
diremo che
{\em spostamenti} e {\em velocit\`a} appartengono esclusivamente ai punti dei
corpi mobili, e sono pertinenti (ma questo si sapeva gi\`a) soltanto agli
istanti di tempo ai quali vengono riferiti. 
Nel dire tutto questo, il linguaggio della matematica e della geometria
ci pu\`o venire in grande aiuto.
La figura \ref{fig:f12} mostra un {\em corpo rigido} {\bf C} che ha sub\`ito
uno spostamento rigido e {\em piano}.
Quest'ultimo aggettivo aggiunge un'informazione
importantissima e restrittiva: tutti i {\em vettori spostamento} degli infiniti
punti di tale corpo saranno contenuti in piani fra loro paralleli.
Per esempio, limitando la presente introduzione
ai soli {\em moti piani}\index{moti piani}, il concetto di {\em velocit\`a angolare}, che introdurremo
a breve, sar\`a  pi\`u semplice rispetto al caso generale 
di moto rigido tridimensionale, di conseguenza lo sar\`a l'intera trattazione dei
moti piani.
Questa scelta limita non poco la portata del presente lavoro e non sarebbe 
giustificata dalla sola semplificazione dei concetti cinematici
che ne consegue,
se 
la maggior parte dei moti che si verificano nelle macchine non fosse
di questa specie, cio\`e {\em moti piani}.
Si possono comunque rilevare molte eccezioni: le viti di azionamento e di manovra, le
ruote sterzanti (e non) dei veicoli, ecc. Ma i casi pi\`u importanti di moti 
non piani, attinenti la meccanica delle macchine, si verificano
nell'ambito dei {\em robot}.
Vi \`e per\`o da dire che la cinematica tradizionale
che considera i moti
tridimensionali, riportata comunque a grandi linee e in ristrettissimo
compendio nel capitolo \ref{3d},
si mostra piuttosto inefficiente per lo studio della 
{\em Meccanica dei Robot}. Infatti, in quest'ambito \`e d'uopo
 una trattazione algebrica {\em ad hoc},
basata su specifiche matrici $4\times4$, \cite{denavit},
e questo ci fa sentire meno in colpa.
\begin{wrapfigure}{r}{0.45\textwidth}
      \begin{center}
      \includegraphics[width=0.42\textwidth]{part1/cinematica/FIG/f12.pdf}
     \end{center}
\begin{picture}(50,0)(-40,0)
\scriptsize{
	\put(31,140){$2)$}
	\put(100,98){$1)$}
\put(12,63){$\Delta \theta$}
\put(-14,98){\bf C}
\put(47,83){\rotatebox{-43}{\bf C}}
\put(53,107){$\Delta{\bm{s}}$}
	\put(4,122){$P_{(2)}$}
	\put(82,73){$P_{(1)}$}
}
\end{picture}
\vskip -7mm
	\caption{\em Spostamento piano di un corpo rigido nel tempo $\Delta t$.}
     \label{fig:f12}
\end{wrapfigure}

\noindent Nella figura \ref{fig:f12} \`e rappresentato lo spostamento {\em piano} del corpo rigido
{\bf C}. Assumiamo che lo spostamento $\Delta{\bm{s}}$,
dalla posizione 1) alla posizione 2), avvenga 
nell'intervallo di tempo $\Delta t$.
Il vettore $\Delta \bm{s}$, che congiunge i punti $P_{(1)}$ e $P_{(2)}$, si chiama
{\em spostamento} \index{spostamento} del punto $P$ nel 
tempo $\Delta t$.
Nello stesso intervallo temporale $\Delta t$ il corpo avr\`a anche mutato orientamento.
Chiamiamo
{\em rotazione} \index{rotazione} $\Delta \theta$ l'angolo compreso
tra una qualsiasi direzione distinguibile sul corpo {\bf C} nella posizione
1) (come pu\`o
essere uno spigolo o un segmento tracciato sul corpo) e lo stesso riferimento,
quando il corpo si trova nella posizione 2).
I rapporti
\begin{equation}
	\lim_{\Delta t \to 0}\;{\Delta {\bm{s}}\over {\Delta t}} = {\bm{v}}\;\;\;\;\;\; \textrm{e}
	\;\;\;\;\;
	\lim_{\Delta t \to 0}\;{{\Delta \theta }\over {\Delta t}} = {\omega}\,, \label{e12}
\end{equation}	
\noindent si chiamano rispettivamente {\em velocit\`a}\index{velocit\`a}
del punto $P$ e {\em velocit\`a angolare}\index{velocit\`a!angolare}
del corpo {\bf C}.
Il lettore con qualche reminiscenza di cinematica
potrebbe chiedersi come mai
non sia stato introdotto alcun {\em sistema di
riferimento}\index{sistema di
riferimento}. 
\`E un fatto che nella maggior parte dei libri moderni di cinematica
si usi spesso riferire gli spostamenti a ``ragionevoli sistemi di riferimento'',
il pi\`u delle volte rappresentati da piccoli {\em assi cartesiani}
disegnati nelle figure nell'angolo in basso a sinistra, come anche a noi accadr\`a di fare pi\`u avanti.
Riteniamo qui superflua la presenza di assi coordinati e, come
abbiamo mostrato, siamo in grado di identificare spostamenti, rotazioni e velocit\`a
senza basarci su di essi.
Ricordiamo per\`o che un sistema di riferimento, almeno implicito, deve
esistere: per noi pu\`o essere la pagina stessa del libro.
A un osservatore (privo di massa) seduto sul corpo {\bf C},  isolato da
tutto il resto dell'universo, sarebbe infatti impossibile rendersi
conto di qualsiasi movimento del corpo stesso.

\begin{figure}[h]
\centering
\includegraphics[width=0.9\textwidth]{part1/cinematica/FIG/f13.pdf}
\begin{picture}(0,0)(290,13)
\scriptsize{
\put(14,126){$\Delta \theta$}
\put(64,122){$A'$}
\put(-15,95){${\Delta{\bm s}}_{\scriptscriptstyle{A}}$}
\put(33,88){$B'$}
\put(-37,81){$A$}
\put(106,80){$\Delta \theta$}
\put(15,66){${\Delta{\bm s}}_{\scriptscriptstyle{B}}$}
\color{gray}
\put(56.1,76){$C$}
\color{black}
\put(-14,55){$1)$}
\put(60,55){$2)$}
\put(-9,32){$B$}
\put(70,25){$C'$}
\put(30,15){$a)$}
}
\end{picture}
\begin{picture}(0,0)(117,13)
\scriptsize{
\put(16,131){$\Delta \theta$}
\put(65,116){$A'$}
\put(-15,95){${\Delta{\bm s}}_{\scriptscriptstyle{A}}$}
\put(33,88){$B'$}
\put(-37,81){$A$}
\put(94,89){$\Delta \theta$}
\put(14,66){${\Delta{\bm s}}_{\scriptscriptstyle{B}}$}
\color{gray}
\put(56.3,76){$C$}
\color{black}
\put(-14,55){$1)$}
\put(60,55){$2)$}
\put(-9,32){$B$}
\put(70,25){$C'$}
\put(30,15){$b)$}
}
\end{picture}
	\caption{\em
I due casi riportano la riduzione del medesimo spostamento, sub\`ito
da un corpo rigido,
alla combinazione di due diverse traslazioni con la stessa rotazione
attorno a due diversi poli.
}
\label{fig:f13}
\end{figure}

\section{Teorema di Chasles}\index{Chasles, teorema di}\index{teorema!di Chasles}
La seguente osservazione \`e di importanza fondamentale: lo spostamento
di un corpo {\bf C} \`e inequivocabilmente individuato
dallo spostamento di un qualsiasi suo punto $P$
unitamente alla rotazione dello stesso corpo
{\bf C}. 
Per convincerci di ci\`o consideriamo la 
figura \ref{fig:f13} a) dove il corpo rigido {\bf C} \`e stato sostituito,
per comodit\`a di esposizione,
dal triangolo $\triangle{ABC}$. Il suo spostamento
dalla posizione 1) alla posizione 2), che muove il triangolo $\triangle ABC$
in $\triangle {A'B'C'}$, pu\`o
essere ottenuto mediante la ``somma'' di due successivi spostamenti elementari.

\begin{wrapfigure}{r}{0.42\textwidth}
      \begin{center}
      \includegraphics[width=0.35\textwidth]{part1/cinematica/FIG/f14.pdf}
     \end{center}
\begin{picture}(50,0)(-40,0)
\scriptsize{
	\put(0,25){$1)$}
	\put(76,25){$2)$}
\put(17,133){$\Delta \theta$}
%\put(-12,75){\bf C}
%\put(60,100){\rotatebox{-78}{\bf C}}
\put(-28,88){$A$}
\put(63,136){$A'$}
\put(0,48){$B$}
\put(32,105){$B'$}
\put(79,40){$C'$}
\put(40,58){$O$}
\put(47.6,64.2){$.$}
\put(-15,105){$\Delta {\bm s}_{\scriptscriptstyle{A}}$}
\put(1,74.5){$\Delta {\bm s}_{\scriptscriptstyle{B}}$}
}
\end{picture}
\vskip -5mm
	\caption{\em Riduzione generica di uno spostamento piano a una
	rotazione.}
     \label{fig:f14}
\end{wrapfigure}

\noindent Supponiamo dapprima di spostare
il vertice $A$ in $A'$ mediante una {\em traslazione}\index{traslazione}
del triangolo $\triangle{ABC}$, cio\`e mediante uno spostamento che mantenga inalterato
l'orientamento primitivo dei tre lati.  Successivamente ruotiamo il triangolo
attorno al punto $A'$ di un angolo pari a $\Delta \theta$, portandolo a coincidere
col triangolo $\triangle A'B'C'$.
Risulta subito chiaro che la scelta del punto $A$ come base per 
la traslazione del triangolo \`e arbitraria e che in sua vece si sarebbe potuto scegliere $B$, $C$, o qualsiasi
altro punto, dentro o fuori dal perimetro del triangolo, purch\'e legato a muoversi
rigidamente col triangolo stesso. In figura \ref{fig:f13} b) viene rappresentato lo stesso spostamento sub\`ito
dal triangolo $\triangle ABC$ assumendo come base per la traslazione il punto $B$.
Naturalmente, attribuendo 
la traslazione a punti diversi si avranno per essi spostamenti
$\Delta {\bm s}$ diversi, come si evince peraltro dalla figura.
Quando la rotazione $\Delta \theta$ non \`e nulla
\`e possibile
individuare un punto, che potrebbe
non essere incluso nel perimetro del triangolo bens\`i appartenere a un {\em piano mobile}\index{piano mobile}
idealmente incollato su di esso, per il quale lo spostamento $\Delta {\bm{s}}$ \`e nullo.
Tale punto, che indichiamo con $O$, si individua
tracciando dalla mezzeria dei segmenti orientati che
rappresentano gli spostamenti 
{$\Delta {\bm s}_{\scriptscriptstyle{A}}$} e
{$\Delta {\bm s}_{\scriptscriptstyle{B}}$} le loro perpendicolari,
come in figura \ref{fig:f14}: il loro punto di incontro 
\`e il {\em centro dello spostamento rigido piano}.
\`E facile dimostrare (e pertanto lasciato al lettore come esercizio\footnote{
	Si noti che  i due triangoli $\triangle{ABO}$ e $\triangle{A'B'O}$
sono uguali.}) che 
\begin{equation}
	\widehat{AOA'} = \widehat{BOB'} = \widehat{COC'} =\Delta \theta \,.
	\label{e13}
\end{equation}	
\noindent Cos\`i, una rotazione rigida
di un angolo $\Delta \theta$ attorno al punto $O$
porter\`a il triangolo dalla posizione 1) alla posizione 2).
Dalla precedente affermazione si pu\`o dedurre la seguente restrizione
al caso di spostamenti rigidi piani del {\em teorema di 
Chasles}: {\em ogni spostamento rigido piano
si pu\`o ridurre in modo univoco
o a una traslazione o a una rotazione}\footnote{
Ricordiamo questo teorema
come riportato in \cite{finzi}:
``Ogni spostamento rigido, non traslatorio,
si riduce in infiniti modi a uno
spostamento roto-traslatorio, sempre per\`o con la stessa componente di
rotazione. Si riduce poi in modo
univoco a uno spostamento elicoidale''.
Torneremo brevemente nel capitolo \ref{3d}
agli spostamenti tridimensionali, col tentativo
di giustificare, anche in questi casi,
l'esistenza di una direzione principale che \`e l'asse
dell'elica.}.  \label{chasles}

\begin{wrapfigure}{r}{0.54\textwidth}
      \begin{center}
      \includegraphics[width=0.35\textwidth]{part1/cinematica/FIG/f15.pdf}
     \end{center}
\begin{picture}(0,0)(-63,1)
\scriptsize{
        \put(-16,178){${\rm d} \theta$}
\put(-17,131){$P$}
\put(40,40){$Q$}
\put(86,61){$O$}
\put(87.5,69.7){$.$}
        \put(-13,153){${\rm d}{\bm s}_{\scriptscriptstyle{P}}$}
        \put(49,77){${\rm d}{\bm s}_{\scriptscriptstyle{P}}$}
        \put(20,54.5){${\rm d}{\bm s}_{\scriptscriptstyle{Q}}$}
        \put(5,26){${\rm d}{\bm \theta} \times \overrightarrow{PQ}$}
}
\end{picture}
        \caption{\em Composizione di spostamenti infinitesimali.}
     \label{fig:f15}
\end{wrapfigure}

\noindent Quanto sin qui riferito circa gli spostamenti piani di un corpo rigido
risulta essere indipendente dalla entit\`a degli spostamenti stessi. Perci\`o le medesime conclusioni resteranno valide anche qualora avessimo a che fare con
spostamenti infinitesimali anzich\'e finiti\footnote{Non \`e vero il contrario: si incorre spesso in errore
trasportando i risultati dedotti considerando spostamenti infinitesimali nel
campo degli spostamenti finiti.}.
Riferendoci alla figura \ref{fig:f15}, possiamo quindi affermare che lo spostamento infinitesimale di un punto $Q$ del triangolo mobile,
{${\rm d} \bm{s}_{\scriptscriptstyle{Q}}$}, 
si possa ottenere sommando lo spostamento infinitesimale
del punto $P$,
{${\rm d} \bm{s}_{\scriptscriptstyle{P}}$}, 
all'effetto che la  rotazione ${\rm d} \theta$, anch'essa infinitesimale,
esercita sul segmento $Q - P$.
E qual \`e questo effetto? Cio\`e: la rotazione ${\rm d} \theta$ applicata
al {\em segmento orientato}\index{segmento orientato} $Q - P$ quale
spostamento produce? Conviene qui fare ricorso
a un ``arnese della matematica'', forse un po' arcigno, 
ma che  risulta troppo comodo
per essere lasciato in disparte: il {\em prodotto vettoriale}\index{prodotto
vettoriale}. Consideriamo quindi la rotazione infinitesimale
${\rm d} \bm{\theta}$\footnote
{
Il vettore ${\rm d} \bm{\theta}$ non \`e, nei casi di moto tridimensionale, il
differenziale di alcun angolo finito di rotazione,
come invece lo \`e nel piano. Rimandiamo
il lettore agli approfondimenti del capitolo \ref{3d} per ulteriori chiarimenti.
}
 come un 
vettore ortogonale al piano dove avviene lo spostamento, senza precisarne
la {\em giacitura}\footnote{Autori importanti, \cite{whittaker} e
\cite{finzi} per esempio, specialmente quando parlano di questo vettore
col senso di velocit\`a angolare (come faremo anche noi tra poco) 
gli  attribuiscono una precisa giacitura, coincidente con
l'{\em asse istantaneo di rotazione}. Ritenendo questa posizione
sterile se non fuorviante, bench\'e non scorretta, preferiamo lasciare tale
vettore
libero di giacere dove vuole.}, col verso dato dalla cosiddetta {\em regola
della mano destra}\index{regola della mano destra} secondo cui il pollice
indica il verso del vettore quando le dita concordano con il senso della
rotazione.
Se al segmento orientato $Q - P$ diamo valore di vettore con 
verso da $P$ a $Q$ si ricava lo spostamento infinitesimale del punto $Q$ come
\begin{equation}
	{\rm d} \bm{s}_{\scriptscriptstyle{Q}} =
	{\rm d} \bm{s}_{\scriptscriptstyle{P}}
	+{\rm d} \bm{\theta} \times \overrightarrow{PQ}\,.
	\label{e14}
\end{equation}	

\vskip 3mm
\noindent Dividendo tutti i termini della \ref{e14} per ${\rm d} t$ otteniamo la seguente 
equazione 

\begin{equation}
	{{ {\rm d} \bm{s}_{\scriptscriptstyle{Q}}}\over{{\rm d} t}} =
	{{ {\rm d} \bm{s}_{\scriptscriptstyle{P}}}\over{{\rm d} t}} +
	{{{\rm d} \bm{\theta}}\over{{\rm d}t}} \times \overrightarrow{PQ}\,,
	\label{e15}
\end{equation}	

\vskip 3mm
\noindent che, con ovvio significato dei simboli, diventa 

\begin{equation}
	{  \bm{v}_{\scriptscriptstyle{Q}}} =
	{  \bm{v}_{\scriptscriptstyle{P}}} +
	{ \bm{\bm \omega}} \times \overrightarrow{PQ}\,.
	\label{e16}
\end{equation}	
\vskip 3mm
\noindent La conoscenza della velocit\`a di un punto e della velocit\`a angolare del corpo,
oppure della velocit\`a angolare e del centro $O$, ci permettono di 
individuare le velocit\`a di tutti i punti del corpo che ne costituiscono il 
suo {\em atto di moto}\index{atto di moto}.
Il centro $O$, nel caso di spostamenti infinitesimali, si chiama
{\em centro istantaneo di rotazione}\index{centro!istantaneo di rotazione} o
{\em centro d'istantanea rotazione}\index{centro!di istantanea rotazione}\footnote{
Se tenessimo traccia delle posizioni via via occupate dal centro d'istantanea
rotazione col procedere del movimento, otterremmo due diverse curve a seconda 
che tale traccia sia lasciata sul piano fisso (il foglio del disegno per
intenderci), oppure sul piano mobile (un piano idealmente incollato col nostro
corpo rigido). Questi due {\em luoghi geometrici} si chiamano rispettivamente
{\em polare fissa}\index{polare!fissa} e {\em polare mobile}\index{polare!mobile}. Esse godono di propriet\`a interessantissime, soprattutto
dal punto di vista della nostra materia, e risultano indispensabili
in alcune analisi cinematiche, come accade nello studio dei profili dei denti
delle ruote dentate. Riteniamo opportuno
rimandare il breve studio che faremo delle polari del moto proprio
al capitolo \ref{cap_ruote_ev},
dove esporremo i principi di progettazione delle ruote dentate.
Segnaliamo agli interessati, che volessero approfondire l'argomento,
 \cite{finzi}, pagg. 168--176
e \cite{sesini1},  pagg. 31--33.}, ed \`e in generale mobile.

\noindent Abbiamo usato la rotazione infinitesimale
${\rm d}{\bm \theta}$ e di conseguenza la velocit\`a angolare $\bm \omega$ come 
vettori ortogonali al piano del moto per i quali il loro {\em prodotto vettoriale} 
con un segmento orientato del piano del moto fornisce uno
spostamento infinitesimale oppure una velocit\`a.

\noindent Vogliamo aggiungere
che la rappresentazione vettoriale delle velocit\`a angolari
non si limita al caso di moti piani, anzi, essa ha origine proprio
nella cinematica dei moti tridimensionali\footnote{Nel caso di moto piano, 
il prodotto vettoriale che coinvolge le rotazioni avviene sempre tra vettori
ortogonali tra loro,
mentre nel caso generale questo non si verifica. Le rotazioni infinitesime
e le velocit\`a
angolari  sono per\`o vettori
particolari. Il loro verso dipende dalla scelta che si fa ``destrorsa''
o ``sinistrorsa'' nella convenzione del prodotto vettoriale e mal
sopportano deformazioni dello spazio che alterino la sua metrica. In pi\`u, la 
possibilit\`a di rappresentare una velocit\`a angolare con un vettore \`e
limitata allo spazio tridimensionale.
},
come esposto in breve nella conclusione del capitolo \ref{3d}.

\chapter{Accelerazioni}

\section{Derivazione della Velocit\`a}

\begin{wrapfigure}{r}{0.43\textwidth}
      \begin{center}
      \includegraphics[width=0.3\textwidth]{part1/cinematica/FIG/f16.pdf}
     \end{center}
\begin{picture}(0,0)(-30,-25)
\scriptsize{
	\put(60,32){${\rm d}\theta$}
	\put(5,84){${\rm d}\bm \theta$}
	\put(32,29){$\bm u$}
	\put(35,48){${\bm u}+{\rm d}{\bm u}$}
	\put(93,22){${\rm d}{\bm u}={\rm d}{\bm \theta} \times {\bm u}$}
}
\end{picture}
	\caption{\em Rotazione di un vettore costante.}
     \label{fig:f16}
\end{wrapfigure}
 
L'accelerazione del punto $Q$ di figura \ref{fig:f15} si pu\`o ottenere
derivando la sua velocit\`a rispetto al tempo.
Derivando perci\`o la \ref{e16} otteniamo

\begin{equation}
	{  \bm{a}_{\scriptscriptstyle{Q}}} =
	{  \bm{a}_{\scriptscriptstyle{P}}} +
	\dot{\bm \omega} \times \overrightarrow{PQ}+
	{\bm \omega} \times \dot{{\overrightarrow{PQ}}}\,.
	\label{e17}
\end{equation}

\noindent Concentriamoci sull'ultimo termine
a destra della \ref{e17}: si tratta di eseguire il prodotto
vettoriale tra $\bm \omega$ e il vettore ``derivata rispetto
al tempo di $\overrightarrow{PQ}$''.
Mediante l'ausilio della figura \ref{fig:f16}, la quale vorrebbe essere un
tentativo (forse maldestro) di mostrare il vettore rotazione infinitesimale
${\rm d}{\bm \theta}$ ortogonale al piano dove avvengono gli spostamenti,
affermiamo che la  derivata di un vettore $\bm u$ di modulo costante, come lo \`e anche 
$\overrightarrow{PQ}$,
risulta dalla sola variazione della sua posizione angolare:

\begin{equation}
{{{\rm d}{\bm u}}\over{{\rm d}t}}=
{{\rm d}{\bm \theta} \times 
{\bm u}\over{{\rm d}t}}=
{{\rm d}{\bm \theta}
\over{{\rm d}t}} \times {\bm u}\,.
	\label{e18}
\end{equation}

\noindent Pertanto la \ref{e17} diventa

\begin{equation}
{{\bm a}_{\scriptscriptstyle{Q}}} =
{{\bm a}_{\scriptscriptstyle{P}}} +
\dot{\bm \omega} \times \overrightarrow{PQ}+
{\bm \omega}\times({\bm \omega} \times {{\overrightarrow{PQ}}})\,,
	\label{e19}
\end{equation}
\noindent relazione di grande importanza 
nella cinematica piana\footnote
{
Troveremo di nuovo la \ref{e19} 
e la precedente espressione delle velocit\`a \ref{e16}
nel paragrafo che tratta i {\em moti relativi}. Di fatto queste
espressioni di velocit\`a e accelerazione del punto $Q$ che si ottengono una
volta note la velocit\`a e l'accelerazione di un punto del sistema
mobile, unitamente alla
velocit\`a e all'accelerazione angolari, permettono
il passaggio dalle grandezze cinematiche ``relative'' a quelle ``assolute''.
 Preferiamo per\`o non
mescolare i due problemi, anche se la differenza riguarda la forma e non la sostanza. Per inciso,
le due espressioni citate sopra costituiscono il
{\em teorema di Rivals},\index{Rivals, teorema di}\index{teorema!di Rivals} \cite{sesini1}, pag. 12.
}.	

\begin{wrapfigure}{r}{0.50\textwidth}
      \begin{center}
      \includegraphics[width=0.25\textwidth]{part1/cinematica/FIG/f17.pdf}
     \end{center}
\begin{picture}(0,0)(-48,-30)
\scriptsize{
\put(6,180){${\bm a}_{\scriptscriptstyle{P}}$}
\put(2,145){$P$}
\put(30,85){${\bm a}_{\scriptscriptstyle{Q}}$}
\put(-3,77){${\bm a}_{\scriptscriptstyle{P}}$}
\put(53,83){${\bm a}_{\scriptscriptstyle{{n}_{Q(P)}}}$}
\put(78,17){$Q$}
\put(41,-2){${\bm a}_{\scriptscriptstyle{{t}_{Q(P)}}}$}
}
\end{picture}
\vskip -7.3mm
	\caption{\em Le componenti dell'accelerazione del punto $Q$.}
     \label{fig:f17}
\end{wrapfigure}
\noindent Il termine
${\bm a}_{\scriptscriptstyle{P}}$
non \`e altro che l'accelerazione del punto $P$.
Il secondo termine,

\begin{equation}
{\bm a}_{\scriptscriptstyle{{t}_{Q(P)}}}=
{\dot{\bm \omega}} \times {{\overrightarrow{PQ}}}\,,
\label{e110}
\end{equation}

\noindent si chiama
{\em accelerazione tangenziale}\index{accelerazione!tangenziale}
del moto di $Q$ attorno a $P$;
il terzo termine,

\begin{equation}
{\bm a}_{\scriptscriptstyle{{n}_{Q(P)}}}=
{\bm \omega}\times({\bm \omega} \times {{\overrightarrow{PQ}}})\,,
\label{e111}
\end{equation}

\noindent  rappresenta l'{\em accelerazione normale}\index{accelerazione!normale}
(a quella tangenziale)
del moto di $Q$ attorno a $P$.
La figura \ref{fig:f17}
dovrebbe chiarire i nomi di questi ultimi due termini.
Tramite la {\em scomposizione del doppio prodotto vettoriale}\index{doppio prodotto vettoriale}\footnote{\label{fig:foot110}
Siano ${\bm u}$, ${\bm v}$, e ${\bm w}$ tre 
vettori in $\mathbb{R}^3$, vale la relazione
\begin{equation}
{\bm u}\times({\bm v}\times{\bm w})=
{\bm v}({\bm u}{\bm w})-
{\bm w}({\bm u}{\bm v})\,.
\label{e2f}
\end{equation}
} 
della \ref{e111} possiamo riscrivere il valore della componente 
normale dell'accelerazione nel modo seguente,
valida nel caso di moti piani
\begin{equation}
{\bm a}_{\scriptscriptstyle{{n}_{Q(P)}}}= -\omega^2 {{\overrightarrow{PQ}}}\,.
\label{e112}
\end{equation}
\noindent La \ref{e19} pu\`o quindi essere
cos\`i riscritta\footnote
{
Ci possiamo
porre la domanda seguente: esiste un punto del piano mobile
con accelerazione nulla?
Uguagliando a zero l'accelerazione
di un generico punto $K$,
dalla \ref{e115}
otteniamo l'equazione

\begin{equation}
{{\bm a}_{\scriptscriptstyle{K}}} = 0 =
{{\bm a}_{\scriptscriptstyle{P}}} +
\dot{\bm \omega} \times \overrightarrow{PK}-
\omega^2 {{\overrightarrow{PK}}}\,,
\label{e1f}
\end{equation}

\noindent che pu\`o essere risolta rispetto al vettore $\overrightarrow{PK}$.
Il punto $K$ cos\`i individuato si chiama {\em centro delle accelerazioni}\index{centro!delle accelerazioni}.
}

\begin{equation}
{{\bm a}_{\scriptscriptstyle{Q}}} =
{{\bm a}_{\scriptscriptstyle{P}}} +
\dot{\bm \omega} \times \overrightarrow{PQ}-
\omega^2 {{\overrightarrow{PQ}}}\,.
	\label{e115}
\end{equation}

\noindent Qualora poi il punto $P$ appartenesse all'asse fisso di rotazione del corpo rigido, 
asse presente in ruote dentate, manovelle,
dischi, volani, ecc.,
la \ref{e115} diventa

\begin{equation}
{{\bm a}_{\scriptscriptstyle{Q}}} =
\dot{\bm \omega} \times \overrightarrow{PQ}-
\omega^2 {{\overrightarrow{PQ}}}\,,
	\label{e116}
\end{equation}

\noindent e i suoi due termini a destra dell'uguale  si chiamano rispettivamente
{\em accelerazione tangenziale}\index{accelerazione!tangenziale}
e {\em accelerazione centripeta}\index{accelerazione!centripeta} del punto $Q$\footnote
{Nel caso in cui il punto $P$ fosse centro d'istantanea rotazione (e non di una rotazione
finita),
la \ref{e116} e i nomi di accelerazione tangenziale e centripeta del punto $Q$ rimarrebbero validi, in quel momento, indipendentemente dall'esistenza o meno
 di un asse fisso di rotazione.
}.

\section{Punto di Vista della Traiettoria}
\noindent Un altro approccio al calcolo dell'accelerazione del punto $Q$ si basa
\begin{wrapfigure}{r}{0.58\textwidth}
      \begin{center}

      \includegraphics[width=0.45\textwidth]{part1/cinematica/FIG/f18.pdf}
     \end{center}
\begin{picture}(-0,0)(-47,-30)
\scriptsize{
\put(81,22){$Q$}
\put(60,46){$Q'$}
\put(124,58){$O$}
\put(13,80){$-{\bm v}_{\scriptscriptstyle{Q}}$}
\put(51,90){${\bm v}_{\scriptscriptstyle{Q}}$}
\put(-28,75){${\bm v}_{\scriptscriptstyle{Q'}}$}
\put(66,120){${\bm v}_{\scriptscriptstyle{Q'}}$}
\put(-25,105){${\rm d}\theta$}
\put(110,28){${\rm d}\theta$}
\put(124,42){${\rm d}{\bm \theta}$}
\put(-5,30){${\rm d}{\bm v}_{\scriptscriptstyle{Q}_{\scriptscriptstyle{n}}}$}
\put(23,60){${\rm d}{\bm v}_{\scriptscriptstyle{Q}_{\scriptscriptstyle{t}}}$}
}
\end{picture}
\vskip -5.8mm
	\caption{\em Accelerazioni di un corpo rigido: punto di vista intrinseco.}
     \label{fig:f18}
\end{wrapfigure}
sulla valutazione diretta
della derivata di ${{\bm v}_{\scriptscriptstyle{Q}}}$.
\noindent Dopo un tempo infinitesimale ${\rm d}t$ il punto $Q$ si sar\`a portato in $Q'$,
come mostrato nella figura \ref{fig:f18}.
Anche le velocit\`a
${\bm v}_{\scriptscriptstyle{Q}}$ e 
${\bm v}_{\scriptscriptstyle{Q'}}$
saranno in generale diverse 
e possiamo quindi evidenziare la loro differenza infinitesimale.
Con l'aiuto delle figure \ref{fig:f16} e \ref{fig:f18}
abbiamo


\begin{equation}
{\rm d}{\bm v}_{\scriptscriptstyle{Q}_{\scriptscriptstyle{t}}}=
{{\bm v}_{\scriptscriptstyle{Q}}\over{
|{\bm v}_{\scriptscriptstyle{Q}}|}}
\dot{|{\bm v}_{\scriptscriptstyle{Q}}|}{\rm dt}\,,
	\label{e117}
\end{equation}

\begin{equation}
{\rm d}{\bm v}_{\scriptscriptstyle{Q}_{\scriptscriptstyle{n}}}=
 {\rm d}{\bm \theta}\times 
{\bm v}_{\scriptscriptstyle{Q}}\,. 
	\label{e118}
\end{equation}


\noindent Il vettore ${\rm d}{\bm \theta}$ che, come abbiamo
precedentemente chiarito, \`e
ortogonale al piano del disegno, non viene rappresentato in figura se non tramite il proprio simbolo.
A questo punto conviene per\`o introdurre un sistema di
tre versori ortogonali tra loro $\hat{\bm t}$, $\hat{\bm n}$ e $\hat{\bm b}$.
Poniamo il {\em versore tangente}\index{versore!tangente}

\begin{equation}
 \hat{\bm t}= {{{\bm v}_{\scriptscriptstyle{Q}}\over{
|{\bm v}_{\scriptscriptstyle{Q}}|}}}\,,
\label{e119}
\end{equation}

\noindent e imponiamo al {\em versore normale}\index{versore!normale} $\hat{\bm n}$
di essere appunto normale a $\hat{\bm t}$ e al vettore ${\rm d}{\bm \theta}$, inoltre, di essere
diretto come $\overrightarrow{QO}$;
infine imponiamo a $\hat{\bm b}$, chiamato
{\em versore binormale}\index{versore!binormale},
di avere la medesima direzione e il medesimo verso di ${\rm d}{\bm \theta}$.
\noindent Le relazioni \ref{e117} e \ref{e118} si possono ora scrivere, aiutandoci con
i versori della {\em terna intrinseca}\index{terna!intrinseca} test\'e introdotta,
nel seguente modo:
\begin{equation}
{\rm d}{\bm v}_{\scriptscriptstyle{Q}_{\scriptscriptstyle{t}}} = {\hat{\bm t}}
\dot{|{\bm v}_{\scriptscriptstyle{Q}}|}{\rm d t}\,,
	\label{e120}
\end{equation}
\begin{equation}
{\rm d}{\bm v}_{\scriptscriptstyle{Q}_{\scriptscriptstyle{n}}} =
{\rm d}{\theta}
({\hat{{\bm b}}}\times 
{\bm v}_{\scriptscriptstyle{Q}})\,. 
	\label{e121}
\end{equation}
\noindent L'angolo ${\rm d}\theta$ si pu\`o ottenere dal rapporto tra arco
$|(Q'-Q)|=|{\bm v}_{\scriptscriptstyle{Q}}|{\rm d}t$
e il {\em raggio di curvatura della traiettoria}, individuato da $\rho=|(O-Q)|$ o alternativamente
da $\rho=|(O-Q')|$, cio\`e
\begin{equation}
{\rm d}{\theta}=
{|{\bm v}_{\scriptscriptstyle{Q}}|{\rm d}t 
\over{\rho}}\,.
\label{e122}
\end{equation}
\noindent Dividendo tutte le espressioni per ${\rm d}t$ e sostituendo ${\rm d}{\theta}$ con l'espressione fornita dalla
\ref{e122}, le \ref{e120} e \ref{e121} diventano
\begin{equation}
{\bm a}_{\scriptscriptstyle{Q}_{\scriptscriptstyle{t}}}=
{\hat{\bm t}}
\dot{|{\bm v}_{\scriptscriptstyle{Q}}|}\,,
	\label{e123}
\end{equation}
\begin{equation}
{\bm a}_{\scriptscriptstyle{Q}_{\scriptscriptstyle{n}}}=
{|{\bm v}_{\scriptscriptstyle{Q}}| 
\over{\rho}}
({\hat{{\bm b}}}\times 
{\bm v}_{\scriptscriptstyle{Q}})\,. 
	\label{e124}
\end{equation}
\noindent Ma il versore 
${\hat{\bm b}}$, che \`e concorde a ${{\rm d}\bm \theta}$ sia come direzione sia come verso,
risulta ortogonale a ${\bm v}_{\scriptscriptstyle{Q}}$ e al versore
${\hat{\bm n}}$, pertanto 
\begin{equation}
({\hat{{\bm b}}}\times 
{\bm v}_{\scriptscriptstyle{Q}})= 
{\hat{\bm n}}
|{\bm v}_{\scriptscriptstyle{Q}}|\,,
	\label{e125}
\end{equation}
\noindent e la \ref{e124} si pu\`o scrivere nel modo seguente
\begin{equation}
{\bm a}_{\scriptscriptstyle{Q}_{\scriptscriptstyle{n}}}=
{\hat{\bm n}}
{|{\bm v}_{\scriptscriptstyle{Q}}|^2 
\over{\rho}}\,.
	\label{e126}
\end{equation}
\noindent Riassumendo, abbiamo per la accelerazione del punto $Q$
\begin{equation}
{{\bm a}_{\scriptscriptstyle{Q}}} =
{\hat{\bm t}}
\dot{|{\bm v}_{\scriptscriptstyle{Q}}|}+
{\hat{\bm n}}
{|{\bm v}_{\scriptscriptstyle{Q}}|^2 
\over{\rho}}\,.
	\label{e127}
\end{equation}
\noindent Ci si pu\`o chiedere quale sia la differenza tra l'espressione \ref{e115},
ottenuta poc'anzi, e la formula alternativa \ref{e127}.
Si nota immediatamente che le informazioni necessarie
all'individuazione dell'accelerazione del punto $Q$, mediante i due approcci,
sono diverse.
Da una parte l'espressione \ref{e115} richiede di sapere l'accelerazione del
punto $P$, unitamente alla velocit\`a e all'accelerazione angolari del corpo
rigido, nonch\'e la posizione di $Q$ rispetto a $P$.
Dall'altra dobbiamo sapere cosa succede alla velocit\`a del punto $Q$ in due
posizioni temporali infinitamente vicine, il che equivale a conoscere anche
la geometria della {\em traiettoria}\index{traiettoria} del punto $Q$. 
L'espressione \ref{e126} risulta molto comoda qualora si conoscano la
velocit\`a e la {\em curvatura della traiettoria}\index{curvatura!della traiettoria} percorsa da $Q$. 
Immaginiamo di percorrere con un'automobile una curva e di trovarci in un
punto dove il suo raggio, in generale variabile, sia noto\footnote
{
Non \`e facile avere a disposizione una tale informazione,
anche se forse si potrebbe 
ottenere dalle mappe, oggigiorno disponibili su vari dispositivi.
}
e di conoscere la nostra velocit\`a
(il suo modulo) mediante la lettura del tachimetro: ecco che tramite la \ref{e126}
risulta immediato calcolare la nostra accelerazione centripeta.
Inoltre, la \ref{e127} \`e equivalente alla \ref{e116}
 (il lettore pu\`o verificarlo tramite
le relazioni che intercorrono tra le varie quantit\`a) quando il moto ha un asse fisso.
Vale la pena aggiungere che la \ref{e127}
\`e stata ricavata senza 
fare uso della planarit\`a del moto, quindi la sua validit\`a si estende
anche ai moti tridimensionali.
\newpage
\thispagestyle{empty}
\null

\endinput 



\chapter{I Moti Relativi}

\section{Introduzione}

Agli spostamenti, alle velocit\`a, con le quali abbiamo
or ora fatto conoscenza, e, con prudenza, alle accelerazioni calza spesso a pennello
l'epiteto ``relativo''.
Stando seduti sul sedile di un treno che si muove per partire
dalla stazione in molti hanno sperimentato
la percezione che un altro treno, vicino di binario al nostro,
si stesse muovendo verso il fondo delle rotaie, in procinto
di scendere i gradini verso Piazza Duca d'Aosta\footnote
{
Per chi parte dalla Stazione Centrale di Milano.
}.
Un tale dubbio non potrebbe mai sorgere
nella mente
di un passeggero di quel treno che invece ci osserva tranquillamente partire.  
Come pure, rimanendo in ambito ferroviario, il passeggero che cammina placido
verso la testa del treno possiede, rispetto ad un passeggero a terra, che aspetta 
sul marciapiede,
la propria velocit\`a sommata a quella del treno.


\begin{figure}[ht]
	\begin{center}
      		\includegraphics{part1/relativi/FIG/f19.pdf}
	\end{center}
\begin{picture}(0,0)(-55,-50)
\scriptsize{
	\put(0,66){$\scriptscriptstyle{\rm {Terra}}$}
	\put(162,-10){$\scriptscriptstyle{\rm {Sole}}$}
	\put(220,98){\rotatebox{35}{$\scriptscriptstyle{\rm {Marte}}$}}
	\put(110,72){$\bm{d}$}
	\put(80,26){$\bm{s}_{\scriptscriptstyle{T}}$}
	\put(188,36){$\bm{s}_{\scriptscriptstyle{M}}$}
	\put(135,125){$\bm{\Omega}$}
}
\end{picture}
\vskip -6mm
	\caption{\em Osservazione del Sole da due pianeti diversi.}
     \label{fig:f19}
\end{figure}

\noindent Insomma, gli spostamenti di un oggetto,
come anche le sue velocit\`a e le sue accelerazioni, appaiono
diversi ad osservatori tra loro in movimento. Anzi tutte le grandezze cinematiche sono
legate (relative) all'osservatore che le misura. La domanda che ci si pu\`o porre, 
la quale costituisce il
{\em problema dei moti relativi}\index{problema dei moti relativi},
\`e la seguente: risulta possibile a due diversi osservatori scambiare i dati
provenienti dalle loro indagini cinematiche e confrontare le rispettive osservazioni
eseguite sopra oggetti comuni? La risposta \`e s\`i, stando a certi patti e a certe condizioni.

\section{Posizione del Problema}
Premettendo che una trattazione del {\em teorema di Coriolis}, rigorosa e completa, si trova
magistralmente esposta nel gi\`a citato \cite{finzi}, pag. 219, riportiamo qui un esempio che 
riteniamo un inquadramento piuttosto generale del {\em problema dei moti relativi}.
Un terrestre e un marziano,
di professione astronomi, sono appassionati di osservazioni solari:
entrambi misurano con facilit\`a posizione, velocit\`a e accelerazione della nostra 
stella. Nessuna delle quantit\`a misurate dal terrestre \`e in accordo con la
corrispondente del marziano\footnote{Un caro amico, leggendo le bozze di questo lavoro,
mi ha fatto notare come, pur lasciando da parte il sapore vagamente tolemaico
delle misure di velocit\`a e accelerazione del Sole da parte dei due astronomi,
 non sia comunque possibile il confronto tra  osservazioni
eseguite da due distinti punti dell'universo separati da immense distanze, in quanto
la luce del Sole impiega tempi rilevanti e diversi per giungere sulla Terra e su Marte.
Rispondiamo a questa giusta critica affermando che,
 fuori dalla sfera dei meccanismi, dove gli esempi di moti relativi abbondano,
purtroppo non ci vengono in mente
altre situazioni da proporre, sia pure di fantasia.
Per questo ci teniamo gli astronomi tolemaici,
 supponendo in tutto questo che la velocit\`a della luce sia infinita.}.
 Il terrestre, forse per protagonismo, decide di misurare
alcune quantit\`a cinematiche di Marte, {\em in primis} la sua posizione $\bm d$ e le 
corrispondenti
derivate ma,
incuriosito dalla rotazione del suo corrispondente (dispongono di cannocchiali potentissimi), decide di misurare anche
$\bm \Omega$, la velocit\`a angolare di Marte rispetto alla Terra.
Mediante lo schema fornito dalla figura
\ref{fig:f19}, impostano la seguente relazione
\begin{equation}
	\bm{s}_{\scriptscriptstyle{T}}=
	\bm{d}+
	\bm{s}_{\scriptscriptstyle{M}}\,,
\label{e128}
\end{equation}
\noindent che infonde ai due astronomi (sono ovviamente in contatto tra loro)
un pizzico di ottimismo circa la loro relazione, perch\'e la trovano soddisfatta.
Il terrestre ha a disposizione anche la velocit\`a del punto di Marte dove lavora il suo
amico extraterrestre,
$\dot{\bm d}$, e il marziano gli fornisce gentilmente la sua misura della
velocit\`a del Sole,
$\dot{\bm{s}}_{\scriptscriptstyle{M}}$.
Ma scrivendo la seguente
\begin{equation}
\xcancel{\dot{\bm{s}}_{\scriptscriptstyle{T}}=
	\dot{\bm{d}}+
\dot{\bm{s}}_{\scriptscriptstyle{M}}}\,,
\label{e129}
\end{equation}
\noindent in apparente accordo con la \ref{e128}, e sostituendo dati in loro
possesso, incontrano una cocente delusione.
Con lo scopo di rendere pi\`u chiare le interazioni future tra i due osservatori
conviene attribuire loro epiteti che ricordino le loro ``prerogative''. A ben 
considerare, essi sono completamente intercambiabili: Marte e Terra sono
due pianeti a pieno titolo e nessuno dei due osservatori
possiede caratteristiche tali da
implicare alcuna sorta di privilegio. Chiameremo tuttavia
{\em osservatore assoluto}\index{osservatore!assoluto} colui che
misura, oltre alle grandezze cinematiche del Sole,
anche le grandezze cinematiche dell'altro osservatore, che sono il
vettore $\bm d$, il vettore $\bm \Omega$ e le loro derivate; per noi costui \`e,
ad arbitrio, l'osservatore terrestre.
Chiameremo invece
il marziano {\em osservatore relativo}\index{osservatore!relativo}.
Risulta chiaro da quanto ora esposto che i due ruoli sono invertibili:
basterebbe all'osservatore marziano essere in possesso di posizione,
velocit\`a e accelerazione, grandezze angolari comprese, della
Terra per potere diventare 
l'osservatore assoluto.
\noindent Il terrestre pensa e ripensa alla \ref{e129} e, cercando di capire dove
si nasconde l'errore, per stanarlo conduce il seguente esperimento mentale: immagina una mosca
aliena ferma sulla lente del cannocchiale marziano e, con lo scopo
di conoscerne la velocit\`a, prova ad applicare a
essa l'errata 
formula \ref{e129}. La {\em velocit\`a relativa}\index{velocit\`a!relativa},
cio\`e la velocit\`a della mosca misurata dal marziano, quella che nella 
\ref{e129} \`e indicata con 
$\dot{\bm{s}}_{\scriptscriptstyle{M}}$,
 \`e in questo caso sicuramente nulla.
Per tale ragione, l'osservatore assoluto (il terrestre) si dovrebbe accontentare di ottenere
per la velocit\`a della mosca, come per le velocit\`a di tutti gli oggetti fermi
 su Marte, un valore uguale a
$\dot{\bm{d}}$. Ebbene questo sarebbe possibile solo se tali oggetti (i
punti considerati di tali oggetti) si trovassero giusto
nel centro della terna del sistema di riferimento marziano, oppure se tutto il pianeta,
e di conseguenza il sistema di riferimento a esso collegato, traslasse.
\noindent Addebitando alla rotazione del sistema relativo la causa delle differenti velocit\`a assolute
degli oggetti fermi su Marte cominciano a intravedere uno
spiraglio di luce verso la soluzione dell'enigma. La formula
\ref{e129} non funziona perch\'e il sistema di riferimento relativo ruota rispetto
al sistema di riferimento assoluto con velocit\`a angolare $\bm \Omega$. Per  stimare
 l'errore
che l'osservatore relativo commette nell'effettuare ``le proprie derivate'' l'osservatore
assoluto chiede al marziano di comunicargli
la derivata di un vettore $\bm a$
di modulo costante e fisso con il sistema assoluto.
Il marziano vede tale vettore
ruotare con velocit\`a angolare $-\bm \Omega$ e ne
determina una variazione temporale (derivata rispetto al tempo)
 pari a $\dot{\bm a}= - {\bm \Omega}\times {\bm a}$.
\noindent Il calcolo di questo termine \`e gi\`a stato svolto in precedenza, nel capitolo
in cui si descrivono gli spostamenti infinitesimali, e in
quel contesto \`e limitato al caso
piano, equazioni \ref{e14}, \ref{e15}, \ref{e16}.
A questo punto l'osservatore assoluto conosce il termine correttivo da introdurre 
nella \ref{e129} e la riscrive correttamente\footnote
{
Per coloro i quali volessero approfondire il calcolo delle
derivate di un vettore rotante o, ci\`o che \`e lo stesso, la relazione
tra la rotazione del sistema di riferimento e la variazione (apparente)
di un vettore, rimandiamo a \cite{finzi}, pag. 163.
}
\begin{equation}
\dot{\bm{s}}_{\scriptscriptstyle{T}}=
	\dot{\bm{d}}+
\dot{\bm{s}}_{\scriptscriptstyle{M}}+
 {\bm \Omega}\times 
	\bm{s}_{\scriptscriptstyle{M}}\,.
\label{e130}
\end{equation}

\noindent Le velocit\`a dei punti che sono fermi su Marte, come la mosca sulla lente del
cannocchiale, assumono il nome speciale di {\em velocit\`a di trascinamento}\index{velocit\`a!di trascinamento}
\begin{equation}
{\bm v}_{\scriptscriptstyle{{\rm tr}}}=
	\dot{\bm{d}}+
 {\bm \Omega}\times 
	\bm{s}_{\scriptscriptstyle{M}}\,,
\label{e131}
\end{equation}
ovviamente variabile da punto a punto del sistema relativo.
\noindent La rimanente componente si chiama {\em velocit\`a relativa}\index{velocit\`a!relativa}
\begin{equation}
{\bm v}_{\scriptscriptstyle{{\rm rel}}}=
\dot{\bm{s}}_{\scriptscriptstyle{M}}\,,
\label{e132}
\end{equation}
\noindent che, sommata alla precedente, permette di ottenere
la {\em velocit\`a assoluta}\index{velocit\`a!assoluta}
del Sole 
\begin{equation}
{\bm v}_{\scriptscriptstyle{{\rm ass}}}=
{\bm v}_{\scriptscriptstyle{{\rm rel}}}+
{\bm v}_{\scriptscriptstyle{{\rm tr}}}\,.
\label{e133}
\end{equation}
\section{L'Accelerazione di Coriolis}

\noindent Giunti sin qui, sapendo che ogni derivata che provenga dal sistema relativo va depurata
aggiungendo il termine opportuno, i due astronomi provano a sostituire i loro
dati di accelerazione del Sole nella derivata della \ref{e130}
\begin{equation}
\ddot{\bm{s}}_{\scriptscriptstyle{T}}=
	\ddot{\bm{d}}+
\ddot{\bm{s}}_{\scriptscriptstyle{M}}+
 {\bm \Omega}\times 
\dot{\bm{s}}_{\scriptscriptstyle{M}}+
\dot{\bm \Omega}\times 
	\bm{s}_{\scriptscriptstyle{M}}+
 {\bm \Omega}\times 
\dot{\bm{s}}_{\scriptscriptstyle{M}}+
 {\bm \Omega}\times( 
 {\bm \Omega}\times 
	\bm{s}_{\scriptscriptstyle{M}})\,;
\label{e134}
\end{equation}
\noindent la quale sostituzione riempie di gioia i due astronomi perch\'e la
\ref{e134} viene soddisfatta: essi battezzano questa formula 
{\em teorema di Coriolis}\index{teorema!di Coriolis}. D'altra parte, una volta imparato a correggere le 
derivate provenienti dall'osservatore relativo, la \ref{e134} si ottiene
meccanicamente e senza impiego di gran ragionamenti derivando la \ref{e130}.
Anche in questo caso, come abbiamo fatto per le velocit\`a,
 conviene dare un nome speciale all'accelerazione 
dei punti fermi costantemente rispetto al sistema relativo (la mosca di 
poco fa ferma sulla lente del cannocchiale), e chiameremo 
{\em accelerazione di trascinamento}\index{accelerazione!di trascinamento} la seguente
\begin{equation}
{\bm a}_{\scriptscriptstyle{{\rm tr}}}=
	\ddot{\bm{d}}+
\dot{\bm \Omega}\times 
	\bm{s}_{\scriptscriptstyle{M}}+
 {\bm \Omega}\times( 
 {\bm \Omega}\times 
        \bm{s}_{\scriptscriptstyle{M}})\,.
\label{e135}
\end{equation}
\noindent Nella \ref{e135} compaiono solamente i termini della \ref{e130} che non contengono
grandezze cinematiche relative, n\'e velocit\`a n\'e accelerazioni (la mosca \`e ferma e rimane ferma).
L'{\em accelerazione relativa}\index{accelerazione!relativa} \`e naturalmente
\begin{equation}
{\bm a}_{\scriptscriptstyle{{\rm rel}}}=
\ddot{\bm{s}}_{\scriptscriptstyle{M}}\,,
\label{e136}
\end{equation}
\noindent ed \`e l'unica che l'osservatore relativo \`e in grado di misurare.
Della \ref{e130} restano due monomi, uguali, contenenti grandezze relative e grandezze
assolute  mescolate. L'espressione
\begin{equation}
{\bm a}_{\scriptscriptstyle{{\rm cor}}}=
2 {\bm \Omega}\times 
\dot{\bm{s}}_{\scriptscriptstyle{M}}=
2 {\bm \Omega}\times 
{\bm v}_{\scriptscriptstyle{{\rm rel}}}
\label{e137}
\end{equation}
\noindent si chiama {\em accelerazione di Coriolis}\index{accelerazione!di Coriolis}, in onore al matematico francese
{\em Gaspard-Gustave de Coriolis} che l'ha scoperta.
A questo punto, possiamo finalmente scrivere l'espressione
dell'{\em accelerazione assoluta}\index{accelerazione!assoluta}
nel modo seguente
\begin{equation}
{\bm a}_{\scriptscriptstyle{{\rm ass}}}=
{\bm a}_{\scriptscriptstyle{{\rm rel}}}+
{\bm a}_{\scriptscriptstyle{{\rm tr}}}+
{\bm a}_{\scriptscriptstyle{{\rm cor}}}\,.
\label{e138}
\end{equation}

\section{Riepilogo}

\noindent  Quando due osservatori,
in moto tra loro, indagano sulla cinematica di un terzo oggetto ne traggono
quantit\`a distinte. Per potere confrontare tali quantit\`a occorre che uno degli
osservatori (assoluto) misuri anche le quantit\`a cinematiche assolute
attribuibili al sistema relativo cio\`e
$\bm d$, $\bm \Omega$ e le loro derivate.
Alla velocit\`a di una formica (punto osservato), misurata dall'osservatore
relativo, l'osservatore assoluto dovr\`a aggiungere la componente di trascinamento
che \`e la velocit\`a assoluta del punto del sistema relativo dove la formica appoggia le 
zampe in quel momento.
Con le accelerazioni le cose vanno diversamente: non \`e sufficiente sommare la componente
del trascinamento all'accelerazione relativa per ottenere l'accelerazione assoluta. La
velocit\`a relativa provoca nei sistemi rotanti un'altra accelerazione,
detta di {\em Coriolis}, da aggiungere a quella relativa e di trascinamento, della quale
daremo nel seguente capitolo qualche ulteriore delucidazione.
\newpage
\thispagestyle{empty}
\null

\chapter{Esempi di Moti Relativi}

\section{Precisazioni sull'Accelerazione di Coriolis} 

\begin{figure}[h]
\centering
\begin{minipage}[b]{0.50\textwidth}
\centering
\includegraphics[width=0.9\textwidth]{part1/relativi/FIG/f110.jpg}
\caption{\em La giostra dei cavalli.}
 \label{fig:f110}
\end{minipage}\hfill
\begin{minipage}[b]{0.50\textwidth}
\centering
      \includegraphics[width=0.9\textwidth]{part1/relativi/FIG/f111.pdf}
\begin{picture}(0,0)(130,0)
\scriptsize{
\put(24,131){${\rm d}\theta$}
\put(-19,116){$\omega$}
\put(-14,30){${\bm d}$}
\put(-29,32){$y$}
\put(-8,10){$x$}
\put(-3,86){${{\bm v}_{\scriptscriptstyle{{\rm rel}}}}(t+{\rm d}t)$}
\put(51,108){${{\bm v}_{\scriptscriptstyle{{\rm rel}}}}(t)$}
\put(60,66){$r-{\rm d}r$}
\put(70,85){$r$}
}
\end{picture}
        \caption{\em Versione schematica.}
     \label{fig:f111}
\end{minipage}
\end{figure}

\noindent La ``giostra dei cavalli'' sembra un laboratorio appositamente costruito per
fare esercizio coi moti relativi.
Essa ci riconduce innanzitutto ai nostri moti piani, dove l'intuito pu\`o meglio padroneggiare
la situazione.
Divertimento per bambini (e non) molto conosciuto, questa giostra si presenta come un grande
piatto rotante, normalmente pavimentato in legno, dove un paio di ranghi
di sedili, cavalli, carrozze e altro ospitano i turisti portandoli in rotazione. Credo che
i modelli pi\`u lussuosi muovano i cavalli anche verticalmente, ma questo non 
importa al nostro scopo, tant'\`e che ci\`o di cui misureremo le quantit\`a cinematiche,
cio\`e il nostro punto osservato, sar\`a una formica intenzionata a esplorare 
la giostra. L'osservatore assoluto potrebbe essere
il bigliettaio ed \`e rappresentato dal riferimento $(x,y)$, coincidente con la sua edicola.
\`E appena il caso di chiarire che il vettore
 ${\bm d}$ \`e costante perch\'e il centro della giostra
\`e fermo rispetto alla piazza, per questo motivo ${\bm d}$
non giocher\`a alcun ruolo nei ragionamenti che seguiranno.
Un turista seduto sulla giostra \`e nei panni dell'osservatore
relativo e, indagando il movimento della formica, riuscir\`a a misurare
la velocit\`a e l'accelerazione relative dell'insetto.
La velocit\`a e l'accelerazione del punto del pavimento dove si trova la formica
forniscono subito
le diverse componenti del trascinamento: tali quantit\`a cinematiche
appartengono ai punti del disco rotante.
Supponiamo ora che la formica si incammini lungo una fessura radiale
della pavimentazione, comune nelle giostre col {\em parquet}, come mostrato in figura \ref{fig:f111}.
Nell'istante di tempo considerato, sia ${{\bm v}_{\scriptscriptstyle{{\rm rel}}}}(t)$
la velocit\`a relativa della formica, cio\`e la velocit\`a con cui sta percorrendo l'interstizio
rettilineo. Si pu\`o osservare nella stessa figura che la circonferenza del
piatto della giostra dove si trover\`a la formica dopo un intervallo infinitesimale
di tempo ${\rm d}t$ avr\`a
un raggio minore. \`E naturale pertanto
pensare che, con riferimento ai due istanti considerati,
si avr\`a una differenza tra le due
velocit\`a di trascinamento cui sar\`a soggetta la formica.

\begin{figure}[hbt]
\centering
\begin{minipage}[b]{0.57\textwidth}
\centering
\includegraphics[width=0.8\textwidth]{part1/relativi/FIG/f112.pdf}
\begin{picture}(0,0)(130,0)
\scriptsize{
\put(59,108){${\rm d}\theta=\omega{\rm d}t$}
\put(-10,132){${{\bm v}_{\scriptscriptstyle{{\rm tr}}}}(t)$}
\put(-40,59){${{\bm v}_{\scriptscriptstyle{{\rm tr}}}}(t+{\rm d}t)$}
\put(101,147){$r$}
\put(118,100){$r-{\rm d}r$}
\put(-36,41){$-{{\bm v}_{\scriptscriptstyle{{\rm tr}}}}(t+{\rm d}t)$}
\put(10,12){${{\bm v}_{\scriptscriptstyle{{\rm tr}}}}(t)$}
\put(22,70){${{\rm d}{{\bm v}_{\scriptscriptstyle{{\rm tr}}}}}_{\scriptscriptstyle{t}}$}
\put(62,57){${\rm d}{{\bm v}_{\scriptscriptstyle{{\rm tr}}}}_{\scriptscriptstyle{n}}$}
}
\end{picture}
\caption{\em Effetto della differenza radiale.}
 \label{fig:f112}
\end{minipage}\hfill
\begin{minipage}[b]{0.38\textwidth}
\centering
      \includegraphics[width=0.8\textwidth]{part1/relativi/FIG/f113.pdf}
\begin{picture}(0,0)(110,0)
\scriptsize{
\put(40,135){${\rm d}\theta$}
\put(3,98){\rotatebox{28}{${\rm d}{\bm v}_{\scriptscriptstyle{{\rm rel}}}$}}
\put(38,65){\rotatebox{-72}{${-{\bm v}_{\scriptscriptstyle{{\rm rel}}}}(t)$}}
\put(0,75){\rotatebox{-58}{${{\bm v}_{\scriptscriptstyle{{\rm rel}}}}(t+{\rm d}t)$}}
\put(70,130){\rotatebox{-72}{${{\bm v}_{\scriptscriptstyle{{\rm rel}}}}(t)$}}
\put(40,115){\rotatebox{-58}{${{\bm v}_{\scriptscriptstyle{{\rm rel}}}}(t+{\rm d}t)$}}
}
\end{picture}
        \caption{\em Cambiamento di direzione della velocit\`a relativa.}
     \label{fig:f113}
\end{minipage}
\end{figure}

\noindent Nella figura \ref{fig:f112} sono rappresentate le due componenti infinitesimali frutto della
differenza tra la velocit\`a di trascinamento al tempo $t+{\rm d}t$ e al tempo $t$.
Non ci occupiamo della componente normale
${\rm d}{{\bm v}_{\scriptscriptstyle{{\rm tr}}}}_{\scriptscriptstyle{n}}$, in quanto a noi gi\`a
nota.
Il lettore pi\`u diligente pu\`o tornare alla figura \ref{fig:f18} e ricavare anche in questo caso
la formula della accelerazione centripeta \ref{e126}.
La velocit\`a infinitesimale 
${{\rm d}{{\bm v}_{\scriptscriptstyle{{\rm tr}}}}}_{\scriptscriptstyle{t}}$ \`e invece
di nostro interesse: tale variazione della velocit\`a tangenziale si esprime come
\begin{equation}
{{\rm d}{{\bm v}_{\scriptscriptstyle{{\rm tr}}}}}_{\scriptscriptstyle{t}}=
\omega[(r -{\rm d}r) - (r)]\hat{{\bm t}}\,,
\label{e139}
\end{equation}
\noindent dove $\hat{\bm t}$ \`e il versore tangente, gi\`a altre volte incontrato (equazione \ref{e119}). Sapendo poi che che la differenza tra i raggi dopo 
un intervallo di tempo ${\rm d}t$ sar\`a ${\rm d} r= |{\bm v}_{\scriptscriptstyle{{\rm rel}}}|
{\rm d}t$, possiamo scrivere 
\begin{equation}
{{\rm d}{{\bm v}_{\scriptscriptstyle{{\rm tr}}}}}_{\scriptscriptstyle{t}}=
-\omega
(|{\bm v}_{\scriptscriptstyle{{\rm rel}}}|{\rm d}t)\,
\hat{{\bm t}}\,,
\label{e140}
\end{equation}
\noindent o anche
\begin{equation}
{{{{\rm d}{{\bm v}_{\scriptscriptstyle{{\rm tr}}}}}_{\scriptscriptstyle{t}}}\over{{\rm d}t}}=
{\bm \omega} \times
{\bm v}_{\scriptscriptstyle{{\rm rel}}}=
{{\bm a}_{\scriptscriptstyle{{\rm cor}}}\over{2}}\,.
\label{e141}
\end{equation}
\noindent Abbiamo cos\`i individuato la natura di met\`a dell'accelerazione
di {\em Coriolis}: essa deriva dalla costrizione subita dal nostro punto osservato a
cambiare raggio sul disco del pavimento, e di conseguenza 
velocit\`a di trascinamento. Questo d\`a conto
di uno dei monomi dell'espressione dell'accelerazione di {\em Coriolis}.
Veniamo ora alla figura \ref{fig:f113} e consideriamo che la velocit\`a relativa, nell'intervallo di
tempo ${\rm d}t$, cambia direzione, ruotando di un angolo pari
a ${\rm d} \theta=\omega{\rm d}t$. 
La variazione della velocit\`a relativa risulta pertanto essere
\begin{equation}
{\rm d}{\bm v}_{\scriptscriptstyle{{\rm rel}}}=
-\omega {\rm d}t 
|{\bm v}_{\scriptscriptstyle{{\rm rel}}}|\hat{\bm t}\,,
\label{e142}
\end{equation}
\noindent quindi 
\begin{equation}
{{\rm d}{\bm v}_{\scriptscriptstyle{{\rm rel}}}\over{\rm d}t}=
{\bm \omega}\times 
{\bm v}_{\scriptscriptstyle{{\rm rel}}}=
{{\bm a}_{\scriptscriptstyle{{\rm cor}}}\over{2}}\,,
\label{e143}
\end{equation}
\noindent che \`e il secondo (o il primo?) monomio dell'accelerazione di {\em Coriolis}.
Cos\`i tale accelerazione, dapprima introdotta
mediante un calcolo formale (derivando la \ref{e130} si \`e ottenuta
la \ref{e134}), si ricava ora mediante le osservazioni test\'e svolte e si
intravede la giustificazione dei suoi due monomi identici: essi sono
da ricondurre a due diverse conseguenze provocate e subite dalla velocit\`a
relativa durante la rotazione del sistema di riferimento relativo.
Possiamo aggiungere che tale chiarimento \`e stato molto
caro all'autore una volta scovato, stampato con tipi minori e con stile
incomparabilmente pi\`u elegante del nostro, su \cite{sesini1}, pag. 11.
Ma anche queste delucidazioni possono lasciare un po' di amaro in bocca per via di ulteriori dubbi che possono sorgere:
cosa succede ai due pluricitati monomi quando la velocit\`a relativa \`e tangente alle circonferenze
del piatto della giostra anzich\'e essere radiale? Che fine fa in tal caso il monomio che
esiste in virt\`u della diminuzione o dell'aumento delle circonferenze percorse nel sistema
relativo? Si riduce forse a zero? Proprio cos\`i! E l'altro monomio, quello dovuto al
cambiamento di direzione della velocit\`a relativa, raddoppia: la dimostrazione di
questa affermazione \`e lasciata come esercizio al lettore.
E cosa succede nelle situazioni intermedie? Cio\`e quando
la velocit\`a relativa non \`e n\'e radiale n\'e tangenziale? Conviene qui
tagliar corto: i procedimenti
matematici servono proprio a questo: la \ref{e134} \`e stata ricavata correttamente
dalla \ref{e130} perci\`o risulta sempre valida. Le altre considerazioni, meno rigorose,
basate su circostanze particolari, servono a stimolare e svegliare l'intuito, ci
fanno conoscere gli attori, solitamente austeri e scontrosi, della matematica con
maggior confidenza.

\section{Ulteriori Precisazioni}

\noindent Ancora qualche parola su questa particolare grandezza cinematica. Anche se
non \`e qui il luogo per trattare argomenti di {\em dinamica}, tutti conosciamo
gli effetti delle accelerazioni sulle masse: esse fanno nascere le {\em forze di inerzia}.
Bench\'e i fisici chiamino queste forze, e non senza ragione,
 ``apparenti'', i loro effetti sono ben noti
ai passeggeri che stanno spesso
in piedi, anzich\'e seduti sugli autobus. Bene, le forze di inerzia legate alle accelerazioni di {\em Coriolis}
provocano particolari azioni dinamiche: sono quelle che al ritrarre o allo stendere
delle braccia del pattinatore ne aumentano o ne diminuiscono
la sua velocit\`a angolare quando piroetta;
esse provocano la circolazione dei venti sul nostro pianeta,
i vortici dei cicloni, e tanti altri effetti analoghi.

\section{Due Esempi Significativi}
\label{par:esempi_moti_rel}

\noindent Abbiamo visto che la soluzione del problema dei
moti relativi si trova in due equazioni che riguardano rispettivamente velocit\`a
e accelerazioni, \ref{e133} e \ref{e138}, ciascuna delle quali mette d'accordo le osservazioni
svolte nei rispettivi sistemi di riferimento, quello assoluto e quello relativo.
Tutte le volte che disponiamo di relazioni di uguaglianza 
siamo autorizzati
a trasformarle in altrettante equazioni
nelle quali qualche termine potr\`a essere incognito.
D'altra parte sappiamo che, dato un sistema di
equazioni tra loro indipendenti, possiamo ricavare il valore delle incognite
se il loro numero uguaglia quello delle relazioni che le legano tra loro.
Desideriamo provare a sfruttare questa attitudine delle uguaglianze \ref{e133} e \ref{e138} per
ricavare alcune grandezze cinematiche incognite
nei due problemi raffigurati nelle figure
\ref{fig:f114} e \ref{fig:f115}.
Andiamo con ordine e cominciamo dalla scala a pioli appoggiata a un muro
che scivola sul pavimento con velocit\`a e accelerazione date,
${{\bm v}_{\scriptscriptstyle{A}}}$ e ${{\bm a}_{\scriptscriptstyle{A}}}$, come rappresentato in figura \ref{fig:f114}.


\begin{figure}[hbt]
\centering
\begin{minipage}[b]{0.48\textwidth}
\centering
\includegraphics[width=0.8\textwidth]{part1/relativi/FIG/f114.pdf}
\begin{picture}(0,0)(130,0)
\scriptsize{
\put(-9,133){$y$}
\put(-2,123){$B$}
\put(1,92){${{\bm a}_{\scriptscriptstyle{B}}}$}
\put(-18,90){${{\bm v}_{\scriptscriptstyle{B}}}$}
\put(36,66){$\omega$}
\put(91,17){${{\bm v}_{\scriptscriptstyle{A}}}$}
\put(109,15){${{\bm a}_{\scriptscriptstyle{A}}}$}
\put(62,1){$A$}
\put(-9,1){$O$}
\put(68,28){$y'$}
\put(81.5,3){$x'$}
\put(105,3){$x$}
}
\end{picture}
\caption{\em Slittamento di una scala appoggiata a un muro.}
 \label{fig:f114}
\end{minipage}\hfill
\begin{minipage}[b]{0.48\textwidth}
\centering
      \includegraphics[width=0.8\textwidth]{part1/relativi/FIG/f115.pdf}
\begin{picture}(0,0)(129,0)
\scriptsize{
\put(-2,123){$B$}
\put(2,87){${{\bm a}_{\scriptscriptstyle{B}}}$}
\put(-25,90){${{\bm v}_{\scriptscriptstyle{B}}}$}
\put(24,88){${{\bm v}_{\scriptscriptstyle{P}}}_{\scriptscriptstyle{{\rm rel}}}$}
\put(10,105){${{\bm a}_{\scriptscriptstyle{P}}}_{\scriptscriptstyle{{\rm rel}}}$}
\put(11,74){$P$}
\put(36,66){$\omega$}
\put(11,40){$y$}
\put(92,17){${{\bm v}_{\scriptscriptstyle{A}}}$}
\put(109,15){${{\bm a}_{\scriptscriptstyle{A}}}$}
\put(11,1){$O$}
\put(62,1){$A$}
\put(60,28){$y'$}
\put(83,22){$x'$}
\put(105,3){$x$}
}
\end{picture}
        \caption{\em Slittamento di una scala sopra la sommit\`a di un muro.}
     \label{fig:f115}
\end{minipage}
\end{figure}

\noindent Ci poniamo il problema di trovare velocit\`a e accelerazione del punto $B$, cio\`e
${{\bm v}_{\scriptscriptstyle{B}}}$ e 
${{\bm a}_{\scriptscriptstyle{B}}}$, conoscendo 
${{\bm v}_{\scriptscriptstyle{A}}}$ e 
${{\bm a}_{\scriptscriptstyle{A}}}$, ossia le grandezze cinematiche del
punto $A$.
Esprimiamo le grandezze cinematiche assolute di $B$ tramite le grandezze
 relative 
appartenenti al sistema di riferimento $(x',y')$. Assumiamo in questo caso
(ma su questa scelta torneremo) che il sistema relativo trasli e si mantenga centrato in $A$.
Nel sistema relativo $(x',y')$ si osserva $B$ 
percorrere degli archi di circonferenza con centro $A$. Riprendendo la \ref{e133} possiamo scrivere
\begin{equation}
{\bm v}_{\scriptscriptstyle{{\rm ass}}}=
{\bm v}_{\scriptscriptstyle{{\rm rel}}}+
{\bm v}_{\scriptscriptstyle{{\rm tr}}}\,,
\label{e144}
\end{equation}
\noindent cio\`e
\begin{equation}
{{\bm v}_{\scriptscriptstyle{B}}}= 
{{\bm v}_{\scriptscriptstyle{B}}}_{\scriptscriptstyle{{\rm rel}}}+ 
{{\bm v}_{\scriptscriptstyle{A}}}\,. 
\label{e145}
\end{equation}
\noindent L'equazione \ref{e145} contiene due incognite, i vettori
${{\bm v}_{\scriptscriptstyle{B}}}$ e 
${{\bm v}_{\scriptscriptstyle{B}}}_{\scriptscriptstyle{{\rm rel}}}$. 
Ma tali grandezze non sono
completamente sconosciute in quanto la loro direzione \`e nota. Siamo pertanto di fronte
a un'equazione vettoriale (che nel piano vale due equazioni scalari) che contiene due incognite, i moduli dei vettori che non conosciamo.
Stabilito quindi che il numero delle incognite non eccede il numero delle equazioni, potremmo
lasciare al lettore il compito di individuare 
uno dei vari metodi per risolvere la \ref{e145}.
Preferiamo per\`o dare un cenno del metodo grafico che si potrebbe seguire,
metodo forse non molto efficiente ma a nostro parere molto istruttivo.
\begin{wrapfigure}{r}{0.35\textwidth}
      \begin{center}
      \includegraphics[width=0.30\textwidth]{part1/relativi/FIG/f116.pdf}
     \end{center}
\begin{picture}(0,0)(0,0)
\scriptsize{
\put(40,91){${{\bm v}_{\scriptscriptstyle{A}}}$}
\put(6,55){${{\bm v}_{\scriptscriptstyle{B}}}$}
\put(36,44){${{\bm v}_{\scriptscriptstyle{B}}\scriptscriptstyle{{\rm rel}}}$}
\put(61,59){\rotatebox{30}{${\scriptstyle{{\rm dir}}}{\scriptscriptstyle{\perp AB}}$}}
\put(13,30){\rotatebox{-90}{${\scriptstyle{{\rm dir}}}{\scriptscriptstyle{\parallel OB}}$}}
}
\end{picture}
        \caption{\em Soluzione grafica di un'equazione vettoriale.}
     \label{fig:f116}
\end{wrapfigure}
\noindent Con riferimento alla figura \ref{fig:f116} tracciamo dapprima il vettore noto 
${{\bm v}_{\scriptscriptstyle{A}}}$. Si tracciano quindi per i suoi estremi
le due direzioni cospicue indicate nella stessa figura.
Non ha alcuna importanza se la direzione parallela a $OB$ viene riportata sulla
coda del vettore noto oppure sulla sua testa e naturalmente lo stesso vale per l'altra direzione nota.
La soluzione consta nell'ammettere
che l'unica configurazione in cui i vettori incogniti
giacciano sulle direzioni tracciate e soddisfino
la relazione \ref{e145} risulta essere quella
riportata in figura, e la chiusura del triangolo
costituisce la soluzione dell'equazione vettoriale.
Una volta determinate
le velocit\`a incognite
si pu\`o senza fatica ricavare
il modulo della velocit\`a angolare della scala
 dalla relazione
$|{{\bm v}_{\scriptscriptstyle{B}}\scriptscriptstyle{{\rm rel}}}|=\omega\overrightarrow{|{AB}|}$;
tenendo quindi presente che la sua direzione non pu\`o 
che essere ortogonale al piano del moto, il suo
verso dovr\`a soddisfare la relazione
${{\bm v}_{\scriptscriptstyle{B}}\scriptscriptstyle{{\rm rel}}}=\bm{\omega}\times\overrightarrow{AB}$. 


\noindent Passiamo alle accelerazioni. In virt\`u del fatto che il sistema di riferimento relativo trasla, scrivendo la \ref{e138} il termine dell'accelerazione di {\em Coriolis} sar\`a assente. 
Avremo pertanto


\begin{equation}
{{\bm a}_{\scriptscriptstyle{B}}}= 
{{\bm a}_{\scriptscriptstyle{B}}}_{\scriptscriptstyle{{\rm rel}}}+
{{\bm a}_{\scriptscriptstyle{A}}}\,.
\label{e146}
\end{equation}

\noindent Conviene per\`o scrivere il termine
${{\bm a}_{\scriptscriptstyle{B}}}_{\scriptscriptstyle{{\rm rel}}}$ come somma delle sue due
componenti: quella avente direzione parallela alla scala e diretta verso $A$ (componente normale) e quella ortogonale alla scala stessa

\begin{equation}
{{\bm a}_{\scriptscriptstyle{B}}}= 
-\omega^2 \overrightarrow{AB} + \dot{\omega}|\overrightarrow{AB}|\widehat{{\perp{AB}}}+
{{\bm a}_{\scriptscriptstyle{A}}}\,,
\label{e147}
\end{equation}

\noindent dove il versore
$\widehat{{\perp{AB}}}$ sta ad indicare che la direzione del penultimo termine, al contrario del suo
modulo, \`e nota e proviene dal prodotto vettoriale
${{{\bm a}_{\scriptscriptstyle{B}}}_{\scriptscriptstyle{{\rm rel}}}}_{\scriptscriptstyle{(t)}}=
\dot{\bm\omega}\times\overrightarrow{AB}$,
pertanto essa risulta perpendicolare alla nostra scala.
Anche la direzione del vettore incognito ${{\bm a}_{\scriptscriptstyle{B}}}$
\`e nota ed \`e l'unica che il moto di strisciamento sul
muro verticale gli consenta.
Gli altri due termini, ${{\bm a}_{\scriptscriptstyle{A}}}$ e
${{{\bm a}_{\scriptscriptstyle{B}}}_{\scriptscriptstyle{{\rm rel}}}}_{\scriptscriptstyle{(n)}}=-\omega^2 \overrightarrow{AB}$, sono
invece completamente noti.
Ricadiamo cos\`i ancora nel caso di un'equazione vettoriale dove le incognite sono i moduli di due vettori; in questo frangente,
numerando i termini senza badare al segno di uguaglianza, risultano essere il primo e il terzo.
Una possibile soluzione \`e quindi di nuovo
lo schema grafico di figura \ref{fig:f116}, magari con l'avvertenza di produrre in
anticipo la somma dei due vettori noti. 
Occorre solo accennare che una volta ricavato
il terzo termine, 
${{{\bm a}_{\scriptscriptstyle{B}}}_{\scriptscriptstyle{{\rm rel}}}}_{\scriptscriptstyle{(t)}}=
\dot{\omega}|\overrightarrow{AB}|\widehat{{\perp{AB}}}$, risulter\`a noto anche
il modulo  dell'accelerazione 
angolare della scala $\dot{\omega}$; la sua direzione sar\`a ortogonale
al piano del moto e il suo verso dovr\`a
essere tale da produrre il
valore corretto di 
${{{\bm a}_{\scriptscriptstyle{B}}}_{\scriptscriptstyle{{\rm rel}}}}_{\scriptscriptstyle{(t)}}$
tramite il prodotto vettoriale con $\overrightarrow{AB}$.


\noindent Esaminiamo ora la scala della figura \ref{fig:f115} che, anzich\'e appoggiarsi al muro verticale
mediante il suo estremo $B$, appoggia in un suo punto imprecisato sulla sommit\`a del muro nel
punto (del muro) $P$. Abbiamo scelto di relazionarci con un osservatore relativo
che, oltre a traslare con il punto $A$, ruota mantenendo uno dei suoi assi diretto
lungo la scala.
Anche qui, le incognite sono velocit\`a e accelerazione del punto $B$, note la velocit\`a
e l'accelerazione del punto $A$\footnote
{
Il lettore, qualora fosse interessato, pu\`o provare
a scambiare le grandezze incognite con quelle note e viceversa, per rendersi
conto che, non solo la soluzione non cambia a livello concettuale, ma persino
nelle equazioni e nei disegni le grandezze coinvolte
rimarranno invariate nella sostanza.
}.
Ancora una volta partiamo dalle velocit\`a ma, anzich\'e tentare di coinvolgere subito
la grandezza incognita  
${{\bm v}_{\scriptscriptstyle{B}}}$,
proviamo a calcolare la velocit\`a del punto $P$, che deve essere nulla essendo
$P$ un punto appartenente al muro.
In tal caso, la \ref{e133} diventa

\begin{equation}
{{\bm v}_{\scriptscriptstyle{P}}}=0= 
{{\bm v}_{\scriptscriptstyle{P}}}_{\scriptscriptstyle{{\rm rel}}}+ 
{{\bm v}_{\scriptscriptstyle{A}}}+
{\omega}|\overrightarrow{AP}|\widehat{{\perp{AP}}}\,,
\label{e148}
\end{equation}

\noindent dove 
${{\bm v}_{\scriptscriptstyle{P}}}_{\scriptscriptstyle{{\rm rel}}}$ ha modulo incognito e direzione lungo la scala,
${{\bm v}_{\scriptscriptstyle{A}}}$ \`e un vettore noto e l'altra componente di trascinamento
${{\bm v}_{\scriptscriptstyle{{\rm tr}}}}_{\scriptscriptstyle{{\rm (n)}}}=
{\omega}|\overrightarrow{AP}|\widehat{{\perp{AP}}}$
ha modulo incognito e direzione ortogonale la scala.
Di nuovo una equazione vettoriale dove le direzioni sono tutte note mentre sono incogniti due dei tre moduli. Si possono quindi ricavare sia la velocit\`a relativa di $P$,
${{\bm v}_{\scriptscriptstyle{P}}}_{\scriptscriptstyle{{\rm rel}}}$, sia il
modulo della velocit\`a angolare
$\omega$; il vettore $\bm \omega$ sar\`a come al solito ortogonale
al piano del moto e il suo verso dovr\`a soddisfare la relazione vettoriale
${{\bm v}_{\scriptscriptstyle{{\rm tr}}}}_{\scriptscriptstyle{{\rm (n)}}}=
{\bm\omega}\times\overrightarrow{AP}$.
La velocit\`a del punto
$B$ sar\`a 
\begin{equation}
{{\bm v}_{\scriptscriptstyle{B}}}= 
{{\bm v}_{\scriptscriptstyle{A}}}+
\bm\omega\times\overrightarrow{AB}\,.
\label{e149}
\end{equation}

\noindent Passando alle accelerazioni dovremo in questo caso fare i conti anche con l'accelerazione di {\em Coriolis}, in quanto il sistema relativo ruota con
velocit\`a angolare $\bm\omega$, ormai nota. Al fine di riportare
le espressioni delle accelerazioni ci avvarremo, anche in questo caso,
delle loro componenti, quando queste saranno pi\`u semplici da calcolare.
Dunque, con riferimento alla \ref{e139},
avremo per l'accelerazione del punto $P$
\begin{equation}
{{\bm a}_{\scriptscriptstyle{P}}}=0= 
{{\bm a}_{\scriptscriptstyle{P}}}_{\scriptscriptstyle{{\rm rel}}}+ 
{{\bm a}_{\scriptscriptstyle{A}}}
-\omega^2 \overrightarrow{AP} + \dot{\omega}|\overrightarrow{AP}|\widehat{{\perp{AP}}}+
2{\bm \omega}\times
{{\bm v}_{\scriptscriptstyle{P}}}_{\scriptscriptstyle{{\rm rel}}}\,,
\label{e150}
\end{equation}

\noindent dove il prodotto vettoriale del
termine di {\em Coriolis} ha il solo scopo di indicarne la direzione e il verso. 
Quando la rotazione $\omega$ segue 
le dita della mano destra, il pollice d\`a la direzione e il verso
di $\bm \omega$. Diamo  ora all'indice direzione e verso della velocit\`a
relativa: tenendo le tre dita reciprocamente ortogonali tra
loro, il dito medio fornir\`a la direzione e il verso
dell'accelerazione di {\em Coriolis}. Il procedimento ora indicato
si chiama {\em regola della mano destra}\index{regola della mano destra}
per il prodotto vettoriale.
Dei cinque termini della \ref{e150} tre sono completamente noti
\begin{equation}
{\rm termine\, noto}\equiv
{{\bm a}_{\scriptscriptstyle{A}}}
-\omega^2 \overrightarrow{AP} + 
2{\bm \omega}\times
{{\bm v}_{\scriptscriptstyle{P}}}_{\scriptscriptstyle{{\rm rel}}}\,,
\label{e151}
\end{equation}

\noindent mentre l'accelerazione relativa di $P$ e la componente tangenziale dell'accelerazione di trascinamento presentano modulo incognito
e direzione nota.
Ancora una volta, il riferimento per la soluzione
pu\`o essere lo schema di figura \ref{fig:f116},
che permette di ricavare i moduli delle
due accelerazioni incognite e di conseguenza ci consentir\`a di conoscere
il modulo dell'accelerazione angolare della scala $\dot{\omega}$ prima e poi
$\dot{\bm\omega}$.
Infine si pu\`o ottenere l'accelerazione dell'estremo 
$B$ della scala dalla seguente
\begin{equation}
{{\bm a}_{\scriptscriptstyle{B}}}= 
{{\bm a}_{\scriptscriptstyle{A}}}
-\omega^2 \overrightarrow{AB} + \dot{\bm\omega}\times\overrightarrow{AB}\,.
\label{e152}
\end{equation}


\noindent A conclusione di questo breve esempio di applicazione del teorema di {\em Coriolis} sorge spontanea
una domanda: gli esercizi di figura \ref{fig:f114} e di figura \ref{fig:f115}
appaiono molto simili; come abbiamo dunque fatto ad affibbiare un sistema
relativo traslante al primo, mentre nel secondo caso scegliamo un sistema
relativo traslante-rotante?
La decisione da prendere in tal senso sta forse alla nostra discrezione?
Purtroppo no. Se utilizzassimo
un sistema relativo rotante anche nel primo caso, ci accorgeremmo subito
che la sua rotazione non sarebbe di utilit\`a alcuna. 
Nel secondo caso, introducendo un riferimento relativo che trasla senza ruotare,
non sarebbe possibile giungere alla soluzione.
Una linea guida per scegliere tra il sistema relativo semplicemente traslante e quello
che ruota ``incollato'' a qualche corpo rigido pu\`o essere la 
seguente: tutte le volte che le grandezze relative non possono
in alcun modo essere rappresentate come prodotte dal trascinamento (un classico esempio
\`e la velocit\`a radiale della formica sul piatto della giostra)  \`e 
indispensabile l'introduzione di un sistema di riferimento relativo che, oltre eventualmente a traslare, ruoti.

\section{Ulteriori Considerazioni sui Moti Relativi}

Nelle applicazioni pratiche \`e uso comune riferire i moti a sistemi che si
qualificano come assoluti quando sono in qualche modo fissati a qualche
cosa di molto grande, in generale il nostro pianeta, mentre 
si qualificano come relativi sistemi di riferimento che sono in moto
rispetto alla Terra. Al fine di contrastare (in modo pacifico) tale punto 
di vista,
oltre a tirare in ballo il valido argomento, che quasi tutti conoscono,
che la Terra si muove, ci preme sottolineare che gli unici
moti sperimentalmente rintracciabili sono quelli relativi.
Un insieme di corpi rigidi \`e un teatro nel quale (quando i vincoli
lo permettono) i vari attori posseggono campi di velocit\`a misurabili soltanto
rispetto agli altri corpi.

\begin{figure}[ht]
	\begin{center}
      		\includegraphics[width=0.9\textwidth]{part1/relativi/FIG/arnold_kennedy.pdf}
	\end{center}
\begin{picture}(0,0)(-55,-50)
\scriptsize{
	\put(0,14){corpo $1$}
	\put(234,72){\rotatebox{53}{corpo $3$}}
	\put(108,22){\rotatebox{24}{corpo $2$}}
	\put(115,48){$\omega_{21}$}
	\put(87,14){$A$}
	\put(198,28){$B$}
	\put(261,138){$P$}
	\put(107,7){$C$}
	\put(230,129){$\omega_{32}$}
	\put(222,166){\rotatebox{-38}{${\bm\omega}_{32}\times\overrightarrow{BP}$}}
	\put(235,189){${{\bm v}_{\scriptscriptstyle{P}}}_{\scriptscriptstyle 1}$}
	\put(162,117){${v_{\scriptscriptstyle B}}_{\scriptscriptstyle 1}$}
}
\end{picture}
	\caption{\em Allineamento dei tre centri di rotazione.}
     \label{fig:arnold_kennedy}
\end{figure}

\noindent Cosicch\'e, considerando la figura \ref{fig:arnold_kennedy}, 
possiamo riferire il moto di ciascuno dei tre corpi a ciascuno dei due
rimanenti. In figura si nota che il corpo $1)$ e il corpo $2)$ sono
tra loro uniti da un giunto rotoidale e cos\`i accade per il corpo $2)$ e
il corpo $3)$. Sappiamo quindi che il moto di $2)$ relativamente a $1)$
sar\`a rotatorio con centro $A$ (e naturalmente lo stesso vale per il 
moto di $1)$ rispetto a $2)$), cos\`i come il moto di $3)$ attorno
a $2)$ sar\`a anch'esso rotatorio con centro $B$. 
Volendo descrivere il moto di $3)$ rispetto a $1)$ richiamiamo la regola della
composizione delle velocit\`a (\ref{e133}):
\begin{equation}
{\bm v}_{\rm ass} = {\bm v}_{\rm tr} + {\bm v}_{\rm rel}\,.
\nonumber
\end{equation}

\noindent Pensando come assoluta la velocit\`a
di $P$, punto  appartenente al corpo  $3)$,
relativa al corpo $1)$, ${{\bm v}_{\scriptscriptstyle P}}_{\scriptscriptstyle 1}$, con ovvio significato dei 
simboli avremo  
\begin{equation}
{{\bm v}_{\scriptscriptstyle P}}_{\scriptscriptstyle 1} = {{\bm v}_{\scriptscriptstyle B}}_{\scriptscriptstyle 1} + {{\bm v}_{\scriptscriptstyle P}}_{\scriptscriptstyle 2}\,.
\label{eq:vel_p}
\end{equation}

\noindent La \ref{eq:vel_p} ci dice che la stessa 
${{\bm v}_B}_1$, componente di trascinamento, \`e anche la velocit\`a del
particolare punto del corpo $3)$ coincidente con $B$.
Questa informazione, da sola, ci d\`a la certezza che il centro di istantanea
rotazione del corpo $3)$ nel moto relativo a $1)$, dovendo stare
sulla normale a tale vettore,
giacer\`a da
qualche parte sulla retta che passa per i punti $A$ e $B$.
Se desideriamo individuare il centro di istantanea rotazione $C$ del
corpo $3)$ rispetto al corpo $1)$ (e viceversa),
sempre tramite la \ref{eq:vel_p} calcoliamo la velocit\`a
del generico punto $P$ del corpo $3)$ rispetto al corpo $1)$.
Una volta identificato il vettore
${{\bm v}_{\scriptscriptstyle P}}_{\scriptscriptstyle 1}$, 
come in figura, baster\`a condurre la perpendicolare al suo piede e
$C$ risulter\`a determinato dall'intersezione con
la retta che passa per i punti $A$ e $B$.
Il punto $C$ \`e il centro di rotazione istantaneo
del moto relativo tra $1)$ e $3)$\footnote{Risulta essere 
{\em istantaneo} in quanto non abbiamo premesso alcuna ipotesi circa la
costanza o meno della velocit\`a angolare 
$\omega_{32}$, per questo la posizione di $C$
relativa ai tre corpi pu\`o variare nel tempo. Se tale velocit\`a angolare fosse costante,
il centro $C$ non si sposterebbe.}.
In questa analisi,
l'esistenza dei vincoli rotoidali $A$ e $B$
\`e stata sfruttata solamente per scrivere 
la \ref{eq:vel_p}, la quale sarebbe corretta anche se $A$
fosse il centro istantaneo di rotazione di $2)$ 
rispetto a $1)$ e $B$ il centro istantaneo del moto relativo tra $2)$ e $3)$.
Questo notevole risultato, che vede i tre centri istantanei dei moti relativi di
tre corpi sempre allineati tra loro, costituisce il {\em teorema di
Aronhold-Kennedy}\index{teorema!di Kennedy}. 
La tesi di questo teorema viene spesso utilizzata nell'analisi cinematica
che gli ingegneri civili eseguono sulle strutture:
se tre corpi, tra loro connessi,
 presentano centri istantanei di rotazione relativi allineati,
tali corpi hanno una possibilit\`a (almeno in quell'istante) di muoversi
l'uno rispetto all'altro, situazione questa quasi sempre da evitarsi
nelle costruzioni.
\newpage
\thispagestyle{empty}
\null

\endinput


\part{Trasmissione e Trasformazione del Movimento }
\chapter{Breve Introduzione}

La movimentazione delle macchine, dei macchinari e degli impianti avviene, nella stragrande maggioranza dei casi,
tramite i {\em motori}\index{motori}. La situazione di gran lunga pi\`u diffusa
vede la sorgente primaria di moto in uno o pi\`u
motori che presentano, in generale,  un albero cilindrico tramite il quale
mettono a disposizione dell'utilizzatore
la potenza meccanica necessaria sotto forma di rotazione di tale albero e di un momento torcente ad esso applicato.
Mi preme tranquillizzare subito quei lettori che iniziassero a fremere
pensando a eccezioni 
illustri ai motori rotativi, e dico loro che il periodo  sopra riportato ci
fornisce soltanto un sistema di
riferimento per inquadrare il problema, e nulla toglie alla sua generalit\`a.
Presenteremo infatti in questa parte del volume
la cinematica di alcuni dispositivi che permettono  di
mettere in relazione reciproca movimenti geometricamente
simili, ovvero geometricamente eterogenei tra loro.
Partiremo  quindi
considerando i casi pi\`u semplici di trasformazione di
movimenti rotatori in movimenti ancora rotatori, e ci spingeremo fino a
studiare il comportamento cinematico
di dispositivi che, sfruttando un movimento di ingresso rotatorio
(spesso a velocit\`a costante), producono in uscita leggi di movimento o traiettorie
molto particolari.
Tali leggi di moto e tali particolari traiettorie saranno ottenute tramite
{\em sistemi articolati}\index{sistemi!articolati}
ed {\em eccentrici}\index{eccentrici} che, in gergo, 
prendono il nome di {\em camme}\index{camme}.

\noindent Pi\`u nello specifico
si parler\`a di
{\em trasmissioni}\index{trasmissione!del moto}
nei casi frequentissimi di generazione di un moto rotatorio da uno simile.
Si annoverano in questa famiglia di dispositivi meccanici
le {\em ruote dentate}\index{ruote!dentate}, delle quali un esemplare
\`e mostrato in figura \ref{fig:f2a}, e gli {\em ingranaggi}\index{ingranaggi}
che esse costituiscono,
le {\em cinghie di trasmissione}\index{cinghie di trasmissione} con le
relative {\em pulegge}, le
{\em catene}, i {\em giunti} omocinetici e non ecc. ma anche, di nuovo,
quadrilateri e altri sistemi articolati.
Il che mostra che la trasmissione del moto e la generazione di particolari
movimenti costituiscono, in fondo,
un'unica problematica della Meccanica Applicata alle
Macchine, anche se i due temi, quello delle trasmissioni e quello
dei meccanismi che generano movimenti peculiari,
posseggono specificit\`a cos\`i
profonde da esser tratti, nei libri, quasi sempre in capitoli disgiunti,
quando non occupano trattazioni specialistiche distinte.
Siccome questo volume \`e dedicato alla Cinematica, per quanto
riguarda la trasmissione del moto rotatorio studieremo
soltanto le ruote dentate, sorvolando sulle cinghie di trasmissione che,
bench\'e rappresentino un dispositivo meccanico molto frequentemente
impiegato e manifestino interessanti problematiche dinamiche,
dal punto di vista cinematico non offrono spunti di particolare
spicco.
\begin{wrapfigure}{r}{0.5\textwidth}
     \begin{center}
     \includegraphics[width=0.42\textwidth]{part2/ruote/FIG/f2a.jpeg}
     \end{center}
        \caption{\em Ruota dentata.}
     \label{fig:f2a}
\end{wrapfigure}
\noindent Quando parliamo di ``generazione di particolari movimenti''
intendiamo, in generale, la possibilit\`a di
produrre moti le cui leggi temporali dello spostamento o delle
sue derivate abbiano determinate caratteristiche, oppure di generare
opportune traiettorie a partire da moti rotatori.
I meccanismi  atti a questi particolari scopi verranno
quindi analizzati in ambito cinematico e tali studi, stante lo scopo di questi
appunti e lo spazio limitato, non avranno alcuna pretesa di completezza n\'e di
profondit\`a.
Introdurremo cos\`i i {\em sistemi articolati}\index{sistemi!articolati}, in particolare i quadrilateri, che possono trasformare 
un moto rotatorio con velocit\`a angolare costante in moti
con velocit\`a angolare variabile, oppure condurre qualche loro
punto, normalmente appartenente alla biella, che chiameremo pertanto {\em punto
di biella}\index{biella!punto di}, a percorrere traiettorie
speciali, le quali possono, ad esempio, approssimare
spostamenti rettilinei: traiettorie
queste molto pregiate nell'ambito delle macchine automatiche 
per impiego manifatturiero. 
Vi sono poi i {\em manovellismi}\index{manovellismo} (casi degeneri dei
quadrilateri) che trasformano il moto rettilineo in  moto
rotatorio e viceversa. Essi imperversano da pi\`u di cent'anni nel
settore dei motori a combustione interna, e tuttora tengono banco essendo
presenti (in genere a quaterne) nei motori delle nostre attuali automobili.
\noindent Esiste, infine, la sofisticata possibilit\`a di
provvedere un {\em cedente}\index{cedente} di precise
{\em leggi di moto}\index{leggi di moto}.
Spesso tale cedente consiste in un'asta guidata da due collari,
azionata tramite dispositivi meccanici
eccentrici, opportunamente profilati: le {\em camme}.


\noindent Ciascuna delle famiglie di dispositivi appena
citati merita, e di fatto possiede, trattazioni
specialistiche monotematiche e il loro approfondimento in un corso di Meccanica
Applicata di base risulterebbe eccessivamente tecnico, perci\`o fuori luogo.
La nostra intenzione \`e quella di mostrare, in questa parte del presente tomo,
alcuni (pochi) membri appartenenti alle tipologie citate di meccanismi,
passandoli in rassegna e tentando di fornire,
ove possibile, qualche spunto per la loro analisi cinematica e per 
la loro sintesi.

\chapter{Trasmissione del Moto mediante Ruote Dentate}\index{ruote!dentate}\label{cap_ruote_ev}

\section{Le Ruote di Frizione}
\noindent Credo che a ben pochi lettori sfugga, almeno da un punto di vista
meramente visivo, cosa sia una {\em ruota dentata}.

\begin{wrapfigure}{r}{0.48\textwidth}
      \begin{center}
      \includegraphics[width=0.43\textwidth]{part2/ruote/FIG/f22.pdf}
     \end{center}
\begin{picture}(0,0)(-63,1)
\scriptsize{
\put(21,235){$\bm F$}
\put(-2,190){$r_1$}
\put(4,108){$r_2$}
\put(-55,159){$\omega_1$}
\put(-55,40){$\omega_2$}
\put(21,40){$\bm F$}
\put(65,89){$v_1=v_2$}
}
\end{picture}
        \caption{\em Ruote di frizione.}
     \label{fig:f22}
\end{wrapfigure}

\noindent Le ruote dentate sono comparse alla maggior parte di noi in contesti
eterogenei che vanno dai giochi per bambini agli elettrodomestici, ad alcuni
strumenti da cartoleria, senza parlare dell'emblema della Repubblica Italiana.
Qualche privilegiato ha potuto vedere le ruote dentate impiegate nella
meccanica delle automobili e dei motoveicoli. Tutti possono agevolmente 
immaginare (qualora non ne avessimo contezza diretta) il relativo ingranamento di una ruota dentata con una sua simile.  

\noindent Lasciamo per il momento da parte la caratteristica fondamentale di queste
ruote, che si manifesta tramite una serie di sporgenze opportune, chiamate denti.
Fissiamo invece 
le idee su due cilindri, che presentino raggi adeguati ai nostri
scopi, pensando di calettarli sui
rispettivi alberi, paralleli e liberi di ruotare attorno ai propri assi.
Possiamo fare in modo che le superfici cilindriche siano tra loro
a contatto e che tale contatto sia mantenuto da una data forza $\bm F$.
La figura \ref{fig:f22} rappresenta i due cilindri coi rispettivi alberi.
Ipotizzando che il moto relativo di una ruota rispetto all'altra sia di puro
rotolamento, possiamo affermare che le velocit\`a angolari dei due alberi
dovranno essere tali da produrre velocit\`a periferiche uguali
tra loro sulle due ruote.
Cio\`e dovr\`a valere $v_1=\omega_1 r_1 = v_2 = \omega_2 r_2$, 
relazione che permette di ricavare subito il {\em rapporto di trasmissione}\index{rapporto di trasmissione} della coppia di 
{\em ruote di frizione}\index{ruote!di frizione},
\begin{equation}
\tau={\omega_1\over{\omega_2}}=
{{r_2}\over{r_1}}\,.
\label{eq:tau}
\end{equation}
\noindent I due alberi ruotano pertanto con velocit\`a angolari differenti tra
loro e in strettissima relazione coi raggi delle due ruote che trasmettono
il moto. Nonostante le ruote di frizione siano realmente esistite nel
mondo della meccanica\footnote
{
Un curioso esempio di applicazione: il FERMIAC era un calcolatore analogico basato su ruote di frizione.
}, e si trovino ancora in applicazioni dove le
potenze in gioco sono sia molto grandi sia molto piccole, credo che il
loro pi\`u grave inconveniente potenziale, cio\`e
la possibilit\`a che il loro moto relativo di puro rotolamento venga
contaminato dallo strisciamento, sia evidente anche ad occhi non
particolarmente esperti. Non riteniamo opportuno chiarire
in questa sede il  meccanismo dell'attrito (in prevalenza statico)
necessario a trasmettere correttamente il moto in determinate condizioni,
n\'e illustreremo come possono essere realizzate, nella pratica, le coppie
di ruote di frizione, che tuttavia si fabbricano ancora, come
abbiamo appena affermato, anche
nel contesto di rilevanti potenze da trasmettere.
In questo luogo, ci basta osservare che
la trasmissione corretta del moto, tramite questi dispositivi,
\`e condizionata dal valore della forza $\bm F$ mediante la quale
le ruote sono costrette a mantenere il contatto tra loro.
In cascata, il valore efficace di tale forza \`e subordinato
al coefficiente di attrito tra le due superfici, a sua volta soggetto
a variazioni non leggere del suo valore,
a causa dell'usura e delle condizioni, probabilmente mutevoli, di
lubrificazione delle ruote stesse.
La garanzia del rapporto di trasmissione teorico, derivante dal rapporto tra i
diametri delle ruote, \`e perci\`o negata dai probabili scorrimenti
tra le due primitive nel punto di contatto.


\section{Percorso Intuitivo verso le Ruote Dentate} \label{percorso-ruote}

\noindent Riprendendo il filo dal precedente paragrafo,
siamo naturalmente tentati di impedire lo slittamento della trasmissione
ponendo su una delle ruote
alcune sporgenze e ricavando sulla seconda ruota degli incavi, tali da accogliere
le suddette sporgenze.
Inizieremo a breve un cammino che possiamo definire {\em induttivo}, a valle
del quale riteniamo che la soluzione,
adottata nella pratica come {\em standard} per la creazione delle
ruote dentate, dovrebbe essere ricevuta senza perplessit\`a. 
Premettiamo per\`o alcune considerazioni che riguardano il taglio che abbiamo
desiderato dare all'esposizione del presente argomento.
Le ruote dentate rappresentano nella meccanica gli organi maggiormente
diffusi nell'ambito della trasmissione e della trasformazione del moto, in
generale rotatorio. Si comprende perci\`o facilmente la loro
centralit\`a
in quasi tutte le applicazioni industriali. Crediamo di non esagerare dicendo
che esse rappresentano l'oggetto simbolo della meccanica e del funzionamento
delle macchine.
Pertanto, anche se nella Meccanica Applicata alle Macchine esse godono di una
posizione privilegiata, il loro studio si incontra anche in altre discipline
dell'ingegneria: nel Disegno di Macchine, nella Costruzione di Macchine e nella
Tecnologia Meccanica. Naturalmente, ci\`o accade anche per altri argomenti
di particolare interesse nella nostra 
materia. Tuttavia, ci \`e sembrato di notare 
che le ruote dentate vengano esposte anche altrove
con molta attenzione persino ai loro aspetti cinematici
e geometrici che ne determinano un
buono o cattivo funzionamento: cosa strana, questa, perch\'e l'aspetto
cinematico legato all'ingranamento \`e un argomento piuttosto intricato ma,
forse per questo, si ritiene che non sia mai chiarito a sufficienza.
Cercheremo, come sempre in queste note,
di rimanere circoscritti ai soli aspetti cinematico-geometrici
di questi dispositivi, divagando solamente quando qualche altra loro
caratteristica
(resistenza, facilit\`a di realizzazione e possibilit\`a di impiego) 
sar\`a determinante nel condizionarne la forma. Per tale motivo, questo volume,
che tratta esclusivamente la cinematica, non include
una porzione massiccia di particolarit\`a e di connotazioni delle
ruote dentate, alcune anche molto comuni, altre maggiormente 
tecniche, fondamentali per il loro
effettivo impiego. Un esempio che crediamo significativo:
non tratteremo le ruote con dentatura elicoidale, di grande importanza
nella pratica ma cinematicamente non dissimili dalle loro sorelle
a denti diritti.
Assicuriamo coloro i quali volessero
investigare le ruote dentate nei loro
aspetti pi\`u profondi e nella loro variet\`a di tipologie e applicazioni
che testi molto pi\`u completi, autorevoli e dettagliati
del nostro non mancano. A nostro parere, ad esempio,
questo argomento \`e esposto in modo rigoroso ed esaustivo in
\cite{pellicano}; indichiamo anche un manuale, a livello di scuola superiore,
che riporta un numero notevole di esempi utili per 
il progettista: \cite{punzi}. La ``Bibbia'' in questo campo, ad oggi, rimane \cite{henriot}.

\noindent Come abbiamo gi\`a accennato, riteniamo che
da un punto di vista didattico
sia efficace seguire una sorta di sentiero, percorrendo il quale
saremo costretti ad ammettere
che la possibilit\`a di proseguire sia condizionata dall'accettazione
di alcuni ``precetti'', pi\`u o meno tassativi, dettati da particolari
e stringenti circostanze.
Lungo tale percorso incontreremo pertanto
precetti forti e precetti deboli, e queste due qualit\`a saranno indicate tra
parentesi. Un precetto forte non potr\`a ammettere eccezioni, mentre uno
debole potr\`a essere in qualche caso disatteso, e quando lo sar\`a,
cercheremo di chiarire il perch\'e.
Inizieremo il nostro cammino dalle
ruote di frizione introdotte poco sopra, proponendo
un tentativo per la loro modifica,
in modo da assicurarle dal pericolo di
scorrere l'una sull'altra.
L'esperimento (mentale) che si pu\`o condurre \`e quello
di dotare una delle due ruote di alcune sporgenze che, inserendosi
in opportuni incavi ricavati sull'altra ruota, riescano a trasmettere
la componente tangenziale della forza, che le due ruote si scambiano,
senza basarsi sul fenomeno dell'attrito. Potremmo 
pensare alla ruota che deve essere incavata come fatta di plastilina
o di un qualunque materiale plastico e tenero, che possa poi indurire.
In questa materia plastica la ruota recante le asperit\`a potr\`a
pertanto operare un'azione
di incisione. Questo procedimento di taglio deve essere
condotto pensando che le due circonferenze delle ruote di frizione, 
che chiameremo {\em circonferenze primitive}\index{circonferenza!primitiva},
rotolino l'una sull'altra senza strisciare. Si intuisce che,
i due profili, quello della ruota creatrice
unitamente a quello tagliato nella plastilina,
una volta indurito, saranno
in grado di trasmettere il moto allo stesso modo delle due primitive
rotolanti, cio\`e con rapporto di trasmissione
costante.

\begin{figure}[hbt]
\begin{center}
\includegraphics[width=0.8\textwidth]{part2/ruote/FIG/ruote/belbon.pdf}
\end{center}
\begin{picture}(0,0)(130,0)
\scriptsize{
\put(229,101){$\alpha=180^{\circ}$}
\put(214,34){$\beta=270^{\circ}$}
\put(350,63){ruota 1 tagliata}
\put(250,50){ruota 2 creatrice}
}
\end{picture}
\vskip -7mm
      \caption{\em Dentatura ingenua, 
ottenuta mediante l'applicazione di sporgenze triangolari alla ruota di
sinistra. $\alpha$ \`e l'arco di dentatura della ruota azzurra,
$\beta$ la rotazione della ruota rosa.}
 \label{fig:belbon}
\end{figure}

\noindent La figura \ref{fig:belbon} rappresenta una ruota creatrice,
di colore rosa, recante delle
sporgenze a triangolo isoscele all'esterno della propria circonferenza primitiva, 
esterne cio\`e alla circonferenza della ruota di frizione.
Tale ruota, di colore rosa,``intaglia'', tramite le sue sporgenze,
i denti sulla ruota azzurra.
Il primo precetto che emerge (forte) esige
un lieto fine per la dentatura della ruota tagliata (azzurra).
In altre parole, risulta fondamentale che,
una volta completato il taglio sui $360^{\circ}$ di tale ruota,
essa presenti
un numero intero di denti. Riteniamo inutile commentare questa
indicazione. Aggiungiamo solamente che, se il numero
di denti sulla ruota azzurra risultasse frazionario,
la sua dentatura verrebbe
distrutta da ulteriori passaggi della ruota creatrice.
Allo scopo di prevenire il citato inconveniente e di ottemperare
a questo precetto, la progettazione delle 
ruote dentate si basa principalmente sul loro numero di denti e
non sui loro diametri,
rendendo cos\`i di fatto impossibili rapporti di trasmissione rappresentati
da frazioni stravaganti o, peggio, da numeri irrazionali. Per fortuna
il mondo della tecnica non richiede tali prestazioni!
Scelta quindi una coppia di numeri di denti che meglio si adatta al
rapporto di trasmissione che vogliamo realizzare, giungeremo poi
alla reale dimensione delle ruote introducendo un numero dimensionale
che chiamiamo {\em modulo}\index{modulo}, $m$, misurato in millimetri,
legato al diametro della ruota nel modo seguente: $D= z\, m$, con $z$
numero di denti della ruota e $D$  diametro primitivo (quello
della corrispondente ruota di frizione). Va da s\'e che il
{\em passo}\index{passo}, cio\`e
la misura dell'arco di circonferenza primitiva che separa i due fianchi
attigui e omologhi, sinistri o destri, di due denti vale $p=\pi m$.

\noindent Un secondo precetto (debole) \`e la richiesta, per le asperit\`a del
creatore, di simmetria speculare rispetto a un raggio.
In effetti, le nostre sporgenze sono triangoli isosceli e ci\`o permette di
ottenere condizioni geometriche pari nei due sensi di marcia. Questa
\`e una restrizione che viene, in particolarissimi e rari casi, disattesa.
Ci accorgiamo per\`o, a questo punto, che la ruota azzurra ora
tagliata non potr\`a facilmente generare altri {\em partner}.
La via pi\`u breve per renderci conto di ci\`o consiste nel tentativo
di tagliare una gemella della ruota azzurra
(magari con ugual numero di denti)
utilizzando lei stessa come
creatore.
Purtroppo, trovandosi in questo caso le sporgenze della ruota creatrice
(la azzurra di figura
\ref{fig:belbon}) al di sotto della primitiva,
essa non sar\`a in grado di incidere alcunch\'e sul disco grezzo, lasciandoci
perplessi circa la nostra scelta.
Tuttavia ci accorgiamo che tagliare la {\em partner} della ruota azzurra di
figura \ref{fig:belbon}
tramite lei stessa non \`e il solo metodo che
abbiamo a disposizione.
\begin{figure}[hbt]
\begin{center}
\includegraphics[width=0.72\textwidth]{part2/ruote/FIG/ruote/belbon_interno.pdf}
\end{center}
\begin{picture}(0,0)(0,0)
\scriptsize{
\put(96,95){$\alpha=180^{\circ}$}
\put(175,174){$\beta=90^{\circ}$}
\put(106,62){ruota 1 tagliata}
\put(300,50){ruota 2 creatrice}
}
\end{picture}
\vskip -7mm
      \caption{\em 
Tentativo di dentatura complementare con dentiera a curvatura negativa. La dentatura risulta impossibile.
}
 \label{fig:belbon_interno}
\end{figure}
La ruota rosa (creatrice) di figura \ref{fig:belbon}
possiede un utensile ``duale'' che consiste nella
stessa ruota con curvatura opposta.
Questa inversione, necessaria, della curvatura del creatore non pu\`o
essere chiarita a dovere in questo luogo. Rassicuriamo il lettore che, nel
paragrafo dove si tratta l'ottenimento dei profili coniugati di assortimento,
ci\`o risulter\`a perfettamente illustrato. Qui basta dire che questo
utensile ``duale'', rappresentato in rosa in figura \ref{fig:belbon_interno},
sarebbe perfettamente in grado,
come si intuisce dal suo aspetto, di creare la ruota 
rosa di figura \ref{fig:belbon}.
Ma anche in questo caso, rappresentato appunto
in figura \ref{fig:belbon_interno},
si nota che la ruota creatrice non trova materiale da incidere sullo
stesso disco grezzo dalla quale abbiamo ottenuto la ruota azzurra.
In tale figura viene riportato, con colore tenue, il sovra-metallo necessario
per potere tagliare questa benedetta {\em partner}. Ma \`e chiaro che l'obbligo
di impiegare dischi grezzi
differenti prelude alla necessit\`a di tollerare l'esistenza di
due famiglie di ruote, i membri di una delle quali, probabilmente,
ingraneranno coi membri dell'altra.
Si giunge cos\`i a un terzo precetto (forte) che \`e quello
di potere creare ruote, magari di due generi diversi che chiamiamo
a) e b), in
cui tutte le ruote appartenenti alla famiglia a) possano ingranare con quelle
appartenenti alla famiglia b)\footnote{
Questo requisito minimo per quanto riguarda l'assortibilit\`a delle ruote
dentate viene quasi sempre abbondantemente rispettato, anzi, nelle dentature
ad evolvente, come vedremo, le due famiglie si riducono a una.
}.
L'idea che ci viene in mente per ottemperare a tale precetto, 
quello cio\`e di far s\`i che esistano due famiglie di ruote,
che chiameremo di {\em assortimento}\index{ruote!di assortimento}
 e che siano simmetriche dal punto di vista 
del sovra-metallo che il disco grezzo deve presentare all'esterno
della circonferenza primitiva, consta nel creare sulla dentatrice
un profilo
che stia a cavallo della sua stessa primitiva.
\begin{figure}[hbt]
\begin{center}
\includegraphics[width=0.85\textwidth]{part2/ruote/FIG/ruote/squalo.pdf}
\end{center}
\begin{picture}(0,0)(0,0)
\scriptsize{
\put(290,115){$\alpha=360^{\circ}$}
\put(220,133){$\beta=540^{\circ}$}
\put(85,145){famiglia a)}
\put(85,80){ruota 1 tagliata}
\put(230,50){ruota 2 creatrice}
}
\end{picture}
\vskip -7mm
      \caption{
\em Dentatura ottenuta da sporgenze a ``dente di squalo''. I denti creatori
misurano, sulla primitiva, quanto gli spazi vuoti.
I numeri di denti valgono $z_1=30$ e $z_2=60$.
}
 \label{fig:squalo}
\end{figure}
Il secondo tentativo,
rappresentato in figura \ref{fig:squalo}, mostra una dentatrice a ``denti di 
pescecane'' che taglia la ruota azzurra.
La dentatura \`e in questo caso presentata
completa, sottolineando cos\`i il rispetto del precetto che
impone numeri
di denti interi sulle due ruote, $z_1$ e $z_2$. Anche se confessiamo
 che le stesse figure
\ref{fig:belbon} e \ref{fig:belbon_interno} presenterebbero, sia
sulla ruota creatrice sia su quella tagliata,
un numero intero di denti
qualora tali ruote fossero rappresentate complete:
 non abbiamo avuto il coraggio di scrivere
una procedura che ammettesse le frazioni di dente.
\begin{figure}[hbt]
\begin{center}
\includegraphics[width=\textwidth]{part2/ruote/FIG/ruote/squalo_interno.pdf}
\end{center}
\begin{picture}(0,0)(130,0)
\scriptsize{
\put(385,128){$\alpha=360^{\circ}$}
\put(200,130){$\beta=180^{\circ}$}
\put(346,155){famiglia b)}
\put(346,110){ruota 1 tagliata}
\put(230,47){ruota 2 creatrice}
}
\end{picture}
\vskip -7mm
      \caption{\em Dentatura ottenuta dall'utensile della figura \ref{fig:squalo}
con curvatura opposta.
I numeri di denti valgono $z_1=30$ e $z_2=60$.
}
 \label{fig:squalo_interno}
\end{figure}
In figura \ref{fig:squalo_interno} si accenna alla dentatura che
si otterrebbe per la famiglia b) di ruote complementari alla ruota 
azzurra della famiglia a) di figura \ref{fig:squalo}. Ovviamente, la ruota
appartenente a questa seconda famiglia b),
recante numero di denti pari alla ruota
 dentatrice (rosa)
di figura \ref{fig:squalo}, coinciderebbe con la dentatrice stessa e di
questa famiglia risulterebbe la maggiore\footnote{
Naturalmente, per la famiglia di ruote b) si potrebbero ottenere
ruote pi\`u grandi
della dentatrice di figura \ref{fig:squalo_interno} utilizzando come
utensile una qualsiasi ruota
appartenente alla famiglia a). Non presentiamo in queste note le dentature
interne, ma quello ora esposto rappresenta l'{\em escamotage} che consente
di eseguire dentature interne derivanti dalla dentiera creatrice {\em standard},
che \`e una cremagliera, quindi sicuramente di diametro maggiore rispetto
alla dentatura interna da eseguire.
}.
Ripetiamo che non siamo ancora in grado di valutare  se tutte le ruote della
famiglia a) siano o meno le {\em partner} cinematiche ideali delle ruote b),
ma il nostro fiuto ci convince che  siamo sulla buona strada.
Aggiungiamo, a questi procedimenti ingenui di taglio delle ruote dentate, due
precisazioni.
La prima: il taglio
della maggior parte delle ruote dentate avviene, da un punto di vista
concettuale, proprio nel modo in cui noi le abbiamo tagliate
nell'esperimento mentale, con qualche opportuno accorgimento pratico.
Siccome il materiale da tagliare, che \`e nella gran parte dei casi
acciaio al carbonio, non \`e ``morbido'', il processo di taglio deve avvenire
per asportazione di truciolo. Nei nostri due casi si dovrebbe operare
spostando la ruota creatrice, rosa, fuori dal piano,
per poi accostarla a quella da tagliare; a questo punto
riportarla nel piano, asportando in tal modo il materiale in eccesso
sulla ruota azzurra. Questo procedimento va ripetuto molte volte e a 
piccoli passi, cos\`i da simulare, durante il taglio, l'effettivo
processo di inviluppo continuo che avremmo col rotolamento, l'una
sull'altra, delle due primitive (le circonferenze delle ruote di frizione
riportate in colore rosso e blu). Da qui deriva la seconda precisazione.
Il movimento alternato dell'utensile, che \`e costretto,
uscendo dal piano, ad uscire di scena,
ci consente di avere delle fasi oziose dove siamo
liberi di imporre, alla ruota creatrice, spostamenti non basati sul
rotolamento delle primitive, ma che risultano utili per i nostri fini di taglio.
Infatti, normalmente, una volta guadagnato sulla
ruota creatrice un angolo $\Delta \beta$
corrispondente a $360 /z_2$, si torna indietro con questa 
ruota di tale quantit\`a. Operando in tale maniera, non siamo
costretti ad avere un utensile creatore con un numero di denti elevato,
anzi in genere basta un numero di denti modesto per creare ruote in
un vasto intervallo di numeri di denti. Anche noi ci siamo comportati
in questo modo per disegnare numericamente le nostre figure,
assicurandoci cos\`i di non rendere il codice
che le genera troppo pesante. Inoltre, il semplice movimento di creazione
per inviluppo, che consiste nel rotolamento delle due primitive, pu\`o
essere sostituito da qualsiasi movimento che rispetti tale moto relativo:
nelle nostre figure si \`e scelto di tenere ferma la ruota azzurra e di
movimentare la ruota creatrice di conseguenza. Gli angoli rappresentati
nelle figure \ref{fig:belbon}, \ref{fig:belbon_interno}, \ref{fig:squalo} e
\ref{fig:squalo_interno} si riferiscono esclusivamente al movimento
del creatore, in particolare alla sua rotazione $\beta$ e all'angolo
della porzione di ruota tagliata $\alpha$.
 
\noindent  Non \`e necessaria un'esperienza da meccanico consumato
per constatare  che i denti tagliati sulla ruota 1) di figura \ref{fig:squalo}
sono estremamente scavati alla base.
Notiamo anche che i contatti tra ruota tagliata e creatore e presumibilmente 
tra ruote a) e ruote b) si manifestano sia su superfici (le ruote
dentate possiedono uno spessore) sia su segmenti singolari, come mostrato
in figura \ref{fig:squalo}.
Tali contatti singolari sono in grado, da un punto di vista teorico,
di garantire il corretto e costante rapporto di trasmissione, \index{trasmissione!omocinetica}
che \`e un altro precetto (forte), il quarto.
Detti contatti non sono per\`o idonei a trasferire forze da un elemento
all'altro;
veniamo dunque al quinto precetto (forte):
nella trasmissione del movimento devono necessariamente
essere implicati contatti di strisciamento tra superfici
aventi raggio di curvatura del loro profilo, nel punto di contatto,
 ben diverso da zero. Vedremo
a breve che questo precetto ci porter\`a a considerare utili,
al fine della costruzione delle ruote dentate, esclusivamente opportune
coppie di profili
coniugati, la cui definizione viene rimandata al prossimo paragrafo.

\noindent Altre due intuizioni potrebbero balenarci in mente:
per prima cosa, ci sembra probabile che una riduzione
dell'altezza dei ``denti di squalo'' nella dentatrice possa essere un rimedio
per la forma dei denti
troppo sotto-tagliata alla base della ruota 1); la seconda osservazione
riguarda il raggio della ruota creatrice.
Infatti, qualora utilizzassimo una dentatrice
che rimanesse inalterata
una volta cambiata di segno la
sua curvatura, non avremmo pi\`u due famiglie di
ruote, ma una soltanto, le quali potrebbero (forse) ingranare tutte
quante tra loro e questo, quello cio\`e di ridurre le famiglie di ruote da due
a una, lo diamo come sesto precetto (debole)\footnote{Chi vuole
pu\`o trovare nel capitolo \ref{ruotecy} 
la descrizione di particolari ruote a profilo
cicloidale per le quali si tollera (o si richiede)
la presenza di due famiglie di ruote.}.
Tale simmetria  dell'utensile creatore si chiama
{\em auto-complementarit\`a}\index{auto-complementarit\`a} e
ci conduce  a un creatore a
dentiera rettilinea coi fianchi dei denti pure costituiti da segmenti di retta.
Ma tutte queste intuizioni, che vengono in mente per lo pi\`u a chi le ruote
dentate gi\`a le conosce, 
andrebbero poste in ordine. Per questo siamo costretti
 a toglierci dal nostro comodo sentiero intuitivo
e addentrarci, per poi uscirne il prima possibile, nel pi\`u impervio cammino
che porta alla costruzione dei cosiddetti {\em profili coniugati di
assortimento}\index{profili!di assortimento}.

\section{Costruzione di Profili Coniugati di Assortimento} \label{prof_con}

\noindent Questo paragrafo tratta un problema di cinematica riportato
sulla maggior parte dei
 libri di Meccanica Applicata alle Macchine, ma anche su
alcuni testi di Meccanica Razionale. Prima di 
impostare la definizione dei profili coniugati e di presentare una soluzione
per il loro ottenimento, una volta date le primitive del moto, avvertiamo
il lettore che di tutto questo studio si utilizzeranno soltanto un paio di
conclusioni, legate a due casi molto particolari,
e che tali risultati saranno 
esposti, in riassunto, all'inizio del prossimo paragrafo, al quale rimandiamo
i lettori che non sentono particolare inclinazione per questo argomento.

\noindent Ma cosa sono due {\em profili coniugati}\index{profili!coniugati}?
Due curve regolari qualsiasi possono essere messe tra loro in una
infinit\`a di relazioni reciproche le quali, determinando il loro
posizionamento e il loro movimento relativo,
le rendono tra loro coniugate. 
Come vedremo a breve, ciascuna di queste infinite possibili relazioni, di
natura matematica, che determina il movimento di una curva rispetto all'altra,
deve per\`o possedere alcune caratteristiche (poche) particolari.
Siano $\sigma_1(t)\equiv [\sigma_1(t)_x,\sigma_1(t)_y]^T$ e
$\sigma_2(\tau)\equiv [\sigma_2(\tau)_x,\sigma_2(\tau)_y]^T$,
nel piano cartesiano $(x,y)$, due curve continue e prive di singolarit\`a
in due dati intervalli dei loro parametri, $t$ e $\tau$.
Si scelga ora una funzione, $\tau=f(t)$ definita negli intervalli di interesse,
che mette in corrispondenza biunivoca due intervalli finiti
di coppie di punti appartenenti alle due curve. La relazione tra i due
parametri, $f()$, lega pertanto tra loro i punti
$\sigma_1(t)$ e $\sigma_2(f(t))$ che saranno i candidati
alla coniugazione andando, dopo opportune trasformazioni, a coincidere.
Inoltre, la stessa funzione $f()$ metter\`a naturalmente in relazione anche
le inclinazioni delle tangenti alle curve stesse, che
indicheremo con $\alpha_1(t)$ e $\alpha_2(f(t))$:
anche questi due valori dovranno coincidere nel punto
di coniugazione tra i profili.
Per ora, richiediamo soltanto che $f()$ sia continua e vedremo in seguito
se a tale funzione siano da richiedere o meno altre restrizioni.
Utilizziamo la relazione tra i due parametri per eseguire
alcune {\em trasformazioni piane}\index{trasformazioni piane} di una delle due
curve. In questa sede si \`e deciso ad arbitrio di operare su $\sigma_2$, per
ottenere da tali trasformazioni (continue) la sua coniugata $\sigma_1$.
Procediamo nel modo
seguente; disegnata $\sigma_1(t)$, proponiamo la famiglia di curve
${\sigma_2(\tau)}_t$, che invilupper\`a $\sigma_1$,
mediante la seguente relazione che consta di quattro trasformazioni piane,
le quali legano tra loro le due curve
\begin{equation}
{\sigma_2(\tau)}_t = 
 {\bm T}[\sigma_1(t)]\cdot {\bm R}[\alpha_1(t)]
\cdot{\bm R}[-\alpha_2(f(t))]\cdot
 {\bm T} [-\sigma_2(f(t))]\cdot \sigma_2(\tau)\,.
\label{eq:inviluppo}
\end{equation}

\begin{figure}[hbt]
\begin{center}
\includegraphics[width=0.8\textwidth]{part2/ruote/FIG/ruote/profili_coniugati_definizione_steps.pdf}
\end{center}
\begin{picture}(0,0)(0,0)
\scriptsize{
\put(318,27){$\sigma_1(t)$}
\put(259,84){$\sigma_1(c)$}
\put(209,186){$\sigma_2(\tau)$}
\put(197,182){$a)$}
\put(176,152){$b)$}
\put(215,75){$d)$}
\put(210,95){$c)$}
\put(315,68){$e)$}
\put(192,125){$\sigma_2(f(c))$}
\put(161,68){$0,0$}
\put(159,74){\rotatebox{120}{$\longrightarrow$}}
\put(75,40){$\sigma_2'(\tau)= {\bm T} [-\sigma_2(f(c))]\cdot \sigma_2(\tau)$}
\put(20,156){$\sigma_2''(\tau)={\bm R}[-\alpha_2(f(c))]\cdot \sigma_2'(\tau)$}
\put(132,150){{$\longrightarrow$}}
\put(180,53){$\sigma_2'''(\tau)={\bm R}[\alpha_1(c)]]\cdot \sigma_2''(\tau)$}
\put(215,62){\rotatebox{130}{$\longrightarrow$}}
\put(236,147){$ \sigma_2(\tau)_c={\bm T}[\sigma_1(t)]\cdot \sigma_2'''(\tau)$}
}
\end{picture}
\vskip -7mm
      \caption{\em
Le quattro trasformazioni della formula \ref{eq:inviluppo}.
}
 \label{fig:profili_coniugati_definizione_steps}
\end{figure}
\noindent Le trasformazioni ${\bm T}$ e ${\bm R}$ sono rispettivamente
le traslazioni e le rotazioni nel piano. Esse agiscono sui vettori colonna
dei punti di $\sigma_2$. In particolare, le traslazioni ${\bm T}$ accetteranno
due parametri, dati dai punti $[x,y]^T$ della
curva che contengono come argomento, mentre le rotazioni ${\bm R}$ ne
accetteranno uno solo: tale parametro
\`e l'angolo compreso tra la tangente alla curva, cui l'argomento della
rotazione si riferisce, e una direzione arbitraria e fissa (nelle
nostre figure, quella verticale).
Allo scopo di chiarire la trasformazione operata in \ref{eq:inviluppo}
illustriamo in figura 
\ref{fig:profili_coniugati_definizione_steps} un esempio dove 
la curva $\sigma_1$, riportata in blu, \`e, nella fattispecie, una parabola di
quarto grado rovesciata con vertice nell'origine, mentre la $\sigma_2$, verde,
consiste in una parabola di secondo grado, sempre col vertice nell'origine
e asse verticale.
Infine, la funzione che lega i due parametri \`e semplicemente data
da $\tau=0.3t$. Rimarchiamo che, nonostante tutte le scelte arbitrarie
operate al fine di ottenere la figura
\ref{fig:profili_coniugati_definizione_steps} riflettano una decisa semplicit\`a
(parabole centrate, funzione lineare che lega i due parametri), ci\`o non
inficia la generalit\`a dell'esempio, come vedremo pi\`u avanti, dove lo stesso
codice che genera la figura appena citata viene impiegato per
la generazione di profili coniugati pi\`u complessi. 
Consideriamo la trasformazione \ref{eq:inviluppo} per un valore del parametro
$t=c$.  Tale parametro individua i punti $\sigma_1(c)$ e $\sigma_2(f(c))$,
in rosso nella figura \ref{fig:profili_coniugati_definizione_steps}.
Analizziamo la \ref{eq:inviluppo} partendo da
destra e mettendo in fila le operazioni necessarie per avere le due
curve coniugate nel punto di $\sigma_1(c)\equiv\sigma_2(f(c))$.
Il termine pi\`u a destra, $\sigma_2(\tau)$,
sar\`a semplicemente la curva mobile a), di colore verde, gi\`a citata poc'anzi.
La trasformazione $ {\bm T} [-\sigma_2(f(c))] $ \`e una traslazione che porta il 
punto di $\sigma_2$ destinato ad essere coniugato (col punto ``c-esimo'' di
$\sigma_1$) nell'origine degli assi. Il risultato \`e rappresentato dalla 
curva di colore violetto b) e nome $\sigma_2'(\tau)$.
L'allineamento degli spazi tangenti alle due curve, nei
punti di coniugazione (di parametri $c$ e $f(c)$), avviene in due fasi.
Applichiamo alla $\sigma_2'(\tau)$ una prima rotazione 
di segno contrario alla tangente a $\sigma_2'$ in $\tau=f(c)$, ottenendo in tal
modo $\sigma_2''(\tau)={\bm R}[-\alpha_2(f(c))]\sigma_2'(\tau)$,
parabola c), riportata in colore azzurro chiaro. Applichiamo quindi una seconda
rotazione, pari all'inclinazione della tangente a $\sigma_1$ in $c$, ottenendo
cos\`i $\sigma_2'''(\tau)={\bm R}[\alpha_1(c)]\cdot\sigma_2''(\tau)$,
che \`e la parabola in colore grigio d).
\begin{figure}[hbt]
\begin{center}
\includegraphics[width=0.8\textwidth]{part2/ruote/FIG/ruote/profili_coniugati_definizione_a.pdf}
\end{center}
\begin{picture}(0,0)(0,0)
\scriptsize{
\put(315,42){$\sigma_1$}
\put(206,189){$\sigma_2$}
\put(67,66){$\lambda_1$}
\put(251,176){$\lambda_2$}
}
\end{picture}
\vskip -7mm
      \caption{\em
Profili coniugati: inviluppo e relative polari.
}
 \label{fig:profili_coniugati_definizione_a}
\end{figure}
Finalmente, non rimane
che traslare la curva $\sigma_2'''(\tau)$, 
la quale si trova ora col suo ``c-esimo'' punto in (0,0) e con la tangente in
quel punto inclinata come la tangente a $\sigma_1(c)$,
sul punto di coniugazione con $\sigma_1$, e questo si esegue
mediante la traslazione $\sigma_2(\tau)_c={\bm T}[\sigma_1(c)]\sigma_2'''(\tau)$,
ottenendo in tal modo la parabola e).
Il risultato di queste operazioni, al variare del punto
di coniugazione, cio\`e al variare di $c$,  \`e quello mostrato in figura
\ref{fig:profili_coniugati_definizione_a}. La parabola di colore
rosso, che crea l'inviluppo di $\sigma_1$, viene rappresentata,
onde evitare un eccessivo
appesantimento dell'immagine, soltanto in una ventina di posizioni,
rendendo comunque perfettamente l'idea del processo di inviluppo.
Da un punto di vista matematico, le condizioni necessarie e sufficienti
affinch\'e la $\sigma_2$, durante il suo movimento, inviluppi la curva
$\sigma_1$ sono due, e precisamente \`e richiesto che i punti omologhi
delle due curve coincidano per tutti i valori omologhi dei due parametri (nei
due intervalli di interesse) e che, in tali punti, anche le tangenti
alle due curve siano la stessa retta\footnote
{
L'inviluppo di una curva, espressa
come $\phi(x,y,\lambda)=0$ dove $\lambda$ \`e il parametro che la muove nel
piano, si ottiene eliminando il parametro stesso
dal sistema
\begin{equation}
\begin{cases}
\phi(x,y,\lambda)=0 \\
\displaystyle \frac{\partial \phi(x,y,\lambda)} {\partial\lambda}=0\,,
\end{cases}
\nonumber
\end{equation}

\noindent il che \`e del tutto equivalente a quanto da noi richiesto.
}.

\vskip .2cm
\footnotesize
\noindent Si intuisce facilmente che entrambe queste condizioni sono assicurate
dal procedimento descritto dalla \ref{eq:inviluppo} e illustrato in figura
\ref{fig:profili_coniugati_definizione_a} ma, rischiando di passare per
pedanti, ne riportiamo comunque la dimostrazione.
Quanto alla prima condizione, cio\`e la prescrizione che le due curve
$\sigma_2$ e l'inviluppo creato dalla \ref{eq:inviluppo}, coincidano,
essa deriva con semplicit\`a da ci\`o che
abbiamo discusso circa la  struttura della \ref{eq:inviluppo} stessa.
Considerando infatti un valore del
parametro $t=c$ \hspace{.1 cm} e, di conseguenza, $\tau=f(c)$, avremo
\begin{equation}
{\sigma_2(f(c))}_c = 
 {\bm T}[\sigma_1(c)]\cdot {\bm R}[\alpha_1(c)]\cdot{\bm R}[-\alpha_2(f(c))]\cdot
 {\bm T} [-\sigma_2(f(c))]\cdot \sigma_2(f(c))\,.
\label{eq:dimost_x}
\end{equation}
\noindent Ma considerando la prima (si parte sempre da destra) delle
trasformazioni della \ref{eq:dimost_x} abbiamo
\begin{equation}
{\bm T} [-\sigma_2(f(c))]\cdot \sigma_2(f(c))=(0,0)^T.
\label{eq:dimost_x1}
\end{equation}
\noindent Pertanto, tenendo presente che le rotazioni non hanno effetto sul vettore
nullo, avremo
\begin{equation}
{\sigma_2(f(c))}_t = 
 {\bm T}[\sigma_1(c)]\cdot {\bm R}[\alpha_1(c)]\cdot{\bm R}[-\alpha_2(f(c))]\cdot
 \left(\begin{matrix} 0 \\ 0 \\ \end{matrix}\right)= 
 {\bm T}[\sigma_1(c)]\cdot 
 \left(\begin{matrix} 0 \\ 0 \\ \end{matrix}\right)= 
{\sigma_1(c)}\,. 
\label{eq:dimost_p}
\end{equation}
\noindent Poco pi\`u laboriosa \`e la dimostrazione
della coincidenza delle tangenti alle due curve nel punto di coniugazione.
Analizziamo, a tale proposito, in che modo la \ref{eq:dimost_x}
trasforma il punto infinitamente vicino a quello di parametro $c$
della curva $\sigma_2$. Tale punto risulta individuato dalla seguente espressione
differenziale
\begin{equation}
\sigma_2(f(c)+{\rm d}\tau))=\sigma_2(f(c))+
{{\rm d}{\sigma_2}\over{{\rm d}\tau}}\bigr|_{(f(c))} 
{\rm d} \tau \,, 
\label{eq:dimost_t1}
\end{equation}
\noindent dove il termine ${{\rm d}{\sigma_2}\over{{\rm d}\tau}}\bigr|_{(f(c))}
{\rm d} \tau$ rappresenta un vettore infinitesimo
diretto come la tangente alla $\sigma_2$ nel punto $\sigma_2(f(c))$.
Chiamiamo con $\bm \delta_c$ il vettore infinitesimo 
\begin{equation}
{\bm \delta_c}= \sigma_2(f(c)+{\rm d}\tau)_c- \sigma_2(f(c))_c\,,
\label{eq:vett_inf}
\end{equation}
\noindent intendendo
con i due addendi a destra del segno di uguaglianza l'applicazione della
\ref{eq:dimost_x} sia al punto $\sigma_2(f(c)+{\rm d}\tau))$, sia
al punto $\sigma_2(f(c))$. Data la linearit\`a degli operatori
presenti nella
\ref{eq:dimost_x} vale, per gli operandi, la propriet\`a distributiva
\begin{equation}
{\bm T}()\cdot{\bm R}()\cdot {\bm R}()\cdot{\bm T}()\cdot({\bm a}+{\bm b})=
{\bm T}()\cdot{\bm R}()\cdot {\bm R}()\cdot{\bm T}()\cdot{\bm a}+
{\bm T}()\cdot{\bm R}()\cdot {\bm R}()\cdot{\bm T}()\cdot{\bm b}\,.
\end{equation}

\noindent Ci\`o che rimane dall'applicazione della \ref{eq:dimost_x}
risulta quindi essere
\begin{equation}
{\bm \delta_c} =
 {\bm T}[\sigma_1(c)]\cdot {\bm R}[\alpha_1(c)]\cdot{\bm R}[-\alpha_2(f(c))]
\cdot 
 {\bm T} [-\sigma_2(f(c))]
	\cdot (
{{\rm d}{\sigma_2}\over{{\rm d}\tau}}\bigr|_{(f(c))} 
{\rm d} \tau )\,.
\label{eq:dimost_x3}
\end{equation}
\noindent Sarebbe tedioso ripercorrere per gradi le operazioni
illustrate in figura \ref{fig:profili_coniugati_definizione_steps}, pertanto ci
limitiamo qui a seguire velocemente le trasformazioni del vettore tangente infinitesimo
${{\rm d}{\sigma_2}\over{{\rm d}\tau}}\bigr|_{(f(c))}
{\rm d} \tau$. Ricordando che si parte sempre da destra, esso si muove, mediante
il primo operatore di traslazione,
nell'origine degli assi. Quindi, tramite le due rotazioni, tale vettore
si dispone come la
tangente alla curva $\sigma_1$ nel punto $t=c$. Infine,
mediante l'ultima traslazione $\bm \delta_c$, esso si ritrover\`a
proprio in tale punto
di $\sigma_1$. Risulta cos\`i provato che anche i vettori tangenti alla
$\sigma_2$ diventano, dopo il procedimento di coniugazione, altrettanti
vettori tangenti alla curva $\sigma_1$.



\vskip .2cm
\normalsize
\noindent Abbiamo gi\`a ammesso che le due dimostrazioni, test\'e sviluppate e
stampate con tipi minori, si
presentano come leggermente superflue, essendo il procedimento contenuto
nella \ref{eq:inviluppo}, ed esposto graficamente in
figura \ref{fig:profili_coniugati_definizione_steps}, intrinsecamente volto
ad ottenere la coniugazione dei profili. Tali dimostrazioni ci sono per\`o
utili per meglio circoscrivere i vincoli matematici ai quali $\sigma_1$,
$\sigma_2$ e $f()$ devono sottostare. Affinch\'e si possa manifestare
la coincidenza nei punti di coniugazione, il che equivale alla possibilit\`a
di scrivere la \ref{eq:dimost_x}, la funzione $\tau=f(t)$, che lega i
parametri delle due curve, deve essere definita per tutti i valori di $t$
nell'intervallo di
interesse. Solo questo? S\`i, soltanto questo. Ad essa non \`e chiesto n\'e
di essere monotona n\'e di sottostare
a restrizioni sulle sue derivate, anzi non \`e richiesta neppure
la sua derivabilit\`a. Deve essere chiaro per\`o che ragionando
in questo modo, cio\`e ammettendo che $f()$ possa manifestare qualsiasi
bizzarria, \`e probabile che non si otterranno profili coniugati con
caratteristiche adatte all'impiego nella
meccanica delle macchine. Per quanto riguarda la possibilit\`a
di scrivere la \ref{eq:dimost_t1}, che costituisce la base della seconda 
dimostrazione, circa la coincidenza delle tangenti ai profili, \`e richiesto
che $\sigma_2$, e di conseguenza $\sigma_1$, siano differenziabili, e nulla pi\`u.
Certo, si intuisce con facilit\`a
che aggiungendo altre restrizioni sia ai profili (immaginandoli
ad esempio curve regolari e lisce), sia
alla funzione che lega tra loro i due parametri (che si potrebbe ipotizzare
continua e monotona), si intuisce, dicevamo, che le speranze di ottenere
profili coniugati lisci, privi di singolari\`a e interferenze 
sarebbe superiore.
Ma, come vedremo, tutte queste apprensioni circa la bont\`a dei profili e della
relazione cinematica che li fa scorrere l'uno sull'altro sono fuori luogo nella
meccanica applicata, dove saremo costretti ad ammettere
che i (il?) soli profili interessanti sono anche estremamente semplici.  
Una volta nota, tramite la \ref{eq:inviluppo}, la traiettoria
di due punti qualsiasi appartenenti alla curva $\sigma_2$, chiamiamoli
$\sigma_2(\tau_1)$ e $\sigma_2(\tau_2)$, mentre essa inviluppa $\sigma_1$,
possiamo individuare, al variare del parametro $t$,
le normali a tali traiettorie. Per ciascun valore di $t$,
dall'intersezione di dette normali, otterremo
il centro istantaneo di rotazione di $\sigma_2$. Il luogo geometrico
delle tracce dei centri istantanei di rotazione di $\sigma_2$
si chiama {\em polare fissa}\index{polare!fissa} e verr\`a indicata con $\lambda_1$.
La {\em polare mobile}\index{polare!mobile}, $\lambda_2$,
si ottiene invece lasciando che, una volta individuato il
centro istantaneo di rotazione per un dato valore di
$t$, tale punto venga trascinato
dal piano mobile. \`E sottinteso che, come sempre accade coi moti relativi,
le parti possono essere 
invertite tenendo ferma $\sigma_2$ e facendo muovere $\sigma_1$. Durante tale moto
le polari verranno a scambiarsi tra loro ottenendo, in questo caso, la 
precedente polare mobile come attuale polare fissa.
Polare fissa e polare mobile sono a contatto
tra loro nel centro istantaneo di rotazione corrispondente ad un 
dato valore del parametro $t$ e il loro movimento relativo \`e di puro rotolamento.
In figura \ref{fig:profili_coniugati_definizione_a} sono rappresentate
le polari del moto di $\sigma_2$, fissa e mobile, mediante un centinaio
di punti scelti dal migliaio a nostra disposizione tramite
il codice che genera la stessa figura
\ref{fig:profili_coniugati_definizione_a} e le altre
\ref{fig:profili_coniugati_definizione_b} e
\ref{fig:profili_coniugati_definizione_b1} e che,
qualora fossero tutti disegnati, avrebbero reso continue tali curve.
La rappresentazione
di un minor numero di punti consente di rendere l'idea,
tramite le distanze reciproche tra i punti stessi,
della equivalenza delle lunghezze di tratti
omologhi delle due polari, equivalenza implicata dal rotolamento.
In generale le polari sono curve poco intuitive e possono presentare
discontinuit\`a e contenere punti impropri.

\begin{figure}[hbt]
\centering
\begin{minipage}[b]{0.45\textwidth}
\centering
\includegraphics[width=0.9\textwidth]{part2/ruote/FIG/ruote/profili_coniugati_definizione_b.pdf}
\begin{picture}(0,0)(130,0)
\scriptsize{
\put(123,33){$\sigma_1$}
\put(46,83){$\sigma_2$}
\put(-3,26){$\lambda_1$}
\put(114,68){$\lambda_2$}
}
\end{picture}
      \caption{\em
Curva $\sigma_2$ in scala ridotta e nuove polari.
}
 \label{fig:profili_coniugati_definizione_b}
\end{minipage}\hfill
\begin{minipage}[b]{0.45\textwidth}
\includegraphics[width=0.9\textwidth]{part2/ruote/FIG/ruote/profili_coniugati_definizione_b1.pdf}
\begin{picture}(0,0)(130,0)
\scriptsize{
\put(124,5){$\sigma_1$}
\put(30,23){$\sigma_2$}
\put(-19,8){\rotatebox{-90}{$\longrightarrow$}}
\put(-13,-3){polare fissa e polare mobile all'infinito}
}
\end{picture}
      \caption{\em
Profili coniugati traslanti.
      }
 \label{fig:profili_coniugati_definizione_b1}
\end{minipage}
\end{figure}


\noindent Una situazione non molto dissimile da quella di figura 
\ref{fig:profili_coniugati_definizione_a} \`e quella rappresentata in
\ref{fig:profili_coniugati_definizione_b}, dove la curva $\sigma_2$ \`e
riportata in scala minore e $\tau=0.07t$, allo scopo di mostrare quale
cambiamento macroscopico subiscono le due polari del moto a fronte di un
cambiamento, apparentemente leggero, di uno dei due profili coniugati.
In figura \ref{fig:profili_coniugati_definizione_b1} viene
invece mostrata
una curva $\sigma_2$ che inviluppa, tramite un moto esclusivamente 
traslatorio, la curva $\sigma_1$, sua simile ma in scala maggiore. Per questo
esempio \`e stata infatti scelta una funzione $f()$, che lega i due parametri,
tale da portare a coniugazione i punti delle due curve aventi, a priori,
 la stessa
tangente. \`E proprio questo il motivo che rende traslatorio il moto di
inviluppo e che manda le polari del moto all'infinito rendendole
due segmenti di rette improprie la cui direzione coincide con quella di
spostamento del profilo  $\sigma_2$. I profili coniugati che abbiamo finora analizzato sono
stati ottenuti tramite curve $\sigma_2(\tau)$ molto semplici,
in pratica delle parabole.  Altrettanto semplici sono
i legami $\tau=f(t)$ che definiscono, tramite la $\sigma_1(t)$,
le leggi degli spostamenti di $\sigma_2$.  In questo modo,
gli inviluppi ottenuti somigliano a oggetti effettivamente
utilizzabili nella meccanica delle macchine.
\begin{figure}[hbt]
\begin{center}
\includegraphics[width=0.8\textwidth]{part2/ruote/FIG/ruote/profili_coniugati_definizione_c.pdf}
\end{center}
\begin{picture}(0,0)(0,0)
\scriptsize{
\put(292,123){$\sigma_1$}
\put(198,87){$\sigma_2$}
\put(26,228){$\lambda_1$}
\put(62,92){$\lambda_1$}
\put(269,98){$\lambda_1$}
\put(103,313){$\lambda_2$}
\put(201,33){$\lambda_2$}
\put(129,50){$\lambda_2$}
}
\end{picture}
\vskip -5mm
      \caption{\em
Profili coniugati che si intersecano e polari discontinue.
}

 \label{fig:profili_coniugati_definizione_c}
\end{figure}
Come abbiamo per\`o ripetuto altre volte, la teoria dei profili coniugati
non mette restrizioni, se non
molto blande, alle tre entit\`a matematiche test\'e menzionate.
Per esempio, le curvature di $\sigma_1$ e  $\sigma_2$ possono essere
tali che, combinate alla funzione $f()$, producano inviluppi che 
si intersecano con le curve che li generano. In questi casi,
le polari del moto presentano spesso discontinuit\`a e punti
all'infinito,
come mostra la figura \ref{fig:profili_coniugati_definizione_c}.
Nel caso riportato in figura,
le curve $\sigma_1$ e $\sigma_2$ sono ancora parabole,
quindi curve semplici e regolari, e la funzione che lega i parametri delle due
curve \`e ancora $\tau=0.3t$. Ma le curvature dei due profili
risultano tra loro ``incompatibili''. Ben inteso, da un punto di
vista puramente matematico tutto fila liscio come sempre.
Nella tecnica, per\`o, tali profili non possono essere utilizzati
per trasmettere il  movimento o per altre funzioni concrete
quindi, per noi ingegneri,
rimangono curiosit\`a che appartengono alla teoria.
Nella figura
si possono anche vedere tutti gli spezzoni delle due polari per le quali
si intravedono gli asintoti che indicano le direzioni
dei loro punti impropri.
Capovolgiamo il problema. Supponiamo di conoscere le polari
del moto relativo, e investighiamo la possibilit\`a
di trovare due profili tra loro coniugati. Dettando le polari
la legge geometrica del moto relativo
del  piano mobile rispetto a quello fisso, ad ogni curva appartenente
al primo corrisponder\`a il relativo inviluppo sul secondo,
il quale, del profilo mobile, sar\`a anche il profilo coniugato.
\begin{wrapfigure}{r}{0.55\textwidth}
      \begin{center}
      \includegraphics[width=0.50\textwidth]{part2/ruote/FIG/ruote/profili_coniugati_invil_polar.pdf}
     \end{center}
\begin{picture}(0,0)(0,0)
\scriptsize{
\put(99.6,65){$\eta_0$}
\put(94.2,71.3){\rotatebox{120}{$\longrightarrow$}}
\put(60,100){$\lambda_2$}
\put(43,45){$\lambda_1$}
\put(170,145){$\sigma_1$}
\put(78,135){$\sigma_2$}
\put(-4,163){$\sigma_1(f(\eta_0))\equiv\sigma_2(g(\eta_0))$}
\put(36,162){\rotatebox{-64}{$\longrightarrow$}}
}
\end{picture}
\vskip -6.3mm
        \caption{\em
Profilo $\sigma_1$ coniugato a $\sigma_2$ ottenuto dall'inviluppo di $\sigma_2$,
trascinata dalla polare mobile $\lambda_2$ che rotola su $\lambda_1$.
}
     \label{fig:profili_coniugati_invil_polar}
\end{wrapfigure}

\noindent Notiamo infatti che le trasformazioni contenute nella
\ref{eq:inviluppo} possono essere riferite a qualsiasi altra coppia di
profili coniugati compatibili con il moto relativo
esplicitato dalla stessa formula. Tra queste coppie di profili
coniugati vi \`e, come caso particolare, quella formata alle polari,
per le quali la legge che lega i due parametri \`e semplicemente
$\tau=t$.
Indicando cos\`i le polari del moto con $\lambda_1$ e $\lambda_2$,
figura \ref{fig:profili_coniugati_invil_polar},
la relazione \ref{eq:inviluppo} \`e perfettamente equivalente a
\begin{equation}
{\sigma_2(\tau)}_\eta = 
 {\bm T}[\lambda_1(\eta)]\cdot {\bm R}[\beta_1(\eta)]\cdot{\bm R}[-\beta_2(\eta)]\cdot
 {\bm T} [-\lambda_2(\eta)]\cdot \sigma_2(\tau)\,,
\label{eq:inviluppo_polar}
\end{equation}

\noindent con ovvio significato degli angoli $\beta_1$ e $\beta_2$.
L'espressione \ref{eq:inviluppo_polar} indica la famiglia di 
curve $\sigma_2$ rappresentata in figura \ref{fig:profili_coniugati_invil_polar}
in colore violetto.
Basta, in questo caso, un solo parametro $\eta$ a determinare in quale
dei loro punti le polari sono a contatto, in quanto
esse rotolano l'una sull'altra senza strisciare.
Riteniamo superfluo dimostrare che la famiglia di curve 
${\sigma_2(\tau)}_\eta$ inviluppa il profilo $\sigma_1$, che nella figura
\ref{fig:profili_coniugati_invil_polar} non viene tracciato in modo esplicito,
 in quanto dovremmo nuovamente
ripetere la dimostrazione svolta tramite le formule da \ref{eq:dimost_x} fino a
\ref{eq:dimost_x3} e i relativi commenti.
I parametri di $\sigma_1(t)$ e di $\sigma_2(\tau)$ saranno legati al parametro
$\eta$ da due relazioni imposte dal processo di inviluppo, $t=f(\eta)$ e
$\tau=g(\eta)$, le quali naturalmente richiederanno che per qualunque
valore del parametro $\eta$ si abbia $\sigma_1(f(\eta))=\sigma_2(g(\eta))$.
Nella figura \ref{fig:profili_coniugati_invil_polar} i due profili $\sigma_1$ e $\sigma_2$
si toccano in $\sigma_1(f(\eta_0))=\sigma_2(g(\eta_0))$, e, ad arbitrio,
abbiamo imposto $\lambda_1(\eta_0)=\lambda_2(\eta_0)=(0,0)^T$, cio\`e che 
le due polari si tocchino, in quel caso,
nell'origine delle coordinate cartesiane del piano fisso.
Le trasformazioni contenute nella \ref{eq:inviluppo_polar} dovrebbero
ormai risultare
famigliari; partendo, come le altre volte,  da destra: 
imponiamo una traslazione del punto della polare $\lambda_2$
di parametro $\eta$ in $(0,0)$, segue una rotazione della stessa di un
angolo pari all'opposto dell'inclinazione della $\lambda_2$ in $\eta$,
$-\beta_2(\eta)$, ancora una rotazione pari all'inclinazione
della $\lambda_1$, sempre nel punto di
parametro $\eta$, $\beta_1(\eta)$. Infine, la traslazione 
nel punto di $\lambda_1(\eta)$ completer\`a la trasformazione;
rimandiamo il lettore alla figura
\ref{fig:profili_coniugati_definizione_steps} per maggior chiarezza.
Rimane sottinteso che, in tutti questi spostamenti e rotazioni, la polare
mobile trascina con s\'e il piano mobile e quindi $\sigma_2$.
\begin{figure}[hbt]
\begin{center}
\includegraphics[width=0.8\textwidth]{part2/ruote/FIG/ruote/profili_coniugati_invil.pdf}
\end{center}
\begin{picture}(0,0)(0,0)
\scriptsize{
\put(177,110){$\eta_0$}
\put(172,116){\rotatebox{120}{$\longrightarrow$}}
\put(329,210){$\sigma_a$}
\put(312,270){$\sigma_b$}
\put(144,233){$\sigma_2$}
\put(78,64){$\lambda_{1a}$}
\put(49,119){$\lambda_{1b}$}
\put(133,160){$\lambda_2$}
}
\end{picture}
\vskip -6mm
      \caption{\em
Profili coniugati inviluppati dalla curva $\sigma_2$ trascinata dalla
polare $\lambda_2$ che rotola sia su $\lambda_{1a}$, sia su $\lambda_{1b}$.
}
 \label{fig:profili_coniugati_invil}
\end{figure}
\noindent Un caso di grande interesse, anzi, il caso che giustifica ampiamente
il procedere, che ammettiamo essere piuttosto laborioso,
di questo paragrafo, \`e quello in cui la
polare $\lambda_2$ rotola su due diverse polari:
identifichiamole coi nomi $\lambda_{1a}$ e $\lambda_{1b}$ e coi
colori verde e blu, come rappresentato
in figura \ref{fig:profili_coniugati_invil}. Anticipiamo, tanta \`e l'importanza
che attribuiamo a questo passo, che una volta ottenuti i due inviluppi, mediante
il rotolamento di $\lambda_2$ su $\lambda_{1a}$ e $\lambda_{1b}$,
sar\`a $\lambda_{1a}$ a rotolare su $\lambda_{1b}$, permettendoci cos\`i di
osservare che i due inviluppi,
precedentemente ottenuti, sono coniugati tra loro.
Le relazioni che descrivono le famiglie delle $\sigma_2$
che inviluppano $\sigma_a$ e $\sigma_b$, seguendo \ref{eq:inviluppo_polar},
saranno date dalla seguente coppia di relazioni
\begin{equation}
{\sigma_{2a}(\tau)}_\eta = 
 {\bm T}[\lambda_{1a}(\eta)]\cdot {\bm R}[\beta_{1a}(\eta)]\cdot{\bm R}[-\beta_2(\eta)]\cdot
 {\bm T} [-\lambda_2(\eta)]\cdot \sigma_2(\tau)\,,
\label{eq:inviluppo_polara}
\end{equation}
\begin{equation}
{\sigma_{2b}(\tau)}_\eta = 
 {\bm T}[\lambda_{1b}(\eta)]\cdot {\bm R}[\beta_{1b}(\eta)]\cdot{\bm R}[-\beta_2(\eta)]\cdot
 {\bm T} [-\lambda_2(\eta)]\cdot \sigma_2(\tau)\,.
\label{eq:inviluppo_polarb}
\end{equation}
\noindent Come premesso, procederemo ora in questo modo: anzich\'e fare rotolare 
la $\lambda_2$ sulle altre due primitive, facciamo rotolare la polare
$\lambda_{1a}$ sulla $\lambda_{1b}$. Il rotolamento di $\lambda_{1a}$
trasciner\`a naturalmente le
$\sigma_{2a}(\tau)_\eta$, cio\`e le curve che stanno a primo membro nella
\ref{eq:inviluppo_polara} e che definiscono l'inviluppo di $\sigma_a$.
Desideriamo dimostrare che tali curve, trascinate
dal movimento di $\lambda_{1a}$, inviluppano anch'esse la $\sigma_{b}$,
ovvero che la $\sigma_a$ inviluppa la $\sigma_b$ durante il rotolamento di 
$\lambda_{1a}$ su $\lambda_{1b}$.
Tale inviluppo sar\`a indicato con ${{\sigma_{2ab}(\tau)}_\eta}_\mu$
e, non avendo direttamente a disposizione l'espressione 
di $\sigma_a$, utilizzeremo la famiglia di curve che ne definisce l'inviluppo,
esplicitate dalla \ref{eq:inviluppo_polara}, ottenendo
\begin{multline}
{{\sigma_{2ab}(\tau)}_\eta}_\mu = 
 {\bm T}[\lambda_{1b}(\mu)]\cdot {\bm R}[\beta_{1b}(\mu)]\cdot{\bm R}[-\beta_{1a}(\mu)]\cdot
 {\bm T} [-\lambda_{1a}(\mu)]\cdot {\sigma_{2a}(\tau)}_\eta = \\
 = 
 {\bm T}[\lambda_{1b}(\mu)]\cdot\color{red} {\bm R}[\beta_{1b}(\mu)]\cdot{\bm R}[-\beta_{1a}(\mu)]\cdot
 {\bm T} [-\lambda_{1a}(\mu)]\cdot \\
\color{red} {\bm T}[\lambda_{1a}(\eta)]\cdot {\bm R}[\beta_{1a}(\eta)]\cdot{\bm R}[-\beta_2(\eta)] \color{black}\cdot
 {\bm T} [-\lambda_2(\eta)]\cdot \sigma_2(\tau)\,.
\label{eq:inviluppo_polarb1}
\end{multline}
\noindent Ricordando che le polari rotolano l'una sull'altra senza strisciare,
selezioniamo da questa doppia famiglia di curve quella di parametro
$\mu=\eta$ e componiamo le trasformazioni, scritte in colore rosso,
contenute nella \ref{eq:inviluppo_polarb1} nel seguente modo
\begin{multline}
\color{red} {\bm R}[\beta_{1b}(\eta)]\cdot{\bm R}[-\beta_{1a}(\eta)]\cdot
 {\bm T} [-\lambda_{1a}(\eta)]\cdot 
 {\bm T}[\lambda_{1a}(\eta)]\cdot {\bm R}[\beta_{1a}(\eta)]\cdot{\bm R}[-\beta_2(\eta)]=\\
={\bm R}[\beta_{1b}(\eta)]\cdot{\bm R}[-\beta_2(\eta)]\,, 
\label{eq:id}
\end{multline}
\noindent da cui si ricava immediatamente
\begin{equation}
{{\sigma_{2ab}(\tau)}_\eta}_\eta = 
 {\bm T}[\lambda_{1b}(\eta)]\cdot {\bm R}[\beta_{1b}(\eta)]\cdot
{\bm R}[-\beta_2(\eta)] \cdot
 {\bm T} [-\lambda_2(\eta)]\cdot \sigma_2(\tau)\,.
\label{eq:inviluppo_polarb2}
\end{equation}
\noindent Le espressioni del secondo membro delle \ref{eq:inviluppo_polarb2}
e \ref{eq:inviluppo_polarb} si equivalgono, pertanto
la famiglia di curve ora ottenuta coincide con quella
che si trova facendo rotolare la $\lambda_2$ sulla $\lambda_{1b}$:
${{\sigma_{2ab}(\tau)}_\eta}_\eta\equiv {\sigma_{2b}(\tau)}_\eta$,
come si evince anche dalla figura
\ref{fig:profili_coniugati_invil_lagrange}, che
rappresenta la situazione in cui la polare $\lambda_{1a}$ rotola sulla
$\lambda_{1b}$ toccandola nel punto
$\lambda_{1a}(\eta_c)\equiv \lambda_{1b}(\eta_c)$\footnote
{
Tutte le figure che rappresentano inviluppi sono state ottenute numericamente
mediante l'applicazione diretta, nel codice, delle relative formule.  Ad esempio,
la curva $\sigma_a$ di figura \ref{fig:profili_coniugati_invil_lagrange} \`e stata
ottenuta dalla implementazione numerica della relazione
\ref{eq:inviluppo_polarb2}. L'aspetto diverso di tale figura
rispetto alla precedente
\ref{fig:profili_coniugati_invil_polar}, dove sono rappresentate tutte le curve
inviluppanti, dipende dal seguente stratagemma.
Oltre alle curve che creano l'inviluppo, disegnate col colore opportuno,
viene tracciato un secondo insieme di curve, con lo stesso procedimento, in
scala leggermente minore e di
colore bianco, cio\`e quello del foglio. In tal modo si ottiene, visivamente,
il solo inviluppo. Questo artificio potrebbe essere percepito da un occhio esperto.}.
\begin{figure}[hbt]
\begin{center}
\includegraphics[width=0.8\textwidth]{part2/ruote/FIG/ruote/profili_coniugati_invil_lagrange.pdf}
\end{center}
\begin{picture}(0,0)(0,0)
\scriptsize{
\put(172,145){$\eta_0$}
\put(166,143){\rotatebox{-120}{$\longrightarrow$}}
\put(203.5,107.4){$\eta_c$}
\put(198,114){\rotatebox{120}{$\longrightarrow$}}
\put(317,242){$\sigma_a$}
\put(303,264){$\sigma_b$}
\put(142,227){$\sigma_2$}
\put(97,36){$\lambda_{1a}$}
\put(63,124){$\lambda_{1b}$}
\put(129,160){$\lambda_2$}
}
\end{picture}
\vskip -6mm
      \caption{\em
I profili $\sigma_a$ e $\sigma_b$ risultano a loro volta coniugati quando la
polare $\lambda_{1a}$ rotola sulla $\lambda_{1b}$ toccandola in $\mu_c=\eta_c$.
}
 \label{fig:profili_coniugati_invil_lagrange}
\end{figure}
\noindent Nella stessa figura sono riportate, con colori meno marcati,
anche la $\sigma_2(\tau_i)$ e la
polare $\lambda_2$ che rimangono al  proprio posto.
\noindent Possiamo quindi affermare che, dati un insieme di polari
$\lambda_i$ e un profilo $\sigma_0$
collegato ad una di queste, diciamo $\lambda_0$,
i profili coniugati $\sigma_{i}$,
che otterremo dagli inviluppi generati da $\sigma_0$ quando $\lambda_0$
rotola sull'insieme delle
$\lambda_i$\footnote{Anche se quanto affermiamo non presenta
ricadute pratiche, possiamo lasciare,
nell'insieme di polari sulle quali rotola $\lambda_0$, anche essa stessa.
 Cos\`i operando $\sigma_0$ invilupper\`a s\'e stessa.}, 
risultano tutti coniugati tra loro a coppie rispetto al moto
generato dal rotolamento delle due polari associate a ciascuna di tali
coppie.  Questo importantissimo risultato costituisce lo scopo del
presente paragrafo.
Da esso discende la possibili\`a di creare {\em profili coniugati di 
assortimento}\index{profili!di assortimento}, questione di estrema
importanza nella fabbricazione di famiglie di ruote dentate che,
a primitive date, possano ingranare l'una con l'altra.


\noindent Prendiamo di nuovo in considerazione
la trasformazione \ref{eq:inviluppo} e la figura
\ref{fig:profili_coniugati_definizione_steps}.
Se omettessimo le ultime due delle quattro trasformazioni presenti
(si parte sempre da destra), cio\`e se ci fermassimo alla parabola c),
rappresentata in azzurro chiaro,
$\sigma''_2(\tau)$,
potremmo ridisegnare la stessa figura applicando
alla $\sigma_1$, con segno opposto, le
trasformazioni non effettuate sulla $\sigma_2$. Ovviamente 
la nuova figura risulterebbe traslata verso sinistra della 
quantit\`a ${\bm T}[-\sigma_1(t)]$ e la $\sigma_1$ sarebbe
ruotata di ${\bm R}[-\alpha_1(t)]$, ma la coniugazione
tra i due profili risulterebbe perfettamente equivalente.
Pi\`u in generale, riconsiderando la famiglia di polari $\lambda_i$
e i relativi profili coniugati $\sigma_i$ possiamo scrivere,
per un dato parametro $\eta_c$,
\begin{equation}
{\sigma_i''(\tau)}_{\eta_c} = 
 {\bm R}[-\beta_i(\eta_c)]\cdot
 {\bm T} [-\lambda_i(\eta_c)]\cdot \sigma_i(\tau)\,,
\label{eq:inviluppo_polar_euler}
\end{equation}

\begin{figure}[hbt]
\begin{center}
\includegraphics[width=0.8\textwidth]{part2/ruote/FIG/ruote/profili_coniugati_invil_euler.pdf}
\end{center}
\begin{picture}(0,0)(0,0)
\scriptsize{
\put(180,153){$\eta_0$}
\put(173,151){\rotatebox{-120}{$\longrightarrow$}}
\put(207,107.4){$\eta_c$}
\put(202,114){\rotatebox{120}{$\longrightarrow$}}
\put(318,244){$\sigma_a$}
\put(305,264){$\sigma_b$}
\put(222,223){$\sigma_2$}
\put(100,35){$\lambda_{1a}$}
\put(67,122){$\lambda_{1b}$}
\put(151,170){$\lambda_2$}
}
\end{picture}
\vskip -6mm
      \caption{\em
Le primitive del moto  rotolano l'una sull'altra in modo tale da mantenere
 fisso, rispetto all'osservatore assoluto, il loro punto di contatto.
}
 \label{fig:profili_coniugati_invil_euler}
\end{figure}

\noindent dove si riconosce una generalizzazione della
\ref{eq:inviluppo_polar} priva per\`o
delle ultime due trasformazioni a sinistra.
Il vantaggio offerto da questo punto di vista \`e quello di trattare
tutte le polari e i relativi profili coniugati mediante trasformazioni 
aventi la medesima struttura, le quali contengono
soltanto informazioni circa quella stessa polare o profilo coniugato.
Agendo in questo modo, il punto di contatto tra tutte le polari rimane fisso
rispetto a un osservatore assoluto e perde di senso parlare di
polari fisse e mobili. 
In figura \ref{fig:profili_coniugati_invil_euler} sono riportati
i profili coniugati e le rispettive polari, coincidenti
nei punti di parametro $\eta_c$ tramite il punto di vista
test\'e descritto e rappresentato dalla relazione \ref{eq:inviluppo_polar_euler}.
Il risultato vede ancora i profili tra loro coniugati (tutti e tre),
mentre il punto di contatto tra tutte le polari, come abbiamo
anticipato, rimane inalterato
rispetto a un osservatore assoluto. Le tre polari si muovono
in modo tale da evocare il funzionamento di un laminatoio.
In questa rappresentazione del moto relativo le polari cambiano nome
e diventano le {\em primitive del moto}\index{primitiva!del moto}.
\begin{figure}[hbt]
\begin{center}
\includegraphics[width=0.8\textwidth]{part2/ruote/FIG/ruote/profili_coniugati_punto.pdf}
\end{center}
\begin{picture}(0,0)(0,0)
\scriptsize{
\put(318,180){$\sigma_a$}
\put(284,250){$\sigma_b$}
\put(124,223){$P$}
\put(48,75){$\lambda_{1a}$}
\put(25,135){$\lambda_{1b}$}
\put(84,238){$\lambda_2$}
}
\end{picture}
\vskip -6mm
      \caption{\em
Profili coniugati inviluppati come traiettorie di un punto $P$ trascinato dalla
polare $\lambda_2$ che rotola sia su $\lambda_{1a}$, sia su $\lambda_{1b}$.
}
 \label{fig:profili_coniugati_punto}
\end{figure}
\noindent Nelle figure \ref{fig:profili_coniugati_invil},
 \ref{fig:profili_coniugati_invil_lagrange} e
 \ref{fig:profili_coniugati_invil_euler} abbiamo scelto  $\sigma_2$ 
che genera gli altri profili coniugati, di forma ellittica. Se tale curva,
rimpicciolendosi, degenerasse in un punto tutto l'impianto della generazione
dei profili coniugati
rimarrebbe inalterato. Le figure
\ref{fig:profili_coniugati_punto},
\ref{fig:profili_coniugati_punto_lagrange} e
\ref{fig:profili_coniugati_punto_euler}
sono quelle che in tal caso
si ottengono. Le riportiamo
prive di commento, in quanto dovremmo ripetere i medesimi
ragionamenti che abbiamo svolto poc'anzi senza alcuna aggiunta.
Possiamo quindi affermare che anche i profili generati
come traiettoria di un punto solidale a una delle polari o primitiva del moto,
durante il suo rotolamento sulle rimanenti,
sono tra loro coniugati. 
I profili coniugati di assortimento ottenuti come inviluppo
di una linea saranno utilizzati nel prossimo paragrafo per disegnare
le ruote dentate con denti profilati a evolvente di cerchio,
le quali costituiscono la quasi totalit\`a delle ruote dentate di
uso comune nelle applicazioni tecniche.
I profili provenienti invece dalla traiettoria di un punto possono
generare ruote con denti a profilo
cicloidale. Queste, unitamente ad alcune loro applicazioni,
verranno descritte brevemente, come gi\`a segnalato, nel capitolo
\ref{ruotecy}.
A conclusione di questo laborioso paragrafo possiamo affermare
che le figure \ref{fig:profili_coniugati_invil} e
\ref{fig:profili_coniugati_punto}, insieme alle loro evoluzioni cinematiche,
sono riportate su quasi
tutti i testi di Meccanica Applicata alle Macchine. Citando il solito
\cite{sesini1}, esse si trovano tra le pagine 71 e 74. Ma tali figure
risultano reperibili in un gran numero di testi
dove campeggiano con aspetti spesso molto simili tra loro. Anche le dimostrazioni
associate, le quali assicurano la possibilit\`a di
 ottenere profili coniugati di 
assortimento, si somigliano tutte e probabilmente sono maggiormente
 valide ed eleganti, confrontate con quella che 
abbiamo proposto. Di sicuro sono molto pi\`u sintetiche.

\begin{figure}[hbt]
\centering
\begin{minipage}[b]{0.45\textwidth}
\centering
\includegraphics[width=0.9\textwidth]{part2/ruote/FIG/ruote/profili_coniugati_punto_lagrange.pdf}
\begin{picture}(0,0)(170,12)
\scriptsize{
\put(88,87){$\eta_0$}
\put(82,85){\rotatebox{-120}{$\longrightarrow$}}
\put(106,54){$\eta_c$}
\put(101,59){\rotatebox{120}{$\longrightarrow$}}
\put(150,114){$\sigma_a$}
\put(145,124){$\sigma_b$}
\put(60,113){$P_0$}
\put(94,120){$P_c$}
\put(25,35){$\lambda_{1a}$}
\put(9,67){$\lambda_{1b}$}
\put(43,118){$\lambda_2$}
}
\end{picture}
      \caption{\em
I profili $\sigma_a$ e $\sigma_b$ risultano a loro volta coniugati quando la
polare $\lambda_{1a}$ rotola sulla $\lambda_{1b}$ toccandola in un
punto con parametro $\mu=\eta=c$.
}
 \label{fig:profili_coniugati_punto_lagrange}
\end{minipage}\hfill
\begin{minipage}[b]{0.45\textwidth}
\includegraphics[width=0.9\textwidth]{part2/ruote/FIG/ruote/profili_coniugati_punto_euler.pdf}
\begin{picture}(0,0)(180,10)
\scriptsize{
\put(88,92){$\eta_0$}
\put(82,90){\rotatebox{-120}{$\longrightarrow$}}
\put(105,54){$\eta_c$}
\put(100,59){\rotatebox{120}{$\longrightarrow$}}
\put(150,115){$\sigma_a$}
\put(145,126){$\sigma_b$}
\put(94,121){$P_c$}
\put(20,33){$\lambda_{1a}$}
\put(14,67){$\lambda_{1b}$}
\put(84,140){$\lambda_2$}
}
\end{picture}
      \caption{\em
Le primitive del moto  rotolano l'una sull'altra in modo tale da mantenere
 fisso, rispetto all'osservatore assoluto, il loro punto di contatto.
      }
 \label{fig:profili_coniugati_punto_euler}
\end{minipage}
\end{figure}



\section{Profili Coniugati a Evolvente}

\noindent Torniamo ora volentieri
al percorso intrapreso, dal quale qualche perplessit\`a di troppo ci
aveva costretto a una deviazione,
per andare dritti allo studio della cinematica
delle ruote dentate.
\begin{figure}[b]
\centering
\includegraphics[width=0.8\textwidth]{part2/ruote/FIG/ruote/evol1.pdf}
\begin{picture}(0,0)(150,0)
\scriptsize{
\put(42,217){$\sigma_2$}
\put(90,198){$T'$}
\put(117,195){$\sigma'_2$}
\put(36,192){$P'$}
\put(88,160){$\mu'$}
\put(23,156){$\alpha=\alpha'$}
\put(-41,148){$P\equiv T\equiv M$}
\put(148,143){$\lambda_2$}
\put(54,119){$H$}
\put(-32,110){$h$}
\put(-12,100){$\theta$}
\put(-30,90){$F$}
\put(-28,183){$F'$}
\put(104,94){$M'$}
\put(-30,63){$r_b$}
\put(147,59){$\mu$}
\put(-55,19){$r$}
\put(146,26){$\lambda'_2$}
\put(110,29){$H'$}
\put(4,13){$O$}
\put(120,-7){$\lambda_1$}
}
\end{picture}
\vskip 3mm
      \caption{\em
Impianto geometrico per la costruzione dell'evolvente (i).
}
 \label{fig:evol1}
\end{figure}
Abbiamo lasciato il sentiero quando i dubbi circa
la bont\`a delle soluzioni che ci proponeva via via
 l'intuito hanno avvolto un punto
chiave: le ruote delle due famiglie, a) e b), rappresentate nelle figure
\ref{fig:squalo} e \ref{fig:squalo_interno}, di colore azzurro, ingraneranno
tra loro correttamente? Le notizie contenute nel precedente paragrafo ci
assicurano che i profili dei denti delle ruote a),
 ottenuti dall'inviluppo di quelli della dentatrice di figura \ref{fig:squalo}
e quelli dei denti delle b), creati invece dalla dentatrice a curvatura
opposta di figura \ref{fig:squalo_interno}, sono tutti tra loro coniugati.
Del resto, essi sono ottenuti come inviluppo di un segmento di retta
(un fianco del dente della dentatrice) che viene
trascinato dalla stessa primitiva mentre rotola sia su una 
primitiva qualsiasi delle ruote a) sia su una appartenente alle  b), inviluppando cos\`i i profili dei loro denti.
Tutte le ruote a) saranno pertanto
ingranabili con le b), con l'ovvio limite, del quale abbiamo gi\`a
fatto cenno in precedenza, che
la dimensione di queste ultime non pu\`o superare quella della ruota 
dentatrice. Tale limite ci spinge a considerare, per il diametro
della ruota di colore rosa di figura \ref{fig:squalo}, valori molto grandi:
e se pensassimo a una dentatrice di diametro infinito? In questo modo, oltre
il problema di dentare le ruote b) di grandi dimensioni, si risolverebbe
felicemente, come vedremo a breve,
 anche la noiosa questione delle due famiglie. Ricordiamo che,
a suo tempo, tra i vari precetti da rispettare nella costruzione delle ruote
dentate era emerso, in forma debole, quello di avere una sola famiglia
di ruote di assortimento.
Ci si rende facilmente conto che,
trasformando la dentatrice in una 
{\em cremagliera}\index{cremagliera}, con denti a fianco rettilineo,
tale precetto viene automaticamente
rispettato. Infatti, una circonferenza di raggio infinito presenta
curvatura pari a zero e, stante il profilo rettilineo dei suoi denti,
la dentatrice con curvatura opposta risulterebbe essere
sempre la cremagliera di partenza. Tale caratteristica \`e conosciuta
come {\em auto-complementarit\`a}\index{auto-complementarit\`a}
(dell'utensile).
Ma andiamo per gradi e consideriamo la figura \ref{fig:evol1}.
Sia $\lambda_1$ la primitiva circolare di raggio $r$ della ruota che
desideriamo costruire.
Su di essa facciamo rotolare, rigorosamente in assenza di strisciamento,
la primitiva $\lambda_2$, che \`e la retta disegnata in verde,
alla quale \`e solidale la
``curva'' $\sigma_2$, il segmento marrone inclinato di $\theta$
rispetto ad una retta ortogonale a $\lambda_2$\footnote{
La primitiva $\lambda_2$, indicata qui seguendo la tradizione come una
retta, sar\`a, per ovvie motivazioni tecniche, un segmento di lunghezza opportuna,
sufficiente ad ottenere il tratto di inviluppo che former\`a il fianco
del dente. Il segmento $\sigma_2$ viene tradizionalmente trattato come una
semiretta la quale partendo dall'estremo $F$
si estende all'infinito, permettendo in questo modo di disegnare l'intera
evolvente. Anche in questo caso ci accontentiamo di un segmento, avendo gi\`a in mente
che esso costituir\`a il fianco dei denti dell'utensile creatore.}:
questo segmento, trascinato dalla primitiva $\lambda_2$ mentre
rotola sulla circonferenza $\lambda_1$, 
creer\`a l'inviluppo che stiamo cercando.
La figura stabilisce che, per un valore dell'angolo
 $\alpha =0$, la retta  $\lambda_2$
si trovi in posizione orizzontale e sia tangente alla
 circonferenza primitiva $\lambda_1$ nel
punto $P$.
 Il processo di inviluppo, dal quale si otterr\`a il profilo coniugato
a $\sigma_2$, consta nel fare appunto rotolare la $\lambda_2$ sulla $\lambda_1$.
A una rotazione generica $\alpha=\alpha'$ corrisponderanno
la retta $\lambda'_2$ e il segmento
$\sigma'_2$.
\`E interessante individuare, su tale segmento, il punto
attivo nella creazione del profilo coniugato, cio\`e il punto $T'$
in cui il segmento
risulter\`a tangente all'inviluppo che stiamo creando.
$T'$ sar\`a perci\`o determinato
dalla intersezione della retta che, passando dal centro istantaneo di rotazione
$M'$, taglia ortogonalmente il segmento $\sigma'_2$, imponendo cos\`i che $T'$
possieda velocit\`a
parallela al segmento $\sigma'_2$ stesso.
L'individuazione di tale punto porta con s\'e
conseguenze profonde. I due triangoli $ \triangle{P' T' M'}$ e
 $\triangle{M'H'O}$
sono simili, pertanto il secondo, al di l\`a della
sua posizione, rimane lo stesso per ogni valore di $\alpha$.
Ci\`o significa che la lunghezza del segmento $(H'M')$ (indicheremo d'ora in poi
la lunghezza dei segmenti tramite le lettere che li individuano poste tra
parentesi tonde) risulter\`a
costante al variare di $\alpha$ e precisamente avremo $(H'M')= r\sin(\theta)$.
La retta $\mu'$ risulta perci\`o tangente alla
circonferenza di raggio $r_b= r \cos(\theta)$ per qualsiasi valore di $\alpha$,
dove $r_b$, grandezza di importanza rilevantissima, \`e il raggio della
{\em circonferenza di base}\index{circonferenza!di base}.
Inoltre, la lunghezza $(H'T')$ sar\`a data dalla quantit\`a 
$(M'T')$, variabile con $\alpha$, sommata alla quntit\`a fissa $(H'M')$.
Per  $\alpha=0$ abbiamo $P\equiv T\equiv M$
e naturalmente  $(H'M')=(HM)$. 
Si osserva che risulta
sempre  $(M'T')=(M'P')\cos(\theta)$ quindi, dato il rotolamento di $\lambda_2$ 
su $\lambda_1$, $(M'P')=r\alpha$ e $(M'T')=r\alpha\cos(\theta)=r_b \alpha$.
Questo indica altres\`i
che la lunghezza dell'arco $H'H$,  che sottende un angolo pari ad $\alpha$,
equivale a
$(M'T')$ e che il segmento $H'T'$ della retta $\mu'$ si comporta, al variare di
$\alpha$, come un filo che si svolge dalla circonferenza di base.

\begin{figure}[hbt] 
\centering
\includegraphics[width=0.8\textwidth]{part2/ruote/FIG/ruote/evol2.pdf}
\begin{picture}(0,0)(150,0)
\scriptsize{
\put(-11,218){$\sigma_2^o$}
\put(74,218){$\lambda_2^o$}
\put(42,217){$\sigma_2$}
%\put(-86,212){$P''$}
\put(90,198){$T'$}
\put(118,195){$\sigma'_2$}
\put(27,179){$P'$}
%\put(-58,181){$\sigma''_2$}
\put(88,160){$\mu'$}
%\put(-96,165){$\lambda''_2$}
\put(-2,160){$P^o$}
%\put(-30,160){$T''$}
\put(31,165){$\alpha$}
\put(6,147){$P\equiv T\equiv M$}
\put(148,143){$\lambda_2$}
\put(148,134){$\mu^o$}
\put(-76,127){$M^o$}
\put(-3,118){$T^o\equiv H^o$}
\put(56,116){$H$}
\put(-184,107){$\alpha^o=-\tan(\theta)$}
%\put(-2.2,105.5){$\delta$}
\put(104,95){$M'$}
\put(13,92){$\gamma$}
%\put(-79,83){$H''$}
\put(-30,63){$r_b$}
\put(147,59){$\mu$}
\put(34,50){$\beta$}
%\put(-133,42){$M''$}
\put(110,29){$H'$}
\put(146,26){$\lambda'_2$}
\put(-55,19){$r$}
\put(120,-7){$\lambda_1$}
%\put(-144,10){$\mu''$}
\put(4,13){$O$}
\put(-30,90){$F$}
\put(-28,183){$F'$}
\put(-13,102){$F^o$}
%
}
\end{picture}
      \caption{\em
Impianto geometrico per la costruzione dell'evolvente (ii).
      }
 \label{fig:evol2}
\end{figure}

\noindent In questo intreccio di angoli, rette e segmenti riteniamo opportuno,
riferendoci ora alla figura \ref{fig:evol2},
individuare il valore dell'angolo $\beta$ dato un generico valore di $\alpha$.
Da semplici considerazioni trigonometriche otteniamo subito
\begin{equation}
\tan(\beta) = \frac{(H'T')}{(OH')}=\
\frac{(H'M')+(M'T')}{(OH')}=\
\frac{r\sin(\theta)+\alpha  r\cos(\theta)}{r\cos(\theta)}=\
\tan(\theta) + \alpha\,.
\label{eq:beta}
\end{equation}
\vskip 1mm
\noindent Il capo del filo che si svolge dalla circonferenza di base
si individua ponendo $(H^oT^o)=0$, ottenendo in questo modo anche
il valore dell'angolo $\alpha^o$,
tale da fare coincidere $H^o\equiv T^o$:
$r\sin(\theta)+\alpha  r\cos(\theta) = 0 \rightarrow \alpha^o= -\tan(\theta)$.
Consideriamo ora la lunghezza del filo svolto dalla posizione $H^o$ alla
posizione $H'$, la lunghezza dell'arco $H'H^o$ sar\`a
\begin{equation}
\overset{\large\frown}{(H'H^o)} = (H'T') \rightarrow (\beta + \gamma) r_b= r_b \tan(\beta)\,,
\label{eq:gamma0}
\end{equation}
\noindent ottenendo pertanto
\begin{equation}
 \gamma(\beta) = \tan(\beta)-\beta \,.
\label{eq:gamma}
\end{equation}
\noindent La quantit\`a ora ricavata, unitamente alla distanza
\begin{equation}
\rho(\beta) = (OT') =  \frac{r_b}{\cos(\beta)}\,,
\label{eq:rho}
\end{equation}
\noindent forma la coppia delle coordinate polari del  punto generico
 $T'$ dell'inviluppo
che risulta essere l'{\em evolvente}\index{evolvente}
 della circonferenza di base con inizio
nel punto $H^o$. Questa curva \`e riportata in vari manuali di Matematica,
e pu\`o essere definita proprio come luogo dei punti dell'estremo di un filo
che si svolge da una determinata circonferenza a partire da un punto iniziale.
L'angolo $\beta$ riveste i panni di
parametro delle
coordinate polari di questa curva ed \`e
legato al nostro angolo di rotolamento
$\alpha$ dalla relazione \ref{eq:beta}. Tale formula permetterebbe,
invertendo la tangente, di sostituire $\beta$ con $\alpha$ nelle
\ref{eq:gamma} e \ref{eq:rho}.
Ma questa operazione
non viene in genere eseguita, forse perch\'e, come vedremo a breve,
in quel poco (o molto) utilizzo che si fa delle coordinate polari dell'evolvente,
l'angolo $\beta$ si comporta egregiamente.

\noindent Riepiloghiamo. La costruzione di profili di assortimento per i denti
ci porta a una primitiva comune a tutti gli inviluppi
che \`e una retta, $\lambda_2$, che rotoler\`a sulle primitive delle ruote da
tagliare.  Ulteriori 
considerazioni, legate alla simmetria, ci indicano come curva inviluppante
un segmento di retta, $\sigma_2$, solidale a $\lambda_2$ e inclinato
rispetto a questa di un angolo opportuno, $\pi/2 -\theta$. L'angolo $\alpha$,
mediante il quale la $\lambda_2$ rotola sulla circonferenza di raggio
$r$, governa il processo di inviluppo. Tramite i ragionamenti sopra esposti
perveniamo alla definizione di tale inviluppo che risulta essere
un'evolvente della circonferenza di base, di raggio $r_b=r \cos(\theta)$.
Pur riconoscendo l'elevato valore culturale 
di questa scoperta, del fatto cio\`e che il profilo dei denti, ottenuti tramite
il taglio con una dentiera rettilinea a fianchi pure rettilinei, sia un'evolvente
della circonferenza
di base, non dedicheremo nel presente lavoro molto spazio a questa affascinante
curva geometrica. Come riportato in \cite{guiggiani},
dal quale confessiamo di avere preso diversi spunti: ``Il fatto che il
profilo attivo sia una evolvente di cerchio \`e, in fondo, una fortunata
coincidenza''.
Ci troviamo d'accordo con questa affermazione, riservandoci
 di valutare meglio l'epiteto. Torneremo a breve
su alcune (poche) applicazioni delle formule relative all'evolvente,
nel paragrafo che tratta la correzione delle ruote mediante lo spostamento
dei profili, in particolare quando opereremo con variazione di interasse.
\noindent Prima di passare agli inviluppi veri e propri, che definiranno
la forma dei denti, \`e bene fare ancora qualche considerazione sulla
figura \ref{fig:evol2}. Quando la retta $\lambda_2$ si trova nella
posizione $\lambda^o_2$ il punto $T^o\equiv H^o$ cade sulla circonferenza
di base. Tale posizione rappresenta, sul segmento $\sigma_2$,
l'estremo oltre il quale quest'ultimo risulta inefficace nel creare
l'inviluppo. Pertanto, una posizione dell'estremo $F$ esterna a tale punto,
come quella riportata nelle  nostre figure, \`e  disutile
 in quanto la frazione in eccesso, non partecipando all'inviluppo, avrebbe come
unico effetto quello di distruggere una porzione 
del materiale da tagliare: una situazione analoga a quella
di figura \ref{fig:profili_coniugati_definizione_c}, innocua nella teoria,
inaccettabile nella pratica. La distanza $(T^oP^o)$ ha perci\`o
un significato tecnico importante e vale
\begin{equation}
(T^oP^o)=\frac{r }{\cos(\theta)}\sin^2(\theta)\,,
\label{eq:maxt}
\end{equation}
come si deduce facilmente considerando i due triangoli simili
$\triangle{OP^oM^o}$ e $\triangle{M^oP^oT^o}$.
In figura \ref{fig:evol1} \`e stata indicata con $h$ la lunghezza della
proiezione di $FP$ su una direzione ortogonale all'epiciclo $\lambda_2$. Tale
lunghezza rappresenta, come vedremo a breve, una porzione
dell'altezza dei denti dell'utensile che stiamo cercando. 

\begin{figure}[hbt] 
\centering
\includegraphics[width=0.7\textwidth]{part2/ruote/FIG/ruote/inviluppo_evolvente.pdf}
\begin{picture}(0,0)(150,0)
\scriptsize{
\put(14,234){$\sigma_2^o$}
\put(17,75){$T^o\equiv H^o$}
\put(-66,19){$r$}
\put(84,19){$r_b$}
\put(150,81){$\lambda_1$}
%
}
\end{picture}
      \caption{\em
Inviluppo dell'evolvente e relativo arco di trocoide. Sporgenza di $\sigma_2$
uguale ad $h_{\rm max}$, $\theta=20^{\circ}$.
      }
 \label{fig:inviluppo_evolvente}
\end{figure}
\noindent Risulta perci\`o ancora pi\`u rilevante
specificare la lunghezza della proiezione di $(T^oP^o)$ sulla direzione
ortogonale a $\lambda^o_2$, che fornir\`a
la massima altezza utile, a partire da $\lambda_2$,
dei denti della nostra dentiera generatrice:
\begin{equation}
h_{\rm max}=r  \sin^2(\theta)\,.
\label{eq:maxh}
\end{equation}
\noindent La figura \ref{fig:inviluppo_evolvente} mostra un fianco di dente,
inviluppato dal segmento $\sigma_2$, proprio in questa condizione limite,
quando cio\`e $(FP)=(T^oP^o)$.
Si pu\`o notare che, affinch\'e l'evolvente sia completa, partendo dalla
circonferenza di base, il fianco della dentiera scava ulteriormente
al di sotto di questa, costruendo un raccordo alla base del fianco del dente 
stesso. La medesima figura ci mostra che, al
variare di $\alpha$, l'estremo del segmento generatore non sconfina mai
al di l\`a del prolungamento ideale di $\sigma_2^o$ verso il centro della ruota,
essendo $\sigma_2^o$ l'ultima posizione utile nell'inviluppo dell'evolvente
stessa. Le ulteriori posizioni di $\sigma_2$, riportate in colore marrone,
rappresentano l'uscita di scena del profilo creatore. Sull'estremo inferiore
di quest'ultimo abbiamo evidenziato un minuscolo cerchio, rispettando la coerenza
del suo colore e quello del segmento che lo porta. Tali cerchi definiscono graficamente
l'arco di {\em trocoide}\index{trocoide}\footnote{
Si chiamano trocoidi le curve disegnate da punti esterni all'epiciclo
e solidali con esso durante il
rotolamento di quest'ultimo sulla circonferenza che percorre. Nel nostro caso
l'epiciclo (o ipociclo: nella teoria \`e possibile far rotolare circonferenze
pi\`u grandi all'interno di quelle pi\`u piccole) 
\`e la retta $\lambda_2$ che rotola sulla primitiva $\lambda_1$.
} che, come si vede, delinea un raccordo del profilo del dente con una
possibile circonferenza interna a quella di base: la {\em circonferenza
di piede}\index{circonferenza!di piede}.

\begin{figure}[h]
\centering
\begin{minipage}[b]{0.45\textwidth}
\centering
\includegraphics[width=0.9\textwidth]{part2/ruote/FIG/ruote/inviluppo_evolvente_sottotaglio.pdf}
\begin{picture}(0,0)(130,0)
\scriptsize{
\put(49.2,137){$\sigma_2^o$}
%\put(52,19){$T^o$}
\put(-8,19){$r$}
\put(89,19){$r_b$}
\put(130,45){$\lambda_1$}
}
\end{picture}
      \caption{\em
Inviluppo con $\theta=20^{\circ}$ e $(FP)=2(T^oP^o)$.
}
 \label{fig:inviluppo_evolvente_sottotaglio}
\end{minipage}\hfill
\begin{minipage}[b]{0.45\textwidth}
\includegraphics[width=0.9\textwidth]{part2/ruote/FIG/ruote/inviluppo_evolvente_corto.pdf}
\begin{picture}(0,0)(130,0)
\scriptsize{
\put(49.2,137){$\sigma_2^o$}
%\put(53,54){$T^o$}
\put(-8,19){$r$}
\put(89,19){$r_b$}
\put(68,19){$r_e$}
\put(130,45){$\lambda_1$}
}
\end{picture}
      \caption{\em
Inviluppo con $\theta=20^{\circ}$ e $(FP)=0.66(T^oP^o)$.
      }
 \label{fig:inviluppo_evolvente_corto}
\end{minipage}
\end{figure}

\noindent Il raggio della circonferenza di piede si deduce considerando che
la proiezione del segmento $FP$ sulla direzione ortogonale a $\lambda_2$
\`e la massima profondit\`a, in direzione radiale, alla quale arriver\`a
la trocoide e vale
\begin{equation}
r_p = r - (FP)\cos\theta = r - h\,.
\label{eq:r_p}
\end{equation}
\noindent La figura \ref{fig:inviluppo_evolvente_sottotaglio} rappresenta
l'inviluppo creato da un segmento $\sigma_2$ di lunghezza eccessiva (verso
la ruota da tagliare). L'evolvente risulta completa e tocca
la circonferenza di base, ma il fianco del dente viene scavato alla base e
una piccola parte della evolvente stessa viene distrutta: 
situazione questa che viene indicata col nome di {\em sottotaglio dei
denti}\index{sottotaglio dei denti}. 
La figura \ref{fig:inviluppo_evolvente_corto} rappresenta invece il
profilo ottenuto quando il segmento $\sigma_2$ ha una lunghezza insufficiente
a inviluppare tutta l'evolvente, alla quale mancher\`a un arco in prossimit\`a
della circonferenza di base.  La coordinata
polare radiale del punto dell'evolvente $T'$, gi\`a riportata
nella formula \ref{eq:rho}, si pu\`o anche ricavare, 
con riferimento alla figura  \ref{fig:evol2}, applicando
il teorema di Pitagora al triangolo $\triangle OT'H'$. 
Al fine quindi di individuare il raggio della circonferenza da cui partir\`a
l'evolvente, sostituiamo il
punto $T'$ con l'estremo di $F$ di $\sigma_2$ prestando attenzione
al segno opportuno per la porzione di segmento $PF$.
Stanti queste premesse, l'evolvente partir\`a da una circonferenza di raggio 
\begin{equation}
r_e = \sqrt{r_b^2 +\left(\frac{r_b\sin\theta}{\cos\theta}-\frac{(FP)\cos\theta}{\sin\theta}\right)^2}\,,
\label{eq:ini_evo}
\end{equation}
\noindent e in termini di sporgenza dei denti della dentiera (utensile) $h$
\begin{equation}
r_e = \sqrt{r_b^2 +\left(\frac{r_b\sin\theta}{\cos\theta}-\frac{h}{\sin\theta}\right)^2}\,.
\label{eq:ini_evoh}
\end{equation}

\begin{figure}[hbt]
\centering
\begin{minipage}[b]{0.45\textwidth}
\centering
\includegraphics[width=0.9\textwidth]{part2/ruote/FIG/ruote/inviluppo_evolvente_sottotaglio35.pdf}
\begin{picture}(0,0)(130,0)
\scriptsize{
\put(43.5,150){$\sigma_2^o$}
%\put(52,19){$T^o$}
\put(20,19){$r$}
\put(88,19){$r_b$}
\put(128,98){$\lambda_1$}
}
\end{picture}
      \caption{\em
Inviluppo con $\theta=35^{\circ}$ e $(FP)=2(T^oP^o)$.
}
 \label{fig:inviluppo_evolvente_sottotaglio35}
\end{minipage}\hfill
\begin{minipage}[b]{0.45\textwidth}
\includegraphics[width=0.9\textwidth]{part2/ruote/FIG/ruote/inviluppo_evolvente_corto35.pdf}
\begin{picture}(0,0)(130,0)
\scriptsize{
\put(43.5,150){$\sigma_2^o$}
%\put(53,54){$T^o$}
\put(20,19){$r$}
\put(88,19){$r_b$}
\put(73,19){$r_e$}
\put(128,98){$\lambda_1$}
}
\end{picture}
      \caption{\em
Inviluppo con $\theta=35^{\circ}$ e $(FP)=0.66(T^oP^o)$.
      }
 \label{fig:inviluppo_evolvente_corto35}
\end{minipage}
\end{figure}

\noindent Giusto per dare un'idea dell'influenza che il valore dell'angolo
 di pressione $\theta$ esercita su quello che sar\`a il profilo dei denti
 riportiamo le figure \ref{fig:inviluppo_evolvente_sottotaglio35} e 
\ref{fig:inviluppo_evolvente_corto35} per le quali, anzich\'e
assumere $\theta=20^{\circ}$, che \`e il valore {\em standard} dell'angolo di pressione,
si \`e posto $\theta=35^{\circ}$.
La differenza pi\`u rilevante tra queste ultime due figure e le loro precedenti
omologhe consta in una significativa riduzione del raggio di base. Ci\`o
permette, aumentando l'angolo di pressione, di poter considerare
altezze utili della dentiera pi\`u elevate, come del resto \`e evidente
anche dalla formula \ref{eq:maxh}.
Ci si pu\`o chiedere se due profili, generati con angoli $\theta$ diversi,
siano tra loro coniugati. Formulata in questo modo la domanda e di conseguenza
la risposta necessitano di qualche precisazione. La risposta \`e negativa
qualora intendessimo far rotolare,
l'una sull'altra, le due circonferenze primitive (assegnate)
e ottenere i due profili tramite due diversi epicicli.
\`E la teoria dei profili coniugati a escluderlo:
deve necessariamente essere lo stesso epiciclo a rotolare
su tutte le primitive per poter ottenere i cosiddetti {\em profili di assortimento}\index{profili!di assortimento}.
Qualora invece ci si interroghi circa la possibilit\`a di coniugare i due
profili a evolvente tramite circonferenze primitive diverse da quelle di
partenza la
risposta \`e affermativa. Infatti le due curve saranno le evolventi di 
due circonferenze di base di raggi $r_{b_1}$ e $r_{b_2}$. Se tali circonferenze
vengono affacciate, mantenendo i loro centri a distanza $i>r_{b_1}+r_{b_2}$,
si avranno due possibili raggi primitivi e un possibile angolo di pressione
dati da
\begin{equation}
r_1=\frac{r_{b_1}i}{r_{b_1}+r_{b_2}}\,,\;\;\;
r_2=\frac{r_{b_2}i}{r_{b_1}+r_{b_2}}\; {\rm e}\;\;\;
\theta=\arccos\frac{r_{b_1}+r_{b_2}}{i}\,.
\label{eq:nuove_primitive}
\end{equation}
\noindent Lasciamo la verifica di queste 
formule alla volont\`a di qualche studente.
Anticipiamo solamente che i valori dei due raggi primitivi
presentano scarso interesse applicativo, mentre quello dell'angolo di pressione
per un interasse dato risulta molto rilevante per gli ingranaggi con ruote
corrette, come vedremo, in succinto, in un prossimo paragrafo.
Il fatto che i profili a evolvente siano, entro limiti ragionevoli,
tutti fra loro potenzialmente coniugati mediante primitive circolari
risulta di grandissima importanza. Da ci\`o deriva infatti
la facolt\`a di variare, sempre entro i limiti menzionati,
 l'interasse di progetto
nella fase di montaggio delle ruote che portano tali profili:
vantaggio notevole, che permette di assorbire eventuali errori
di fabbricazione e messa in opera. Un risultato ancora pi\`u importante
della flessibilit\`a derivante dalla \ref{eq:nuove_primitive} \`e
la possibilit\`a di tagliare
le cosiddette ruote a profilo corretto nelle quali, sacrificando
senza alcun rimorso la geometria ``normale''  di tali ruote e spostando
opportunamente l'utensile creatore, si ovvia al grave problema del sottotaglio,
chiamato spesso {\em interferenza} che, come vedremo in un prossimo
paragrafo, pone un limite inferiore al numero di denti dei pignoni; ma
soprattutto offre al progettista l'opportunit\`a di individuare una strategia
di ottimizzazione degli ingranaggi, quasi sempre basata sull'utilizzo di
interassi diversi da quelli che imporrebbe il taglio di ruote ``normali''.


\section{Proporzionamento Modulare delle Ruote}

\noindent Si \`e gi\`a accennato altrove che, in ottemperanza al primo
precetto che abbiamo ammesso, cio\`e il vincolo  che impone un
numero intero di denti per le ruote dentate, nel progetto di queste ultime
si parte dal loro numero di denti, e non dai loro diametri o dai loro raggi.
Il numero di denti $z$ di una ruota rappresenta per\`o solamente una
``situazione'' geometrica angolare. In pratica, $z$ impone il
numero di spicchi mediante il quale si desidera
dividere la torta. La dimensione effettiva della
ruota deve pertanto essere legata a un fattore di scala dimensionale.
Tale parametro si chiama {\em modulo}\index{modulo}, $m$, ed \`e
a sua volta legato al numero di denti e al raggio della circonferenza primitiva
 dalla relazione
\begin{equation}
2 r =mz\,,
\label{eq:mod_z}
\end{equation}
\noindent gi\`a riportata in precedenza, ma di importanza tale da non
farci sentire in imbarazzo ripetendola. Il modulo $m$ viene espresso,
nei paesi in cui si usa il sistema metrico decimale, in ${\rm mm}$. Abbiamo 
gi\`a fatto cenno anche al {\em passo}\index{passo} della ruota che \`e
la lunghezza dell'arco di primitiva che separa due fianchi, destri o sinistri,
attigui:
\begin{equation}
p=\pi m\,.
\label{eq:passo}
\end{equation}
\noindent Dalla \ref{eq:passo} si deduce immediatamente che la possibilit\`a
di ingranare di due ruote dentate \`e condizionata dal loro modulo:
soltanto ruote con moduli uguali tra loro potranno formare ingranaggi.
L'unificazione impone per il modulo $m$ una certa successione di valori
normalizzati, della quale riportiamo i pi\`u usuali
$m=1,\;1.25,\;1.5,\;2,\;2.5,\;3,\;4,\;5,\;6\; {\rm mm}$.
Si riduce, in questo modo, la variet\`a di scelta di questo fondamentale
parametro a fronte
di notevoli vantaggi, evidenti nella fase di progettazione,
ma ancora di pi\`u se si pensa alla fase di taglio.
A ciascun valore del modulo corrisponde infatti un certo numero di utensili,
{\em in primis} la dentiera tagliente (o l'utensile creatore)
 e a seguire le mole di rettifica,
ove necessario.
Per passare dai profili rappresentati nelle precedenti figure
agli effettivi fianchi dei denti delle ruote da
tagliare occorre estendere l'impianto teorico della primitiva $\lambda_2$
equipaggiata di un solo segmento $\sigma_2$ aggiungendo altre parti e
apportando qualche leggera modifica a tali elementi.
\begin{figure}[hbt]
\begin{center}
\includegraphics[width=0.94\textwidth]{part2/ruote/FIG/ruote/dentiera_riferimento.pdf}
\end{center}
\begin{picture}(0,0)(0,9.6)
\scriptsize{
\put(15,110){\rotatebox{90}{$2.5\, m$}}
\put(38,120.5){\rotatebox{90}{$1.25\, m$}}
\put(38,95){\rotatebox{90}{$1.25\, m$}}
\put(167,52){$\theta$}
\put(170,86){$F$}
\put(267,43){$p=\pi m$}
\put(264,123){$\lambda_2$}
\put(264,113.5){$\lambda_t$}
\put(76,135){$\sigma_2^-$}
\put(143,135){$\sigma_2^-$}
\put(210,135){$\sigma_2^-$}
\put(277,135){$\sigma_2^-$}
\put(121,135){$\sigma_2$}
\put(188,135){$\bm \sigma_2$}
\put(255,135){$\sigma_2$}
\put(322,135){$\sigma_2$}
\put(293,62){$p/2$}
\put(332,114){$c=mx$}
}
\end{picture}
\vskip -10mm
      \caption{\em
Dentiera normalizzata di riferimento.\index{dentiera!di riferimento}
}
 \label{fig:dentiera_riferimento}
\end{figure}

\noindent In figura \ref{fig:dentiera_riferimento} si possono notare
 la primitiva $\lambda_2$ e il segmento $\sigma_2$ contornati da
 altri oggetti
dei quali definiremo, in breve, natura e scopo.
Innanzitutto, notiamo che, a parte il numero
 di denti dell'utensile\footnote{
Il numero di denti $z_c$ dell'utensile creatore ({\em Maag}) \`e normalmente contenuto
tra i limiti $1\le z_c \le 8$, riservando il valore maggiore,
in generale, al taglio di ruote con numero di denti $z>60$.
}
e l'angolo $\theta$, tutte
le altre grandezze sono facilmente rapportabili al modulo $m$.
Il profilo $\sigma_2$, che nelle figure \ref{fig:evol1} e \ref{fig:evol2}
terminava in basso col punto $F$, e che veniva volutamente lasciato di
lunghezza indefinita all'altro capo, viene ora troncato in modo speculare
(non invertito)
rispetto a $\lambda_2$, sapendo gi\`a che i denti, dovendo avere altezza
finita, non beneficerebbero del prolungamento dell'evolvente oltre quella che,
con nome significativo, si chiama {\em circonferenza di testa}\index{circonferenza!di testa}.
Tale segmento viene poi ripetuto diverse volte e precisamente
il doppio del numero di denti di cui si vuole dotare l'utensile creatore.
La met\`a di tali segmenti, segnati coi nomi $\sigma_2^-$, 
\`e orientata in modo da creare un angolo
$-\theta$ con una linea perpendicolare  $\lambda_2$. Questi fianchi
dell'utensile sono necessari
per costruire, in una sola passata, entrambi i fianchi dei denti: le ruote
dentate trasmettono il moto in tutte e due le direzioni possibili e per
tali direzioni il loro comportamento deve essere equivalente\footnote{
Ricordiamo i precetti introdotti: il secondo (debole) prescrive
proprio questa simmetria. Esso viene talvolta disatteso, sia per particolari
applicazioni che richiedono asimmetria di comportamento dell'ingranaggio
rispetto alle direzioni del flusso di potenza, sia per ottenere
vantaggi circa la riduzione di rumorosit\`a della trasmissione, come 
testimoniato da alcune specifiche e recenti ricerche.
Stando alla nostra modesta esperienza,
confessiamo per\`o di non aver mai toccato con mano
 {\em ruote asimmetriche}\index{ruote!asimmetriche}.}.
Per una {\em dentiera normalizzata}\index{dentiera!normalizzata}
valgono le seguenti relazioni, una parte delle quali \`e riportata 
anche nella figura:
l'altezza dei denti \`e legata al modulo e vale $2.5\,m$, ugualmente
ripartita tra {\em addendum}\index{addendum} e {\em dedendum},\index{dedendum}
che pertanto valgono entrambi $1.25\, m$; i denti e i vani hanno la 
medesima forma e il passo tra due fianchi consecutivi e omologhi,
misurato su una
qualunque retta orizzontale che intercetti tali fianchi, vale $p=\pi m$.
Il valore dell'angolo di pressione che, in Europa, si ritiene normalizzato
vale $\theta=20^{\circ}$.
Oltre alla primitiva $\lambda_2$, sempre in figura 
\ref{fig:dentiera_riferimento}, troviamo la retta orizzontale $\lambda_t$
che chiameremo {\em primitiva di taglio}\index{primitiva!di taglio},
e che individua l'effettiva primitiva usata durante il moto di taglio:
essa, rotolando sulla primitiva della ruota da tagliare,
crea, tramite i profili $\sigma$, i profili dei denti.
Tale retta, rappresentata a distanza $c=mx$ 
da $\lambda_2$, si trover\`a a coincidere con questa ($c=0$) nel
taglio di {\em ruote normali non corrette}\index{ruote!non corrette}:
$x$ si chiama {\em fattore di correzione adimensionale}\index{correzione!adimensionale}, del quale daremo qualche ulteriore notizia in
un prossimo paragrafo.

\section{Taglio delle Ruote Dentate per Inviluppo}

\noindent Siamo ormai dotati dell'utensile opportuno per poter
eseguire di nuovo le dentature delle figure \ref{fig:squalo} e
 \ref{fig:squalo_interno}. Chiediamo al lettore di pazientare ancora
un momento, in modo da poter ritoccare il nostro utensile e renderlo somigliante
a quelli che realmente si usano nel processo industriale del taglio delle 
ruote dentate.
\noindent Chiariamo subito che la dentiera rappresentata in figura \ref{fig:dentiera_maag}
corrisponde solo teoricamente a quelle che si usano nel processo di taglio
per inviluppo delle ruote dentate e che portano il nome {\em Maag}\index{dentiera!Maag}.
\begin{figure}[t]
\begin{center}
\includegraphics[width=0.98\textwidth]{part2/ruote/FIG/ruote/dentiera_maag.pdf}
\end{center}
\begin{picture}(0,0)(0,9.6)
\scriptsize{
\put(5,99){\rotatebox{90}{$2.5\, m$}}
\put(33,116){\rotatebox{90}{$=$}}
\put(33,93){\rotatebox{90}{$=$}}
\put(181,55){$r_r=0.4m$}
\put(132,36){$\theta=20^{\circ}$}
\put(90,69){\rotatebox{-20}{$p_b=p \cos\theta$}}
\put(287,42){$p=\pi m$}
\put(277,110){$\lambda_2$}
\put(277,102){$\lambda_t$}
\put(307,59){$p/2$}
\put(339,102){$c=mx$}
}
\end{picture}
\vskip -6mm
      \caption{\em
Utensile Maag.
}
\vskip -3mm
\label{fig:dentiera_maag}
\end{figure}
Infatti, tali utensili a forma di pettine
devono possedere, oltre alle qualit\`a geometriche che interessano la nostra
disciplina, molti altri accorgimenti, come quelli che permettano il loro
afferraggio sulla macchina dentatrice
e altri che facilitino il processo tecnologico del taglio mediante asportazione
di truciolo: i taglienti affilati e
gli angoli di spoglia, irrinunciabili in tali processi\footnote{
Nella pratica \`e molto frequente mettere diversi pettini come quelli di
figura \ref{fig:dentiera_maag} adagiati sulle generatrici di un cilindro, ottenendo
cos\`i quello che in gergo si chiama {\em creatore}\index{creatore}. Il creatore taglia i denti 
tramite un movimento rotatorio e traslatorio e il moto di rotolamento tra le due
primitive avviene con continuit\`a. Per quest'ultimo motivo, l'asse
dei denti di ciascun pettine deve essere un'elica, e l'insieme di queste
schiere di denti conferiscono all'utensile l'aspetto di una vite.
} e non facilmente rappresentabili nel piano dove giacciono i profili dei denti.
La differenza pi\`u evidente tra la dentiera di riferimento di figura
\ref{fig:dentiera_riferimento} e l'utensile Maag, \`e la forma arrotondata
delle teste e del fondo dei denti. Questo accorgimento si rende necessario
per due ordini di difficolt\`a che l'uso di una dentiera
con spigoli vivi comporterebbe.
La prima di tali difficolt\`a risiede in ambito tecnologico
e si pu\`o riassumere dicendo che i taglienti spigolosi (bench\'e ammissibili)
sono estremamente pi\`u fragili rispetto a quelli che presentano spigoli
arrotondati. Il secondo vantaggio offerto dallo smusso dei denti riguarda
la possibilit\`a di costruire pignoni recanti un minor numero di denti,
cio\`e di ovviare in parte a quello che tradizionalmente si
chiama {\em problema dell'interferenza}, alla soluzione del
 quale dedicheremo, appena pi\`u avanti, uno spazio adeguato.
\begin{figure}[hbt]
\centering
\begin{minipage}[b]{0.48\textwidth}
\centering
\includegraphics[width=0.98\textwidth]{part2/ruote/FIG/ruote/17denti_4mod_20gradi_0corr_0p01rr_inv.pdf}
\begin{picture}(0,0)(130,0)
\scriptsize{
\put(7,45){\rotatebox{25}{$\longrightarrow$}}
\put(9,38){\rotatebox{35}{$\longrightarrow$}}
\put(14,30){\rotatebox{45}{$\longrightarrow$}}
\put(27,21){\rotatebox{75}{$\longrightarrow$}}
\put(0,13){punti della trocoide}
}
\end{picture}
      \caption{\em
Profilo dei denti generati dalla dentiera di riferimento: $z=17$, $\theta=20^{\circ}$, $r_r=0$.
}
 \label{fig:17denti_4mod_20gradi_0corr_0p01rr_inv}
\end{minipage}\hfill
\begin{minipage}[b]{0.48\textwidth}
\includegraphics[width=0.98\textwidth]{part2/ruote/FIG/ruote/17denti_4mod_20gradi_0corr_0p4rr_inv.pdf}
\begin{picture}(0,0)(130,0)
\scriptsize{
}
\end{picture}
      \caption{\em
Profilo dei denti generati da dentiera {\em Maag}: $z=17$, $\theta=20^{\circ}$, $r_r=0.4m$.
      }
 \label{fig:17denti_4mod_20gradi_0corr_0p4rr_inv}
\end{minipage}
\vskip -3mm
\end{figure}
Crediamo sia istruttivo fare un confronto tra gli inviluppi
prodotti dalla dentiera di riferimento
e dall'utensile {\em Maag} sullo stesso pignone con $z=17$, che, come
vedremo, presenta il minimo numero di denti possibili, al di sotto del quale
si manifesta il sottotaglio alla base dei denti. 
Come mostrato in figura 
\ref{fig:17denti_4mod_20gradi_0corr_0p01rr_inv}, la quale riporta
l'inviluppo ottenuto da una dentatrice con spigoli vivi,
lo sconfinamento verso l'asse del dente dei
 punti del raccordo di piede (l'arco di trocoide) rende la base del 
dente stesso sottotagliata. Questo comporta una sua diminuzione della
resistenza a flessione, comportandosi il dente, dal punto di vista
strutturale, come una mensola
sollecitata da un carico collocato a distanza variabile dall'incastro.
 Del resto il sottotaglio era prevedibile in quanto,
 stando al proporzionamento modulare della dentiera di riferimento,
abbiamo per l'altezza dei suoi denti $h=1.25 m$.  Sostituendo questo valore nella
\ref{eq:maxh} otteniamo
\begin{equation}
z_{\rm min}=\frac{2.5}{\sin^2(\theta)}\,,
\label{eq:zminrif}
\end{equation}
\noindent e per $\theta= 20^{\circ}$ si  ottiene
$z_{{\rm min}_{20^{\circ}}}=22$\index{numero minimo di denti}. Gi\`a in questo luogo, desideriamo
mettere in evidenza, alla buona, cio\`e mediante il calcolo
diretto della \ref{eq:zminrif} per
$\theta=25^{\circ}$, la sensibilit\`a del risultato di
 tale formula alla variazione dell'angolo di 
pressione: abbiamo, ad esempio, $z_{{\rm min}_{25^{\circ}}}=14$. La figura
\ref{fig:17denti_4mod_20gradi_0corr_0p4rr_inv} riporta invece il
profilo dei denti di un pignone con $z=17$ e $\theta = 20^{\circ}$ ottenuto per\`o
tramite utensile {\em Maag}, a spigoli smussati, con raggio di raccordo delle creste
dei denti $r_r=0.4m$.
In quest'ultimo caso, si nota che le teste dei denti
dell'utensile sono molto vicine al limite del sottotaglio\footnote{
Quanto qui asserito ($z_{{\rm min}_{20^{\circ}}}=17$) si ricava facilmente
osservando che nella formula \ref{eq:zminrif}, nel caso di utensile con denti
 smussati, il numeratore passerebbe da $2.5$ a $2$.
},
pertanto si ammette in generale,
e nelle condizioni ora menzionate, che il numero minimo  di denti per
un pignone (senza correzione) \`e $z_{{\rm min}_{20^{\circ}}}=17$.
Nelle due figure ora citate riportiamo anche un arco della circonferenza di
testa o di troncatura della ruota da tagliare, circonferenza che delimita
l'estremit\`a esterna dei denti.
Notiamo che in questo caso
la presenza di uno smusso alla base dei denti della dentiera non
incide sulla forma della testa dei denti della ruota:
 tale raccordo \`e un mero accorgimento costruttivo,
\begin{figure}[t]
\centering
\begin{minipage}[b]{0.48\textwidth}
\centering
\includegraphics[width=0.98\textwidth]{part2/ruote/FIG/ruote/17denti_4mod_20gradi_0corr_0p4rr.pdf}
\vskip 5mm
\begin{picture}(0,0)(130,0)
\scriptsize{
}
\end{picture}
      \caption{\em
Pignone con $z=17$, $\theta=20^{\circ}$: assenza di sottotaglio.
}
 \label{fig:17denti}
\end{minipage}\hfill
\begin{minipage}[b]{0.48\textwidth}
\includegraphics[width=0.98\textwidth]{part2/ruote/FIG/ruote/60denti_4mod_20gradi_0corr_0p4rr.pdf}
\begin{picture}(0,0)(130,0)
\scriptsize{
}
\end{picture}
      \caption{\em
Ruota dentata con $z=60$, $\theta=20^{\circ}$: fianchi dei denti somiglianti a quelli di una cremagliera.
      }
 \label{fig:60denti}
\end{minipage}
\vskip -3mm
\end{figure}
\noindent In figura \ref{fig:17denti} riportiamo finalmente ci\`o che si ottiene
tagliando una ruota con la dentiera {\em Maag} quando le due primitive, quella
della dentiera e quella della ruota, coincidono. Il pignone raffigurato ha
numero di denti $z=17$, angolo di pressione $\theta=20^{\circ}$
e, come si pu\`o notare visivamente,
 ci troviamo al limite del sottotaglio, anche se il numero di denti \`e
inferiore a 22 che, ripetiamo, rappresenta il minimo teorico per un
utensile con spigoli vivi.
La figura \ref{fig:60denti} mostra invece una ruota con $z=60$ e angolo
di pressione sempre $\theta=20^{\circ}$. In questo caso si nota che,
trovandoci molto lontani dal numero minimo di denti, la circostanza
del sottotaglio \`e
ben lungi dal presentarsi e i denti appaiono molto robusti alla base.
Si pu\`o osservare inoltre che i fianchi
dei denti, pur essendo (ovviamente) degli archi di evolvente,
assomigliano molto a quelli rettilinei della cremagliera.
Crediamo sia istruttivo rappresentare l'ingranamento di queste
due ruote, che naturalmente
le riconduce nella stessa scala (figura \ref{fig:1760}).

\noindent Siamo finalmente giunti alla fine del nostro percorso che, rispettando
i vincoli ``naturali'' che a mano a mano si sono palesati, ha portato a una
soluzione che appare quasi unica. Rimane ancora una verifica da effettuare,
che sottintende un precetto forte (l'ultimo): l'ingranaggio di due ruote
dentate deve essere in grado di trasmettere il moto con continuit\`a. Questo
significa che le coppie di denti a contatto devono sempre essere in numero
 maggiore
dell'unit\`a.
\begin{figure}[hbt]
\begin{center}
\includegraphics[width=1.0\textwidth]{part2/ruote/FIG/ruote/1760.pdf}
\begin{picture}(0,0)(200,0)
\scriptsize{
\put(180,197){$A$}
\color{white}
\put(280,160){$z_1=17$}
\put(50,160){$z_2=60$}
\put(331,139){$O_1$}
\put(151,145){$P$}
\put(134,97){$B$}
\put(122,83){$p_b$}
\put(126,73){$\mu$}
\put(100,46){$r_{b_2}$}
\color{black}
\put(146,46){$r_{p_2}$}
\put(198,31){$r_{p_1}$}
\put(216,31){$r_{b_1}$}
\put(158,23){$r_{t_2}$}
\put(181,23){$r_{t_1}$}
}
\end{picture}
\end{center}
\vskip -5mm
\caption{ \em
Ingranaggio con $z_1=17$, $z_2=60$, $\theta=20^{\circ}$.
} 
\vskip -3mm
\label{fig:1760}
\end{figure}
\noindent Con riferimento alla figura \ref{fig:1760}, ricordando che 
i punti di contatto (coniugati) tra due fianchi di evolvente giacciono
sulla perpendicolare alle evolventi stesse in tali punti, possiamo affermare
che i contatti tra i denti avvengono solo sulla retta $\mu$, tangente
alle due circonferenze di base di raggi $r_{b_1}$ e $r_{b_2}$, e chiamata
{\em retta delle pressioni}\index{retta delle pressioni}.
Pi\`u in dettaglio, le posizioni dei contatti saranno interne alle due 
circonferenze di testa delle ruote, le quali delimitano il tratto $AB$ che,
nella figura \ref{fig:1760}, \`e rappresentato in rosso. Affinch\'e
la trasmissione del moto avvenga con continuit\`a, sul segmento $AB$ devono
essere presenti uno o pi\`u punti di contatto. Sulla retta delle pressioni
i profili dei denti si succedono col passo che si misura sulle circonferenze
di base, dalle quali ``nascono'' le evolventi, destre o sinistre,
 e che naturalmente coincide con il passo della  successione dei fianchi, 
destri o sinistri, dei denti dell'utensile, misurato in direzione normale ai
fianchi stessi. Tale misura, indicata in figura \ref{fig:dentiera_maag}, vale
$p_b=\pi m \cos \theta$. Anche questa lunghezza \`e riportata in 
figura \ref{fig:1760}, giusto per
un confronto visivo, mediante un segmento verde.
Riepilogando, la trasmissione del moto sar\`a continua se il 
{\em fattore di ricoprimento}\index{fattore di ricoprimento} $f_c$, dato
dal rapporto tra la lunghezza del segmento $AB$ e il passo di base, \`e
maggiore dell'unit\`a: $f_c=(AB)/p_b>1$.
 La lunghezza del segmento $PB$ (come del resto quella del segmento
$PA$) si determina applicando il teorema del coseno al triangolo $\triangle{O_1BP}$ (oppure $\triangle{O_2AP}$). Per $(PB)$ si ha
\begin{equation}
(PB)^2  +2r_{p_1}\sin( \theta)(PB)+r_{p_1}^2 - r_{t_1}^2=0\,.
\label{eq:segpb2}
\end{equation}
\noindent Pertanto, indicando le lunghezze $(PA)$ e $(PB)$ semplicemente
con $\eta$ e togliendo i pedici che identificano una delle due ruote, possiamo 
scrivere
\begin{equation}
\eta= -r_p\sin( \theta)\pm \sqrt{r_p^2\sin^2(\theta)-r_p^2 + r_t^2}\,.
\label{eq:segpb3}
\end{equation}
\noindent Introducendo le grandezze legate al proporzionamento normale 
e facendo il rapporto tra $\eta$ e il passo base $p_b$ abbiamo
\begin{equation}
\frac{\eta}{p_b}= \frac{-(z/2)\sin( \theta)\pm \sqrt{(z/2)^2\sin^2(\theta)+ z+1}}{\pi \cos (\theta)}\,.
\label{eq:segpb3_1}
\end{equation}
\noindent La formula \ref{eq:segpb3_1} fornisce,
con  $z=17$, un fattore di ricoprimento,
nel tratto di competenza $PB$, $\underline f_c=\eta/p_b=0.75$. La qual cosa  significa
che, nel caso peggiore (ammissibile per non avere sottotaglio),
cio\`e quando l'ingranaggio \`e costituito da due pignoni da diciassette denti,
essi possono ingranare tra loro con continuit\`a essendo, in questo caso,
$f_c=1.5$.

\section{Ruote Corrette mediante Spostamento del Profilo I} \label{ruote-corr1}

Le proporzioni della dentiera di riferimento di figura
\ref{fig:dentiera_riferimento} risultano sicuramente ragionevoli ma, 
rappresentando un compromesso esteso a tutte le ruote, qualsiasi
sia il loro numero di denti e quello della ruota accoppiata,
si intuisce facilmente che esse potrebbero non rappresentare
la scelta ottimale.
\begin{figure}[b]
\begin{center}
\includegraphics[width=1.0\textwidth]{part2/ruote/FIG/ruote/1160.pdf}
\begin{picture}(0,0)(200,0)
\scriptsize{
\put(180,197){$A$}
\color{white}
\put(280,155){$z_1=11$}
\put(50,155){$z_2=60$}
\put(285,138){$O_1$}
\put(151,145){$P$}
\put(123,82){$p_b$}
\put(124,71){$\mu$}
\put(99,46){$r_{b_2}$}
\color{black}
\put(135.5,103.5){$B$}
\put(146,46){$r_{p_2}$}
\put(213,44){$r_{p_1}$}
\put(224,50){$r_{b_1}$}
\put(158,23){$r_{t_2}$}
\put(190,23){$r_{t_1}$}
}
\end{picture}
\end{center}
\vskip -5mm
\caption{\em
Ingranaggio con $z_1=11$, $z_2=60$, $\theta=20^{\circ}$ senza correzioni.
} 
\vskip -3mm
\label{fig:1160}
\end{figure}
\noindent In figura \ref{fig:1160} \`e riportato l'ingranaggio di un pignone
con $z_1=11$ denti e $z_2=60$ denti. Si nota subito che i denti del pignone
si presentano scavati alla base. Ci\`o era ampiamente prevedibile in quanto
la \ref{eq:zminrif}, pur considerando l'arrotondamento dei taglienti
dell'utensile creatore, non permette di scendere sotto il numero 
$z_{\rm min}=17$, senza che il sottotaglio si manifesti. La scarnitura del
dente alla base \`e una circostanza molto grave? Certamente la resistenza
a flessione del dente stesso ne risente. E dal punto di vista della nostra
disciplina? Apparentemente, in quest'ambito, le cose vanno meglio: la cinematica
dell'ingranaggio \`e comunque corretta. Ma l'impegno dell'evolvente di cerchio
fino alla sua radice fa sorgere un'altra preoccupazione, dovuta al fatto che
in quel punto, com'\`e noto, il raggio di curvatura del profilo \`e pari a zero,
e da ci\`o possono emergere altre problematiche di carattere strutturale.
Per\`o, la grande produttivit\`a delle dentatrici di ruote a evolvente,
la relativa semplicit\`a di costruzione degli utensili necessari,
la possibilit\`a di affilare tali utensili senza compromettere la loro
forma ``operativa'' ci convincono, in un ragionamento un po' fantastico,
che il mondo della meccanica avrebbe accettato tranquillamente anche le
ruote sottotagliate qualora, a tale problema, non ci fosse stato rimedio
 e anzi, in talune applicazioni, cui faremo cenno alla
fine del prossimo paragrafo, tali ruote si accettano appunto.
Ma il rimedio esiste, ed \`e uno di quei ``toccasana'' che non comporta
spese, al contrario, oltre a determinare la scomparsa del sottotaglio,
le {\em correzioni tramite
lo spostamento del profilo}\index{correzione!profilo} portano una serie
di vantaggi tali da essere impiegate anche quando il sottotaglio non si
manifesta affatto.
In aggiunta, mentre da un punto di vista didattico le dentature
corrette presentano per lo studente un ulteriore (l'ultimo) piccolo
grattacapo, nella pratica esse si progettano e si eseguono con estrema
disinvoltura: se si dovesse progettare l'ingranaggio di figura \ref{fig:1760}
e il pignone di $z_1=17$ denti fosse la ruota motrice
ben difficilmente esso non verrebbe corretto.
\begin{figure}[hbt]
\begin{center}
\includegraphics[width=1.0\textwidth]{part2/ruote/FIG/ruote/1160c_i0.pdf}
\begin{picture}(0,0)(200,0)
\scriptsize{
\color{white}
\put(170,172){$A$}
\put(280,155){$z_1=11$}
\put(50,155){$z_2=60$}
\put(285,138){$O_1$}
\put(148,144){$P$}
\put(123,82){$p_b$}
\put(124,71){$\mu$}
\put(97,46){$r_{b_2}$}
\put(138,86){$B$}
\color{black}
\put(146,46){$r_{p_2}$}
\put(213,38){$r_{p_1}$}
\put(225,49){$r_{b_1}$}
\put(157,23){$r_{t_2}$}
\put(190,23){$r_{t_1}$}
}
\end{picture}
\end{center}
\vskip -5mm
\caption{\em
Ingranaggio con $z_1=11$, $z_2=60$, $\theta=20^{\circ}$, con correzioni $x_1=0.5$ e $x_2=-0.5$.
} 
\vskip -3mm
\label{fig:1160c_i0}
\end{figure}
Ricordiamo che nelle due figure \ref{fig:dentiera_riferimento} e
\ref{fig:dentiera_maag} abbiamo riportato una retta tratteggiata indicata
con $\lambda_t$, chiamata {\em primitiva di taglio}\index{primitiva!di taglio}, distante dalla primitiva della
dentiera di una quantit\`a $c=m x$. Quando le ruote vengono tagliate facendo
rotolare $\lambda_t$ sulla primitiva teorica delle ruote stesse con $c \ne 0$
esse si dicono {\em corrette}\index{ruote!corrette}. Riferendoci all'utensile {\em Maag} (ma
la questione rimane identica anche per utensili creatori a vite), si tratta
semplicemente di spostare la dentiera di una quantit\`a $c$ verso l'esterno
della ruota da
tagliare se $c >0$, oppure verso l'interno qualora la correzione fosse negativa.
In figura \ref{fig:1160c_i0} rappresentiamo l'ingranaggio
$z_1=11$ e $z_2=60$, le cui ruote sono state realizzate spostando la
primitiva di taglio dei due 
utensili delle quantit\`a adimensionali $x_1=0.5$ e $x_2=-0.5$. Anche i raggi
delle due circonferenze di testa dovranno subire
un incremento $\delta r_{1_t}=0.5 m$
e un decremento  $\delta r_{2_t}=-0.5 m$. 
Notiamo con facilit\`a il cambiamento di forma della base dei denti del 
pignone, che non palesano pi\`u alcun sottotaglio, a fronte della riduzione
modesta e
probabilmente accettabile dello spessore alla base di quelli dell'altra ruota.
Come riportato in figura \ref{fig:1160c_i0},
l'angolo di pressione vale ancora $\theta=20^{\circ}$, mantenendo 
il valore di quello dell'utensile, e l'interasse vale $i=(z_1+z_2)m/2$.
Queste ultime affermazioni sono meno ovvie
rispetto a ci\`o che a prima vista si potrebbe ritenere. In particolare
la seconda (dalla quale per\`o dipende la prima), e cio\`e che le ruote
corrette possano formare ingranaggi senza variazione di interasse rispetto
a quello teorico.
Infatti, quando la somma $c_1+c_2 \ne 0$, l'interasse di funzionamento,
cio\`e la distanza che deve separare i centri delle ruote corrette affinch\'e
ingranino correttamente, non corrisponde al valore di tale somma aggiunta
all'interasse teorico, come dimostreremo a breve. 
Intanto osserviamo ancora che l'ingranaggio di figura \ref{fig:1160c_i0} presenta
fattore di ricoprimento $f_c > 1$, fatto che andrebbe calcolato tramite
la formula \ref{eq:segpb2} e le successive, opportunamente modificate, anche se,
in questa sede, ci accontentiamo di ci\`o che si manifesta visivamente. Appare
anche chiaro che il tratto $PA$, che si chiama {\em accesso},\index{accesso}
risulta alquanto pi\`u corto di $PB$, che si chiama {\em recesso}\index{recesso}.
I due nomi
presuppongono che il pignone funga da motore, altrimenti tali nomi devono essere
invertiti. Senza entrare nel merito delle considerazioni dinamiche che
supportano l'affermazione che formuliamo, affermiamo che \`e sempre conveniente
avere un tratto di accesso ridotto rispetto a quello di recesso.
L'ingranaggio di figura \ref{fig:1160c_i0} \`e pertanto formato da due ruote
corrette mediante lo spostamento dei loro profili in due sensi opposti,
e il filo logico sottinteso alle correzioni
potrebbe essere cos\`i riassunto: si corregge il pignone, in senso positivo, 
fino alla eliminazione del sottotaglio; lo stesso valore della correzione,
cambiato di segno, si applica all'altra ruota, sperando di non spostare su
quest'ultima il sottotaglio e mantenendo invariato l'interasse teorico.
Anche se questa soluzione appare elegante, di solito non viene scelta.
Come chiarito anche in \cite{righettini}, il ragionamento 
sviluppato appena sopra sembra presupporre che il proporzionamento normale
rifletta condizioni di funzionamento ottimali e pertanto sia opportuno
derogare alle sue leggi il meno possibile. Gli spostamenti dei profili si
dimostrano invece vantaggiosi a prescindere dal problema del sottotaglio.
\`E tutt'altro che infrequente progettare le ruote dentate sacrificando
l'interasse teorico a favore dell'ottenimento
di altre caratteristiche, oppure imponendo un 
interasse funzionale a geometrie particolari. Quest'ultimo
caso si presenta spesso
per gli assi paralleli dei cambi delle automobili sui quali si devono
calettare quattro o pi\`u coppie di ruote dentate, sempre in presa,
aventi somme
dei denti leggermente diverse tra loro e caratteristiche geometriche tali
(si tratta quasi sempre di ruote elicoidali)
da non poter garantire un interasse comune senza essere corrette.
 Non entreremo nell'ambito dei criteri
di scelta delle correzioni che si apportano per ottimizzare, sotto svariati
aspetti, il funzionamento degli ingranaggi. Diciamo solamente che il caso
di figura \ref{fig:1160c_i0} potrebbe essere approcciato da un progettista 
correggendo solo il pignone e lasciando la ruota da sessanta denti con geometria
normale. Di seguito analizzeremo appunto questo caso, dove la correzione sar\`a
apportata solamente sul pignone, con la conseguenza di dover 
determinare un nuovo angolo di
pressione e un nuovo interasse, diversi da quelli normali.


\section{Ruote Corrette mediante Spostamento del Profilo II}\index{correzione!con variazione interasse}

\noindent Ricordiamo che, durante il taglio delle ruote corrette,
nonostante lo spostamento della primitiva della dentiera, tale primitiva
dovr\`a rotolare sulla primitiva teorica della ruota da tagliare, inviluppando,
in questo modo, evolventi di circonferenze di base ancora di raggio
$r_{1_b}= z_1 m /2 \cos \theta$ e $r_{2_b}= z_1 m /2 \cos \theta$. Di
conseguenza, qualsiasi siano le correzioni dei profili, il
{\em passo base}\index{passo!base} $p_b$, gi\`a riportato in figura
\ref{fig:dentiera_maag}, non cambia.
La figura \ref{fig:1160c_ix} riporta l'ingranaggio di due ruote
dentate, sempre con $z_1=11$ denti e $z_2=60$ denti. Come si legge nella didascalia,
la correzione positiva \`e stata applicata solamente al pignone, con $x=0.5$,
mentre la seconda ruota non \`e corretta. Prima di rispondere alla legittima
domanda: come sono stati determinati l'interasse effettivo di funzionamento
e il nuovo angolo di pressione? Riteniamo opportuno ribadire che
la teoria, e in particolare le premesse alle formule \ref{eq:nuove_primitive},
ci assicurano che i nostri profili a evolvente sono ancora coniugati.
In quest'ambito, quello cio\`e delle correzioni con variazione di
 interasse, si hanno
due casi emblematici. 
\begin{figure}[t]
\begin{center}
\includegraphics[width=1.0\textwidth]{part2/ruote/FIG/ruote/1160c_ix.pdf}
\begin{picture}(0,0)(200,0)
\scriptsize{
\color{white}
\put(270,40){$z_1=11$}
\put(40,120){$z_2=60$}
\put(270,18){$O_1$}
\put(39,28){$r_{b_2}$}
\put(96,78){$r_{{p_2}_f}$}
\put(182,185){$l_1$}
\put(98,140){$l_2$}
\put(167,178){$\mu$}
\put(258,187){$\zeta$}
\put(262,175){$\gamma(\theta_0)$}
\put(264,165){$\gamma(\theta_f)$}
\put(217,83){$\zeta$}
\put(220,54){$\gamma(\theta_0)$}
\put(227.5,27){$\gamma(\theta_f)$}
\put(291,106){$\delta_1$}
\color{black}
\put(232,185){$r_{b_1}$}
\put(120,210){$r_{t_2}$}
\put(97,167){$p$}
\put(162,219){$r_{t_1}$}
\put(147,149){$p$}
\put(354.5,189){$r_{{p_1}_f}$}
\put(329.5,199.5){$r_{{p_1}_0}$}
}
\end{picture}
\end{center}
\vskip -5mm
\caption{\em
Ingranaggio con $z_1=11$, $z_2=60$, $\theta_0=20^{\circ}$, $\theta_f=22^{\circ}$, $i_0=142$, $i_f=143.9$, con le correzioni $x_1=0.5$ e $x_2=0$.
} 
\vskip -3mm
\label{fig:1160c_ix}
\end{figure}
Esporremo per primo quello che porta 
all'individuazione dell'interasse
di funzionamento $i_f$ e del corrispondente angolo di pressione $\theta_f$,
date le due correzioni arbitrarie delle ruote. 
Il secondo caso fa invece riferimento
all'individuazione di quelle correzioni necessarie a produrre un interasse
di funzionamento dato.
Riprendiamo dalla figura \ref{fig:1160c_ix} e cerchiamo di individuare
un criterio, date le correzioni delle due ruote, per potere stabilire
l'effettivo interasse di funzionamento e l'effettivo angolo di pressione.
Da un punto di vista cinematico, l'ingranaggio della figura appena citata
funzionerebbe egregiamente anche con l'interasse dato dalla somma dei
raggi delle primitive teoriche pi\`u la somma delle correzioni:
$i_g=(z_1 +z_2) m /2 +(x_1+x_2)m$. Se montassimo in questo modo le ruote,
noteremmo per\`o un certo gioco che, in generale, non pu\`o essere tollerato
nelle trasmissioni di potenza, onde evitare che i denti ``sbattano''
\index{sbattimento} a causa delle eventuali (frequentissime nelle varie
applicazioni) inversioni della direzione del flusso di
potenza. Ecco quindi il criterio cercato: le ruote devono essere montate senza gioco\footnote{
Un minimo gioco tra le ruote deve sempre sussistere, per evitare ai cuscinetti
di supporto di funzionare in regime di carico gravoso e non previsto,
ma soprattutto per consentire la lubrificazione dei denti. La 
quantificazione di tale gioco \`e una questione delicata che dipende da
un numero rilevante di fattori: condizione di impiego della coppia di ruote,
precisione e finitura 
delle ruote stesse, tipo di lubrificazione, e molti altri. Rimandiamo
gli studenti desiderosi di approfondire alla letteratura specializzata, 
come ad esempio \cite{henriot}.
}.
Riferendoci ancora alla figura \ref{fig:1160c_ix}, quanto appena stabilito
si traduce nella seguente relazione $p=l_1 + l_2$. Questa relazione afferma
che il passo sulle primitive di funzionamento $p$, il quale ovviamente,
affinch\'e le ruote possano ingranare, deve essere lo stesso
su entrambe, \`e dato dalla
somma dello spessore dei denti sulle primitive stesse.
\`E questo uno dei casi in cui scopriamo l'utilit\`a della rappresentazione
analitica dell'evolvente e, in particolare, della \ref{eq:gamma}.
Cominciamo col calcolare la lunghezza dell'arco della primitiva teorica, o
di riferimento,
che attraversa un dente. A questo proposito notiamo che, quando la dentiera viene
spostata della quantit\`a $c=mx$, il vano sulla primitiva di taglio, che sar\`a
pari allo spessore del dente sulla primitiva di riferimento della ruota, vale
\begin{equation}
l_0= m \left( \tfrac{\pi}{2} + 2 x \tan(\theta_0) \right)\,,
\label{eq:l0}
\end{equation}
\noindent dove l'addendo di destra non \`e altro che la base di un triangolo
isoscele di altezza pari a $c=mx$ e angolo di apertura uguale al 
doppio dell'angolo dei profili della dentiera. La \ref{eq:l0} va specializzata
per le due ruote, in quanto le correzioni $x_1$ e $x_2$ saranno,
in generale, 
diverse. La conoscenza di $l_{0_1}$ e $l_{0_2}$ ci fornisce con facilit\`a
il valore degli angoli sottesi da tali archi:
\begin{equation}
\delta_1 = l_{0_1}/r_{{p_1}_0}\;\;\; {\rm e} \;\;\;\;
\delta_2 = l_{0_2}/r_{{p_2}_0}\,.
\label{eq:delta}
\end{equation}
\noindent Le lunghezze degli archi $l_1$ e $l_2$ si otterranno considerando
gli effettivi valori del raggio primitivo di funzionamento delle due ruote 
e l'aumento degli angoli $\delta$ dovuti all'angolo $\zeta$, come riportato
in figura \ref{fig:1160c_ix}.
Quanto a $\zeta$, sempre dalla figura, abbiamo
\begin{equation}
\zeta= \gamma(\theta_0)-\gamma(\theta_f)\,,
\label{eq:zeta}
\end{equation}
\noindent e tale valore\footnote{
Per la coordinata polare dell'evolvente abbiamo mantenuto la nomenclatura
esposta nella formula \ref{eq:gamma}, $\gamma()$, bench\'e
il nome pi\`u comunemente usato nella letteratura sia {\em involute
function}\index{involute function}. Pertanto, con un'abbreviazione a nostro
parere infelice, abbiamo
\begin{equation}
 {\rm inv}(\theta) = \tan(\theta)-\theta \,,
\label{eq:inv}
\end{equation}
come si trova, ad esempio, in \cite{henriot}.
} sar\`a lo stesso per le due ruote. Finalmente ricaviamo
\begin{equation}
l_1= r_{{p_1}_f} (\delta_1-2\zeta)\;\;\; {\rm e} \;\;\; l_2= r_{{p_2}_f} (\delta_2-2\zeta)\,,
\label{eq:l1l2}
\end{equation}
\noindent quindi, per il passo $p=l_1+l_2$, misurato sulle primitive di 
funzionamento possiamo scrivere
\begin{multline}
p= l_1+l_2= 
2 r_{{p_1}_f}[ \frac{1}{z_1} (\pi/2 + 2 x_1 \tan(\theta_0)) + \zeta] +\\
+2 r_{{p_2}_f}[ \frac{1}{z_2} (\pi/2 + 2 x_2 \tan(\theta_0)) + \zeta]\,.
\end{multline} 
\noindent A valle di alcune semplici sostituzioni, tenendo presente 
che $r_{{p_1}_f}= z_1 p/\pi$ e $r_{{p_2}_f}= z_2 p/\pi$,
esplicitando $\zeta$ tramite la \ref{eq:zeta}, otteniamo
\begin{equation}
\gamma(\theta_f)=2 \frac{x_1+x_2}{z_1+z_2}\tan(\theta_0) +\gamma(\theta_0)\,.
\label{eq:thetaf}
\end{equation}
\noindent La funzione $\gamma(\theta)$, col nome citato nella \ref{eq:inv}, \`e
riportata, tabulata, in diversi trattati: citiamo \cite{henriot} dove,
dalla pagina 72 alla pagina 79, $\theta$ spazia da $10^{\circ}$ a $50^{\circ}$. Tale funzione \`e
regolare e monotona crescente, pertanto, senza scomodare metodi di ricerca
degli zeri recanti nomi altisonanti, il comune metodo di bisezione
(che abbiamo usato per ottenere le nostre figure) porta,
con estrema rapidit\`a, a soluzioni apprezzabili della \ref{eq:thetaf}. Una
volta individuato il valore di $\theta_f$, che sar\`a l'effettivo angolo di
pressione, ricordando che in ogni caso
il raggio della circonferenza di base non cambia,
come abbiamo gi\`a altrove sottolineato, avremo
\begin{equation}
r_{{p_1}_f} \cos(\theta_f)=r_{{p_1}_0} \cos(\theta_0)\;\;\; {\rm e} \;\;\;
r_{{p_2}_f} \cos(\theta_f)=r_{{p_2}_0} \cos(\theta_0)\,,
\end{equation}
\noindent dalle quali si ottiene
\begin{equation}
i_f= r_{{p_1}_f} + r_{{p_2}_f}=  (r_{{p_2}_0}+ r_{{p_1}_0}) \frac{ \cos(\theta_0)}{\cos(\theta_f)}= m(z_1+z_2)\frac{ \cos(\theta_0)}{2 \cos(\theta_f)}\,.
\label{eq:if}
\end{equation}
\begin{wrapfigure}{r}{0.54\textwidth}
      \begin{center}
      \includegraphics[width=0.52\textwidth]{part2/ruote/FIG/ruote/interasse.pdf}
     \end{center}
\begin{picture}(0,0)(-63,1)
\scriptsize{
        \put(-46,119){$\ni$}
        \put(78,19){$1-i_f/i_0$}
}
\end{picture}
\vskip -4.1mm
        \caption{\em Rapporto tra interasse effettivo $i_f$ e interasse
intuitivo $i_0+m(x_1+x_2)$.}
     \label{fig:interasse}
\end{wrapfigure}
\noindent Ci sembrano opportune le seguenti due considerazioni. La prima:
$\cos(\theta_f)$, e perci\`o anche l'interasse di funzionamento $i_f$, dipende
soltanto dalla somma delle correzioni quindi, come gi\`a ammesso nel
precedente paragrafo, se $x_1 + x_2=0$ risulta $\theta_f=\theta_0$ e $i_f=i_0$,
 cio\`e angolo di pressione e interasse rimangono inalterati.
 La seconda osservazione pu\`o essere elaborata
come risposta alla domanda: l'interasse di funzionamento sar\`a dato da
 $i_f=i_0+m(x_1+x_2)$? La risposta \`e no, e per convincerci di ci\`o
studiamo l'andamento del rapporto tra l'interasse effettivo e la somma delle
correzioni addizionata all'interasse teorico, che chiamiamo $\ni$
\begin{equation}
\ni=\frac{i_f}{i_0+m(x_1+x_2)}\,,
\label{eq:ni}
\end{equation}
\noindent in funzione
della variazione adimensionale dell'interasse $\delta_i=1-i_f/i_0$.
\noindent La figura \ref{fig:interasse} mostra quindi il grafico di $\ni$
al variare di $\delta_i$, che per $0<=\delta_i<=0.1$  
evidenzia una riduzione di $i_f$ rispetto
all'interasse che si sarebbe individuato a intuito $i_0+m(x_1+x_2)$.
Tale situazione dovr\`a, peraltro, risultare compatibile con i giochi tra testa e piede
 dei denti e richiede pertanto una specifica verifica.
\vskip 3mm
\noindent Veniamo ora al secondo dei due casi emblematici
per le correzioni con variazione di interasse, che \`e quello 
in cui le correzioni incognite devono
risultare compatibili con un interasse dato. Stabilito quindi il valore di $i_f$,
dalla \ref{eq:if} ricaviamo $\theta_f$, come gi\`a accennato con la 
formula \ref{eq:nuove_primitive}.
Dalla \ref{eq:thetaf} ricaviamo poi,
con facilit\`a, la somma delle correzioni $x_1+x_2$ ammissibile con 
l'interasse di funzionamento assegnato. 
Il tema della ripartizione delle
correzioni sulle due ruote \`e complesso e rientra nell'ambito dello
studio dei criteri di scelta delle correzioni stesse. Ribadiamo che
decidere di dare al pignone la minima correzione necessaria a scongiurare
il sottotaglio dei denti e quindi regolarci di conseguenza circa la correzione
della seconda ruota, magari con una correzione di segno opposto, cos\`i
da lasciare inalterato l'interasse di progetto,
 \`e molto riduttivo e anacronistico. 
Una guida alla scelta delle correzioni tramite percorsi volti
 all'ottimizzazione degli ingranaggi (percorsi e criteri
 a volte in contesa
tra loro) \`e riportata in \cite{righettini}. Riteniamo per\`o fuori
dallo scopo di queste note approfondire ulteriormente questo tema che,
come abbiamo anticipato, \`e materia decisamente specialistica.
\begin{figure}[hbt]
      \begin{center}
      \includegraphics[width=0.65\textwidth]{part2/ruote/FIG/ruote/forma_dente.pdf}
     \end{center}
\begin{picture}(0,0)(-63,1)
\scriptsize{
\color{white}
        \put(30,53){$x=-0.3$}
        \put(37,65){$x=0$}
        \put(31,79){$x=0.3$}
        \put(31,93){$x=0.6$}
        \put(25,110){$x=1.086$}
}
\end{picture}
\vskip -5mm
        \caption{\em
Variazione della forma dei denti di una ruota con $z=17$ e correzioni $x=-0.3$, $x=0$, $x=0.3$, $x=0.6$ e $x=1.086$.
}
     \label{fig:forma_dente}
\end{figure}

\noindent Riepiloghiamo
l'argomento della correzione del profilo dei denti sottolineando tre punti
che ci sembrano importanti. 
Il primo riguarda la forma del dente, sensibilmente modificata
dallo spostamento del profilo,
come riportato in figura \ref{fig:forma_dente}.
Al crescere della correzione si passa infatti dalla condizione di
 sottotaglio a quella
che vede l'irrobustimento della 
base dei denti, fino all'appuntimento dei denti stessi quando si 
adottano valori di correzione eccessivi.
Nella figura \`e appunto riportato, per una ruota con $z=17$, il
caso di appuntimento del dente, quando le due evolventi si incrociano esattamente
sulla circonferenza di testa. La situazione dei denti a punta, ma anche con
spessore di testa troppo piccolo, \`e da evitarsi per ragioni di resistenza
meccanica.

\noindent Il secondo punto saliente riguarda
 il fattore di ricoprimento $f_c$. Esso, in sostanza, dipende
dalla somma delle correzioni $x_1 + x_2$, e 
diminuisce al crescere di questa somma, mentre non \`e molto sensibile ai valori
delle singole correzioni.
Con $x_1+x_2=0$, il fattore di ricoprimento \`e compreso tra $1.4<f_c<2$, che
ovviamente \`e anche l'intervallo dove spazia tale fattore per le ruote
non corrette.
Adottando correzioni la cui somma \`e negativa si possono
ottenere elevati valori del fattore di ricoprimento, a beneficio della
silenziosit\`a dell'ingranaggio.
Quando invece $x_1 + x_2 > 0$ il fattore di ricoprimento diminuisce e questo
pu\`o arrivare a compromettere
 la continuit\`a della trasmissione. Pertanto, nel caso di correzioni a
somma algebrica maggiore di zero, cio\`e con aumento dell'interasse,
bisogna sempre verificare il valore di $f_c$.

\noindent Il terzo punto riguarda l'allontanamento del segmento dei
 contatti $AB$ dalle ruote con correzione positiva. Ridurre la porzione
di accesso sulla linea dei contatti porta un netto
miglioramento del rendimento degli ingranaggi. Per tale motivo ai pignoni
motori, che sono una buona parte di quelli che si costruiscono, si applicano
volentieri correzioni positive che, eliminando il sottotaglio e irrobustendo
la base del dente a favore della sua resistenza strutturale, ne migliorano anche
il rendimento diminuendo il tratto di accesso, come mostra il confronto
tra i tratti $PA$ e $PB$ della figura \ref{fig:1160c_i0}.
Le cose vanno diversamente per i moltiplicatori di velocit\`a, dove la ruota
motrice non necessiterebbe, di per s\'e, correzioni positive con lo scopo
di eliminare il sosttotaglio. In tali casi, i diagrammi che ci guidano
nel ripartire la somma delle correzioni $x_1+x_2$ sulle due ruote possono
allocare una correzione positiva maggiore alla ruota motrice, cio\`e,
in questo caso, quella con maggiore numero di denti, come mostrato
in \cite{righettini}, pag. 52, e, in alcuni casi, persino 
tollerare un leggero
sottotaglio sulla ruota condotta, quella cio\`e con numero di denti minore.

\section{Note Conclusive sulle Ruote Dentate}

A conclusione di questo lungo capitolo, proponiamo un'analisi
ad elevatissima
 velocit\`a del percorso che, altrove, abbiamo seguito con la giusta
cadenza in modo da poter essere sviluppato senza incertezze. Tale traccia
ci ha condotto, quasi attraverso alcuni passaggi obbligati, alle ruote
con denti profilati a evolvente che dominano
interamente il panorama degli ingranaggi. Ad esempio, anche se non
ne abbiamo fatto cenno esplicito in questo lavoro, le ruote con denti ad asse
elicoidale, cio\`e le {\em ruote elicoidali}\index{ruote!elicoidali},
 pur possedendo caratteristiche specifiche, che un tecnico
deve conoscere perfettamente per poterle inserire correttamente 
nel progetto di una macchina, non
presentano, dal punto di vista della loro cinematica, nulla di diverso
dalle ruote che abbiamo studiato. Anzi, la loro fabbricazione si ottiene
semplicemente inclinando di un certo angolo l'utensile {\em Maag} di cui abbiamo
parlato diffusamente (oppure l'asse della ruota da tagliare).
Quindi profilo a evolvente ovunque, sulle ruote
coniche, che possono trasmettere il moto tra alberi concorrenti, sulle dentature
interne, e cos\`i via. Tale profilo si ottiene ``automaticamente'', come abbiamo visto,
dal processo di inviluppo tramite il taglio con utensili molto semplici.
In pratica, l'utilizzo dell'utensile pi\`u semplice da realizzare e da
affilare, cio\`e della dentiera riportata in figura
\ref{fig:dentiera_maag},
 ci fornisce direttamente e senza sforzo le ruote dentate a evolvente,
assortibili, facilmente adattabili a modeste variazioni di interasse,
facili da correggere per il raggiungimento di determinati scopi, come
l'eliminazione del sottotaglio dei denti.

\noindent E questa \`e
la conclusione alla quale ci premeva arrivare. Il successo di una 
tecnologia \`e spesso legato alla sua semplicit\`a e
alla sua attuabilit\`a.
Per questo tipo di  ruote non ha gran senso neppure lo studio numerico
del loro profilo come invece risulta d'obbligo,
da quando i calcolatori sono entrati a far parte degli strumenti
del progettista,  nella profilatura 
delle camme. Infatti, come abbiamo ormai troppe volte ripetuto,
l'ottenimento del loro sofisticato profilo
segue un processo
indiretto rappresentato dall'inviluppo creato dall'utensile dentiera,
processo che abbiamo usato anche noi per le nostre figure.
Anni fa, soprattutto
in fase prototipale, alcune
ruote dentate si lavoravano mediante l'impiego di frese a disco
che, scavando nel cilindro grezzo della futura ruota dentata
il vuoto tra un dente e il successivo,
modellavano i fianchi dei denti stessi.
Pu\`o darsi che qualche meccanico affezionato
alla {\em fresatura di forma}\index{fresatura di forma}
dei denti si possa ancora trovare presso
le officine che eseguono riparazioni e restauro di oggetti particolari.
La forma delle frese impiegate in questo procedimento
dipende naturalmente dal modulo della dentatura che si desidera ottenere,
ma anche, come \`e ovvio, dal numero dei denti della ruota.
Anzi, a ben vedere, sarebbe necessario impiegare 
una fresa per ogni ben preciso numero di denti.
Nella pratica per\`o bastano una decina di
utensili (di un certo modulo) per avere risultati soddisfacenti con
$z$ che spazia in un largo intervallo. Tagliando in
questo modo  il profilo dei denti
si perde per\`o completamente il controllo circa la 
capacit\`a della dentatura ottenuta di ospitare correttamente i denti
dell'altra ruota. \`E qui infatti che nasceva il problema dell'interferenza
propriamente detta,
la quale invece si manifesta come sottotaglio
dei denti quando questi si ottengono da un inviluppo. Rispetto alla
fresatura di forma,
il processo di taglio con creatore o dentiera {\em Maag} \`e
talmente pi\`u comodo, pi\`u preciso, pi\`u flessibile (si pensi
alle correzioni) da avere dato, dopo la sua comparsa,
una propulsione violenta a questo settore e allo sviluppo
della meccanica dello scorso secolo. Ma
cosa ci riserva il futuro? Si potrebbe pensare di fabbricare le ruote
dentate tagliandone direttamente il profilo con macchine
a controllo numerico, come potrebbe fare un'elettroerosione a filo?
Oppure, si potrebbe concepire di ottenerle tramite deposizione controllata
di materiale?  \`E chiaro che se si considerano queste opportunit\`a,
l'ottenimento dell'evolvente non \`e automatico e dovremmo
basarci su opportuni codici di calcolo che ci forniscano la geometria del
dente stesso. Su questo campo da gioco, poco conosciuto,
sospettiamo per\`o che il profilo a evolvente
avrebbe pi\`u di un concorrente.

\endinput

\chapter{Cenni ai Sistemi Articolati}

\section{Introduzione}

\noindent Due o pi\`u corpi rigidi possono essere collegati tra loro mediante
{\em coppie rotoidali}\index{coppia!rotoidale}\footnote{In questi appunti non si descrivono le coppie rigide 
elementari; quale esempio di tali scorrimenti di superfici di corpi rigidi l'una
sull'altra citiamo la coppia rotoidale che collega le due lame delle comuni forbici
mediante un perno, precisando che le superfici combacianti sono quella del cilindro
del perno e le due superfici cilindriche dei fori che lo ospitano. Rimandiamo  il
lettore che desideri approfondire l'argomento 
a \cite{sesini1}, pagg. 22-24, o anche \cite{martin}, pagg. 9-10, o ancora
\cite{sesini3} che tratta numerosi casi di notevole interesse.}.
\noindent Le catene cinematiche che si ottengono in questo modo prendono
il nome di {\em sistemi articolati}\index{sistemi!articolati} e tale denominazione
viene mantenuta anche nei casi in cui una o pi\`u coppie rotoidali sia sostituita
da una sua degenerazione che, portando l'asse della coppia all'infinito, diventa una
{\em coppia prismatica}\index{coppia!prismatica}.
Questa vasta famiglia di meccanismi trova estese e notevolissime applicazioni:
ogni motore a combustione interna contiene almeno un sistema articolato
(il manovellismo ordinario) e svariate
sono le applicazioni dei {\em quadrilateri articolati}\index{quadrilateri!articolati} 
nelle macchine automatiche, nelle presse, nelle macchine per cucire; l'elenco
completo delle applicazioni risulterebbe molto lungo.
Esistono sistemi articolati con coppie rotoidali i cui assi sono orientati
in direzioni sghembe tra loro oppure concorrenti in un  punto 
e qualcuno di questi {\em sistemi articolati spaziali}
ha notevoli applicazioni, come il {\em giunto di Cardano}. 
La stragrande maggioranza dei sistemi articolati presenta invece coppie rotoidali
con assi tra loro paralleli. I movimenti relativi tra i vari membri si 
realizzeranno perci\`o solo in piani paralleli tra loro e ortogonali
a tali assi e questo ci permette
di identificarli come {\em sistemi articolati piani}\index{sistemi!articolati piani},
gli unici di cui accenniamo lo studio.

\noindent Dato un sistema articolato piano, senza perdita di generalit\`a, possiamo
considerare i moti piani dei suoi membri, che in realt\`a hanno sempre luogo
in piani paralleli tra loro, ma distinti, come contenuti in un solo piano. Le tracce degli assi delle coppie rotoidali che intersecano questo piano costituiscono perci\`o
un insieme di $n$ punti, tanti quanti sono gli snodi, e tali punti sono tra loro
collegati tramite $m$ corpi rigidi. Ogni punto nel piano gode di due gradi di
libert\`a di movimento, mentre i collegamenti rigidi, che mantengono invariata
la distanza tra questi punti, assorbono un grado di libert\`a ciascuno.
Ne consegue che, quando non si presentano situazioni anomale nelle quali una o 
pi\`u parti del meccanismo risulti sovra-vincolata e magari altre parti labili,
il numero di gradi di libert\`a residuo $g$ per un sistema articolato piano 
risulta essere $g=2n - m$.

\begin{figure}[hbt]
\centering
\begin{minipage}[b]{0.30\textwidth}
\centering
\includegraphics[width=0.8\textwidth]{part2/quadri/FIG/reticolare_triangolo.pdf}
\begin{picture}(0,0)(130,0)
\scriptsize{
\put(84,85){$n=3$}
\put(84,75){$m=3$}
\put(84,65){$g=3$}
}
\end{picture}
      \caption{\em Sistema articolato con tre coppie rotoidali e tre membri.}
 \label{fig:f_triangolo}
\end{minipage}\hfill
\begin{minipage}[b]{0.30\textwidth}
\centering
\includegraphics[width=0.8\textwidth]{part2/quadri/FIG/pentalatero.pdf}
\begin{picture}(0,0)(129,0)
        \scriptsize{
\put(82,85){$n=4$, $m=5$}
\put(82,75){$g=3$}
}
\end{picture}
        \caption{\em Sistema articolato con quattro coppie rotoidali e cinque membri.}
     \label{fig:f_pentalatero}
\end{minipage}\hfill
\begin{minipage}[b]{0.30\textwidth}
\centering
\includegraphics[width=0.8\textwidth]{part2/quadri/FIG/quadrilatero.pdf}
\begin{picture}(0,0)(129,0)
        \scriptsize{
\put(82,85){$n=4$, $m=5$}
\put(82,75){$g=4$}
}
\end{picture}
        \caption{\em Sistema articolato con quattro coppie rotoidali e quattro membri.}
     \label{fig:f_quadrilatero}
\end{minipage}
\end{figure}

\noindent In figura \ref{fig:f_triangolo} \`e riportato il pi\`u semplice sistema
articolato non banale. I gradi di libert\`a residua per tale sistema sono
$g=3\times 2- 3=3$. Tre gradi di libert\`a, nel piano, sono quelli posseduti da
qualsiasi corpo rigido svincolato, e il nostro sistema comportandosi come
un corpo indeformabile  viene classificato
e studiato nella  scienza delle costruzioni come sistema reticolare. Anche il sistema
articolato riportato in figura \ref{fig:f_pentalatero} presenta un numero di gradi
di libert\`a pari a tre. Anch'esso si comporta come un corpo rigido e fa quindi 
parte delle {\em strutture reticolari}\index{strutture reticolari}.
Il sistema riportato in figura \ref{fig:f_quadrilatero} presenta invece  $g=4$,
quindi conserva un grado di deformabilit\`a interna e per questo fa parte dei
meccanismi articolati.


\section{Quadrilateri Articolati: Velocit\`a e Accelerazioni}

\noindent Il {\em quadrilatero articolato piano}\index{quadrilateri!articolati piani}
sar\`a il solo oggetto di questa breve trattazione, in questo capitolo.
La condizione pi\`u comune \`e quella che vede
uno dei suoi quattro membri fungere da telaio: in figura
\ref{fig:f_quad_telaio} il membro $AD$ \`e il telaio. 
Il membro $BC$, sul quale non si trovano coppie rotoidali vincolate
a terra, prende il nome di {\em biella}\index{biella}. Gli altri due membri, $AB$ e $DC$, 
si chiamano {\em manovelle}\index{manovella} oppure
{\em bilancieri}\index{bilanciere} a seconda che la circostanza di potere
compiere l'intero angolo giro sia o meno verificata.

\noindent Fissato il lato che funger\`a da telaio e quindi la biella, la condizione degli
altri due membri viene determinata dalla seguente
{\em regola di Grashof}\index{Grashof, regola di}
che riportiamo pari pari e senza dimostrazione da \cite{sesini1}, pag. 99:
{\em il quadrilatero articolato piano pu\`o essere a doppia manovella o
a manovella-bilanciere,
soltanto se la somma del pi\`u piccolo e del pi\`u grande dei suoi lati non \`e
maggiore della somma degli altri due. In tal caso si ha una doppia manovella,
se l'asta fissa \`e la pi\`u corta, una manovella e un bilanciere se l'asta
fissa \`e una delle contigue alla pi\`u corta (essendo la manovella il lato pi\`u corto).
In tutti gli altri casi il quadrilatero \`e a doppio bilanciere}.

\begin{figure}[hbt]
\centering
\begin{minipage}[b]{0.48\textwidth}
\centering
     \includegraphics[width=0.88\textwidth]{part2/quadri/FIG/quadrilatero_telaio.pdf}
\begin{picture}(0,0)(170,25)
        \scriptsize{
        \put(12,169){$B$}
        \put(155,111){$C$}
        \put(132,73){$D$}
        \put(12,60){$A$}
}
\end{picture}
        \caption{\em Quadrilatero con telaio vincolato a terra.}
     \label{fig:f_quad_telaio}
\end{minipage}\hfill
\begin{minipage}[b]{0.48\textwidth}
\centering
    \includegraphics[width=0.88\textwidth]{part2/quadri/FIG/quadrilatero_schematico.pdf}
\begin{picture}(0,0)(170,25)
        \scriptsize{
        \put(12,169){$B$}
        \put(94,153){$\omega_{\scriptscriptstyle{BC}}$}
        \put(115,145){${\dot{\omega}}_{\scriptscriptstyle{BC}}$}
        \put(46,112){$\omega_{\scriptscriptstyle{AB}}$}
        \put(46,100){${\dot{\omega}}_{\scriptscriptstyle{AB}}$}
        \put(155,111){$C$}
        \put(148,90){$\omega_{\scriptscriptstyle{AD}}$}
        \put(142,78){${\dot{\omega}}_{\scriptscriptstyle{AD}}$}
        \put(128,72){$D$}
        \put(12,60){$A$}
}
\end{picture}
        \caption{\em Rappresentazione schematica del quadrilatero.}
     \label{fig:f_quad_schema}
\end{minipage}
\end{figure}

\noindent Ad esempio nel caso di figura \ref{fig:f_quad_telaio} abbiamo $|\overrightarrow{DC}|+|\overrightarrow{BC}|<|\overrightarrow{AB}|+|\overrightarrow{AD}|$,
quindi avremo tre possibilit\`a: due manovelle scegliendo $DC$ come telaio,
una manovella e un bilanciere quando il telaio \`e $AD$ oppure $BC$, 
infine due bilancieri con telaio $AB$. Nel caso riportato in figura
\ref{fig:f_quad_telaio}
 il telaio risulta essere
un lato contiguo al lato pi\`u corto: il membro $CD$ funger\`a, in questo
caso, da manovella
mentre il lato $AB$ sar\`a un bilanciere.



\noindent La cinematica del quadrilatero articolato si pu\`o studiare mediante le nozioni,
esposte in un capitolo precedente, circa i moti relativi. In particolare, ricordando
l'esempio di figura \ref{fig:f114} e la relativa soluzione, possiamo anche in questo caso
esprimere il moto di un membro rispetto a un riferimento fisso e a uno
mobile e imporre che le quantit\`a cinematiche coinvolte, velocit\`a e accelerazione,
corrispondano tra loro.

\noindent Riferiamoci alla rappresentazione schematica di figura
\ref{fig:f_quad_schema}, dove oltre al quadrilatero sono rappresentate velocit\`a
angolare e accelerazioni angolari della manovella $DC$ che si supporranno note.
Se ci proponessimo di trovare le analoghe grandezze cinematiche per il bilanciere $AB$,
potremmo scrivere la seguente relazione che esprime la velocit\`a del punto $C$,
sia in termini assoluti, sia in termini relativi, nel suo moto attorno al
punto $B$. Qui immaginiamo incardinato il sistema di riferimento relativo che
trasla, mantenendo perci\`o i suoi assi sempre nelle medesime direzioni.
Avremo
\begin{equation}
{\bm v}_{\scriptscriptstyle{{\rm ass}}}=
{\bm v}_{\scriptscriptstyle{{\rm rel}}}+
{\bm v}_{\scriptscriptstyle{{\rm tr}}}\,,
\label{eq:e_v_rivals_q}
\end{equation}
\noindent cio\`e
\begin{equation}
{{\bm v}_{\scriptscriptstyle{C}}}=
{{\bm v}_{\scriptscriptstyle{C}}}_{\scriptscriptstyle{{\rm rel}}}+
{{\bm v}_{\scriptscriptstyle{B}}}\,.
\label{eq:e_v_quad_schematico}
\end{equation}
\noindent Riferendoci alle quantit\`a indicate in figura \ref{fig:f_quad_schema},
possiamo riscrivere la \ref{eq:e_v_quad_schematico} nel modo seguente
\begin{equation}
{{\bm v}_{\scriptscriptstyle{C}}}=
\omega_{\scriptscriptstyle{BC}} |\overrightarrow{BC}| \widehat{{\perp{BC}}}
+\omega_{\scriptscriptstyle{AB}} |\overrightarrow{AB}| \widehat{{\perp{AB}}}\,,
\label{eq:e_v_quad_schematico0}
\end{equation}
\noindent dove i versori
$\widehat{{\perp{BC}}}$ e $\widehat{{\perp{AB}}}$ indicano che la direzione dei
due termini incogniti a destra del segno di uguaglianza \`e nota.
L'equazione \ref{eq:e_v_quad_schematico0} contiene pertanto due 
termini vettoriali di modulo incognito e direzione conosciuta.
Si pu\`o dunque seguire (in linea di principio) lo schema grafico riportato in figura
\ref{fig:f116}, e lo svolgimento riportato nel medesimo contesto che ci permette 
di determinare la velocit\`a angolare del bilanciere e della biella,
$\omega_{\scriptscriptstyle{AB}}$ e $\omega_{\scriptscriptstyle{BC}}$.

\noindent Consideriamo ora le accelerazioni. Dato che il sistema di riferimento relativo trasla, scrivendo la \ref{e138} il termine dell'accelerazione di Coriolis sar\`a assente:
\begin{equation}
{{\bm a}_{\scriptscriptstyle{C}}}=
{{\bm a}_{\scriptscriptstyle{C}}}_{\scriptscriptstyle{{\rm rel}}}+
{{\bm a}_{\scriptscriptstyle{B}}}\,.
\label{eq:e_a_quad_schematico}
\end{equation}
\noindent Risulta conveniente scrivere i termini
${{\bm a}_{\scriptscriptstyle{C}}}_{\scriptscriptstyle{{\rm rel}}}$ e 
${{\bm a}_{\scriptscriptstyle{B}}}$ 
come somma delle loro
componenti: normale e tangenziale
\begin{equation}
{{\bm a}_{\scriptscriptstyle{C}}}=
-\omega^2_{\scriptscriptstyle{BC}} \overrightarrow{BC} + {\dot{\omega}_{\scriptscriptstyle{BC}}}|\overrightarrow{BC}|\widehat{{\perp{BC}}}
-\omega^2_{\scriptscriptstyle{AB}} \overrightarrow{AB} + {\dot{\omega}_{\scriptscriptstyle{AB}}}|\overrightarrow{AB}|\widehat{{\perp{AB}}}\,,
\label{eq:e_a_quad_schematico2}
\end{equation}
\noindent dove, ancora una volta, i versori
$\widehat{{\perp{BC}}}$ e $\widehat{{\perp{AB}}}$ indicano che la direzione
di questi due termini incogniti \`e conosciuta.
Gli altri termini risultano essere tutti noti, quindi
ricadiamo ancora nel caso di una equazione vettoriale dove le incognite sono
i moduli del secondo e del quarto termine a destra del segno di uguaglianza della
\ref{eq:e_a_quad_schematico2} e per la soluzione ci riferiamo, una volta di pi\`u,
al procedimento di figura \ref{fig:f116}, dal quale
otterremo anche
${\dot{\omega}}_{\scriptscriptstyle{AB}}$ e ${\dot{\omega}}_{\scriptscriptstyle{BC}}$.

\section{Applicazioni dei Quadrilateri} \label{q_squilibrio}
\noindent Abbiamo gi\`a ricordato che il quadrilatero articolato si presta a
svariatissime applicazioni. Nella sua versione con doppia manovella esso
pu\`o trasmettere il moto rotatorio
tra due alberi paralleli, nella configurazione in cui
le tracce degli assi di tali alberi siano i centri delle
rispettove coppie rotoidali sul telaio. In questo caso,
se il quadrilatero presenta i lati opposti uguali tra loro, la trasmissione
del moto sar\`a {\em omocinetica}\index{trasmissione!omocinetica} con
rapporto di trasmissione unitario e il meccanismo
prende il nome di {\em parallelogramma articolato}\index{parallelogramma articolato},
famoso per aver collegato, un tempo, le ruote delle locomotive a vapore in
modo tale da renderle tutte motrici.
La trasmissione del moto tramite un quadrilatero generico non sar\`a pi\`u,
in generale, omocinetica
e si potranno avere due casi di particolare interesse: manovella-manovella, oppure
manovella-bilanciere. In entrambi i casi, il progetto (sintesi) del quadrilatero
mira a ottenere un particolare movimento rotatorio o oscillatorio 
sulla manovella o sul bilanciere in uscita, tale da rendere diverse tra loro
 le velocit\`a di
percorrenza ciclica delle due fasi che questo membro attraversa durante il moto
a velocit\`a angolare costante della manovella motrice.

\begin{figure}[hbt]
\centering
\begin{minipage}[b]{0.48\textwidth}
\centering
\includegraphics[width=0.9\textwidth]{part2/quadri/FIG/quadri/doppia_manovella.pdf}
\begin{picture}(0,0)(130,0)
\scriptsize{
\put(39,125){$\alpha_2=81^{\circ}$}
\put(8,119){$B$}
\put(88,114){$C$}
\put(26,53){$A$}
\put(55,53){$D$}
\put(85,26){$\alpha_1=333^{\circ}$}
\put(109,97){$\beta_2=34^{\circ}$}
\put(11,37){$\beta_1=214^{\circ}$}
}
\end{picture}
      \caption{\em Quadrilatero a doppia manovella.}
 \label{fig:doppia_manovella}
\end{minipage}\hfill
\begin{minipage}[b]{0.48\textwidth}
\centering
\includegraphics[width=0.9\textwidth]{part2/quadri/FIG/quadri/vel_doppia_manovella.pdf}
\vspace*{6mm}
\begin{picture}(0,0)(129,0)
\scriptsize{
\put(-5,42){${\omega_{\scalebox{.5}{DC}}}/{\omega_{\scalebox{.5}{AB}}}$}
\put(45,60){${\omega_{\scalebox{.5}{BC}}}/{\omega_{\scalebox{.5}{AB}}}$}
\put(124,2){$\alpha$}
}
\end{picture}
\vskip .5mm
      \caption{\em Velocit\`a angolare della manovella $CD$ e della biella $BC$.}
     \label{fig:vel_doppia_manovella}
\end{minipage}
\end{figure}

\noindent La figura \ref{fig:doppia_manovella} riporta un quadrilatero con doppia
manovella: $AB$ e $CD$ compiono intere rotazioni, mentre l'asta $BC$ funge da biella.
In tale figura sono riportate alcune posizioni occupate dalle cerniere $B$ e $C$ 
durante la rotazione a velocit\`a angolare costante della manovella $AB$.
Analizzando tali posizioni si nota che la spaziatura dei punti
generati dalla cerniera $B$, che \`e l'estremo della manovella motrice, \`e appunto
costante mentre la spaziatura delle tracce relative alla cerniera $C$, che \`e l'estremo
della manovella condotta, \`e variabile.
Ci\`o rende perfettamente conto delle velocit\`a angolari delle due manovelle:
come abbiamo detto, costante per quella motrice e variabile per la manovella cedente.
La figura \ref{fig:vel_doppia_manovella} riporta la velocit\`a angolare della manovella
$CD$ e, per completezza-curiosit\`a, quella della biella $BC$, rapportate alla
velocit\`a angolare della manovella motrice. 

\noindent La velocit\`a angolare variabile della manovella $DC$ pu\`o essere a
sua volta utilizzata per azionare un secondo meccanismo come, ad esempio, il 
manovellismo ordinario di una pressa, al quale sarebbe in tal modo possibile
un funzionamento con {\em ritorno rapido}\index{ritorno rapido}.
Consideriamo infatti l'ipotesi di muovere quest'ultimo manovellismo mediante la rotazione della
manovella $DC$ (che fisicamente potrebbe coincidere in parte con quella del manovellismo
stesso). Scegliendo opportunamente la posizione del punto morto di tale manovellismo,
cio\`e il punto di partenza per gli angoli di manovella, potremmo avere, per il
cursore del manovellismo mosso, una corsa di andata pi\`u lenta rispetto a quella di
ritorno o viceversa.
Con riferimento alla figura \ref{fig:doppia_manovella} scegliamo come angolo di
partenza $\beta_1$. Il cursore del manovellismo invertir\`a il suo moto in $\beta_2$,
distante $180^{\circ}$ da $\beta_1$, a valle di una rotazione della
manovella $AB$ di un angolo di ``andata'' pari a
$\alpha_a=\alpha_2-\alpha_1$. Parimenti, il movimento di ritorno si estender\`a
sull'arco percorso sempre di $AB$ pari a
$\alpha_r=360-(\alpha_2-\alpha_1)$.
I tempi di andata e di ritorno, $t_a$ e $t_r$, saranno perci\`o diversi tra loro
e proporzionali ai due angoli percorsi dalla manovella motrice durante le
rispettive fasi:
\begin{equation}
\alpha_a = 360^{\circ} {t_a \over{t_a+t_s}}\,,\;\;\;
\alpha_r = 360^{\circ} {t_s \over{t_a+t_s}}\,.
\label{eq:angoli_andata_ritorno}
\end{equation}
\noindent Il rapporto $s={{\alpha_a}\over{\alpha_r}}$ 
viene talvolta indicato come {\em squilibrio del quadrilatero}\index{quadrilatero!squilibrio}
e, come abbiamo gi\`a accennato, dipende anche dalla scelta della
 posizione iniziale
$\alpha_1$. Lo squilibrio si ritiene tanto pi\`u elevato quanto pi\`u
esso si scosta dall'unit\`a.
In pratica, accade raramente che 
il rapporto tra gli angoli di andata e ritorno, cio\`e lo squilibrio, sia maggiore di
tre o minore di un terzo.
In figura \ref{fig:doppia_manovella} sono indicate le coppie di angoli
$\alpha_1$, $\beta_1$, $\alpha_2$, $\beta_2$ che per questo quadrilatero
realizzano il {\em massimo squilibrio}:
\begin{equation}
s={{333^{\circ}-81^{\circ}}\over{360^{\circ} -333^{\circ} +81^{\circ}}} = 2.3\,.
\label{eq:esempio_squilibrio}
\end{equation}
\noindent Evidenziamo, sperando di contribuire in modo positivo alla chiarezza dell'argomento,
che esiste sempre la
possibilit\`a di scegliere l'inizio delle rotazioni della manovella $DC$, $\beta_1$,
in modo tale da rendere lo squilibrio minimo, cio\`e pari a uno. Con dubbia 
utilit\`a pratica, vedremo in tal caso compiere
alla manovella cedente un mezzo giro dopo che la motrice ha compiuto anch'essa $180^{\circ}$,
rimanendo tuttavia variabile il rapporto di trasmissione tra le due manovelle.
\begin{figure}[hbt]
\centering
\begin{minipage}[b]{0.48\textwidth}
\centering
\includegraphics[width=0.9\textwidth]{part2/quadri/FIG/quadri/manovella_bilanciere.pdf}
\begin{picture}(0,0)(130,0)
\scriptsize{
\put(110,152){$C$}
\put(35,83){$B$}
\put(46,75){$\alpha_1=51^{\circ}$}
\put(8,34){$A$}
\put(111,34){$D$}
\put(4,-6){$\alpha_2=263^{\circ}$}
}
\end{picture}
\vskip 2mm
      \caption{\em Quadrilatero a manovella e bilanciere.}
 \label{fig:manovella_bilanciere}
\end{minipage}\hfill
\begin{minipage}[b]{0.48\textwidth}
\centering
\includegraphics[width=0.9\textwidth]{part2/quadri/FIG/quadri/vel_manovella_bilanciere.pdf}
\vspace*{6mm}
\begin{picture}(0,0)(129,0)
\scriptsize{
\put(35,75){${\omega_{\scalebox{.5}{DC}}}/{\omega_{\scalebox{.5}{AB}}}$}
\put(127,2){$\alpha$}
}
\end{picture}
\vskip 2mm
      \caption{\em Velocit\`a angolare della manovella $CD$.}
     \label{fig:vel_manovella_bilanciere}
\end{minipage}
\end{figure}
\noindent In figura \ref{fig:manovella_bilanciere} \`e rappresentato un quadrilatero il cui
membro $DC$ non compie intere rotazioni ma oscilla percorrendo archi di $54.2^{\circ}$  che, anche in questo esempio, assumono i nomi di andata e ritorno.
Nella stessa figura sono evidenziate, tramite la traccia della cerniera $C$, 
le differenti velocit\`a con le quali vengono descritti tali archi.
Anche qui i tempi di percorrenza
dell'andata e del ritorno, $t_a$ e $t_b$, saranno proporzionali ai corrispondenti
angoli percorsi,
a velocit\`a costante, dalla manovella $AB$. Perci\`o, appoggiandoci 
alle precedenti notazioni,
lo squilibrio di questo quadrilatero diventa
\begin{equation}
s={{263^{\circ}-51^{\circ}}\over{360^{\circ} -263^{\circ} +51^{\circ}}} = 1.43\,.
\label{eq:squilibrio_manovella_bilanciere}
\end{equation}
\noindent In questo caso, gli angoli sottesi dal bilanciere nei quali avviene
l'inversione della sua corsa
risultano  inequivocabili e lo squilibrio massimo,
riportato nella \ref{eq:squilibrio_manovella_bilanciere},
\`e l'unico squilibrio del quale abbia senso parlare.

\noindent Interessanti applicazioni di meccanismi con tempi di andata e ritorno diversi
tra loro si possono realizzare mediante l'impiego di particolari quadrilateri nei
quali una coppia rotoidale ha asse improprio. Tali meccanismi si chiamano {\em manovellismi}\index{manovellismo} e, 
data la loro importanza e diffusione, verranno trattati nel successivo
capitolo, a loro completamente dedicato.

\noindent Il problema di ottenere particolari quadrilateri, a doppia manovella oppure a manovella e bilanciere, che realizzino un dato squilibrio tra i tempi di andata e 
ritorno pu\`o essere risolto tramite specifici metodi
grafici i quali, se coadiuvati da una buona dose di esperienza del progettista,
portano al risultato desiderato.
A tali metodi di sintesi non facciamo cenno, ritenendoli
fuori luogo in questi appunti e ottimamente collocati in libri specialistici come il
pluricitato \cite{ruggieri}, pagg. 215--220. Aggiungiamo che la sintesi dei quadrilateri
avviene molto frequentemente tramite analisi ripetute, partendo da geometrie iniziali
che solitamente il progettista porta con s\'e nel suo bagaglio di esperienza.

\section{Curve di Biella}

\noindent Fin qui, poco si \`e detto del moto della biella eccezion fatta per la 
figura \ref{fig:vel_doppia_manovella} dove \`e riportata la velocit\`a angolare
della biella $BC$ rapportata a quella della manovella $AB$.
D'altra parte lo studio del moto della biella di un quadrilatero articolato ha
applicazioni pratiche molto rilevanti.
Per un dato quadrilatero, i punti del piano su cui giace la
biella (pensato solidale con essa) percorrono, durante il movimento di questa, una doppia infinit\`a di curve.
Al variare del {\em punto di biella}\index{biella!punto di} considerato, tali
{\em curve di biella}\index{biella!curve di} assumono forme diversissime tra loro
e possono risultare appetibili quando appunto servono movimentazioni
cicliche e peculiari.
Le curve di biella presentano infatti una grande
variet\`a di archi e intrecci
e vengono spesso prese in considerazione in diversi contesti
del progetto di macchine automatiche.
In tempi passati, il progettista interessato all'impiego di una curva di biella
si affidava per la sua scelta ai cosiddetti {\em atlanti}\index{atlanti di quadrilateri}.
Tra queste opere, la pi\`u
completa \`e quella di Hrones e Nelson \cite{hrones}, che presenta, tramite
un ingegnoso ordine sistematico, pi\`u di settemila curve di biella.  
\begin{figure}[hbt]
\centering
\includegraphics[width=0.9\textwidth]{part2/quadri/FIG/quadri/curve_di_biella.pdf}
\begin{picture}(0,0)(250,0)
\scriptsize{     
\put(10,352){\rotatebox{-54}{$\longrightarrow$}}
\put(-13,354){punto di biella}
\put(107,298){\rotatebox{-54}{$\longrightarrow$}}
\put(83,301){punto di biella}
\put(125,182){$C$}
\put(22,54){$A$}
\put(-30,54){$B$}
\put(142,54){$D$}
}
\end{picture}
      \caption{\em Alcune curve di biella di un quadrilatero articolato manovella-bilanciere. I punti di biella sono evidenziati.}
 \label{fig:f_cur_biella}
\end{figure}

\noindent Oggigiorno,  non \`e difficile trovare ``pagine di calcolo'' che contengono
analizzatori di 
quadrilateri i quali, a fronte dei dati che identificano il meccanismo da
studiare e della scelta del punto di biella, forniscono immediatamente la 
relativa curva. Il progettista, guidato dall'esperienza e dalle analisi numeriche
ripetute, pu\`o cos\`i
avvicinarsi per passi alla soluzione del proprio problema di progetto.
Naturalmente esistono interi programmi commerciali di calcolo che possono analizzare
i quadrilateri articolati, i quali programmi spesso sono molto generali e pensati
per l'analisi di meccanismi anche molto complessi.
Diciamo per\`o, per esperienza, che un tecnico, qualora si accinga all'utilizzo
delle curve di biella, difficilmente rinuncia a trarre ispirazione dai casi
riportati sugli atlanti, se non altro per generare l'embrione di quello che poi
costituir\`a la soluzione finale. 
\begin{wrapfigure}{r}{0.5\textwidth}
     \begin{center}
     \includegraphics[width=0.42\textwidth]{part2/quadri/FIG/quadri/watt.pdf}
     \end{center}
\begin{picture}(0,0)(0,0)
        \scriptsize{
        \put(77,85){$B$}
        \put(58,18){$B'$}
        \put(90,45){$C$}
        \put(90,20){$C'$}
        \put(89,69){$E$}
        \put(73,19){$E'$}
        \put(155,58){$D$}
        \put(11,73){$A$}
}
\end{picture}
        \caption{\em Quadrilatero di {\em Watt}.}
     \label{fig:watt}
\end{wrapfigure}
\noindent A titolo di esempio riportiamo in figura \ref{fig:f_cur_biella}
le traiettorie di alcuni punti di biella del quadrilatero a manovella-bilanciere
$ABCD$. Anche in questo caso, la distanza
tra le tracce lasciate dal punto di biella d\`a un'idea precisa della 
velocit\`a mediante la quale tale punto percorre la propria traiettoria
(accorgimento grafico, questo,
riportato con precisione sugli atlanti tramite linee a tratteggio di lunghezza
diversificata).
\noindent Non deve sorprendere se particolari curve
di biella, che, da un punto di vista teorico, rettilinee non sono, vengano
invece sfruttate nella pratica come generatori di spostamenti rettilinei
o {\em guide rettilinee
approssimate}\index{guide rettilinee approssimate}.
La figura \ref{fig:watt} mostra una di queste guide approssimata dal
movimento del
punto centrale $E$ della biella $BC$. Il quadrilatero $ABCD$ prende il nome dal suo 
inventore e si chiama quadrilatero di {\em Watt} (talvolta parallelogramma
di {\em Watt}, anche se le applicazioni che lo vedono
in tal veste sono molto rare).
Nella nostra figura, le tracce del punto di biella $E$ non sono equidistanziate,
infatti la figura \`e stata ottenuta
mediante il movimento a velocit\`a costante del bilanciere $AB$;
inoltre, sempre nella figura citata, viene riportato un pezzo
molto esteso della curva che il punto $E$ pu\`o percorrere.
Normalmente per\`o questo meccanismo viene
impiegato sfruttando la guida rettilinea approssimata, fornita dal punto mediano della 
biella, a fronte di un azionamento
direttamente applicato allo stesso punto $E$.
Tale azionamento pu\`o essere l'oscillazione verticale del telaio di un'automobile
quando, ad esempio, il quadrilatero di Watt viene utilizzato nelle sospensioni.
Ma il movimento oscillatorio verticale della scocca di un veicolo
\`e limitato nella propria corsa dai vincoli geometrici della sospensione
stessa, pertanto va da s\'e che il punto di biella percorrer\`a
soltanto una porzione, quella quasi rettilinea, della
 sua traiettoria possibile.
\newpage
\thispagestyle{empty}
\null

\endinput 
\section{Cenni Generali sulla Sintesi dei Quadrilateri}
\noindent Desideriamo concludere questo paragrafo con un brevissimo cenno
riassuntivo alle tecniche di sintesi del quadrilatero articolato.
Per prima cosa distinguiamo da subito due strade possibili: quella {\em diretta}\index{
sintesi diretta dei quadrilateri} e quella {\em indiretta}\index{sintesi
indiretta dei quadrilateri}. La prima di queste strade conduce alla determinazione,
tramite metodi grafici o analitici, dei punti di cerniera del quadrilatero
che soddisfa le richieste di
progetto in modo univoco oppure ad infinit\`a di soluzioni dipendenti da uno o 
pi\`u parametri liberi, come potrebbero essere il fattore di scala o il posizionamento
di una o pi\`u cerniere.
Il metodo indiretto consiste invece nell'analisi numerica ripetuta di diverse
soluzioni che,
guidata dall'intuito e dall'esperienza del progettista, nonch\'e dalle agili
possibilit\`a interattive dei moderni programmi di calcolo, porta via via alla
soluzione ottimale. 
\`E possibile una classificazione sistematica dei vari problemi di sintesi
che possono essere ricondotti a tre problemi paradigmatici.

\noindent Il primo di questi problemi \`e gi\`a stato esposto, con piglio descrittivo,
poco sopra in queste note. Si tratta, con riferimento alle figure
\ref{fig:doppia_manovella} o \ref{fig:manovella_bilanciere}, di correlare le posizioni
angolari della manovella motrice con le rispettive posizioni dell'altra manovella
(oppure bilanciere). Si parla in questi casi di uso del quadrilatero come
{\em generatore di funzioni}\index{generatore di funzioni}. Crediamo sia evidente
che non si possa sperare di ottenere un meccanismo in grado di seguire
pedissequamente una data funzione arbitraria $\beta(\alpha)$. Siccome le incognite,
o gradi di libert\`a,
risultano essere le quattro lunghezze dei membri del quadrilatero \`e ragionevole
aspettarsi che potranno essere soddisfatte puntualmente soltanto quattro valori:
$\beta(\alpha_1)$, $\beta(\alpha_2)$, $\beta(\alpha_3)$ e $\beta(\alpha_4)$.
In questo caso quindi vengono fatte corrispondere cinque posizioni del quadrilatero e 
si parla di problema a dieci gradi di libert\`a.
Anzich\'e essere fatte valere tutte sugli angoli, queste quattro condizioni potrebbero,
in parte, essere  applicate alle
derivate di tali angoli: ad esempio, nella sintesi del quadrilatero di figura
\ref{fig:manovella_bilanciere} dovremmo imporre velocit\`a nulla nelle due posizioni
di {\em punto morto}\index{punto morto} del bilanciere $DC$ e avremo cos\`i assegnate
$\beta(\alpha_1)$, ${\dot\beta}(\alpha_1)=0$, $\beta(\alpha_2)$ e ${\dot\beta}(\alpha_2)=0$. 
\noindent Un altro tipico problema di sintesi pu\`o essere ricondotto alla ricerca 
di traiettorie desiderabili per un dato punto del {\em piano di biella}\index{biella!piano di}, traiettorie
cio\`e percorse da un punto di biella. Ci siamo gi\`a occupati di queste traiettorie,
in modi descrittivo, in un paragrafo precedente e potremmo, in linea con la maggior parte
degli autori che trattano questo argomento, chiamare questa funzione richiesta al
quadrilatero {\em generazione di traiettorie}\index{generazione!di traiettorie}.
Da un punto di vista molto formale si potrebbe asserire che, in questo caso,
le incognite
saranno le quattro coordinate delle cerniere del telaio, le tre lunghezze delle
aste rimanenti, le due lunghezze dei rimanenti lati del triangolo che individua (tramite
uno dei suoi vertici) il punto di biella e infine la posizione iniziale della 
manovella movente: totale 10 gradi di libert\`a. A questi gradi di libert\`a,
sempre molto formalmente, corrisponderanno cinque posizioni (due coordinate ciascuna)
che potranno essere soddisfatte, oppure quattro posizioni e una tangente alla curva
di biella in uno dei quattro punti, oppure tre punti e due tangenti.
 
\noindent Un terzo problema di sintesi riguarda la possibilit\`a di fare guidare alla
biella un oggetto da posizionare nel piano. Questa funzione si pu\`o chiamare 
{\em generazione di moti piani}\index{generazione!di moti piani}. Anche qui, 
\cite{ruggieri}, pag. 279, i gradi di libert\`a sono dieci e si potrebbero soddisfare
cinque posizioni per l'oggetto movimentato, che saturerebbero tutte le incognite.

\noindent I tre problemi di sintesi dimensionale del quadrilatero sono per\`o tutti
riconducibili l'uno all'altro come riportato in \cite{ruggieri} pag 279.
Ci sembra opportuno dare solo un brevissimo cenno alla soluzione al terzo di questi
problemi a 5 gradi di libert\`a (o dieci se si considera che ogni punto \`e individuato
da due coordinate. La generazione di particolari posizioni della biella \`e stato
formulato, in modo assai generale, da Burmester quindi tale problema ne prende il
nome.

\noindent Nella sua formulazione pi\`u generale, il {\em problema di Burmester}
consiste nell'individuare un quadrilatero che,  tramite un punto di biella, tocchi
alcune posizioni precise.
Doverebbe essere chiaro da quanto fin qui esposto che il numero massimo di tali punti
di precisione non pu\`o superare il numero cinque. 
Difficilmente per\`o ci si pu\`o spingere alla saturazione di tutti i gradi di
libert\`a disponibili, in quanto si potrebbero trovare soluzioni non praticabili.
Ci si accontenta spesso di progettare quadrilateri in cui un punto di biella 
passa per tre punti di precisione, lasciando liberi due gradi di libert\`a in modo
da individuare, nel novero delle soluzioni ottenute,
meccanismi effettivamente realizzabili.

\noindent Sebbene gli studi del Burmester siano molto interessanti e
indispensabili laddove venga
richiesto al punto di biella il passaggio da quattro o da cinque punti di
precisione,
facciamo semplicemente notare che la soluzione per tre soli punti di precisione
\`e talmente a portata di mano che sarebbe un peccato non descriverla.

\noindent Sia dunque da realizzare un quadrilatero in cui un punto di biella
passi da tre punti stabiliti. Per tali punti passer\`a una e una sola circonferenza.
Il moto della biella avr\`a quindi come centro di spostamento finito il centro
di tale circonferenza.
(da concludere)


\chapter{Manovellismi}

\section{Introduzione}

\noindent I meccanismi che si identificano col nome di 
manovellismi non sono altro, da un punto di vista cinematico, che
casi particolari di quadrilateri articolati. Essi sarebbero potuti entrare
a pieno titolo nel capitolo precedente, visto che le tecniche per
il loro studio, siano queste grafiche, analitiche o numeriche, non
differiscono da quelle che si  impiegano nell'analisi dei loro
progenitori. Il presente capitolo, dedicato esclusivamente allo studio dei
manovellismi, si giustifica valutando l'importanza industriale che 
tali meccanismi hanno
avuto e quindi anche la profondit\`a e la mole di studi che
li riguarda.
Affermiamo, senza timore di esagerare, che il
{\em manovellismo ordinario}\index{manovellismo!ordinario} o
{\em di spinta}\index{manovellismo!di spinta} \`e stato il cuore della
Rivoluzione Industriale, essendo un componente essenziale del motore a vapore.
Anche ai giorni nostri,  nel settore del trasporto su strada,
tale meccanismo entra, con numero
variabile di esemplari, praticamente in tutti i motori endotermici, tralasciando
pochissime eccezioni.

\noindent Richiamiamo la nomenclatura delle quattro cerniere
di un quadrilatero ricordando la figura \ref{fig:f_quad_schema} ovvero
prendendo visione della figura \ref{fig:quad_D_infinito}.
Ciascuna delle quattro cerniere $A$, $B$, $C$, e $D$ pu\`o essere mandata all'infinito
seguendo una precisa direzione e divenendo una di quelle
entit\`a che la geometria
definisce come {\em punti impropri} del piano.
\`E scontato che, affinch\'e una (alla volta, per il momento)
delle cerniere si trovi
infinitamente distante dalle altre tre, \`e necessario che le due aste
che vi concorrono siano a loro volta di lunghezza infinita.
\noindent Data la ovvia equivalenza cinematica dei punti $A$ e $D$ e dei punti
$B$ e $C$ si pu\`o intuire che pensare situato all'infinito l'uno o l'altro dei punti di
tali coppie produca, a livello cinematico, lo stesso effetto,
e cos\`i \`e.

\noindent Affermiamo che allo spostarsi all'infinito delle cerniere
$A$ oppure $D$ il quadrilatero d\`a origine ai {\em manovellismi ordinari}\index{manovellismo!ordinario}, naturalmente dopo avere mutato la cerniera 
``fuori uso'' con un altro tipo di vincolo, in generale costituito da un 
{\em cursore}\index{cursore} che scorre in un {\em glifo}\index{glifo}.

\noindent Nel secondo caso, quando ad andare all'infinito \`e una delle
cerniere $B$ oppure $C$, avremo i {\em manovellismi
non ordinari}.\index{manovellismo!non ordinario}
\noindent Andiamo con ordine.


\section{Il Manovellismo Ordinario}

\begin{figure}[hbt]
\centering
\begin{minipage}[b]{0.48\textwidth}
\centering
\includegraphics[width=0.95\textwidth]{part2/manovellismi/FIG/manovellismi/quad_D_infinito.pdf}
\begin{picture}(0,0)(155,-12)
\scriptsize{
\put(17,53){$A$}
\put(44,85){$B$}
\put(140,56){$C$}
\put(73,-3){$D$}
}
\end{picture}
\vskip .2mm
      \caption{\em Generico quadrilatero manovella-bilanciere.}
 \label{fig:quad_D_infinito}
\end{minipage}\hfill
\begin{minipage}[b]{0.48\textwidth}
\centering
\includegraphics[width=0.95\textwidth]{part2/manovellismi/FIG/manovellismi/quad_D_infinito_man.pdf}
\begin{picture}(0,0)(155,3)
	\scriptsize{
\put(16,53){$A$}
\put(42,85){$B$}
\put(132,60){$C$}
\put(87,58){$\sigma$}
\put(93,3){$D$}
\put(100,10){\rotatebox{-90}{$\longrightarrow$}}
\put(104,3){$\infty$}
}
\end{picture}
	\caption{\em Manovellismo ordinario centrato ottenuto spostando la cerniera $D$ a distanza ``infinita''.}
     \label{fig:quad_D_infinito_man}
\end{minipage}
\end{figure}

\noindent In figura \ref{fig:quad_D_infinito_man} \`e
riportato il quadrilatero degenere che si ottiene mandando all'infinito la
cerniera $D$, una delle due vincolate al telaio, del quadrilatero di
figura \ref{fig:quad_D_infinito}.
La direzione secondo la 
quale la cerniera $D$ diventa un punto improprio \`e, in questo caso, quella 
ortogonale al segmento $AC$: otteniamo in questo modo un {\em manovellismo
ordinario centrato}\index{manovellismo!ordinario}. Si ottiene
invece un
{\em manovellismo ordinario deviato}\index{manovellismo!ordinario deviato}
scegliendo, per la cerniera $D$, una direzione diversa\footnote{
\noindent Dal punto di vista cinematico, si ottiene un meccanismo
perfettamente equivalente

\begin{wrapfigure}{r}{0.33\textwidth}
     \begin{center}
     \includegraphics[width=0.32\textwidth]{part2/manovellismi/FIG/capsulismo.pdf}
     \end{center}
\begin{picture}(0,0)(0,0)
	\scriptsize{
\put(118,81){$C$}
\put(29,97){$B$}
\put(59,35){$D$}
\put(-3,90){$A$}
\put(-5,92){\rotatebox{-159}{$\longrightarrow$}}
\put(-3,80){$\infty$}
}
\end{picture}
\vskip -5.3mm
        \caption{\em Manovellismo ordinario centrato a glifo mobile.}
     \label{fig:capsulismo}
\end{wrapfigure}

\noindent a quello di figura \ref{fig:quad_D_infinito_man} mandando all'infinito
la cerniera $A$ anzich\'e la $D$.
\noindent Per dare un po' di vitalit\`a a questa nota, 
scegliamo in questo caso di rendere mobile il glifo e tenere ferma
la cerniera $B$.
Il dispositivo pratico che ne risulta \`e piuttosto
diverso nei due frangenti.
Nella stragrande maggioranza dei casi,
il manovellismo ordinario centrato viene progettato e realizzato
tenendo fisso il glifo.
Nell'altro caso si ottiene invece il
sistema articolato che \cite{sesini1} riporta a pag. 107, Fig. IV-4 (e che
noi riportiamo in figura \ref{fig:capsulismo}), 
dichiarando che ``tale meccanismo non trova notevoli applicazioni''; anche
all'autore
non risultano applicazioni rilevanti di questo
manovellismo, neppure ai nostri giorni. La figura riportata su queste note
appare per\`o piuttosto diversa
da quella che si trova in \cite{sesini1}. Questo perch\'e
abbiamo preferito privilegiare, in questo caso, la coerenza con la precedente
figura \ref{fig:quad_D_infinito}, mantenendo perci\`o gli stessi rapporti
dimensionali tra i membri del quadrilatero e gli stessi orientamenti
relativi,
piuttosto che garantire alla manovella $CD$ di potere compiere intere
rotazioni.}.
\noindent Il percorso della cerniera $C$ diventa il segmento $\sigma$,
che pu\`o essere una scanalatura (glifo) in cui scorre un cursore di forma
opportuna. Ad esempio, nei motori endotermici, i cilindri costituiscono i
glifi e gli stantuffi hanno la funzione di cursore. \`E per\`o chiaro che l'elemento che funge
da guida cava pu\`o essere invertito con l'elemento che funge da cursore,
trasformando quest'ultimo in un glifo, e
questa opportunit\`a si presenta per tutti i manovellismi, ordinari e non.
\noindent In ambito industriale il manovellismo si studia
e si progetta, come \`e ovvio, tramite l'ausilio dei programmi di calcolo 
per l'analisi di sistemi articolati.
\noindent Tuttavia il manovellismo ordinario centrato \`e talmente diffuso,
specialmente in ambito {\em automotive} (ogni vettura  a combustione
interna ne ospita in media quattro) e gli
studi che lo hanno riguardato all'inizio del secolo scorso sono cos\`i rilevanti
e fecondi di risultati, da imporre alla coscienza 
di chi insegna Meccanica Applicata l'esposizione, magari in breve,
di quella che si pu\`o chiamare ``trattazione analitica del manovellismo'';
e cos\`i faremo noi.
\begin{figure}[hbt]
     \begin{center}
     \includegraphics[width=0.85\textwidth]{part2/manovellismi/FIG/man_ord.pdf}
\begin{picture}(0,0)(270,-3)
        \scriptsize{
        \put(38,109){$B$}
        \put(53.4,74.3){$\alpha$}
        \put(187,57){$C$}
        \put(5,59){$A$}
\put(44,66.6){\line(1,0){50}}
        \put(50,73.6){\rotatebox{-67}{$\curvearrowleft$}}
}
\end{picture}
\vskip -1mm
        \caption{\em Manovellismo ordinario centrato.}
     \label{fig:man_ord}
	\end{center}
\vskip -3mm
\end{figure}
\noindent Cominciamo col dare uno sguardo al moto del {\em piede di biella}\index{piede di biella} $C$, senza pretese di esattezza, giusto per renderci conto in modo 
qualitativo del suo movimento alternativo. Riportiamo a tal fine
di nuovo il manovellismo di figura
\ref{fig:quad_D_infinito_man} nella nuova figura
\ref{fig:man_ord},
e le quantit\`a cinematiche del punto $C$, ottenute da un'analisi
numerica, in figura
\ref{fig:vel_manovellismo_ordinario_centrato}.
\begin{figure}[hbt]
\begin{center}
    \includegraphics[width=0.55\textwidth]{part2/manovellismi/FIG/manovellismi/vel_manovellismo.pdf}
\begin{picture}(0,0)(170,25)
        \scriptsize{
        \put(80,118){$v_{\scriptscriptstyle C}$}
        \put(120,50){$a_{ \scriptscriptstyle C}$}
        \put(165,25){$\alpha$}
}
\end{picture}
        \caption{\em Velocit\`a e accelerazione nel manovellismo ordinario centrato.}
     \label{fig:vel_manovellismo_ordinario_centrato}
\end{center}
\vskip -3mm
\end{figure}
\noindent  In figura \ref{fig:man_ord}
si possono vedere le tracce
del {\em bottone di manovella}\index{bottone di manovella} $B$
e del piede di biella $C$,
mentre in figura  \ref{fig:vel_manovellismo_ordinario_centrato} sono riportati
i grafici di velocit\`a e accelerazione del punto $C$,
assumendo per la  velocit\`a angolare della manovella $AB$, di
lunghezza unitaria,
$\omega_{\scriptscriptstyle AB}=1$.
Notiamo subito, sia esaminando le tracce del piede di biella,
sia dal grafico della velocit\`a, che il punto $C$ deve rallentare 
nei punti morti (manovella con $\alpha=0$ e $\alpha=180$) fino a
fermarsi: presupposto ovvio affinch\'e il moto di $C$ si possa invertire.

\section{Espressioni Analitiche di Velocit\`a e Accelerazione del Piede
di Biella}

\noindent Con riferimento alla figura \ref{fig:moc_schematico} possiamo
innanzitutto considerare che la distanza $\overline{AC}$ \`e pari a

\begin{equation}
\overline{AC} = r \cos (\alpha) + l \cos (\beta)\,.
\end{equation}

\begin{wrapfigure}{r}{0.5\textwidth}
     \begin{center}
     \includegraphics[width=0.48\textwidth]{part2/manovellismi/FIG/moc_schematico.pdf}
     \end{center}
\begin{picture}(0,0)(32,3)
	\scriptsize{
        \put(64,97){$B$}
        \put(68,64){$\alpha$}
        \put(54,80){$r$}
        \put(120,78){$l$}
        \put(122,64){$\beta$}
        \put(164,65){$C$}
        \put(168,24){$x$}
        \put(168,16){$\dot x$}
        \put(168,8){$\ddot x$}
        \put(34,59){$A$}
\put(54,60){\line(1,0){30}}
\put(115,60){\line(1,0){40}}
}
\end{picture}
\vskip -3.95mm
        \caption{\em Rappresentazione schematica del manovellismo ordinario centrato.}
     \label{fig:moc_schematico}
\end{wrapfigure}

\noindent \`E facile perci\`o scrivere per la coordinata $x$, che indica lo
spostamento del piede di biella ($C$) dal punto morto esterno, 
\begin{equation}
x= r+l - r \cos (\alpha) - l \cos (\beta)\,.
\label{eq:formula_manovellismo}
\end{equation}
\noindent \`E opportuno (e tradizionale) introdurre a questo punto un parametro
che caratterizza fortemente il manovellismo: si tratta del rapporto tra il raggio
di manovella e la lunghezza della biella:
\begin{equation}
\lambda= \frac{r}{l}\,.
\end{equation}
\noindent Bench\'e questo rapporto possa, in teoria, spaziare in un intervallo
amplissimo, esso \`e soggetto a restrizioni quando si desidera una manovella 
libera di compiere intere rotazioni. In tal caso $\lambda< 0.28$\footnote
{Valore di lambda per i manovellismi del motore {\em Alfa Romeo} 1600 cc della
{\em Giulia GTA Sprint}.}\footnote{Al di l\`a degli ovvi problemi che 
nascerebbero aumentando $\lambda$, problemi principalmente dovuti 
alle possibili interferenze tra organi in moto alterno, di solito gli stantuffi,
il rapporto $\lambda$ si tiene comunque piuttosto distante dal limite teorico
$\lambda=0.5$, cio\`e da quello che vedrebbe la biella di lunghezza doppia
della manovella e che consentirebbe la completa rotazione di quest'ultima.
Dalla figura \ref{fig:moc_schematico} risulta $\lambda={\sin
\beta_{\scriptscriptstyle{(\alpha= 90^{\circ})}}}$. Pertanto, nelle posizioni
$\alpha=90^{\circ}$ e $\alpha=270^{\circ}$ le forze che agiscono sullo stantuffo
solleciteranno la biella
moltiplicate per $\frac{1}{\cos \beta_{\scriptscriptstyle {\rm max}}}=\frac{1}{\sqrt{1-\lambda^2}}$; \`e evidente che valori
elevati di $\lambda$ renderebbero
la trasmissione del moto inefficiente, dando luogo a ingenti spinte laterali sul glifo
e pregiudicando dal punto di vista strutturale
gli organi meccanici coinvolti.}
risulta essere un limite ragionevole per i manovellismi ordinari centrati impiegati nei motori a combustione interna e nei compressori.
Tale limitazione non vale per manovellismi dove la manovella si comporta in
realt\`a da bilanciere, come accade nella maggior parte delle presse a
ginocchiera, dove $\lambda$ pu\`o persino superare l'unit\`a.
Osserviamo che $r \sin \alpha = l \sin\beta$, e che
$\cos \beta=\sqrt{1-\lambda^2 \sin^2\alpha}$.
Tramite queste relazioni la \ref{eq:formula_manovellismo} diventa 
\begin{equation}
x= r[1 - \cos \alpha + \frac{1}{\lambda} (1 -  \sqrt{1-\lambda^2 \sin^2\alpha})]\,.
\label{eq:formula_manovellismo_1}
\end{equation}
\noindent Derivare rispetto al tempo una prima, quindi una seconda volta
la \ref{eq:formula_manovellismo_1},
allo scopo di ottenere velocit\`a e accelerazione del piede di biella,
non presenta certo difficolt\`a. Tale calcolo per\`o somiglierebbe
a un esercizio
meramente didattico di studio ``a livello liceale'' delle operazioni di
derivazione di semplici funzioni.
Come \`e noto, infatti, le derivate di radicali portano radicali, e 
le funzioni $\dot x$ e $\ddot x$, ottenute per questa via, dovrebbero
essere graficate con grande sforzo di calcolo,
oppure per via numerica, circostanza quest'ultima che vanificherebbe lo spirito
di questo percorso.


\noindent Si preferisce pertanto partire da una formula pi\`u semplice,
che approssimi lo spostamento del piede di biella, e ricavare da
quest'ultima le quantit\`a cinematiche derivate.
La strada \`e quella di esprimere il radicando della \ref{eq:formula_manovellismo_1} tramite i primi due termini del suo sviluppo in
serie, sviluppo eseguito considerando $\lambda$
(e non gi\`a $\alpha$) come variabile indipendente:
\begin{equation}
\sqrt{1-\lambda^2 \sin^2\alpha}\simeq 1 - \frac{1}{2} \lambda^2\sin^2\alpha \,.\notag
%\label{eq:svil_sqrt}
\end{equation}
\noindent Otteniamo cos\`i l'espressione arrestata al secondo ordine per lo
spostamento del piede di biella,
che risulter\`a tanto pi\`u  approssimato quanto pi\`u piccolo sar\`a
il rapporto $\lambda$:
\begin{equation}
x= r[1 - \cos \alpha + \frac{1}{2} \lambda \sin^2\alpha]\,;
\label{eq:formula_manovellismo_approx0}
\end{equation}
\noindent essendo poi
\begin{equation}
\sin^2\alpha= \frac{1 - \cos 2\alpha}{2}\,,
\label{eq:formula_manovellismo_approx}
\end{equation}
\noindent otteniamo finalmente   
\begin{equation}
x= r[1 +\frac{\lambda}{4} - \cos \alpha - \frac{\lambda}{4} \cos 2\alpha]\,.
\label{eq:formula_manovellismo_approx_1}
\end{equation}
\noindent La velocit\`a ``approssimata al second'ordine'' si
trova derivando rispetto al tempo la \ref{eq:formula_manovellismo_approx_1},
tenendo presente che ${\rm d}\alpha/{\rm d}t=\omega$:
\begin{equation}
\dot x= r\omega[\sin \alpha +\frac{\lambda}{2} \sin 2\alpha]\,,
\label{eq:formula_manovellismo_approx_vel}
\end{equation}
\noindent mentre per l'accelerazione avremo
\begin{equation}
\ddot x= r\omega^2[\cos \alpha +\lambda \cos 2\alpha]\,.
\label{eq:formula_manovellismo_approx_acc}
\end{equation}
\noindent Le espressioni \ref{eq:formula_manovellismo_approx_vel} e \ref{eq:formula_manovellismo_approx_acc}
mostrano chiaramente un andamento delle due grandezze cinematiche
composto da due armoniche.


\noindent Tramite le seguenti due osservazioni,
proviamo a evidenziare il lato a nostro avviso
pi\`u significativo
di questo studio analitico del manovellismo ordinario.
In primo luogo notiamo che per $\lambda$ molto
piccoli (biella molto lunga in rapporto al raggio di manovella), il moto del
piede di biella risulta coincidere sostanzialmente col moto
armonico:
\begin{equation}
\begin{split}
r-x=r\cos \alpha\,,\;\;\; \dot x=\omega r \sin\alpha\,\;\; {\rm e}\;\;\;\ddot x= \omega^2 r \cos\alpha\,. 
\end{split}
\label{eq:formule_primo_o}
\end{equation}
\begin{wrapfigure}{r}{0.5\textwidth}
     \begin{center}
     \includegraphics[width=0.48\textwidth]{part2/manovellismi/FIG/manovellismi/vel_manovellismo_x.pdf}
     \end{center}
\begin{picture}(0,0)(0,0)
	\scriptsize{
	\put(66,100){${\ddot x}(x)$}
	\put(80,168.5){${\dot x}(x)$}
}
\end{picture}
\vskip -5mm
        \caption{\em Velocit\`a e accelerazione del piede di biella in funzione del suo spostamento normalizzato per $r$ ($1^o$ ordine tratteggiato).}
     \label{fig:vel_manovellismo_x}
\end{wrapfigure}

\noindent \`E interessante esprimere
velocit\`a e accelerazione del piede di biella 
in funzione della coordinata $x$
anzich\'e dell'angolo $\alpha$, figura \ref{fig:vel_manovellismo_x}.
In tal caso, posto $\omega = 1$, quadrando e sommando membro a membro le formule che esprimono
velocit\`a e spostamento approssimati al primo ordine abbiamo
\begin{equation}
(r-x)^2 + {\dot x}^2=r^2\,.
\label{eq:cerchio_vel}
\end{equation}
\noindent La \ref{eq:cerchio_vel} rappresenta l'equazione di una circonferenza
di raggio $r$ e centrata a una distanza $r$ dal punto morto esterno: tale
risulter\`a quindi  l'andamento della
velocit\`a, graficata prendendo $x/r$ come ascissa.
Questo andamento verr\`a ovviamente modificato 
da valori di 
$\lambda$ non  trascurabili, quando cio\`e il moto del piede di 
biella si scoster\`a sensibilmente dal moto armonico.
La figura \ref{fig:vel_manovellismo_x} mostra l'andamento effettivo di
$\dot x(x)$, che tiene conto anche delle successive armoniche,
per $\lambda=0.28$.

\noindent La seconda osservazione riguarda l'accelerazione.
Il grafico di $\ddot x(x)$ \`e anch'esso riportato nella figura appena
citata e ha una forma che somiglia a quella di una parabola.
Fermandoci per\`o al primo ordine,
tale grafico sarebbe una retta. Infatti dalle \ref{eq:formule_primo_o}
otteniamo subito $\cos \alpha=(r-x)/r$ e poi $\ddot x = \omega^2(r-x)$.
Ponendo quindi come al solito $\omega=1$, il grafico di $\ddot x(x)$ sar\`a
una retta a quarantacinque gradi, passante dallo zero a distanza $r$ dal
punto morto esterno. In prima approssimazione, possiamo pertanto
attribuire al piede di biella una accelerazione che varia linearmente
tra i due punti morti: considerazione questa che porta con s\'e una decisa
facilit\`a di calcolo.
Un'ultima osservazione riguarda la possibilit\`a di calcolare
i due valori di accelerazione nei punti morti tramite la formula 
\ref{eq:formula_manovellismo_approx_acc}: otteniamo
$\ddot x(0^{\circ})= \ddot x(0)=r \omega^2(1+\lambda)$  e
$\ddot x(180^{\circ})= \ddot x(2r)=r \omega^2(1-\lambda)$. Lasciamo al 
paziente e volenteroso lettore il compito di mostrare che 
questi due valori
dell'accelerazione del piede di biella nei punti morti sono esatti:
essi non
risentono cio\`e dell'approssimazione introdotta nello spostamento dalla
\ref{eq:formula_manovellismo_approx_1}.

\section{Velocit\`a e Accelerazione del Piede di Biella: Analisi
mediante i Moti Relativi}

\noindent Come nello studio del quadrilatero, \`e possibile, anche in questo caso,
esprimere le quantit\`a cinematiche dei diversi punti del manovellismo
mediante le relazioni che legano queste grandezze quando
si considerano i moti relativi che avvengono tra i vari membri.

\begin{wrapfigure}{r}{0.5\textwidth}
     \begin{center}
     \includegraphics[width=0.48\textwidth]{part2/manovellismi/FIG/moc_schematico_graf.pdf}
     \end{center}
\begin{picture}(0,0)(32,7)
	\scriptsize{
        \put(58,94){$B$}
        \put(43,85){${\omega_{\scalebox{.5}{AB}}}$}
        \put(55,51){${\dot\omega_{\scalebox{.5}{AB}}}$}
        \put(102,91){${\omega_{\scalebox{.5}{BC}}}$}
        \put(124,84){${\dot\omega_{\scalebox{.5}{BC}}}$}
        \put(136,58){${{\bm v}_{\scalebox{.5}{C}}}$}
        \put(164,65){$C$}
        \put(34,59){$A$}
}
\end{picture}
\vskip -4mm
        \caption{\em Velocit\`a e accelerazioni nei moti relativi del manovellismo.}
     \label{fig:moc_schematico_graf}
\end{wrapfigure}

\noindent Con riferimento alla figura \ref{fig:moc_schematico_graf},
posizioniamo un sistema di riferimento traslante centrato nel
{\em bottone di manovella}\index{bottone di manovella} $B$. 
Possiamo cos\`i esprimere la velocit\`a del piede di biella come somma delle
due componenti:
quella relativa
e quella di trascinamento
\begin{equation}
{{\bm v}_{\scriptscriptstyle{C}}}= 
{{\bm v}_{\scriptscriptstyle{C}}}_{\scriptscriptstyle{{\rm rel}}}+
{{\bm v}_{\scriptscriptstyle{B}}}\,,
\label{eq:man_vel_graf0}
\end{equation}
\noindent ovvero esprimendo i due addendi di destra tramite le loro relazioni con
le velocit\`a angolari di biella e manovella,
\begin{equation}
{{\bm v}_{\scriptscriptstyle{C}}}= 
{\omega_{\scriptscriptstyle{BC}}}|\overrightarrow{BC}|\widehat{{\perp{BC}}}+
{{\bm v}_{\scriptscriptstyle{B}}}\,.
\label{eq:man_vel_graf1}
\end{equation}
\noindent La \ref{eq:man_vel_graf1} \`e un'equazione vettoriale nella quale
la direzione di tutti e tre i termini \`e nota
mentre $|{{\bm v}_{\scriptscriptstyle{C}}}|$ e 
${\omega_{\scriptscriptstyle{BC}}}$ risulteranno incogniti, supponendo data
la velocit\`a angolare della manovella. Si ripropone quindi
lo schema risolutivo illustrato in figura \ref{fig:f116},
dal quale ricaveremo
il valore della velocit\`a del piede di biella e la velocit\`a angolare
della biella.
Per quanto riguarda le accelerazioni possiamo scrivere
\begin{equation}
{{\bm a}_{\scriptscriptstyle{C}}}= 
{{\bm a}_{\scriptscriptstyle{C}}}_{\scriptscriptstyle{{\rm rel}}}+ 
{{\bm a}_{\scriptscriptstyle{B}}}\,. 
\label{eq:acc_man_graf0}
\end{equation}
\noindent Come abbiamo gi\`a chiarito altrove, conviene scrivere il termine
${{\bm a}_{\scriptscriptstyle{C}}}_{\scriptscriptstyle{{\rm rel}}}$
 come somma delle sue due
componenti: quella avente direzione parallela alla biella e diretta verso $B$
(componente normale) e quella ortogonale alla biella stessa.
Anche il termine ${{\bm a}_{\scriptscriptstyle{B}}}$ si pu\`o scrivere 
mediante la somma delle sue componenti,
quella radiale e quella tangenziale, le quali
risulteranno note
supponendo di conoscere la velocit\`a e la accelerazione angolari della
manovella.  Otteniamo in questo modo
\begin{equation}
{{\bm a}_{\scriptscriptstyle{C}}}= 
-{\omega_{\scriptscriptstyle{BC}}}^2 \overrightarrow{BC}+
{\dot\omega_{\scriptscriptstyle{BC}}}|\overrightarrow{BC}|\widehat{{\perp{BC}}}
-{\omega_{\scriptscriptstyle{AB}}}^2 \overrightarrow{AB}+
{\dot\omega_{\scriptscriptstyle{AB}}}|\overrightarrow{AB}|\widehat{{\perp{AB}}}\,.
\label{eq:acc_man_graf1}
\end{equation}
\noindent Ancora una volta siamo di fronte a una equazione vettoriale
dove le incognite sono
il modulo dell'accelerazione del piede di biella
$|{{\bm a}_{\scriptscriptstyle{C}}}|$ e l'accelerazione angolare della biella
${\dot\omega_{\scriptscriptstyle{BC}}}$, essendo note le direzioni di
tutti i vettori.
Una possibile soluzione \`e quindi, anche qui, 
lo schema grafico di figura \ref{fig:f116}; per inciso, la \ref{eq:acc_man_graf1} \`e formalmente identica alla \ref{eq:e_a_quad_schematico2},
la qual cosa non ci stupisce visto che stiamo di nuovo trattando
la cinematica di un quadrilatero, sia pure degenere.
Riteniamo di dover ribadire che il ruolo dei termini noti
e di quelli incogniti possono essere scambiati a discrezione: se considerassimo
assegnate la velocit\`a e l'accelerazione del piede di biella potremmo ricavare,
dalle \ref{eq:man_vel_graf1} e \ref{eq:acc_man_graf1},
le quantit\`a {${\omega_{\scalebox{.5}{AB}}}$} e
{${\dot\omega_{\scalebox{.5}{AB}}}$} della manovella $AB$.

\section{Manovellismi Non Ordinari}

\noindent Come \`e stato accennato nel paragrafo introduttivo, se ad andare
all'infinito \`e una delle due cerniere mobili di un quadrilatero si ottengono,
sostituendo opportunamente le aste e le cerniere degeneri con
membri appropriati,
i {\em manovellismi non ordinari}\index{manovellismo!non ordinario}. Questi manovellismi, ovviamente {\em cinematicamente
equivalenti}\index{equivalenza cinematica} ai quadrilateri degeneri che li 
originano,
vengono spessissimo chiamati, con ragione,
meccanismi 
{\em a glifo rotante}\index{glifo!rotante} ovvero
{\em a glifo oscillante}\index{glifo!oscillante},
dando per scontato che un preciso membro di tali meccanismi
sia un glifo. Non ci stanchiamo
di ripetere che, a rigore, i due membri dell'accoppiamento glifo-cursore
possono sempre essere scambiati tra loro senza alcuna mutazione
della cinematica del sistema.
Prima di entrare nel vivo di questo paragrafo, desideriamo esporre una breve
precisazione, che per la verit\`a avrebbe potuto precedere anche la
trattazione dei manovellismi ordinari, ma che si rende ancora
pi\`u necessaria in questa sede: cosa
intendiamo, esattamente, con la locuzione
{\em equivalenza cinematica}\index{equivalenza
cinematica}? Ancora pi\`u esplicitamente, quando asseriamo che la cerniera
$D$ e il bilanciere $DC$ del quadrilatero
degenere di figura \ref{fig:quad_D_infinito_man} vengono sostituiti
dalla coppia glifo-cursore $\sigma$, ottenendo in tal modo il
manovellismo ordinario
equivalente, di quale equivalenza parliamo?
Nello studio dei manovellismi si parla di meccanismo equivalente
quando la biella (o un corpo che la sostituisce) sar\`a interessata dallo
stesso moto piano, mantenendo per tutti i punti di biella le stesse
traiettorie, le stesse velocit\`a e le stesse accelerazioni. Nel
caso che abbiamo studiato dei manovellismi ordinari l'equivalenza
del comportamento della biella di figura \ref{fig:quad_D_infinito_man} con
la cerniera $D$ a grande distanza o con la coppia glifo-cursore $\sigma$ 
sembra all'autore sufficientemente evidente.
Qualche attenzione in pi\`u \`e richiesta
nello studio dei manovellismi non ordinari.

\noindent Giusto per fissare qualche idea circa il comportamento di 
un quadrilatero che degenera tramite l'allontanamento di una delle
due cerniere (perfettamente equivalenti) $B$ oppure $C$, eseguiamo un tentativo
di analisi qualitativa tramite i nostri consueti (e rudimentali)
strumenti di indagine numerica.
\begin{figure}[hbt]
\centering
\begin{minipage}[b]{0.52\textwidth}
\centering
\includegraphics[width=0.9\textwidth]{part2/manovellismi/FIG/manovellismi/c_infinito.pdf}
\begin{picture}(0,0)(138,-6)
\scriptsize{
\put(11,65){$B$}
\put(110,107){$C$}
\put(110,100){\rotatebox{42.45}{$\longrightarrow$}}
\put(113,98){$\infty$}
\put(2,32){$A$}
\put(50,23){$D$}
\put(52.5,37.5){$\alpha$}
\put(78,60){$\beta$}
\put(101,71){$\gamma$}
\color{red}
\put(99.3,83.3){\rotatebox{-47.55}{\line(1,0){15}}}
\put(73,60){\rotatebox{-47.55}{\line(1,0){15}}}
\put(48,38){\rotatebox{-47.55}{\line(1,0){15}}}
}
\end{picture}
      \caption{\em Quadrilatero degenere ottenuto mandando la cerniera $C$ a distanza infinita.}
 \label{fig:c_infinito}
\end{minipage}\hfill\hfill
\begin{minipage}[b]{0.40\textwidth}
\centering
\includegraphics[width=0.9\textwidth]{part2/manovellismi/FIG/glifi_c_infinito.pdf}
\vspace*{10mm}
\begin{picture}(0,0)(139,18)
\scriptsize{
\put(13,65){$B$}
\put(110,107){$C$}
\put(110,100){\rotatebox{42.45}{$\longrightarrow$}}
\put(113,98){$\infty$}
\put(0,32){$A$}
\put(54,14){$D$}
\put(50,44){$a)$}
\put(74,65){$b)$}
\put(100,91){$c)$}
}
\end{picture}
      \caption{\em Possibili meccanismi a glifo mobile derivati dalla figura a lato.}
     \label{fig:glifi_c_infinito}
\end{minipage}
\end{figure}
\noindent Abbiamo scelto, con criterio del tutto arbitrario, di rendere
improprio l'asse della cerniera $C$. Come nel caso dei manovellismi ordinari,
la scelta della direzione secondo la quale mandare tale cerniera 
a grande distanza non
\`e indifferente. Quando la cerniera, come in questo caso, \`e mobile
 tale direzione muta durante il 
movimento del meccanismo. Notiamo per\`o che, scelta in un dato istante
la direzione del punto improprio $C$, essa rimarr\`a, durante il movimento,
costante rispetto alla direzione del segmento $BD$;
in altre parole l'angolo tra la direzione del punto improprio
e $BD$ si manterr\`a costante.
Ci\`o accade in quanto il triangolo $\triangle{BCD}$ presenta angolo costante
e nullo in $C$. Ritenendo quindi le lunghezze dei lati
$BC$ e $CD$ costanti, anche se infinite, rimarranno invariati durante
il movimento anche gli altri suoi angoli\footnote
{
La lunghezza del lato $BD$, ovviamente, non si mantiene costante. Essa
per\`o risulta infinitesimale rispetto alle lunghezze degli altri
due lati, pertanto la sua variazione provoca alterazioni infinitesimali
e trascurabili del triangolo $\triangle{BCD}$.
}.
Nel caso delle figure \ref{fig:c_infinito} e \ref{fig:glifi_c_infinito}
la direzione scelta coincide con la perpendicolare al segmento $BD$,
scelta che consente di ottenere {\em manovellismi non ordinari
centrati}\index{manovellismo!non ordinario}.
\noindent In figura \ref{fig:c_infinito} sono riportate tre curve di biella
relative a tre punti che giacciono sul membro $DC$.
Lo scopo di tale figura \`e quello 
di rendere maggiormente evidente il moto del piano di
biella nei pressi del bilanciere  ``improprio''.
Da questa analisi si pu\`o notare che nei punti di intersezione con 
il bilanciere $DC$, le curve di biella sono 
ortogonali ad esso. La figura, rimarcando questa circostanza, facilita
l'intuizione. Tali direzioni
erano peraltro prevedibili a valle di semplici ragionamenti geometrici, osservando che 
tutti i punti di biella che giacciono su $DC$ descrivono, nel loro
moto relativo a tale membro,  delle circonferenze di centro $C$,
le cui tracce sono in figura rappresentate da brevi segmenti:
$\alpha$, $\beta$ e $\gamma$.
La condizione poi che il punto $C$ si trovi a grande distanza ci
autorizza, come abbiamo appena sopra chiarito, a ritenere costante
l'angolo compreso tra $DC$ e il segmento $DB$,
pari nel nostro caso a $90^{\circ}$. 
Queste considerazioni ci consentono di sostituire la cerniera in $C$ con
gli accoppiamenti prismatici che danno origine alla figura
\ref{fig:glifi_c_infinito}.
\noindent I tre meccanismi a), b) e c) sono, da un punto
di vista cinematico, un solo meccanismo, essendo il piano di biella lo stesso
per le tre soluzioni e rendendo ridondante la distinzione della cerniera
in $B$ in tre cerniere separate.
I tre sistemi articolati   
presentano un glifo sul bilanciere $CD$ dove scorre la relativa biella.
Oltre a ripetere (lo abbiamo detto gi\`a troppe volte) che il ruolo
di chi porta la
scanalatura (glifo) e di chi invece fa da cursore possono indifferentemente
essere scambiati, osserviamo che tali
accoppiamenti prismatici possono essere dislocati in qualsiasi punto della
biella equivalente. 
\noindent In particolare, considerando il meccanismo a) di figura \ref{fig:glifi_c_infinito}, che riteniamo maggiormente significativo,
possiamo ottenere i sistemi articolati equivalenti
 di figura \ref{fig:c_infinito_equivalenti}.

\begin{figure}[hbt]
\centering
\includegraphics[width=0.9\textwidth]{part2/manovellismi/FIG/c_infinito_equivalenti.pdf}
\begin{picture}(0,0)(320,0)
\scriptsize{
\put(46,75){\tiny biella equivalente}
\put(45,75){\rotatebox{220}{$\longrightarrow$}}
\put(161,75){\tiny biella equivalente}
\put(159,75){\rotatebox{220}{$\longrightarrow$}}
\put(268,86){\tiny biella equivalente}
\put(254,85){\rotatebox{180}{$\longrightarrow$}}
\put(94,14){\tiny glifo}
\put(90,22){\rotatebox{180}{$\longrightarrow$}}
\put(146,45){\tiny glifo}
\put(156,41){\rotatebox{0}{$\longrightarrow$}}
\put(253,45){\tiny glifo}
\put(264,41){\rotatebox{0}{$\longrightarrow$}}
\put(0,50){$1)$}
\put(113,50){$2)$}
\put(222,50){$3)$}
}
\end{picture}
      \caption{\em Manovellismi non ordinari centrati equivalenti.}
 \label{fig:c_infinito_equivalenti}
\end{figure}

\noindent I tre meccanismi qui riportati presentano glifi
oscillanti attorno alla cerniera $D$, dove essi sono incernierati.
La lunghezza della ``biella'' \`e, dalla configurazione $1)$ alla
configurazione $3)$,  via via decrescente fino ad annullarsi.
Nel meccanismo $3)$ il piano di biella con tutte le sue curve
va pensato infatti solidale al cursore che scorre nel glifo.
Occorre per\`o precisare che, quando si parla di glifi oscillanti e ci si riferisce
appunto alla figura \ref{fig:c_infinito_equivalenti}, caso $3)$, ancora
pi\`u delle curve
dei punti di biella interessa il movimento del glifo stesso,
il quale \`e un bilanciere, e come tale si muove.

\section{Meccanismo a Ritorno Rapido}\index{meccanismo a ritorno rapido}\index{guida di Fairbairn}

La figura \ref{fig:c_infinito_quick} riporta nuovamente il meccanismo
di figura \ref{fig:c_infinito_equivalenti}, caso $3)$, ruotato
di un angolo opportuno, in modo da porre il suo telaio in verticale. 
\begin{figure}[hbt]
\centering
\includegraphics[width=0.9\textwidth]{part2/manovellismi/FIG/manovellismi/c_infinito_quick.pdf}
\begin{picture}(0,0)(250,0)
\scriptsize{
\put(235,245){$V$}
\put(136,125){$B$}
\put(73,158){$A$}
\put(73,62){$D$}
\put(180,50){$C$}
\put(175,50){\rotatebox{-42}{$\longrightarrow$}}
\put(173,40){$\infty$}
}
\end{picture}
      \caption{\em Meccanismo a ritorno rapido.}
 \label{fig:c_infinito_quick}
\end{figure}
\noindent A una rotazione completa della manovella
$AB$, di solito a velocit\`a $\omega$ costante, corrispondono
un'oscillazione di andata del glifo e una di ritorno che si svolgono
impiegando tempi diversi fra loro. In particolare, come abbiamo gi\`a 
visto nello studio dei quadrilateri, pagina
\pageref{eq:angoli_andata_ritorno}, tali tempi
saranno proporzionali ai rispettivi angoli di manovella.
Nella figura poc'anzi citata, ottenuta da un'analisi numerica,
non trova espressione grafica il pattino
che deve permettere al punto $B$ di scorrere
in una apposita scanalatura dell'asta $DV$, che per questo
assume generalmente il nome di glifo. Con la speranza  di porre
rimedio a tale lacuna grafica, riportiamo la prossima illustrazione,
nella quale
si pu\`o anche apprezzare l'utilizzo  concreto di questo manovellismo.
\begin{figure}[hbt]
\centering
\includegraphics[width=0.9\textwidth]{part2/manovellismi/FIG/fairbairn.pdf}
\begin{picture}(0,0)(300,10)
\scriptsize{
\put(42,376){\tiny andata lenta}
\put(40,347){\tiny ritorno rapido}
\put(258,298){$V$}
\put(251,278){$K$}
\put(175,287){\tiny slitta portautensile}
\put(135,107){$B$}
\put(50,211){$\omega=cost$}
\put(66,126){$\alpha$}
\put(70,106){$\beta$}
\put(42,164){$A$}
\put(32,28){$D$}
}
\end{picture}
      \caption{\em Meccanismo di Fairbairn.}
 \label{fig:fairbairn}
\end{figure}
\noindent La {\em guida di Fairbairn}\index{glifo!di Fairbairn}\index{Fairbairn, guida di},
rappresentata in figura \ref{fig:fairbairn},
si ottiene dal meccanismo
di figura \ref{fig:c_infinito_quick} con l'aggiunta di
una bielletta (o biscottino) di
snodo oppure di un opportuno accoppiamento prismatico (\cite{sesini1}, Fig. IV-8,
pag. 108, ma anche \cite{hartog}, pag. 168) in modo tale da permettere alla
sommit\`a
del glifo di azionare una slitta rettilinea. Tale soluzione costituiva,
fra l'altro,
il cuore di una ben nota macchina utensile che necessita di una corsa di andata
``lenta'' e di un ritorno ``rapido''.
Le tracce lasciate dalla sommit\`a del glifo, durante il suo movimento, dovrebbero dare
conto, mediante la loro spaziatura, delle differenti velocit\`a di $V$ durante
un ciclo di manovella. Le stesse spaziature, proiettate in pianta, dovrebbero
altrettanto rendere l'idea della velocit\`a del punto $K$, quindi della velocit\`a
del portautensile.
Anche qui, come nel caso del manovellismo ordinario centrato, \`e possibile
impostare le equazioni che legano l'angolo $\beta$ percorso dal glifo
all'angolo $\alpha$ di manovella;
derivando queste equazioni
rispetto ad $\alpha$ oppure rispetto al tempo si otterranno infine
le equazioni della cinematica del glifo oscillante centrato. Analogamente,
non presenterebbe certo difficolt\`a alcuna mettere
 in relazione l'angolo di oscillazione del glifo con lo spostamento
della slitta portautensile, della quale si otterrebbero finalmente
 velocit\`a e accelerazione, come mostrano
importanti autori, ad esempio \cite{hartog}, pag. 168.
Non ce la sentiamo per\`o di seguirli in tale percorso,
e ci accontentiamo dell'analisi numerica
riportata in modo qualitativo nelle figure
 \ref{fig:c_infinito_quick} e
 \ref{fig:fairbairn}, mediante le tracce del bottone di manovella
 e della estremit\`a del particolare glifo che abbiamo scelto di
proporre come esempio.

\noindent A nostro parere, dalla trattazione analitica del glifo
non emergono importanti spunti didattici,
al di l\`a dei meri passaggi trigonometrici:
i volenterosi e curiosi studenti
sapranno agilmente trovare tali calcoli nella citata bibliografia, se
interessati.
Per contro, colui che si vedr\`a assegnato il compito di progettare realmente
un meccanismo a ritorno rapido si affider\`a ai contemporanei codici numerici
di simulazione di meccanismi che  forniscono, con precisione, tutte le
grandezze cinematiche in gioco.
Inoltre sottolineiamo di nuovo che
 l'importanza di questo meccanismo risiede nella 
possibilit\`a di avere tempi di andata e ritorno sensibilmente diversi fra
loro, in modo da non sprecare tempo prezioso nella fase
oziosa del moto dell'utensile,
pi\`u che nella conoscenza puntuale della velocit\`a o della
accelerazione di quest'ultimo.
Nel nostro caso, invero un po' estremo,  abbiamo $\alpha_a=257^{\circ}$ e $\alpha_r=103^{\circ}$ che permettono
al tempo di andata, corrispondente al tempo di lavoro dell'utensile, di essere
pi\`u che doppio rispetto al tempo ``morto'' del ritorno.
Ai tempi in cui le macchine utensili venivano realizzate tramite pregiate
fusioni in ghisa, il glifo di {\em Fairbairn} era l'azionamento
principale della {\em limatrice}\index{limatrice}.
A proposito di questa applicazione, sottolineiamo che
la velocit\`a massima, raggiunta dalla slitta durante
la fase di andata (lenta), si manifesta quando il punto
$K$ \`e allineato coi punti $A$, $D$ e $B$. Essa \`e facilmente calcolabile
e vale, con $K$ in tale
posizione, $|v_{max}|=|\omega| |\overrightarrow{AB}||\overrightarrow{KD}|/|\overrightarrow{BD}|$; essa
rappresenta la velocit\`a  di taglio nominale dell'utensile.

\noindent Mi si permetta (tanto  tempo \`e trascorso quindi mi sento al sicuro da
rivendicazioni o proteste di carattere commerciale) di citare la ``gloriosa
e indistruttibile'' limatrice
Garavaglia.
Tutte le officine meccaniche rispettabili
ne avevano almeno una a disposizione. L'autore per\`o confessa che, gi\`a
trent'anni or sono, queste macchine beneficiavano di
notevolissimi periodi di riposo.

\section{Cinematica del Meccanismo a Ritorno Rapido: Analisi 
mediante i Moti Relativi}

\noindent La soluzione grafica che abbiamo esposto per l'analisi cinematica
di quadrilateri e manovellismi ordinari si pu\`o naturalmente estendere
anche al caso del meccanismo a ritorno rapido.
\noindent Riferendoci alla figura \ref{fig:soluzione_grafica_glifo}
posizioniamo, in questo caso, un sistema di riferimento relativo (al quale cio\`e
verranno riferite le grandezze relative) che ruota solidalmente con il
glifo ed \`e centrato in $D$.
Siamo in grado ora di esprimere la velocit\`a del punto $B$ (bottone di
manovella) mediante la somma delle due componenti: quella relativa
e quella di trascinamento:
\begin{equation}
{{\bm v}_{\scriptscriptstyle{B}}}_{\scriptscriptstyle{{\rm ass}}}= 
{{\bm v}_{\scriptscriptstyle{B}}}_{\scriptscriptstyle{{\rm rel}}}+
{{\bm v}_{\scriptscriptstyle{B}}}_{\scriptscriptstyle{{\rm tr}}}\,,
\label{eq:glifo_vel_graf0}
\end{equation}
\noindent ovvero, 
\begin{equation}
{{\bm v}_{\scriptscriptstyle{B}}}_{\scriptscriptstyle{{\rm ass}}}= 
{{\bm v}_{\scriptscriptstyle{B}}}_{\scriptscriptstyle{{\rm rel}}}+
{\omega_{\rm g}}|\overrightarrow{DB}|\widehat{{\perp{DB}}}\,.
\label{eq:glifo_vel_graf1}
\end{equation}

\begin{wrapfigure}{r}{0.5\textwidth}
     \begin{center}
     \includegraphics[width=0.48\textwidth]{part2/manovellismi/FIG/soluzione_grafica_glifo.pdf}
     \end{center}
\begin{picture}(0,0)(32,7)
	\scriptsize{
        \put(130,113){$B$}
        \put(41,37){$D$}
        \put(45,162){$A$}
        \put(34,71){$\scriptstyle{y'}$}
        \put(71,72){$\scriptstyle{x'}$}
        \put(62,130){${\omega_{\rm{m}}}$,}
        \put(77,130){${\dot\omega_{\rm{m}}}$}
        \put(63,112.5){${\omega_{\rm{g}}}$,}
        \put(78,112.5){${\dot\omega_{\rm{g}}}$}
        \put(132,157){${{\bm v}_{\scalebox{.5}{B}}}$,}
        \put(143,157){${{\bm a}_{\scalebox{.5}{B}}}$}
}
\end{picture}
\vskip -1.3mm
        \caption{\em Velocit\`a e accelerazioni nei moti relativi del meccanismo a ritorno rapido.}
     \label{fig:soluzione_grafica_glifo}
\end{wrapfigure}
\noindent La \ref{eq:man_vel_graf1} \`e un'equazione vettoriale nella quale
la direzione di tutti e tre i termini \`e nota
mentre
$|{{\bm v}_{\scriptscriptstyle{B}}}_{\scriptscriptstyle{{\rm rel}}}|$
e ${\omega_{\rm g}}$ sono incogniti, una volta data
la velocit\`a angolare della manovella. Anche qui riproponiamo quindi
lo schema risolutivo illustrato in figura \ref{fig:f116},
mediante il quale ricaveremo i valori delle incognite.

\noindent Per quanto riguarda le accelerazioni scriviamo
\begin{equation}
{{\bm a}_{\scriptscriptstyle{B}}}_{\scriptscriptstyle{{\rm ass}}}= 
{{\bm a}_{\scriptscriptstyle{B}}}_{\scriptscriptstyle{{\rm rel}}} 
+{{\bm a}_{\scriptscriptstyle{B}}}_{\scriptscriptstyle{{\rm tr}}} 
+{{\bm a}_{\scriptscriptstyle{B}}}_{\scriptscriptstyle{{\rm cor}}}\,,
\label{eq:acc_glifo_graf0}
\end{equation}
\noindent dove si nota la presenza del termine dell'accelerazione di {\em Coriolis}
dovuta alla rotazione del sistema di riferimento relativo.
Come al solito sostituiamo alle accelerazioni le loro componenti tangenziali
e normali ottenendo
\begin{equation}
{{\bm a}_{\scriptscriptstyle{B}}}_{\scriptscriptstyle{{\rm ass}}}= 
{{\bm a}_{\scriptscriptstyle{B}}}_{\scriptscriptstyle{{\rm rel}}} 
-{\omega_{\rm g}}^2 \overrightarrow{DB} + \dot{\omega}_{\rm g}|\overrightarrow{DB}|\widehat{{\perp{DB}}}+
2{\bm \omega}_{\rm g}\times
{{\bm v}_{\scriptscriptstyle{B}}}_{\scriptscriptstyle{{\rm rel}}}\,.
\label{eq:acc_glifo_graf1}
\end{equation}
\vskip 1mm
\noindent Osserviamo che risultano note le direzioni di tutti i termini
di questa equazione.
Contando poi i termini della \ref{eq:acc_glifo_graf1} da sinistra,
affermiamo che gli unici moduli incogniti
sono quelli dei vettori che occupano il primo e il terzo posto
dopo il degno di uguaglianza. 
Perci\`o otteniamo ancora un'equazione vettoriale
dove le incognite sono
il modulo dell'accelerazione relativa di $B$,
$|{{\bm a}_{\scriptscriptstyle{B}}}_{\scriptscriptstyle{{\rm rel}}}|$,
e il modulo dell'accelerazione angolare del glifo
${\dot\omega_{\rm g}}$.
Il riferimento per una soluzione grafica resta dunque 
il solito, cio\`e lo schema grafico di figura \ref{fig:f116}.
Ammettiamo che, passando dall'esempio di pagina \pageref{fig:f116}
all'approccio grafico alla cinematica dei quadrilateri e infine
alla stessa analisi operata sui manovellismi, ammettiamo
dicevo che la nostra esposizione si 
\`e fatta via via pi\`u succinta, sottacendo di passo in passo
alcune precisazioni presenti nei primi esempi.
Il motivo di questa progressiva reticenza sta nel pudore, forse ingiustificato,
che nasce quando si \`e costretti a ripetere molte volte contenuti non
molto dissimili da quelli espressi poco addietro.
 Siamo certi che lo studente volenteroso
sapr\`a rintracciare in questo percorso tutte le informazioni che
sono state via via tralasciate.


\section{Manovellismi non Ordinari con Due Membri Rotanti}\index{glifo!rotante}

\noindent Nel precedente paragrafo abbiamo analizzato il meccanismo a ritorno
rapido nel quale, stante l'opportunit\`a di realizzare il glifo sull'asta
$DV$, figura \ref{fig:fairbairn}, abbiamo ottenuto un glifo oscillante,
che si comporta cio\`e come un  bilanciere, e abbiamo citato una delle sue
pi\`u famose applicazioni che consiste nella movimentazione di
utensili da taglio.
\begin{wrapfigure}{r}{0.4\textwidth}
     \begin{center}
\includegraphics[width=0.38\textwidth]{part2/manovellismi/FIG/manovellismi/vel_b_infinito.pdf}
\end{center}
\begin{picture}(0,0)(0,0)
\scriptsize{
\put(55,70){${\omega_{\scalebox{.5}{DC}}}/{\omega_{\scalebox{.5}{AK}}}$}
        \put(141,27){$\alpha$}
}
\end{picture}
\vskip -3.2mm
      \caption{\em Velocit\`a del membro $DC$.}
     \label{fig:vel_b_infinito}
\end{wrapfigure}
\noindent Non \`e difficile, adottando
opportune configurazioni, ottenere per tale meccanismo un glifo
 rotante (sempre che
la volont\`a del progettista lo foggi come un glifo). L'analisi
qualitativa di tali meccanismi sar\`a oggetto di questo paragrafo.
\noindent Mandiamo
all'infinito la cerniera $B$ di un opportuno quadrilatero di partenza,
avente geometria tale da poter funzionare (si riveda la regola di {\em Grashof})
 a ``doppia-manovella''.
Abbiamo scelto la cerniera $B$ anzich\'e la $C$ semplicemente per
mostrare al lettore la completa equivalenza, in questa circostanza,
 delle
due cerniere. Nella nostra analisi, che svolgeremo come abbiamo fatto
altre volte tramite l'approccio numerico,
la manovella a velocit\`a costante \`e sempre la $AB$, che lascia
le sue tracce equispaziate; questo ci
consentir\`a di avere sulla manovella $DC$, l'elemento che di solito
\`e costituito da un glifo come in figura \ref{fig:glifo_rotante},
la velocit\`a variabile.
 
\begin{figure}[htb]
\begin{center}
\includegraphics[width=.8\textwidth]{part2/manovellismi/FIG/manovellismi/b_infinito.pdf}
\end{center}
\begin{picture}(0,0)(-35,20)
\scriptsize{
\put(92,330){$B$}
\put(100,330){\rotatebox{67}{$\longrightarrow$}}
\put(104,328){$\infty$}
\put(-2,225){$C$}
\put(21,212){$K$}
\put(88,177){$A$}
\put(142,177){$D$}
\put(154,332){$\beta_2=86^{\circ}$}
\put(135,259){$\alpha_2=68^{\circ}$}
\put(115,100){$\alpha_1=292^{\circ}$}
\put(121,38){$\beta_1=274^{\circ}$}
}
\end{picture}
\vskip -3mm
      \caption{\em Manovellismo non ordinario che permette la rotazione completa
di due membri.}
 \label{fig:b_infinito}
\end{figure}

\noindent In particolare, la figura \ref{fig:vel_b_infinito} riporta l'andamento
del rapporto tra la velocit\`a angolare del glifo, $\omega_{DC}$, e quella
della manovella (presente nella figura \ref{fig:glifo_rotante}), $\omega_{AK}$.
Se volessimo valorizzare il massimo squilibrio che si pu\`o ottenere
tra la rotazione della manovella e quella del glifo,
come abbiamo esposto nel paragrafo {\ref{q_squilibrio}},
dovremmo anche qui trovare una posizione ben
precisa ``di montaggio'' di questo meccanismo.
Quest'ultimo quindi, se opportunamente sfruttato, potr\`a presentare, a fronte di due
semi-rotazioni contigue di $180^{\circ}$ del glifo, due angoli di manovella
sensibilmente diversi tra loro.
\begin{figure}[b]
\begin{center}
\includegraphics[width=.8\textwidth]{part2/manovellismi/FIG/glifo_rotante.pdf}
\end{center}
\begin{picture}(0,0)(-35,20)
\scriptsize{
\put(12,222){$C$}
\put(34,219){$K$}
\put(88,179){$A$}
\put(163,183){$D$}
\put(180,274){$\beta$}
\put(100,230){$\alpha$}
\put(52,142){$\omega={\rm cost}$}
}
\end{picture}
\vskip -5mm
      \caption{\em Manovellismo non ordinario 
 con glifo rotante.}
 \label{fig:glifo_rotante}
\end{figure}
 Riferendoci all'esempio di figura
\ref{fig:b_infinito}, figura sulla quale si 
basa quella successiva, \ref{fig:glifo_rotante}, 
troviamo indicate le coppie di angoli
$\alpha_1$, $\beta_1$, $\alpha_2$, $\beta_2$ che realizzano in questo caso
il {\em massimo squilibrio} pari a
\begin{equation}
s={{292^{\circ}-68^{\circ}}\over{360^{\circ} -292^{\circ} +68^{\circ}}} = 1.6\,.
\label{eq:esempio_squilibrio_glifo}
\end{equation} 
\noindent A una rotazione completa della manovella $AK$, a velocit\`a costante,
corrisponderanno quindi due semi-rotazioni del glifo che si completeranno
in tempi diversi tra loro con $t_a/t_r=1.6$. Si potrebbe quindi impiegare questo
glifo come manovella di
un manovellismo ordinario centrato, per il quale le corse di andata e ritorno
si svolgerebbero in tempi il cui rapporto reciproco \`e stato or ora indicato. 
Questa soluzione \`e stata talvolta adottata (glifo di {\em Whitworth}) in sostituzione
del glifo di {\em Fairbairn}, \cite{sesini1}, pag. 109.

\section{Quadrilateri Doppiamente Degeneri e Altre Considerazioni}

Tratteremo brevemente soltanto il caso in cui, spostando le due cerniere mobili,
$B$ e $C$, in due punti impropri distinti del piano, si d\`a luogo
a un interessante meccanismo di discreta importanza nelle applicazioni pratiche:
il {\em giunto di {\em Oldham}}\index{giunto di Oldham}.
La figura \ref{fig:oldham} mostra un quadrilatero a doppia-manovella
dove il telaio $AD$ misura circa la ventesima parte delle altre aste.
Siamo costretti, si fa per dire, a svelare ora l'arcano che sta alla base
delle figure dei manovellismi riportati in questo capitolo.
\begin{figure}[htb]
\begin{center}
\vspace{-0cm}\hspace{-.10cm}\includegraphics[width=.8\textwidth]{part2/manovellismi/FIG/manovellismi/oldham.pdf}
\hbox{\vspace{-12cm}\hspace{-4cm}\includegraphics[width=0.3\textwidth]{part2/manovellismi/FIG/oldham_particolare.pdf}}
\end{center}
\begin{picture}(0,0)(-55,15)
\scriptsize{
\put(102,333){$B$}
\put(110,330){\rotatebox{90}{$\longrightarrow$}}
\put(114,333){$\infty$}
\put(232,293){$C$}
\put(234,287){\rotatebox{43.4}{$\longrightarrow$}}
\put(232,283){$\infty$}
\put(111,178){$A$}
\put(125,178){$D$}
\put(175,45){$A$}
\put(198,45){$D$}
\put(205,128){$T$}
\put(72,264){$P$}
}
\end{picture}
\vskip -5mm
      \caption{\em Quadrilatero doppiamente degenere e relativo meccanismo equivalente.
}
 \label{fig:oldham}
\end{figure}
Di fatto,
le analisi numeriche che danno origine alle figure \ref{fig:quad_D_infinito_man},
\ref{fig:c_infinito}, \ref{fig:c_infinito_quick} e \ref{fig:b_infinito} sono
state eseguite mediante lo  stesso (semplice e imperfetto) codice che genera
anche le analisi dei quadrilateri del precedente capitolo dedicato
allo studio dei quadrilateri articolati. In ambito numerico,
esiste infatti la possibilit\`a di analizzare, a volte con un certo agio,
anche quelle situazioni limite lasciate
giustamente dalla matematica nella sfera delle idee astratte, come sono in
generale le grandezze che tendono all'infinito.
Inoltre, le scienze ingegneristiche, pur abbracciando senza perplessit\`a il
concetto di distanza infinita, possono in altri casi permettersi il 
lusso  di accontentarsi di qualcosa di meno e magari pi\`u a portata di mano.
\`E il nostro caso, quando pretendiamo di trattare numericamente
le lunghezze, teoricamente infinite, delle aste $AD$ e $CD$ del quadrilatero
di figura \ref{fig:quad_D_infinito_man} sostituendo tali membri con
aste di lunghezza finita, bench\'e 
cinquanta o sessanta volte maggiore rispetto alle altre due.
Del resto, a nessun ingegnere di buon senso verrebbe in mente di obiettare
che la traiettoria del cursore, in figura \ref{fig:quad_D_infinito_man},
indicata con $\sigma$,
sia un arco di circonferenza anzich\'e un segmento di retta.
Torniamo alla figura \ref{fig:oldham} e  notiamo che la manovella $CD$
\`e omocinetica alla $AB$ (nell'approssimazione numerica adottata per le
lunghezze infinite delle tre aste $AB$, $BC$, e $CD$), cio\`e le due manovelle
hanno la stessa velocit\`a angolare.
La curva percorsa dal generico punto di biella $P$ approssimer\`a
tanto meglio una circonferenza quanto pi\`u distante si trover\`a dal telaio
(come nella figura) e anch'essa risulta percorsa con velocit\`a di modulo
costante. Risulta poi evidente che le curve di biella dei punti giacenti
(punti non rappresentati in figura per evitare confusione nel disegno)
 sulle manovelle, taglieranno queste ultime
seguendo direzioni ad esse ortogonali\footnote{\`E facile convincersi
che passando a lunghezze delle tre aste $AB$, $BC$ e $CD$ cento volte
superiori a quelle rappresentate  nella figura, i punti $A$ e $D$ verrebbero
a coincidere nel disegno in un punto $V$. Il meccanismo diventerebbe cos\`i
un triangolo $\triangle VBC$ che ruota attorno al suo vertice trascinando
in questa rotazione, oltre alle due manovelle, tutto il piano di biella.}.
\begin{figure}[hbt]
\centering
\begin{minipage}[b]{0.48\textwidth}
\centering
\includegraphics[width=0.95\textwidth]{part2/manovellismi/FIG/oldham_schema.pdf}
\begin{picture}(0,0)(190,0)
\scriptsize{
\put(36,40){$A$}
\put(48,75){$\scriptstyle{M_1}$}
\put(105,45){$\scriptstyle{M_2}$}
\put(110,110){$T$}
\put(85,57){$D$}
\put(116.5,57.5){$\omega$}
\put(35,51){$\omega$}
\put(33,54){\rotatebox{30}{$\curvearrowleft$}}
\put(110,62){\rotatebox{-60}{$\curvearrowleft$}}
}
\end{picture}
      \caption{\em Rappresentazione funzionale del giunto di Oldham.}
 \label{fig:oldham_schema}
\end{minipage}\hfill
\begin{minipage}[b]{0.48\textwidth}
\centering
\includegraphics[width=0.95\textwidth]{part2/manovellismi/FIG/oldham_commerciale.pdf}
\begin{picture}(0,0)(100,0)
\scriptsize{
\put(30,80){$\scriptstyle{M_1}$}
\put(65,54){$\scriptstyle{M_2}$}
\put(-6,28){$T$}
}
\end{picture}
	\caption{\em Giunto di Oldham commerciale che permette disassamenti tra gli alberi fino a $5\, {\rm mm}$.}
     \label{fig:oldham_commerciale}
\end{minipage}
\end{figure}
\noindent Nel riquadro in basso a destra
della figura \ref{fig:oldham} \`e riportato un meccanismo equivalente
al quadrilatero doppiamente degenere che stiamo analizzando e 
la trave $T$ 
di tale figura riproduce quindi il comportamento della
biella $BC$.
\noindent Si intravede cos\`i la possibilit\`a di trasmettere
omocineticamente il moto
tra due alberi paralleli e tra loro disassati. In figura \ref{fig:oldham_schema}
\`e riportato un {\em giunto di Oldham} che presenta la trave $T$ come
una squadra. Questa soluzione \`e suggerita dal fatto che appare come privo di
senso privilegiare particolari direzioni radiali in un dispositivo il cui
funzionamento spazia su giri completi. D'altra
parte anche la posizione delle due ``chiavi'' pu\`o essere ottimizzata in modo
da diminuire l'ingombro complessivo. La figura \ref{fig:oldham_commerciale}
mostra un piccolo giunto di {\em Oldham} commerciale, dove i glifi sono stati
ricavati sulla trave $T$ e non sulle due manovelle; la figura mostra 
i tre componenti disgiunti, le due manovelle e la trave di collegamento,
nonch\'e il giunto
assemblato che dovrebbe convincere il lettore
della sua semplicit\`a e compattezza.

\noindent Come promesso nel titolo del presente paragrafo, esponiamo
qualche ulteriore considerazione circa i manovellismi e la loro classificazione
in ordinari e non ordinari. Nel percorso che abbiamo seguito in questo lavoro
viene evidenziata la differente genesi delle due tipologie di manovellismo:
quando a divenire un punto improprio \`e una delle due cerniere fisse
si ottiene un manovellismo ordinario, viceversa
si ottiene un manovellismo non ordinario quando a diventare impropria \`e
una delle altre due cerniere.  Mentre nella sfera delle applicazioni
pratiche la scelta del ``membro telaio'' \`e tutt'altro che indifferente (a
ben pochi \`e venuto in mente di costruire un motore dove a ruotare fosse
il blocco dei cilindri),
nell'ambito della cinematica teorica possiamo individuare soltanto
i moti relativi tra
le varie parti di un quadrilatero o di un manovellismo. In altre parole,
considerando di nuovo il manovellismo ordinario di figura \ref{fig:man_ord}
e tirando in ballo la nostra ormai affezionata formica esploratrice
(ve la ricordate? Quella della giostra dei cavalli...), potremmo
affermare che tale insetto vedrebbe un manovellismo ordinario quando
sta fermo sul glifo (il cilindro), un manovellismo non ordinario 
come quello di figura \ref{fig:c_infinito_equivalenti}- $3)$ quando
sta fermo sulla manovella $AB$ e un meccanismo equivalente a quello di
figura  \ref{fig:c_infinito_equivalenti}- $1)$ quando la formica sta ferma 
sulla biella $BC$. 
Questo punto di vista circa i diversi manovellismi,
che riteniamo chiarificatore e
didattico, viene esposto in molti libri che trattano questo argomento come,
ad esempio, \cite{sesini1}, pag. 107.
%\newpage
%\thispagestyle{empty}
%\null
\endinput 

\chapter{Leggi di Moto e Profilatura delle Camme}\index{eccentrici}\index{camme}
\label{camme}

\section{Introduzione}

\noindent Nell'ambito delle macchine automatiche con elevata cadenza di funzionamento,
dedicate alla produzione su vasta scala di manufatti,
si incontrano sovente organi meccanici che eseguono movimenti ripetitivi,
composti da ben precise successioni di salite, soste e discese, nomi che tendono a essere
conservati anche quando i movimenti non sono verticali. La legge temporale di tali
movimenti \`e  solitamente
dettata dal particolare processo tecnologico che la macchina o una parte di essa
\`e chiamata a svolgere.
\begin{figure}[hbt]
\centering
\begin{minipage}[b]{0.600\textwidth}
\centering
\includegraphics[width=0.9\textwidth]{part2/camme/FIG/sagoma_traslazione.pdf}
\begin{picture}(0,0)(130,5)
\scriptsize{
\put(-22,235){$y(t)$}
\put(-25,94){$v$}
\put(-15,70){$T$}
}
\end{picture}
      \caption{\em Movimentazione di un cedente mediante sagoma di traslazione.}
 \label{fig:f_sagoma_traslazione}
\end{minipage}\hfill
\begin{minipage}[b]{0.300\textwidth}
\centering
\includegraphics[width=0.9\textwidth]{part2/camme/FIG/sagoma_avvolta.pdf}
\begin{picture}(0,0)(129,0)
\scriptsize{
\put(72,220){$y(t)$}
\put(50,51) {\color{white}$\omega$}
}
\end{picture}
	\caption{\em Sagoma di traslazione avvolta su un tamburo cilindrico.}
     \label{fig:f_sagoma_avvolta}
\end{minipage}
\end{figure}
\noindent Sono frequentissimi i casi di presse e fustelle montate in cosiddette ``stazioni''
a bordo di una macchina che eseguendo particolari {\em leggi di
moto}\index{leggi di moto} tagliano, sagomano, incollano
i vari pezzi di materiale cosicch\'e, alla fine di un ``ciclo
macchina'', otterremo il manufatto finito. 
Giusto per intenderci, \`e tipico incontrare in questo perimetro
macchine automatiche  che 
producono quantit\`a elevate di stoviglie in carta o alluminio, di semplici
calzature, di tessuti, di contenitori per la cosmetica, di bottiglie
realizzate in svariate materie plastiche vuote o riempite contestualmente
alla loro formatura,
e chi pi\`u ne ha pi\`u ne metta. La cadenza di tali macchine pu\`o raggiungere e superare 
quella di due ``colpi'' al secondo.
La figura \ref{fig:f_sagoma_traslazione} riporta una idea {\em na\"ive} di un cedente 
che, tastando con una punta (o coltello) una sagoma che trasla,  viene azionato verticalmente, quindi
costretto a compiere la legge costituita dal profilo della sagoma stessa. La {\em camma lineare}\index{camme!lineari}
cos\`i ottenuta viene normalmente ``avvolta'', come in figura
\ref{fig:f_sagoma_avvolta}, su un tamburo il quale, ruotando generalmente a velocit\`a
costante, pu\`o riprodurre automaticamente ad ogni giro
la legge di moto senza la necessit\`a di ripeterla molte volte in sequenza,
come sarebbe necessario fare sulla sagoma traslante.  
\begin{wrapfigure}{r}{0.5\textwidth}
     \begin{center}
     \includegraphics[width=0.42\textwidth]{part2/camme/FIG/generic_law/legge_completa.pdf}
     \end{center}
\begin{picture}(0,0)(0,0)
	\scriptsize{
	\put(155,111){$\alpha$}
	\put(6,118){$y(\alpha)$}
	\put(7,105){$y(t)$}
	\put(155,25){$t$}
	\put(130,18){$T$}
}
\end{picture}
        \caption{\em Diagramma delle alzate.}
     \label{fig:f_diagramma_alzate}
\end{wrapfigure}
\noindent Nella stragrande maggioranza delle macchine automatiche la generazione di movimenti opportuni rispettando tempi
precisi \`e di fondamentale importanza. Questo perch\'e i movimenti di fustelle,
presselli, spruzzatori, togli-metti foglio, e altri dispositivi, frequenti in 
questo ambito, necessitano di tempi e luoghi precisi per  soddisfare il processo
tecnologico al quale sono adibiti: processo che pu\`o essere rappresentato da una operazione di imbutitura, di
formatura, di incollaggio ecc. Unitamente a ci\`o, i movimenti di questi
dispositivi devono spesso 
rispettare
sincronizzazioni temporali tra le operazioni stesse e
movimenti di altri organi della medesima macchina.
Questi movimenti ripetitivi possono essere rappresentati in un grafico che prende
il nome di {\em diagramma temporale delle alzate}\index{diagramma delle alzate}, del quale un esempio \`e riportato in
figura \ref{fig:f_diagramma_alzate}.
Abbiamo poc'anzi accennato al fatto che normalmente 
la generazione delle alzate avviene tramite la trasformazione del moto
rotatorio di un albero, il quale ruota
a velocit\`a angolare  $\omega$ costante: ad ogni rotazione
completa dell'{\em albero a camme}\index{albero a camme},
sul quale sar\`a presente la {\em camma}\index{camme}, si completer\`a un intero ciclo del diagramma delle alzate.
Entreremo tra poco nel cuore di una 
descrizione, certamente non esaustiva, dei dispositivi a camma che operano la suddetta trasformazione;
anteponiamo qui solo qualche considerazione generale sulle leggi di movimento.
Tali considerazioni, peraltro, hanno applicazioni ben pi\`u estese rispetto al solo ambito delle camme meccaniche.
Riteniamo tuttavia che i meccanismi a camma offrano un generoso campo di
esempi che aiutano a comprendere l'argomento.

\section{Le Leggi di Moto}

\noindent Fissiamo le idee sulla movimentazione della mazza di una pressa
la quale dovr\`a eseguire gli spostamenti ciclici indicati
in figura \ref{fig:f_diagramma_alzate}
con una alzata massima del cedente pari
a $40\,{\rm mm}$.
Se si conviene di indicare con $T$ il {\em periodo} di tempo entro il 
quale si completa un {\em ciclo macchina}\index{ciclo macchina}, pari
nel nostro caso a $0.8\,{\rm s}$,
la velocit\`a angolare costante dell'albero che genera tale movimento sar\`a
$\omega = {2\pi \over{T}}=7.85\,{\rm rad}/{\rm s}$ che corrisponde a 75
giri al minuto.
Dato il legame tra l'angolo percorso dall'albero a camme e il tempo trascorso,
${\rm d}\alpha= \omega {\rm d}t$,  possiamo sostituire alla variabile temporale
della funzione $y(t)$ la variabile angolare $\alpha$, considerando che
in un ciclo macchina sar\`a $0 \leq \alpha \leq 2\pi$.  
Si chiama {\em diagramma delle alzate}\index{diagramma delle alzate}
il grafico di $y(\alpha)$ che mostra il legame geometrico tra le alzate di una
camma e la sua rotazione.
Si noti che la sovrapposizione del diagramma temporale delle alzate con il diagramma in funzione dell'angolo $\alpha$ \`e possibile
soltanto quando $\omega$ \`e costante, condizione questa, come abbiamo gi\`a affermato,
generalmente soddisfatta. La velocit\`a e l'accelerazione del cedente hanno
quindi un loro corrispondente di natura geometrica. Infatti derivando la 
funzione $y(t)$ rispetto al tempo otteniamo per la velocit\`a
\begin{equation}
{{\rm d}y(t)\over{{\rm d}t}}= \dot y= \omega {{\rm d}y(\alpha)\over{{\rm d}\alpha}}= \omega y'\,,
\end{equation}
\noindent e per l'accelerazione
\begin{equation}
{{\rm d}^2y(t)\over{{\rm d}t^2}}= \ddot y= \dot\omega 
{{\rm d}y(\alpha)\over{{\rm d}\alpha}}+
\omega^2 {{\rm d}^2y(\alpha)\over{{\rm d}\alpha^2}}=\dot\omega y' + \omega^2 y''\,,
\end{equation}
\noindent che, nel caso di velocit\`a angolare costante, si riduce a
\begin{equation}
\ddot y= \omega^2 y''\,.
\end{equation}
\noindent Essendo le {\em quantit\`a cinematiche geometriche}\index{velocit\`a!geometrica}\index{accelerazione!geometrica}, legate esclusivamente alla forma
della camma (pi\`u in generale legate al particolare meccanismo che le produce
utilizzando un moto rotatorio), e indipendenti dal regime di rotazione dell'albero,
esse risultano
estremamente comode per confrontare tra loro
{\em leggi di moto} diverse,
nel processo di sintesi progettuale della camma stessa.
In moltissimi casi, il diagramma delle alzate presenta una certa libert\`a interpretativa.
Spesso, ci\`o che sostanzialmente importa nella movimentazione dell'utensile
sono i valori delle successive alzate e le loro durate temporali, nonch\'e le
durate delle varie soste:
la legge di moto $y(t)$ che collega 
due posizioni di alzata invece non \`e quasi mai fissata in maniera stringente.
Nel caso del diagramma delle alzate di figura \ref{fig:f_diagramma_alzate}, ad esempio,
contiamo tre leggi di moto distinte sugli intervalli di tempo:
$0-0.1$, $0.2-0.3$ e $0.5-0.7$, frammiste a tre 
soste: $0.1-0.2$, $0.3-0.5$ e $0.7-0.8$. 
La mancanza di vincoli rigidi per i tratti di salita e discesa, cui abbiamo test\'e accennato, ci consente  la possibilit\`a di scegliere, a parit\`a
 sostanziale del
diagramma delle alzate, leggi di moto diverse che soddisfacendo tale diagramma 
possono sottostare anche ad alcuni altri criteri che sceglieremo e imporremo in
modo opportuno. Il pi\`u costrittivo di tali criteri \`e quello che permette
al cedente di giungere nei tratti di sosta effettivamente fermo. Indicando 
perci\`o con $t_s$ il tempo in cui si svolge la legge di salita deve 
essere
\begin{equation}
\int_0^{t_s} \ddot y(t) {\rm d}t=0\,.
\label{e_vel_zero}
\end{equation}
\noindent Oppure, considerando la variabile angolare anzich\'e il tempo e
indicando con $\alpha_s$ l'angolo di rotazione della camma
sul quale si estende quel tratto di salita, dovr\`a essere
\begin{equation}
\int_0^{\alpha_s}  y''(\alpha) {\rm d}\alpha=0\,.
\label{e_vel_zero_alpha}
\end{equation}
\noindent Una volta scelta la ``forma'' della
legge di accelerazione che soddisfa le relazioni \ref{e_vel_zero} e
\ref{e_vel_zero_alpha}, dovranno essere verificate le seguenti relazioni
\begin{equation}
\int_0^{t_s} {\rm d}t\int_0^{t_s} \ddot y(t) {\rm d}t\;\;=\;\;
\int_0^{\alpha_s} {\rm d}\alpha\int_0^{\alpha_s} y''(\alpha) {\rm d}\alpha=h\,,
\label{eq:ha}
\end{equation}
\noindent o anche
\begin{equation}
\int_0^{t_s} \dot y(t) {\rm d}t\;\;=\;\;
\int_0^{\alpha_s} y'(\alpha) {\rm d}\alpha=h\,,
\label{eq:hv}
\end{equation}
\noindent in maniera che tale legge produca effettivamente l'alzata richiesta.
Segnaliamo che,
nella parte dedicata agli approfondimenti, al capitolo \ref{cam2},
\`e presente 
una breve nota circa un'ulteriore propriet\`a notevole delle leggi di moto.
\vskip .3cm
\begin{figure}[hbt]
\centering
\begin{minipage}[b]{0.49\textwidth}
\centering
\includegraphics[width=0.9\textwidth]{part2/camme/FIG/generic_law/confronto_alzate_acc_diverse_acc.pdf}
\begin{picture}(0,0)(130,0)
\scriptsize{
\put(-20,96){${{}\ddot{y}_{3}}_{\rm max}=17.5$}
\put(-40,90){$y''$}
\put(-40,80){$\ddot y$}
\put(10,81){${{}\ddot{y}_{2}}_{\rm max}=6$}
\put(45,70){${{}\ddot{y}_{1}}_{\rm max}=4$}
\put(120,1){$\alpha$}
\put(114,1){$t,$}
}
\end{picture}
      \caption{\em Diverse leggi di accelerazione.}
 \label{fig:f_confronto_accelerazioni}
\end{minipage}\hfill
\begin{minipage}[b]{0.49\textwidth}
\centering
\includegraphics[width=0.9\textwidth]{part2/camme/FIG/generic_law/confronto_alzate_acc_diverse_spost.pdf}
\begin{picture}(0,0)(130,0)
	\scriptsize{
\put(-38.5,90){$y(\alpha)$}
\put(-38.5,80){$y(t)$}
\put(114,1){$t,$}
\put(120,1){$\alpha$}
}
\end{picture}
	\caption{\em Corrispondenti diagrammi dell'alzata.}
     \label{fig:f_confronto_alzate}
\end{minipage}
\end{figure}

\noindent In  figura \ref{fig:f_confronto_alzate} sono riportate tre diverse leggi
di spostamento che producono la stessa alzata, $h=1$, durante una rotazione
dell'albero di un radiante, oppure dopo il tempo di un secondo (si sottintende
perci\`o $\omega=1$), terminando la corsa con velocit\`a nulla.
Mentre la discrepanza
tra i diagrammi dello spostamento, $y(t)$ ovvero $y(\alpha)$,
risulta essere modesta, 
si pu\`o osservare che i tre diagrammi di accelerazione,
riportati in figura \ref{fig:f_confronto_accelerazioni}
unitamente ai loro valori massimi, sono molto diversi tra loro.
Ci si pu\`o chiedere: \`e sensato scegliere la
legge di moto 1) che, a parit\`a di alzata, presenta la minima accelerazione
massima? La risposta \`e s\`i e no. Infatti, il processo di sintesi
di una legge di moto per camme deve tenere conto di molte
circostanze, e la scelta di una legge che privilegi le basse accelerazioni pu\`o
presentarsi mediocre sotto altri aspetti. Cercheremo nel seguito
 di dare un 
quadro, peraltro parziale, del complesso di vincoli che guidano
la progettazione di una legge di moto e di una camma.
Riportiamo in figura \ref{fig:f_confronto_velocita}
il confronto tra le velocit\`a associate alle precedenti tre leggi di
moto in cui si nota che la velocit\`a pi\`u elevata viene raggiunta 
dalla legge con accelerazione massima minore
e viceversa.
\begin{figure}[hbt]
\centering
\begin{minipage}[b]{0.49\textwidth}
\centering
\includegraphics[width=0.9\textwidth]{part2/camme/FIG/generic_law/confronto_alzate_acc_diverse_acc.pdf}
\begin{picture}(0,0)(130,0)
\scriptsize{
\put(-20,96){${{}\ddot{y}_{3}}_{\rm max}=17.5$}
\put(-40,90){$y''$}
\put(-40,80){$\ddot y$}
\put(10,81){${{}\ddot{y}_{2}}_{\rm max}=6$}
\put(45,70){${{}\ddot{y}_{1}}_{\rm max}=4$}
\put(120,1){$\alpha$}
\put(114,1){$t,$}
}
\end{picture}
      \caption{\em Diverse leggi di accelerazione.}
 \label{fig:f_confronto_accelerazioni2}
\end{minipage}\hfill
\begin{minipage}[b]{0.49\textwidth}
\centering
\includegraphics[width=0.9\textwidth]{part2/camme/FIG/generic_law/confronto_alzate_acc_diverse_vel.pdf}
\begin{picture}(0,0)(130,0)
\scriptsize{
\put(-48,90){$y'(\alpha)$}
\put(-48,80){$\dot y(t)$}
\put(114,1){$t,$}
\put(120,1){$\alpha$}
\put(22,80){${{}\dot{y}_{1}}_{\rm max}=2$}
\put(62,63){${{}\dot{y}_{2}}_{\rm max}=1.42$}
\put(17,39){${{}\dot{y}_{3}}_{\rm max}=1.11$}
}
\end{picture}
	\caption{\em Corrispondenti diagrammi della velocit\`a.}
     \label{fig:f_confronto_velocita}
\end{minipage}
\end{figure}

\noindent L'importanza di evitare picchi troppo elevati di accelerazione
deriva dal diretto collegamento tra questa grandezza
e l'entit\`a delle forze d'inerzia che insorgono durante il funzionamento della macchina.
Ed \`e questo il motivo che induce il progettista a disegnare la camma partendo
proprio dalle leggi di accelerazione, tenendo cos\`i sotto controllo
tale grandezza.
La legge con {\em accelerazione costante} riportata in
figura \ref{fig:f_legge_acc_cost} \`e la legge di moto pi\`u ``semplice''
che si possa adottare.
In tale figura, come nelle precedenti 
 \ref{fig:f_confronto_accelerazioni}--\ref{fig:f_confronto_velocita},
e nelle seguenti, riguardanti le leggi di moto, le ascisse riportano
sia il tempo $t$, sia l'angolo $\alpha$. Pertanto, queste leggi
possono essere qui interpretate come eseguite in un tempo unitario oppure
in un angolo unitario, ponendo cos\`i $\omega=1\; {\rm rad}/{\rm s}$. Le stesse leggi producono tutte una alzata unitaria.
\begin{figure}[hbt]
\centering
\begin{minipage}[b]{0.49\textwidth}
\centering
\includegraphics[width=0.9\textwidth]{part2/camme/FIG/generic_law/legge_acc_cost.pdf}
\begin{picture}(0,0)(150,-10)
\scriptsize{
\put(-17,75){$\ddot y$}
\put(-17,65){$\dot y$}
\put(-17,55){$y$}
\put(15,65){${{}\ddot{y}}_{\rm max}=4$}
\put(64,59){${{}\dot{y}}_{\rm max}=2$}
\put(140,-9){$\alpha$}
\put(134,-9){$t,$}
}
\end{picture}
\caption{\em Legge di moto con accelerazione costante.}
 \label{fig:f_legge_acc_cost}
\end{minipage}\hfill
\begin{minipage}[b]{0.49\textwidth}
\centering
\includegraphics[width=0.9\textwidth]{part2/camme/FIG/generic_law/legge_acc_cost_tagliata.pdf}
\begin{picture}(0,0)(150,-10)
\scriptsize{
}
\put(-17,75){$\ddot y$}
\put(-17,65){$\dot y$}
\put(-17,55){$y$}
\put(39,70){${{}\ddot{y}}_{\rm max}=4.76$}
\put(48,55){${{}\dot{y}}_{\rm max}=1.43$}
\put(140,-9){$\alpha$}
\put(134,-9){$t,$}
\end{picture}
	\caption{\em Legge di moto con accelerazione costante tagliata.}
     \label{fig:f_acc_cost_tagliata}
\end{minipage}
\end{figure}
La scelta di rappresentare leggi di moto che producono
alzata unitaria in tempi o angoli unitari, leggi che chiameremo {\em unitarie}\index{leggi di moto!unitarie},
agevola il confronto
tra accelerazioni ``di forme diverse'', in quanto tale confronto
pu\`o avvenire semplicemente tra
alcuni dei valori notevoli che le grandezze cinematiche, di tali leggi
unitarie, assumono. Per passare da questi valori, cio\`e quelli
relativi alle grandezze cinematiche delle leggi unitarie e che,
limitatamente a questo capoverso, indicheremo
con pedice $u$, alle grandezze cinematiche effettive (pedice $e$) di una legge
 di uguale forma,
la quale per\`o produce un'alzata pari a $h$ in un tempo $t_s$ o in un angolo
$\alpha_s$, con riferimento agli integrali
\ref{eq:ha} e \ref{eq:hv}, si ha
\begin{equation}
\ddot y_e = \ddot y_u \frac{h}{t_s^2}\,,\;\;\;\;
{y''_e} = {y''_u} \frac{h}{\alpha_s^2}\,,\;\;\;\;
\dot y_e = \dot y_u \frac{h}{t_s}\,,\;\;\;\;
{y'_e} = {y'_u} \frac{h}{\alpha_s}\,.
\end{equation}
\noindent Ad esempio, si pu\`o notare che per la legge
ad accelerazione costante che si compie in un tempo unitario e fornisce
una alzata $h$ unitaria, si ha accelerazione massima ${{}\ddot {y}}_{\rm max}=4$,
come riportato nella figura \ref{fig:f_legge_acc_cost}.
Chiamando questa grandezza, cio\`e la massima accelerazione raggiunta
da una legge unitaria, {\em coefficiente di accelerazione}\index{coefficiente!di accelerazione} $c_a$, avremo per l'effettiva accelerazione
fornita da una legge di moto qualsivoglia
\begin{equation}
{{}\ddot{y}}_{\rm max}= c_a {h\over t_s^2}\;\;\; {\rm e}\;\;\;\;  
{{y''}}_{\rm max}= c_a {h\over \alpha_s^2}\,.  
\label{e_ca}
\end{equation}
\noindent \`E facile osservare che per una legge di moto di forma qualsiasi,
che fornisca
al suo termine $\dot {y}=0$ e che sia simmetrica rispetto alla met\`a dell'angolo sul quale si svolge, risulta $c_a \geq 4$, essendo il valore quattro
riservato proprio al caso di accelerazione costante simmetrica.
Sempre richiamando la figura \ref{fig:f_legge_acc_cost} notiamo che la velocit\`a
massima raggiunta dal cedente \`e $\dot {y}_{\rm max}=2$. Questo valore
\`e il {\em coefficiente di velocit\`a}\index{coefficiente!di velocit\`a}
 $c_v$ di questa legge di moto. 
Analogamente a quanto esposto per le accelerazioni si ha
\begin{equation}
{{}\dot{y}}_{\rm max}= c_v {h\over t_s}\;\;\; {\rm e}\;\;\;\;  
{{y'}}_{\rm max}= c_v {h\over \alpha_s}\,.  
\label{e_cv}
\end{equation}
\noindent Il coefficiente di velocit\`a, che qualifica la legge di moto rispetto
alla velocit\`a massima raggiunta, \`e altrettanto importante nella sintesi
delle camme in quanto, come vedremo tra poco, influenza direttamente
una grandezza di primaria importanza nella trasmissione del movimento.
In figura \ref{fig:f_acc_cost_tagliata} sono riportati i grafici di accelerazione
velocit\`a e spostamento per una legge cosiddetta con {\em accelerazione
costante tagliata}\index{legge!accelerazione costante} dove si nota
che, rispetto alla legge ad accelerazione costante, il coefficiente di
accelerazione cresce, mentre diminuisce quello di velocit\`a.
\begin{figure}[hbt]
\centering
\begin{minipage}[b]{0.49\textwidth}
\centering
\includegraphics[width=0.9\textwidth]{part2/camme/FIG/generic_law/legge_basso_cv.pdf}
\begin{picture}(0,0)(150,-10)
\scriptsize{
}
\put(-17,75){$\ddot y$}
\put(-17,65){$\dot y$}
\put(-17,55){$y$}
\put(2,86){${{}\ddot{y}}_{\rm max}=44.9$}
\put(40,56){${{}\dot{y}}_{\rm max}=1.03$}
\put(140,-9){$\alpha$}
\put(134,-9){$t,$}
\end{picture}
      \caption{\em Legge di moto a basso coefficiente di velocit\`a.}
 \label{fig:f_legge_basso_cv}
\end{minipage}\hfill
\begin{minipage}[b]{0.49\textwidth}
\centering
\includegraphics[width=0.9\textwidth]{part2/camme/FIG/generic_law/legge_acc_sin_tagliata.pdf}
\begin{picture}(0,0)(150,-10)
\scriptsize{
\put(-17,75){$\ddot y$}
\put(-17,65){$\dot y$}
\put(-17,55){$y$}
\put(35,74){${{}\ddot{y}}_{\rm max}=6.54$}
\put(47,56){${{}\dot{y}}_{\rm max}=1.66$}
\put(140,-9){$\alpha$}
\put(134,-9){$t,$}
}
\end{picture}
	\caption{\em Legge di moto con accelerazione sinusoidale tagliata.}
     \label{fig:f_acc_sin_tagliata}
\end{minipage}
\end{figure}
\noindent In figura \ref{fig:f_legge_basso_cv} \`e riportata una legge che presenta 
un bassissimo coefficiente di velocit\`a $c_v = 1.03$ (il limite inferiore teorico \`e uno),  
a discapito dei valori dell'accelerazione che risultano essere elevati, con
coefficiente di accelerazione $c_a=44.9$.
La legge con {\em accelerazione sinusoidale tagliata}\index{legge!accelerazione
sinusoidale} di figura \ref{fig:f_acc_sin_tagliata} soddisfa le esigenze
del progettista che \`e preoccupato delle vibrazioni che possono insorgere
durante il funzionamento della camma e che sono influenzate da cambiamenti repentini
di accelerazione: \cite{ruggieri}, pag. 45.
\begin{wrapfigure}{r}{0.5\textwidth}
     \begin{center}
     \includegraphics[width=0.42\textwidth]{part2/camme/FIG/generic_law/legge_trapezoidale.pdf}
     \end{center}
\begin{picture}(0,0)(-29,-35)
\scriptsize{
}
\put(-17,75){$\ddot y$}
\put(-17,65){$\dot y$}
\put(-17,55){$y$}
\put(-5,81){${{}\ddot{y}}_{\rm max}=7.67$}
\put(40,52){${{}\dot{y}}_{\rm max}=1.53$}
\put(129,-8){$\alpha$}
\put(123,-8){$t,$}
\end{picture}
\vskip -6mm
        \caption{\em Legge trapezoidale tagliata (asimmetrica).}
     \label{fig:f_legge_trapezoidale}
\end{wrapfigure}
\noindent La derivata terza dello spostamento,
grandezza che peraltro compare assai raramente nella meccanica, \`e in certo
modo responsabile delle vibrazioni che insorgono nella catena cinematica
del cedente durante il funzionamento delle camme veloci. La legge sinusoidale
citata poc'anzi presenta bassi valori di {\em jerk}\index{jerk} (cos\`i viene
talvolta chiamata $\dddot{y}(t)$). Purtroppo tale legge presenta, a parit\`a di
altre circostanze, valori di $c_v$ pi\`u elevati della legge ad
accelerazione costante.
\noindent In un'ottica di compromesso si potrebbe modificare la legge ad accelerazione costante rendendo le sue due ``gobbe'' di forma trapezoidale, come in figura 
\ref{fig:f_legge_trapezoidale},  passando
cos\`i da valori molto alti di {\em jerk} (teoricamente infiniti) a valori finiti
e controllabili.
Anche qui per\`o non \`e tutto rose e fiori: le rampe inclinate di questa legge di moto
si traducono in una partenza della salita poco decisa e, quel che \`e peggio,
in una difficolt\`a materiale di realizzazione in officina di queste rampe con
precisione tale da garantire l'esattezza della legge di moto
dato che, per un certo tratto, si discostano troppo poco dalla circonferenza.
\begin{figure}[t]
\centering
\begin{minipage}[b]{0.49\textwidth}
\centering
\includegraphics[width=0.9\textwidth]{part2/camme/FIG/generic_law/legge_7tratti.pdf}
\begin{picture}(0,0)(150,-10)
\scriptsize{
\put(-17,75){$\ddot y$}
\put(-17,65){$\dot y$}
\put(-17,55){$y$}
\put(43,72){${{}\ddot{y}}_{\rm max}=5.08$}
\put(48,57){${{}\dot{y}}_{\rm max}=1.66$}
\put(140,-9){$\alpha$}
\put(134,-9){$t,$}
}
\end{picture}
      \caption{\em Legge di moto con accelerazione trapezoidale modificata.}
 \label{fig:f_legge_7tratti}
\end{minipage}\hfill
\begin{minipage}[b]{0.49\textwidth}
\centering
\includegraphics[width=0.9\textwidth]{part2/camme/FIG/generic_law/legge_7tratti_non_simm.pdf}
\begin{picture}(0,0)(150,-10)
\scriptsize{
\put(-17,75){$\ddot y$}
\put(-17,65){$\dot y$}
\put(-17,55){$y$}
\put(42,72){${{}\ddot{y}}_{\rm max}=4.86$}
\put(48,60){${{}\dot{y}}_{\rm max}=1.59$}
\put(38,4){${{}\ddot{y}}_{\rm min}=-7.6$}
\put(140,-9){$\alpha$}
\put(134,-9){$t,$}
}
\end{picture}
      \caption{\em Legge di moto con accelerazione asimmetrica.}
     \label{fig:f_legge_7tratti_non_sim}
\end{minipage}
\vskip -3mm
\end{figure}
\noindent Una legge di moto raffinata e spesso adottata per la sintesi delle camme, \`e senza dubbio la cosiddetta {\em legge
di moto trapezoidale modificata}\index{legge!trapezoidale modificata}, figura \ref{fig:f_legge_7tratti}.
Si compone di sette tratti: sinusoide, costante, sinusoide, nulla, sinusoide,
costante, sinusoide. Ciascuno di questi intervalli pu\`o essere lungo a
piacere permettendo in tal modo di dare origine eventualmente a leggi non simmetriche, come quella di figura \ref{fig:f_legge_7tratti_non_sim}. 

\definecolor{dark-green}{rgb}{0,0.5,0}
\begin{wrapfigure}{r}{0.6\textwidth}
\begin{center}
\includegraphics[width=0.55\textwidth]{part2/camme/FIG/generic_law/legge_acc_cost_e_7tratti_confronto.pdf}
\end{center}
\begin{picture}(0,0)(0,0)
\put(5,135){\color{red}$\ddot y$}
\put(5,120){\color{dark-green}$\dot y$}
\put(5,105){\color{blue}$y$}
\put(190,26){$\alpha$}
\put(184,26){$t,$}
\end{picture}
\vskip -6.5mm
      \caption{\em Confronto tra legge ad accelerazione costante e legge a sette tratti con lo stesso tempo di taglio.}
 \label{fig:f_confronto_cost_7}
\end{wrapfigure}
\noindent In figura \ref{fig:f_confronto_cost_7} si riporta,
a parit\`a di alzata e di lunghezza del ``taglio'', il confronto tra
una legge ad accelerazione costante e una legge trapezoidale modificata.
Sebbene una tale analisi
sia gi\`a stata illustrata in figura \ref{fig:f_confronto_accelerazioni},
si desidera qui esplicitarla di nuovo tramite l'impiego di due leggi molto
pi\`u usate nella pratica, in quanto la presenza del taglio assicura $c_v$
meno elevati: caratteristica, questa, molto apprezzabile, come vedremo.
\noindent Rimarchiamo ancora una volta la sostanziale equivalenza delle due leggi a livello
del diagramma delle alzate, praticamente indistinguibili nella figura.
L'inconveniente presentato dalla legge trapezoidale modificata
di avere un $c_a$ leggermente
pi\`u elevato \`e ben assorbito dalla maggiore ``dolcezza di movimento'' che essa 
garantisce.
\noindent Infine riportiamo nelle figure
\ref{fig:f_legge_completa_7_a} e
\ref{fig:f_legge_completa_7_h} rispettivamente il diagramma delle accelerazioni completo per progettare la nostra camma e il corrispondente diagramma delle alzate.
Tale diagramma, gi\`a illustrato in figura
\ref{fig:f_diagramma_alzate}, qui ottenuto per\`o
partendo da leggi di accelerazione trapezoidali modificate tagliate \`e
il candidato idoneo a essere ``avvolto'' sul tamburo della camma.
\begin{figure}[hbt]
\centering
\begin{minipage}[b]{0.49\textwidth}
\centering
\includegraphics[width=0.9\textwidth]{part2/camme/FIG/generic_law/legge_completa_7tratti_a.pdf}
\hfill
\begin{picture}(0,0)(130,0)
\scriptsize{
\put(105,2){$\alpha$}
\put(89,97){$T$}
\put(105,89){$t$}
\put(-46,86){$y''$}
}
\end{picture}
      \caption{\em Legge completa con tratti di  accelerazione trapezoidali modificate.}
     \label{fig:f_legge_completa_7_a}
\end{minipage}
\hfill
\begin{minipage}[b]{0.49\textwidth}
\centering
\includegraphics[width=0.9\textwidth]{part2/camme/FIG/generic_law/legge_completa_7tratti_h.pdf}
\begin{picture}(0,0)(116,0)
\scriptsize{
\put(105,2){$\alpha$}
\put(89,97){$T$}
\put(105,89){$t$}
\put(-58,86){$y(\alpha)$}
}
\end{picture}
      \caption{\em Legge completa delle alzate ottenuta da tratti di  accelerazione trapezoidali modificate.}
     \label{fig:f_legge_completa_7_h}
\end{minipage}
\end{figure}
\noindent Un'ultima considerazione circa la progettazione delle leggi di moto
dovrebbe mettere in luce la delicatezza di questo processo e l'atteggiamento
di compromesso necessario. Diamo uno sguardo alla seconda legge di accelerazione
presente nella figura \ref{fig:f_legge_completa_7_a}.
Tale legge produce in un angolo di $45^{\circ}$ un'alzata di $30\, {\rm mm}$ e,
delle tre, \`e quella con le caratteristiche pi\`u esasperate. Essa si presenta
pi\`u
``squadrata'' rispetto alle altre in modo da contenere i valori di $c_{a^+}$
(coefficiente di accelerazione considerando $\ddot y_{\rm max}$)
e di $c_{a^-}$
(coefficiente di accelerazione considerando $\ddot y_{\rm min}$). Tali valori sarebbero ancora inferiori se si eliminasse
il taglio ma questa scelta peggiorerebbe il valore dell'angolo di pressione 
$\theta$ aumentandolo e compromettendo, come vedremo a breve, il buon
funzionamento del meccanismo. 

\noindent Da un diverso punto di vista,
 contenere il valore di $c_{a^-}$ significa aumentare il valore del
minimo raggio di curvatura convessa $\rho_{\rm min}$.
Questo risulta essere un beneficio, come cercheremo di spiegare pi\`u avanti;
tuttavia,
eccedere in tal senso (aumentando il tempo su cui si estende la ``gobba''
negativa dell'accelerazione) produrrebbe il pessimo risultato di
aumentare $c_{a^+}$, parametro che governa
le forze d'inerzia durante l'alzata.

\vskip 1mm
\null
\section{Profilatura delle Camme}\index{camme}

\noindent Come accennato all'inizio di questo capitolo, la legge di spostamento
$y(\alpha)$ 
viene (salvo casi particolari) avvolta su di un tamburo in modo da generare
una camma piana.
Nel processo di sintesi della camma le grandezze  in gioco sono molte e tra
queste il raggio del tamburo test\'e citato, che si chiama {\em raggio di base}\index{raggio
di base} e che verr\`a indicato con $R_b$, gioca il  ruolo pi\`u importante.
All'aumentare del raggio di base scompaiono infatti dalla
camma le problematiche principali di profilatura,
cui accenneremo fra poco.  Di contro,  col crescere di $R_b$,
aumentano le dimensioni della camma stessa rendendola 
scomoda da utilizzare nel progetto organico di una macchina\footnote{Ci preme
precisare 
che il raggio di base pu\`o essere riferito a entit\`a geometriche diverse
dal profilo fisico della camma. In particolare, vedremo che le camme tastate
da una rotella presentano un importantissimo luogo geometrico ``virtuale'', il
{\em profilo del
centro rotella}, che coincide con il profilo di una camma tastata da una punta che esegue le stesse leggi di moto.
In questo caso $R_b$ sar\`a il raggio della circonferenza minima
tangente al tracciato del centro rotella, in quanto difficilmente
risulter\`a utile tirare in ballo il raggio di base del profilo fisico.
Pertanto, nonostante
molteplici trattazioni pi\`u autorevoli della nostra introducano per
il raggio di base del profilo del centro rotella  il simbolo $R_{b0}$, siamo convinti
di non creare equivoci accontentandoci di usare soltanto $R_b$.
}.
Si intravede quindi un percorso di progetto 
che ha come obiettivo quello di trovare una soluzione
che, evitando eventuali criticit\`a 
derivanti da dimensioni geometriche mal scelte (quasi sempre
valore di $R_b$ insufficiente),
soddisfino anche requisiti di compattezza, limitazione
dell'ingombro
 e funzionalit\`a.
La figura \ref{fig:f_ang_press} ci aiuta a riconoscere una delle grandezze
geometriche 
fondamentali di un meccanismo a camma: l'{\em angolo di pressione} $\theta(
\alpha)$, funzione i cui valori, in generale,
variano nell'arco dell'angolo giro.
Notiamo subito che $\theta$ \`e anche l'angolo compreso tra la retta tangente al
profilo nel punto di contatto col cedente e una retta generica
ortogonale al cedente stesso. 
\begin{figure}[bht]
\centering
\includegraphics[width=0.8\textwidth]{part2/camme/FIG/ang_press.pdf}
\begin{picture}(70,0)(130,0)
\scriptsize{
\put(222,202){$\theta$}
\put(208,183.5){${\rm d}y$}
\put(168,167){${\rm d}\alpha (R_b+y)$}
\put(168,111){$\theta$}
\put(55,99){\rotatebox{90}{{$R_b + y$}}}
\put(150,82){$\omega$}
\put(100,33){$R_b$}
}
\end{picture}
      \caption{\em Angolo di pressione $\theta$.}
 \label{fig:f_ang_press}
\end{figure}
La forza $F$  che la camma deve trasmettere al cedente,
soggetto a una spinta assiale $S$, risulta
essere $F=S/\cos(\theta)$, relazione dalla quale si evince l'importanza
di mantenere l'angolo $\theta$ entro stretti limiti.
Angoli $\theta > 30^{\circ}-35^{\circ}$ non sono
in genere accettabili nei tratti di salita, dove cio\`e la camma spinge la
punteria, riservando il valore limite pi\`u elevato a soluzioni
che prevedono guide di scorrimento per la punteria con attriti particolarmente
ridotti, \cite{molian}, pagg. 131-132. \`E frequente, nell'ambito delle macchine
automatiche, l'utilizzo di punterie costituite da colonne 
a sezione circolare, rettificate, che scorrono in manicotti a ricircolo di
sfere con attrito molto limitato.
Pertanto poniamo il massimo angolo di
pressione tollerato al valore $\theta_{\rm lim}=35^{\circ}$, osservando che angoli
di pressione maggiori genererebbero una forza ortogonale al cedente troppo
elevata col pericolo di impuntamento di quest'ultimo, oppure di una trasmissione
del movimento poco efficiente.
Poniamoci dal punto di vista relativo a un osservatore solidale con la camma.
In tale riferimento la punta percorrer\`a, durante una rotazione infinitesima
${\rm d}\alpha$ dell'albero, un tratto infinitesimo della tangente nel punto di contatto.
Questo tratto infinitesimo si scompone facilmente sui cateti 
del triangolo rettangolo
${\rm d}y$ e ${\rm d}\alpha (R_b +y)$, dalla quale scomposizione risulta
\begin{equation}
\tan(\theta)={{\rm d}y\over{(R_b+y)}{\rm d}\alpha}= {y'\over{R_b+y}}\,.
\label{e_ang_press}
\end{equation}
\noindent Affinch\'e il meccanismo funzioni correttamente il valore massimo di questa espressione dovr\`a essere contenuto al di
sotto del valore prescritto. Purtroppo \`e alquanto probabile che il massimo valore di $y'(\alpha)$ e il minimo di $y(\alpha)$ non si verifichino per lo
stesso angolo $\alpha$, quindi 
occorrerebbe cercare il massimo della funzione \ref{e_ang_press}. 
\noindent Ricordiamo che $y'_{\rm max}=c_v h/\alpha_s$, cio\`e il massimo
valore della velocit\`a geometrica \`e dato dal coefficiente di velocit\`a
della legge di moto moltiplicato per l'alzata e diviso per l'angolo di salita.
In via cautelativa si pu\`o asserire che affinch\'e la trasmissione del moto
dalla camma al cedente non sia critica dovr\`a essere 
\begin{equation}
\tan(\theta_{\rm lim})> {c_v h\over{\alpha_s(R_b+y)}}\,.
\label{e_ang_press_cv}
\end{equation}
\noindent Cautelandoci ulteriormente si pu\`o porre  $y(\alpha)=0$ e scrivere
\begin{equation}
\tan(\theta_{\rm lim})> {c_v h\over{\alpha_s R_b}}\,,
\label{e_ang_press_max}
\end{equation}
\noindent ottenendo in questo modo un'indicazione per la scelta del raggio di base:
\begin{equation}
R_b > {c_v h\over{\alpha_s \tan(\theta_{\rm lim})}}\,.
\label{e_R_b}
\end{equation}
\noindent Dalla \ref{e_R_b} emerge appieno l'importanza di utilizzare
leggi di moto a basso valore di coefficiente di velocit\`a: da tale
grandezza dipende infatti linearmente la dimensione e quindi l'ingombro della camma.
Si pu\`o anche osservare come alzate $h$ importanti
basate su angoli di salita $\alpha_s$ modesti comportino
elevati valori del raggio di base. A questo proposito chiariamo
che nell'ambito delle  macchine automatiche \`e tutt'altro che infrequente
lavorare con salite che si svolgono in $20^{\circ}$ o $30^{\circ}$. 
\noindent La figura \ref{fig:f_camma_punteria} riporta una camma con {\em cedente a
punteria}\index{cedente!a punteria}. La punta (o coltello visto che le camme hanno
spessore diverso da zero) tasta, in questo caso, direttamente il profilo della camma e segue
la legge di figura \ref{fig:f_legge_completa_7_h}.
Diciamo per inciso che, tanto la camma di figura \ref{fig:f_camma_punteria} quanto
tutte le altre che seguiranno nelle varie illustrazioni, sono state
ottenute da questa medesima legge di moto che trae origine dal progetto di una 
pressa reale, tuttora funzionante in una fabbrica cartotecnica.
\noindent Molto raramente il profilo della camma viene per\`o tastato 
da una punta come nelle figure \ref{fig:f_ang_press} e \ref{fig:f_camma_punteria}.
Per evidenti questioni tribologiche si preferisce affidare lo scorrimento
del cedente al rotolamento di una rotella sopra un secondo profilo
tracciato in modo che 
il centro di tale rotella segua 
il profilo generato dalla legge delle alzate.
Normalmente, le rotelle sono  cuscinetti a sfere o a rulli appositamente
costruiti per questo scopo.
\vskip 2mm
\begin{figure}[hbt]
\centering
\begin{minipage}[b]{0.49\textwidth}
\centering
{\includegraphics[width=0.9\textwidth]{part2/camme/FIG/camma/camma_punteria.pdf}}
\hfill
%\hspace{\fill}
\begin{picture}(0,0)(150,0)
\scriptsize{
\put(0,150){${\rho_0}_{\rm min}=46\; {\rm mm,\;}\;\alpha=126^{\circ}$}
\put(9,144){\rotatebox{-54}{$\longrightarrow$}}
\put(39,145){\rotatebox{-77}{$\longrightarrow$}}
\put(43,140){$\theta_{\rm max}=35^{\circ},\; \alpha=103^{\circ}$}
}
\end{picture}
      \caption{\em Camma tastata da una punta, $R_b=80\;{\rm mm}$.}
 \label{fig:f_camma_punteria}
\end{minipage}
\hfill
\begin{minipage}[b]{0.46\textwidth}
\centering
\includegraphics[width=0.9\textwidth]{part2/camme/FIG/camma/inviluppo.pdf}
\begin{picture}(0,0)(148,-2)
\scriptsize{
\put(5,135){${\rho}_{\rm min}=26\; {\rm mm,\;}\;\alpha=126^{\circ}$}
\put(49,99){\rotatebox{-234}{$\longrightarrow$}}
\put(41.5,85.2){\rotatebox{-112}{\color{yellow}$\scriptscriptstyle{\rm circ}$}}
\put(40.9,74.9){\rotatebox{-94.5}{\color{yellow}$\scriptscriptstyle{\rm onf}$}}
\put(40.7,65){\rotatebox{-81}{$\color{yellow}\scriptscriptstyle{\rm ere}$}}
\put(42,56.5){\rotatebox{-70}{\color{yellow}$\scriptscriptstyle{\rm nza}$}}
\put(46,46){\rotatebox{-55}{\color{yellow}$\scriptscriptstyle{\rm di}$}}
\put(50.5,39.7){\rotatebox{-39}{\color{yellow}$\scriptscriptstyle{\rm  base}$}}
}
\end{picture}
      \caption{\em Profilo della camma inviluppato dal moto relativo della rotella di raggio  $r_r=20\; {\rm mm}$.}
     \label{fig:f_inviluppo}
\end{minipage}
\end{figure}
\noindent In figura \ref{fig:f_inviluppo} viene riportata la camma con il suo profilo
effettivo ottenuto dall'inviluppo
della rotella il cui centro segue il tracciato blu sul quale \`e posizionata,
cio\`e il {\em profilo del centro rotella}\index{profilo centro rotella}.
\noindent Le figure \ref{fig:f_camma_punteria} e \ref{fig:f_inviluppo} evidenziano un'altra
grandezza di elevato
rilievo nel progetto di una camma: 
la {\em curvatura}\index{curvatura!del profilo} del suo profilo.
Nelle due figure appena citate sono infatti riportati i valori minimi
di $\rho_0(\alpha)$ e $\rho(\alpha)$, che sono rispettivamente il raggio di curvatura
del profilo seguito dal centro della rotella e il raggio di curvatura del profilo fisico
della camma, valori variabili ovviamente da punto a punto.
Come nel caso dell'angolo di pressione, anche il controllo del raggio di curvatura \`e
di importanza fondamentale durante la fase di progetto.
Il profilo inviluppato dalla rotella, cio\`e l'effettivo profilo
della camma, presenter\`a raggi di curvatura  che saranno dati da $\rho=\rho_0-r_r$,
dove $r_r$ rappresenta il raggio della rotella stessa. \`E chiaro quindi che  
il raggio di curvatura del profilo effettivo pu\`o diventare
nullo o persino negativo nel caso in cui il raggio della rotella sia maggiore 
del minimo raggio di curvatura delle zone convesse 
del profilo del centro rotella.
Perci\`o, valori di $\rho_0$ eccessivamente bassi, nei tratti convessi di
tale profilo, possono causare le situazioni di {\em sottotaglio}\index{sottotaglio}
rappresentate in figura \ref{fig:f_sottotaglio},
e contrassegnate con le lettere $A$ e $B$. In tali zone si verifica che l'inviluppo
generato dalla rotella porta a un profilo intrecciato a coda di rondine,
quindi impossibile da seguire con un normale cedente.
\begin{figure}[t]
\hbox{
\vspace{-2cm} \hspace{.6cm}
\includegraphics[width=0.3\textwidth]{part2/camme/FIG/camma/sottotaglio_dettaglio.pdf}
}
%\rule[-4cm]{5pt}{6cm}
\hspace{-6cm}
\includegraphics[width=1.4\textwidth]{part2/camme/FIG/camma/sottotaglio.pdf}

\begin{picture}(0,0)(-100,0)
\scriptsize{
\put(-31,245){A}
\put(64.5,175){A}
\put(137,154){B}
}
\end{picture}
      \caption{\em Camma con evidente sottotaglio, raggio rotella $r_r=58\;{\rm mm}$.}
     \label{fig:f_sottotaglio}
\end{figure}
\noindent Va per\`o detto che occorre tenere il valore di ${\rho}_{\rm min}$
ben lontano anche dallo zero, cio\`e dallo spigolo vivo che si verificherebbe
se il centro della rotella seguisse  una curva di raggio uguale a quello della rotella stessa.
La trasmissione delle forze, a volte ingenti, tra camma e punteria a rotella si basa
sul contatto tra una generatrice del ``cilindroide'' della camma e una del
cilindro della rotella stessa. Le pressioni di contatto che si generano in questi casi
si possono ricavare dalla {\em teoria di Hertz} per i contatti tra solidi elastici.
In particolare, nel nostro caso il contatto avviene tra due cilindri aventi raggio
$\rho$ per la camma e $r_r$ per la rotella, quindi possiamo affermare che la pressione di contatto
$p_c$ sar\`a \cite{timoshenko}, pagg. 418-419\footnote{Questo testo ha rappresentato per  
molto tempo un utile riferimento per tutti i problemi inerenti
il comportamento elastico dei solidi.},
\begin{equation}
p_c \propto \sqrt {{\rho + r_r}\over {\rho r_r}}\,,
\label{eq_hertz}
\end{equation}
\noindent che esplicita la proporzionalit\`a inversa tra tale pressione e la radice
quadrata del prodotto dei due raggi. Risulta quindi chiaro che il desiderio
di non far nascere pressioni troppo elevate, le quali potrebbero compromettere
la forma della
camma deformandola plasticamente, ci porta a desiderare un raggio minimo per il
profilo percorso dalla rotella ben lontano dallo zero
(lo spigolo). Raramente si progettano camme in cui si trovi nelle convessit\`a
del profilo un raggio minore alla met\`a di quello della rotella che lo tasta.
A questo proposito, mostriamo in figura \ref{fig:f_camma_rotella} una camma
correttamente profilata, la quale rappresenta con pi\`u
dettagli la stessa di figura \ref{fig:f_inviluppo}. Si nota che il valore dell'angolo
di pressione massimo non cambia rispetto alla situazione di cedente a punta,
\cite{molian}, pag. 61,
e il raggio minimo di curvatura del profilo risulta
accettabile.
Anche qui, come abbiamo fatto per l'angolo di pressione, ci si potrebbe incamminare alla ricerca
del legame che intercorre tra le leggi di moto, il raggio di base della camma, il raggio della rotella e  
il minimo raggio di curvatura (convessa) che si otterr\`a sul profilo.
Il volume \cite{ruggieri} riporta, a mio avviso, una delle pi\`u ampie e utili
trattazioni della sintesi di meccanismi a camma. Su questo libro si trova,
a pagina 59, una formula che lega il raggio del percorso del centro della rotella
con i valori del raggio di base, della salita $y(\alpha)$ e delle sue derivate, prima e seconda.
Si ha per il raggio del percorso del  ``centro rotella''\index{curvatura!formula}
\begin{equation}
\rho_0={{[{y'}^2+(R_b + y)^2]}^{3/2}\over{(R_b + y)^2-(R_b +y)y'' +2{y'}^2}}\,.
\label{r_cur}
\end{equation}
\noindent Pertanto elevati valori di accelerazione negativa generano raggi di curvatura modesti, come anche 
eccessivi valori per la velocit\`a $y'(\alpha)$, mentre, ancora una volta, a
un aumento del raggio di base $R_b$ corrisponde un aumento del raggio di curvatura
del profilo, ma come sempre anche un poco desiderabile aumento dell'ingombro. 
\begin{figure}[h]
\centering
\includegraphics[width=0.9\textwidth]{part2/camme/FIG/camma/camma_rotella.pdf}
\begin{picture}(0,0)(235,-72)
\scriptsize{
\put(0,150){${\rho}_{\rm min}=26\; {\rm mm,\;}\;\alpha=126^{\circ}$}
\put(24,133){\rotatebox{126}{$\longrightarrow$}}
\put(69,124){\rotatebox{103}{$\longrightarrow$}}
\put(58,112){$\theta_{\rm max}=35^{\circ}, \alpha=103^{\circ}$}
}
\end{picture}
      \caption{\em Camma correttamente profilata, $R_b=80\, {\rm mm}$, $r_r=20\, {\rm mm}$.}
 \label{fig:f_camma_rotella}
\end{figure}
\noindent Nonostante ricavare la \ref{r_cur} non presenti difficolt\`a insormontabili, abbiamo
preferito evitare la sua  dimostrazione che sarebbe risultata lunghetta anzich\'e no,
e avrebbe aggiunto poco alla comprensione del problema che stiamo trattando.
\noindent Oggigiorno la profilatura delle camme si esegue totalmente
con tecniche numeriche e l'individuazione dei punti cospicui, cio\`e quelli che presentano il valore pi\`u elevato dell'angolo
di pressione e il minore raggio di curvatura, risulta essere piuttosto semplice:
normalmente i programmi di calcolo li rendono facilmente rintracciabili
evidenziandoli.
Le figure \ref{fig:ap_camma_rotella} e \ref{fig:cur_camma_rotella} dovrebbero
dare
l'idea della variazione delle due grandezze fondamentali, $\theta$ e $\rho$, per
la camma di figura \ref{fig:f_camma_rotella}. Mentre il grafico di
$\theta$ \`e contenuto entro precisi limiti, quello di $\rho$ pu\`o assumere
valori assoluti molto grandi ($\rho$ \`e infinito nei punti di flesso
del profilo), pertanto una parte di tale grafico \`e escluso dalla figura.
\begin{figure}[b]
\centering
\begin{minipage}[b]{0.48\textwidth}
\centering
\includegraphics[width=0.9\textwidth]{part2/camme/FIG/camma/ck_ang_press_camma_punteria.pdf}
\begin{picture}(0,0)(130,0)
\scriptsize{
}
\put(127,2){$\alpha$}
\put(-31,86){$\theta$}
\end{picture}
      \caption{\em Andamento di $\theta$ per la camma di figura \ref{fig:f_camma_rotella}.}
 \label{fig:ap_camma_rotella}
\end{minipage}\hfill
\begin{minipage}[b]{0.48\textwidth}
\centering
\includegraphics[width=0.9\textwidth]{part2/camme/FIG/camma/ck_cur_camma_rotella.pdf}
\begin{picture}(0,0)(129,0)
\scriptsize{
\put(127,2){$\alpha$}
\put(-38,86){$|\rho|$}
}
\end{picture}
      \caption{\em Andamento di $|\rho|$ per la camma di figura \ref{fig:f_camma_rotella}.}
     \label{fig:cur_camma_rotella}
\end{minipage}
\end{figure}
\noindent Mediante opportune maggiorazioni e minorazioni dei termini della
\ref{r_cur} si pu\`o ottenere, anche in questo caso, 
una formula di progetto per $R_b$, come mostrato in \cite{ruggieri}, pag. 72, formula che non riportiamo.
Aggiungiamo, tuttavia, qualche considerazione
qualitativa circa l'individuazione del raggio di base, mettendolo maggiormente
in luce nella \ref{r_cur}, riscrivendola nel seguente modo
\begin{equation}
\rho_0={{{{[{{y'}^2\over{(R_b + y)^2}}+1]}^{3/2}(R_b+y)}\over{{1-{y''\over{(R_b + y)}}} +{2{y'}^2\over{(R_b + y)^2}}}}}\,.
\label{r_cur1}
\end{equation}
\noindent Notiamo che, a un aumento del raggio di base $R_b$, il valore assoluto del 
termine del denominatore che contiene $y''$
diminuisce, e con esso
cala la possibilit\`a che il denominatore (nel caso di valori negativi di
$y''$) diventi troppo grande e tale da compromettere il 
raggio di curvatura della camma.
Ma soprattutto osserviamo una squisita ``quasi proporzionalit\`a''
tra $\rho_0$ e $R_b$ essendo $\rho_0\approxprop (R_b +y)$, la quale proporzionalit\`a ci garantisce che all'aumentare del raggio di base spariscono
i problemi legati a raggi di curvatura troppo piccoli del profilo.
Perci\`o, aumentando il raggio di base ci togliamo
sia dall'impaccio di angoli di pressione troppo elevati
sia dal pericolo di curvature eccessive del profilo della camma.


\noindent Nonostante la presenza di una formula che pu\`o guidarci
durante il progetto verso un valore di $R_b$ rispettoso dei limiti da imporre
alla curvatura,
la sintesi delle camme, come si \`e detto, si esegue sempre pi\`u spesso 
per tentativi ripetuti, mediante codici dedicati
che rendono veloce ed efficace il confronto tra soluzioni diverse: a un angolo di
pressione massimo troppo grande o a un raggio di curvatura del profilo eccessivamente
ridotto segue un nuovo tentativo con $R_b$ aumentato, oppure anche con modifiche
(tollerabili) alle leggi di moto.

\noindent Anche se in questa sede abbiamo ritenuto 
di fornire soltanto un'infarinatura generale che
renda chiari i concetti e le difficolt\`a che stanno alla base della progettazione
delle camme, in un perimetro di studi pi\`u vasto si trova qualche
altro ``trucco'' che permette
di non aumentare, oltre il ragionevole, l'ingombro di questi eccentrici e allo stesso
tempo di non eccedere coi valori dell'angolo di pressione. 
Uno di tali {\em escamotage} \`e quello di costruire un dispositivo camma-cedente in cui
l'asse di quest'ultimo non passi per il centro della camma stessa.
Il disassamento del cedente \`e in qualche modo equivalente a inclinare il cedente di un dato angolo che 
si sottrarr\`a all'angolo di pressione durante la salita, sommandosi invece a
questo nei tratti di discesa. Una discreta percentuale delle camme
che si realizzano sono di questo tipo: {\em camme a punteria deviata}\index{punteria deviata} oppure, con terminologia anglosassone,
{\em camme offset}\index{camme!offset}.
Un breve cenno alla progettazione di questi dispositivi si trova nel
capitolo di approfondimento \ref{cam2}.

\begin{figure}[b]
\centering
\begin{minipage}[b]{0.68\textwidth}
\centering
\includegraphics[width=0.9\textwidth]{part2/camme/FIG/cambil.pdf}
\begin{picture}(0,0)(150,-20)
\scriptsize{
}
\put(35,143){$\beta(t)$}
\put(-29,60){$\omega$}
\end{picture}
        \caption{\em Camma con cedente a bilanciere.}
     \label{fig:f_cambil}
\end{minipage}\hfill
\begin{minipage}[b]{0.28\textwidth}
\centering
     \includegraphics[width=0.9\textwidth]{part2/camme/FIG/campiat.pdf}
\begin{picture}(0,0)(70,-20)
\scriptsize{
}
\put(7,214){$y(t)$}
\put(10,45){$\omega$}
\end{picture}
        \caption{\em Camma tastata da un piattello.}
     \label{fig:f_campiat}
\end{minipage}\hfill
\end{figure}

\noindent Le camme possono azionare cedenti che non traslano, come mostrato
in figura \ref{fig:f_cambil}.
Sono diffusissimi meccanismi in
cui la camma impone al
cedente a bilanciere un moto rotatorio.
Mutando opportunamente
ci\`o che nel ragionamento si deve mutare e in particolare sostituendo l'alzata $y(t)$ con l'angolo
che descrive il bilanciere nella sua rotazione, $\beta(t)$, tutto quanto esposto circa le precauzioni
nel progetto della camma rimane invariato: le leggi di moto con i loro
 coefficienti di velocit\`a e di accelerazione si scelgono con gli stessi criteri e
nella profilatura della camma si devono fare i soliti conti con curvatura e angolo di 
pressione. Anche qui, rimandiamo il lettore desideroso di approfondire un poco la
progettazione delle camme tastate da cedente a bilanciere al capitolo
\ref{cam2}, oppure a \cite{ruggieri}, pag. 74.


\noindent La trasmissione del movimento tra camma e cedente non avviene sempre tramite il rotolamento
di una rotella sul profilo della camma. Bench\'e tastare una camma con un {\em piattello}\index{piattello}, come in figura \ref{fig:f_campiat}, sia ormai raro nell'ambito delle macchine automatiche, in altri 
domini della meccanica la soluzione camma-piattello \`e molto usata.
Anzi, a dire il vero, vi \`e un vastissimo repertorio di camme tastate esclusivamente
da piattelli che, strisciando sul profilo, in questo caso
rigorosamente convesso di quest'ultime,
aziona cedenti a punteria o a bilanciere: le camme dei motori endotermici a quattro
tempi. Si tratta in genere di camme di dimensioni modeste alle quali calza naturalmente
tutta la teoria qui esposta. 

\noindent Una breve introduzione alla progettazione
delle camme tastate da piattelli si trova nel capitolo di
approfondimento \ref{cam2} e, per esteso, in \cite{ruggieri}, pag. 65.
 I problemi legati alla progettazione di queste camme
sono per\`o di natura cos\`i
diversa da quelli che si incontrano nelle macchine automatiche, la loro progettazione
segue dettami cos\`i intimamente legati alla legge delle alzate, a sua volta legata
al percorso dei gas nei condotti fluidodinamici che esse aprono e chiudono, che
ci sembra meglio lasciare questi piccoli e sofisticati eccentrici al loro
mondo, cio\`e quello motoristico, e tornare alle nostre camme per macchine
automatiche.

\begin{figure}[t]
\centering
\begin{minipage}[b]{0.48\textwidth}
\centering
\includegraphics[width=0.9\textwidth]{part2/camme/FIG/camma/camma_desmo.pdf}
\begin{picture}(0,0)(130,0)
\scriptsize{
}
\end{picture}
      \caption{\em Camma desmodromica.}
 \label{fig:f_camma_desmo}
\end{minipage}\hfill
\begin{minipage}[b]{0.48\textwidth}
\centering
\includegraphics[width=0.9\textwidth]{part2/camme/FIG/camma/desmo_costola.pdf}
\begin{picture}(0,0)(129,0)
\scriptsize{
}
\end{picture}
      \caption{\em Camma a costola.}
     \label{fig:f_desmo_costola}
\end{minipage}
\end{figure}

\section{Camme Coniugate}

\noindent L'accoppiamento, cio\`e il contatto tra camma e cedente, pu\`o essere
 garantito nelle fasi di discesa\footnote{
Ribadiamo che il termine discesa si riferisce al movimento contrario
a quello generato dalla camma mentre compie il suo lavoro funzionale.
Esso non si riferisce
quindi in alcun modo all'avvicinamento alla Terra.
} in due diversi modi:
imponendo che la punteria prema sulla camma durante tutto l'angolo giro, oppure
facendo in modo, mediante una camma accoppiata a quella primaria, che la punteria
sia costretta a seguire il profilo di questa anche nei tratti che presentano $\dot y$ 
negative.
Nel primo caso si parla di {\em accoppiamento di forza}\index{accoppiamento!di forza}
e molto spesso tale forza viene esercitata da una molla (quasi sempre
pre-caricata) che, deformata
durante le fasi di salita, garantisce mediante la sua elasticit\`a il contatto
anche durante la discesa. Questa soluzione \`e raramente impiegata nell'ambito
delle macchine automatiche, l'autore non ne ha mai vedute (ma il mondo
\`e grande...), mentre trova largo impiego nelle piccole camme ad uso motoristico.
Nel secondo caso si impone il contatto mediante un {\em accoppiamento di forma}\index{accoppiamento!di forma}, cio\`e progettando una camma ``negativa'' detta
{\em camma coniugata}\index{camme!coniugata} o {\em camma desmodromica}\index{camme!desmodromica}, riportata in colore rosa
nella figura \ref{fig:f_camma_desmo}.
La generazione di questa camma potrebbe in linea di principio seguire il percorso
per la progettazione della camma ``primaria'' avendo, in questo caso, l'avvertenza di invertire i
segni delle leggi di moto. Di fatto, una volta disegnata la camma primaria,
baster\`a
considerare il prolungamento del cedente fino al luogo ove si desidera tastare 
la camma desmodromica e riportare una seconda rotella.

\begin{figure}[hbt]
\centering
\includegraphics[width=0.85\textwidth]{part2/camme/FIG/camma/rotella_interna.pdf}
\begin{picture}(0,0)(130,0)
\scriptsize{
}
\end{picture}
      \caption{\em Camma tastata dall'anello interno di un cuscinetto.}
 \label{fig:f_rotella_interna}
\end{figure}


\noindent Tale 
rotella, durante il movimento generato dalla camma primaria (sottintendendo il contatto
costante tra cedente e profilo),  ``taglier\`a'' la camma desmodromica o {\em
camma coniugata}\index{camme!coniugata}, e questo modo di 
procedere \`e anche quello utilizzato dal nostro codice
per produrre le illustrazioni \ref{fig:f_camma_desmo} e 
\ref{fig:f_desmo_costola}. Delle due citate figure, la seconda rappresenta
una soluzione molto elegante di accoppiamento
desmodromico, dove le rotelle tastano una costola. Purtroppo, questa
soluzione non \`e
di facilissima realizzazione in officina. Normalmente le camme coniugate
sono pi\`u sottili delle relative primarie e le loro rotelle sono di
minori dimensioni.
Questo \`e dovuto alla funzione ``di servizio'' da loro svolta,
che semplicemente consiste nel mantenere il contatto tra la punteria e la
camma primaria
e nel riportare il cedente nella posizione di partenza: esse
svolgono molto raramente azioni di lavoro. Persino le cautele sull'angolo di
pressione massimo e sul minimo raggio di curvatura convessa,
ammesse per queste camme 
(parliamo sempre delle fasi attive), possono essere leggermente rilassate.

\noindent Anche la fantasia pu\`o partecipare con successo nella progettazione delle camme.
Riportiamo, in figura \ref{fig:f_rotella_interna}, 
una soluzione poco comune di camma-punteria dove
il profilo della camma viene tastato dall'anello interno di un cuscinetto
(a rulli). In tale figura, il meccanismo viene rappresentato nel momento in
cui la camma ha svolto $180^{\circ}$ della sua rotazione e, stanti le leggi
proposte per questa camma, la punteria si trova alla massima estensione 
dell'alzata.
Le due curvature, quella della camma e quella della rotella,
sono, in questo caso, concordi. Pertanto,
nella \ref{eq_hertz} uno dei due
raggi cambia segno e la pressione andrebbe addirittura
ad annullarsi nel caso (da non contemplarsi) di raggi di curvatura
uguali tra loro.
\begin{figure}[hbt]
\centering
\begin{minipage}[b]{0.48\textwidth}
\centering
\includegraphics[width=0.9\textwidth]{part2/camme/FIG/albero.pdf}
\begin{picture}(0,0)(130,0)
\scriptsize{
}
\end{picture}
      \caption{\em Albero a camme per pressa di grandi dimensioni, tastato dall'anello interno di un cuscinetto.}
 \label{fig:f_albero}
\end{minipage}\hfill
\begin{minipage}[b]{0.48\textwidth}
\centering
\includegraphics[width=0.9\textwidth]{part2/camme/FIG/rotella.pdf}
\begin{picture}(0,0)(129,0)
\scriptsize{
}
\end{picture}
      \caption{\em Rotella costituita da cuscinetto a rulli che tasta la camma mediante il suo anello interno.}
     \label{fig:f_rotella}
\end{minipage}
\end{figure}
\noindent La pressione di contatto risulta quindi
notevolmente ridotta in questo tipo di camme che possono essere usate quando le forze da trasmettere 
sono ingenti. Nelle figure \ref{fig:f_albero} e \ref{fig:f_rotella} riportiamo una realizzazione
della tipologia di camme appena descritte e di una delle quattro rotelle,
con il desiderio di mostrare al lettore quali ragguardevoli 
dimensioni possono essere raggiunte da un albero a camme e da una rotella.

\endinput



\part{Approfondimenti}

\chapter{Moto Rigido Tridimensionale}\label{3d}
\section{Il Moto nelle Tre Dimensioni dello Spazio}
\begin{wrapfigure}{r}{0.45\textwidth}
\centering{
      \includegraphics[width=0.40\textwidth]{part3/3d/FIG/f11_app11.pdf}}
\begin{picture}(0,0)(118,23)
\scriptsize{
\put(37.5,285){$P$}
\put(-11,221){$P$}
\put(-18,133){$P$}
\put(14,29){$\bm {G} \equiv \bm \Gamma$}
\put(74,161.5){${\bm \Lambda}$}
\put(94,190){${\bm J}$}
\put(45,100){${{\bm s}}$}
\put(-14,165){${\bm s}_{\scriptscriptstyle{P}}$}
\put(-3,89){${\bm z}_{\scriptscriptstyle{P}}$}
\put(53,249){${\bm z}_{\scriptscriptstyle{P(J)}}$}
\put(20,197){${\bm z}_{\scriptscriptstyle{P(\Lambda)}}$}
\put(1,22){${{\bm e}_1}$}
\put(50,32){${{\bm e}_2}$}
\put(18,74){${{\bm e}_3}$}
\put(84,167){${{\bm e}_1}_{R}$}
\put(85,206){${{\bm e}_2}_{R}$}
\put(45,187){${{\bm e}_3}_{R}$}
}
\end{picture}
	\caption{\em Moto tridimensionale di una clavetta.}
     \label{fig:f11_app11}
\end{wrapfigure}
Il moto di un corpo  nello
spazio tridimensionale rappresenta il
caso generale della {\em cinematica
del corpo rigido}\index{cinematica del corpo rigido}.
Lo si studia con lo scopo,
come avviene nel caso piano, di mettere in relazione tra loro
le varie grandezze cinematiche che vi partecipano,
riferendole, in generale, a un osservatore fisso.
Costui sar\`a nel seguito surrogato
da una {\em terna orientata}\index{terna!orientata} 
$\bm{G}$.
Ci occuperemo, in prima battuta (e in coda), di uno spostamento finito,
del quale sar\`a protagonista
il nostro corpo rigido, che nelle figure \`e rappresentato da una clavetta.
Tale spostamento
far\`a passare la clavetta da una posizione iniziale, nella quale
la terna $\bm \Gamma$, ad essa rigidamente collegata, coincide con
la terna $\bm G$, a una seconda posizione,
nella quale la terna $\bm \Gamma$
si trover\`a in un altro punto dello spazio e avr\`a 
un diverso orientamento: in questa nuova posizione, ci riferiremo ad essa
col nome ${\bm \Lambda}$.

\noindent Consideriamo le terne 
$\bm G$
e $\bm \Gamma$
costituite ciascuna 
da tre versori ortogonali, 
come in figura \ref{fig:f11_app11},
che si possono  
assumere quali {\em vettori di base} per
gli spazi fisso e mobile.
Come abbiamo gi\`a accennato, la terna $\bm{ \Gamma}$, vincolata alla clavetta,
accompagner\`a quest'ultima nei suoi spostamenti e un osservatore
ad essa solidale vedr\`a il punto $P$, appartenente alla clavetta stessa, 
stare fermo.
Evidenziamo che, una volta nota la posizione dell'origine di $\bm \Gamma$,
si potr\`a conoscere la posizione di tutti i punti della clavetta, cio\`e
dello spazio mobile,
conoscendo la posizione di un suo punto generico,
che abbiamo gi\`a indicato con $P$. 
Consideriamo ora lo spostamento generico della clavetta
come idealmente composto da due distinte fasi.
Durante la prima di queste due fasi portiamo semplicemente 
la terna $\bm \Gamma$ a sovrapporsi alla terna $\bm J$.
La seconda fase dello spostamento
vede invece il vettore 
${{\bm z}}_{\scriptscriptstyle{P(J)}}$ andare a coincidere con 
${{\bm z}}_{\scriptscriptstyle{P(\Lambda)}}$, e $\bm J$
stessa sovrapporsi a $\bm \Lambda$, lasciando fisso il centro di $\bm J$.
Pertanto, queste due operazioni
forniscono una ``ricetta''
per conoscere la posizione ``assoluta'' di $P$.
Ma esse presentano difficolt\`a di grado diverso.
Infatti, per portare la terna $\bm\Gamma$ nella terna $\bm J$
baster\`a aggiungere a tutti i punti
dello spazio mobile,
riferiti a $\bm G$, il vettore
${\bm{s}}$.  Invece, per orientare gli assi
di $\bm J$ in modo che essi si sovrappongano a quelli
di $\bm \Lambda$ si intuisce che 
occorrer\`a
operare mediante ``opportune rotazioni'' della stessa terna $\bm J$:
ma quali sono queste rotazioni?
Lasciamo questa domanda in sospeso e rispolveriamo qualche strumento
di cui ci siamo dotati
seguendo il corso di Algebra Lineare.  

\noindent Consideriamo la trasformazione
del vettore 
${{\bm z}}_{\scriptscriptstyle{P(J)}}$ che lo porta a coincidere col vettore
${{\bm z}}_{\scriptscriptstyle{P(\Lambda)}}$.
A tale scopo, chiamiamo con $\bm R$ un operatore di trasformazione dei 
vettori nello spazio tale per cui 
\begin{equation}
{{\bm z}}_{\scriptscriptstyle{P(\Lambda)}}=
{\bm R}{\bm z}_{\scriptscriptstyle{P(J)}}\,.
\label{e11_app11}
\end{equation}
\noindent Affinch\'e $\bm{R}$ rappresenti una rotazione rigida, la sua applicazione a
due vettori qualsiasi dovr\`a 
lasciare inalterato l'angolo tra le direzioni dei due vettori stessi;
inoltre essa non dovr\`a n\'e dilatare
n\'e restringere i loro moduli. Tutto questo, sinteticamente, si esprime dicendo
che la trasformazione $\bm R$
 deve mantenere invariante il {\em prodotto scalare}
tra due vettori qualsiasi, prima e dopo la sua applicazione\footnote
{
Ci siamo riferiti all'operatore $\bm R$ come, appunto, a un'operazione
sui vettori dello spazio tridimensionale. A questo operatore, come anche
ai vettori sui quali esso opera, compete ovviamente
una rappresentazione matriciale
legata ai vettori di base: matrici $1\times3$ saranno le rappresentazioni dei vettori,
mentre la rappresentazione dell'operatore $\bm R$ consister\`a in una matrice
$3\times3$. Speriamo di non ingenerare dubbi nei lettori in conseguenza della
scelta di mantenere inalterati i nomi dei vettori, degli operatori e
delle loro rappresentazioni.
}.
Quindi, posto che, tramite $\bm{R}$,
 $\bm{v} \rightarrow {\bm{v}}_R$ e  $\bm{u} \rightarrow {\bm{u}}_R$,
possiamo scrivere l'uguaglianza (tra le rappresentazioni di questi oggetti)
\begin{equation}
{{\bm{u}}_R}^T{\bm{v}}_R={\bm{u}}^T{\bm{R}}^T{\bm{R}}{\bm{v}}\,.
\label{e12_app11}
\end{equation}
\noindent Ma per quanto appena affermato circa l'invarianza del prodotto scalare tra
due vettori dovr\`a essere
\begin{equation}
\bm{R}^T\bm{R}={\bm{I}}\,.
\label{e13_app11}
\end{equation}
\noindent Quando la trasposta (coniugata, nel caso di matrici complesse) di una matrice si comporta come la sua inversa
sappiamo, sempre dall'Algebra Lineare, che tale
matrice 
\`e {\em unitaria} e che il suo determinante vale uno. Vale poi anche
\begin{equation}
\bm{R}\bm{R}^T={\bm{I}}\,.
\label{e14_app11}
\end{equation}
\noindent Pertanto, se $\bm{R}$ rispetta la \ref{e13_app11}, essa ruota lo spazio a cui si applica
e lascia invariati gli angoli e le dimensioni degli oggetti.
Restringendo le nostre considerazioni a uno spazio con solo due 
dimensioni, giusto per proporre qualcosa di facilmente verificabile
 a quei lettori
poco desiderosi di accordare tutta questa confidenza all'Algebra Lineare,
tale matrice assume la forma
\begin{equation}
{\bm R}=\left( \begin{array}{cc}
\cos \theta & -\sin \theta\\ 
\sin \theta & \cos \theta
\end{array}
\right)
\,,
\label{e15_app11}
\end{equation}
 \noindent dove
$\theta$ rappresenta l'unico possibile angolo di rotazione nel piano.
\noindent La situazione \`e concettualmente analoga, bench\'e 
sia pi\`u complessa, nello spazio
tridimensionale. 
In questo caso infatti la matrice\footnote{
Sarebbe ${\bm R}^{-1}$ ad essere cos\`i composta, ma
${\bm R}^T={\bm R}^{-1}$.
}
${\bm{R}}^T$,
formata dalle colonne della
rappresentazione dei vettori $\left({{\bm e}_1}_{R},{{\bm e}_2}_{R},
{{\bm e}_3}_{R}\right)$ sulla base $\left(\bm{e}_1, \bm{e}_2, \bm{e}_3\right)$,
\`e tale per cui
\begin{equation}
({{\bm e}_1}_{R},{{\bm e}_2}_{R}, {{\bm e}_3}_{R})
= {\bm R}\cdot
(\bm{e}_1, \bm{e}_2, \bm{e}_3)
\,,
\label{eq:rot}
\end{equation}
\noindent e, dato un ben preciso orientamento di ${\bm \Lambda}$,
essa pu\`o essere ottenuta seguendo diversi percorsi.
Fatta salva quindi l'unicit\`a della matrice $\bm R$, esistono diverse
composizioni di rotazioni successive della terna $\bm J$ che la portano
a coincidere con $\bm \Lambda$. Tali procedimenti di rotazione
tengono naturalmente conto degli angoli via via coinvolti e prendono
il nome dal loro primo proponente ovvero dal particolare ambito nel quale il
loro uso \`e prevalente.
Abbiamo cos\`i gli {\em angoli di Eulero}\index{angoli!di Eulero}: {\em
precessione}\index{precessione}, {\em nutazione}\index{nutazione} e 
{\em rotazione}\index{rotazione}, normalmente indicati coi simboli
$\theta$, $\phi$  e $\psi$,
\cite{whittaker}, pagg. 8--12, gli {\em angoli di Cardano}\index{angoli!di Cardano}, meno noti e non molto dissimili da quelli di {\em Eulero} e talvolta
chiamati col nome di {\em angoli nautici}\index{angoli!nautici} e
spesso indicati coi nomi
di {\em imbardata}\index{imbardata}, {\em rollio}\index{rollio} e 
{\em beccheggio}\index{beccheggio}, e altri ancora.

\section{Rotazioni Infinitesimali}

\noindent Se restringiamo il campo delle rotazioni a quelle
infinitesimali, che ormai sappiamo essere
strettissime parenti delle velocit\`a angolari, una volta rapportate 
al tempo infinitesimale durante il quale esse si manifestano,
le cose marciano con maggior semplicit\`a.
Consideriamo la \ref{e11_app11} e differenziamola:
\begin{equation}
{\rm d}{\bm{z}}_{\scriptscriptstyle{P(\Lambda)}}=
\left({\rm d}\bm{R}\right){\bm{z}}_{\scriptscriptstyle{P(J)}}+
\bm{R}\left({\rm d}{\bm{z}}_{\scriptscriptstyle{P(J)}}\right)\,.
\label{e16_app11}
\end{equation}
\noindent Si nota che, in virt\`u della condizione di corpo rigido, 
${\rm d}{\bm{z}}_{\scriptscriptstyle{P(J)}}=0$, perch\'e la posizione del
punto $\bm P$ nel sistema di riferimento mobile con la clavetta \`e costante.
Dobbiamo quindi osservare attentamente il termine 
$\left({\rm d}\bm{R}\right){\bm{z}}_{\scriptscriptstyle{P(J)}}$,
nel quale compare il differenziale della matrice di rotazione $\bm R$. 
Premettiamo che le prossime mosse potrebbero apparire un po' artificiose,
bench\'e il loro
senso sia contenuto nel seguente {\em modus operandi}: costruire una
``matrice derivata'' che operi sul vettore che identifica il punto $P$
{\em gi\`a ruotato}, 
${\bm{z}}_{\scriptscriptstyle{P(\Lambda)}}$,
anzich\'e su
${\bm{z}}_{\scriptscriptstyle{P(J)}}$.
\begin{wrapfigure}{r}{0.45\textwidth}
\centering{
      \includegraphics[width=0.38\textwidth]{part3/3d/FIG/f12_app11.pdf}
}
\begin{picture}(40,25)(18,10)
\scriptsize{
\put(8,87){${\bm{e'}_1}$}
\put(93,114){${\bm{e'}_2}$}
\put(38,185){${\bm{e'}_3}$}
\put(14,134){${\rm d}\phi_{{\scriptscriptstyle{1,2}}}$}
\put(74,137){${\rm d}\phi_{{\scriptscriptstyle{3,1}}}$}
\put(-9.5,69.5){${\rm d}\phi_{{\scriptscriptstyle{1,2}}}x$}
\put(4,50){${\rm d}\phi_{{\scriptscriptstyle{3,1}}}x$}
\put(-33,69){$\bm x$}
}
\end{picture}
\vskip -10mm
	\caption{\em Rotazioni infinitesimali.}
     \label{fig:f12_app11}
\end{wrapfigure}
Per la propriet\`a di $\bm R$ espressa da \ref{e13_app11}, riscriviamo
la \ref{e16_app11} interponendo tra i due fattori la matrice identit\`a
$\bm{R}^T\bm{R}$ ottenendo cos\`i la seguente relazione:
\begin{equation}
{\rm d}{\bm{z}}_{\scriptscriptstyle{P(\Lambda)}}=
\left({\rm d}\bm{R}\right)\bm{R}^T\bm{R}{\bm{z}}_{\scriptscriptstyle{P(J)}}
\,,
\label{e17_app11}
\end{equation}
\noindent cio\`e
\begin{equation}
{\rm d}{\bm{z}}_{\scriptscriptstyle{P(\Lambda)}}=
\left({\rm d}\bm{R}\right)\bm{R}^T{\bm{z}}_{\scriptscriptstyle{P(\Lambda)}}
\,.
\label{e18_app11}
\end{equation}
\noindent Ma derivando la \ref{e14_app11}, notiamo che
\begin{equation}
\left({\rm d}\bm{R}\right)\bm{R}^T+
\bm{R}\left({\rm d}\bm{R}\right)^T=0\,,
\label{e19_app11}
\end{equation}
\noindent o anche

\vskip 1mm
\begin{equation}
\left({\rm d}\bm{R}\right)\bm{R}^T=
-\bm{R}\left({\rm d}\bm{R}\right)^T\,,
\label{e110_app11}
\end{equation}
\noindent La matrice $\left({\rm d}\bm{R}\right)\bm{R}^T$ \`e quindi una matrice
antisimmetrica e infinitesimale (prodotto di una matrice infinitesimale
per una finita) che chiameremo 
${\rm d}\bm{\Phi}=\left({\rm d}\bm{R}\right)\bm{R}^T$. \`E facile convincersi che tale matrice
contiene i tre angoli infinitesimali
di rotazione nei tre piani individuati da ciascuna coppia di vettori di base di
$\bm \Lambda$: 
%$\left(\bm{e'}_1, \bm{e'}_2, \bm{e'}_3\right)$:
\begin{equation}
{\rm d}{\bm\Phi}=\left( \begin{array}{ccc}
0 & -{\rm d}\phi_{{\scriptscriptstyle{1,2}}} & {\rm d}\phi_{{\scriptscriptstyle{3,1}}}\\
{\rm d}\phi_{{\scriptscriptstyle{1,2}}} & 0 & -{\rm d}\phi_{{\scriptscriptstyle{2,3}}}\\
-{\rm d}\phi_{{\scriptscriptstyle{3,2}}} & {\rm d}\phi_{{\scriptscriptstyle{2,3}}} & 0\\
\end{array}
\right)
\,.
\label{e111_app11}
\end{equation}
\noindent Ai suoi elementi abbiamo dato nomi evocativi e segni disparati: conviene
ora fare una prova per renderci conto se tali nomi e tali segni siano o meno
opportuni.
Applichiamo la matrice ${\rm d}\bm \Phi$ al vettore $\bm x$ di modulo costante rappresentato
in figura \ref{fig:f12_app11}, la cui giacitura coincide con quella del
versore $\bm{e'}_1$ con lo scopo di trovare ${\rm d}\left(\bm{x}\right)$.
Avremo
\begin{equation}
{\rm d}
 \left( \begin{array}{c}
x \\ 
0 \\ 
0 \\
\end{array}
\right)=
\left( \begin{array}{ccc}
0 & -{\rm d}\phi_{{\scriptscriptstyle{1,2}}} & {\rm d}\phi_{{\scriptscriptstyle{3,1}}}\\
{\rm d}\phi_{{\scriptscriptstyle{1,2}}} & 0 & -{\rm d}\phi_{{\scriptscriptstyle{2,3}}}\\
-{\rm d}\phi_{{\scriptscriptstyle{3,1}}} & {\rm d}\phi_{{\scriptscriptstyle{2,3}}} & 0\\
\end{array}
\right)
 \left( \begin{array}{c}
x \\ 
0 \\ 
0 \\
\end{array}
\right)=
 \left( \begin{array}{c}
0 \\
{\rm d}\phi_{{\scriptscriptstyle{1,2}}}x \\
-{\rm d}\phi_{{\scriptscriptstyle{3,1}}}x \\
\end{array}
\right)
\,,
\label{e112_app11}
\end{equation}
\vskip 2mm
\noindent che, confrontato con la figura \ref{fig:f12_app11}, ci conforta alquanto.
Il procedimento si pu\`o ripetere (con successo) per i vettori che giacciano
sulle altre due direzioni, cos\`i avremo provato ``alla buona'' che
gli elementi della matrice
${\rm d}\bm \Phi$ sono le rotazioni infinitesime.
Dividendo poi tale matrice per ${\rm d}t$ si ottiene la matrice (finita) delle
velocit\`a angolari
\begin{equation}
{\bm\Omega}=\left( \begin{array}{ccc}
0 & -\omega{{\scriptscriptstyle{1,2}}} & \omega{{\scriptscriptstyle{3,1}}}\\
\omega{{\scriptscriptstyle{1,2}}} & 0 & -\omega{{\scriptscriptstyle{2,3}}}\\
-\omega{{\scriptscriptstyle{3,1}}} & \omega{{\scriptscriptstyle{2,3}}} & 0\\
\end{array}
\right)
\,.
\label{e113_app11}
\end{equation}
\noindent Siamo ora in grado di rispondere al nostro quesito iniziale, cio\`e di conoscere
posizione e velocit\`a del punto $P$
sapendo come si muove la terna $\bm \Lambda$. Per gli spostamenti avremo
\begin{equation}
{\bm s}_{\scriptscriptstyle{P}}=
{{\bm s}}+
{\bm R}{\bm z}_{\scriptscriptstyle{P}}-
{\bm z}_{\scriptscriptstyle{P}}\,,
\label{e114_app11}
\end{equation}
\noindent e per quanto riguarda le velocit\`a
\begin{equation}
{\bm v}_{\scriptscriptstyle{P}}=
{\bm v}_{\scriptscriptstyle\Lambda}+
{\bm \Omega}{\bm z}_{\scriptscriptstyle{P}}\,.
\label{e115_app11}
\end{equation}

\section{Asse del Mozzi e Altre Considerazioni}
\noindent Ci chiediamo a questo punto se, come nel caso piano, vi siano dei punti
o dei luoghi privilegiati che,
a fronte del movimento tridimensionale della clavetta,
non subiscono spostamenti, oppure siano sottoposti
soltanto a movimenti particolarmente semplici.
Come ora sappiamo, l'equazione \ref{e114_app11} fornisce lo spostamento di un
generico punto dello {\em spazio mobile}.
Imponendo che lo spostamento di un punto incognito,
che chiameremo
$Q$, sia nullo
\begin{equation}
0=
{{\bm s}}+
{\bm R}{\bm z}_{\scriptscriptstyle{Q}}-{\bm z}_{\scriptscriptstyle{Q}}\,,
\label{e116_app11}
\end{equation}
\noindent potremmo trovare un punto fisso della trasformazione considerata.
Ma la possibilit\`a di risolvere la \ref{e116_app11} con soddisfazione,
fuori cio\`e da circostanze molto particolari, \`e remota: basti pensare al
caso semplicissimo (che implica una traslazione) dove la matrice di
rotazione $\bm R\equiv \bm I$, caso in cui il sistema di equazioni
summenzionato risulta impossibile per
${{\bm s}}\ne0$.
Osserviamo per\`o che nella \ref{e116_app11} compare una somma notevole:
$ {\bm R}{\bm z}_{\scriptscriptstyle{Q}}-{\bm z}_{\scriptscriptstyle{Q}}$.
Come sappiamo, possono esistere alcuni vettori ${\bm u}_i$, che sono
gli {\em autovettori}\index{autovettore} della
matrice $\bm R$, che soddisfano sempre la relazione 
\begin{equation}
{\bm R}{\bm u}_i-\lambda{\bm u}_i=0\,.
\label{e117_app11}
\end{equation}
\noindent Asseriamo, senza dimostrarlo\footnote
{
Si consideri che, se $\bm u$ e $\lambda$ soddisfano l'equazione
${\bm R}{\bm u}= \lambda {\bm u}$, si ha 
$|| {\bm u}||^2 = {\bm u}^T{\bm R}^T{\bm R}{\bm u}= \lambda^2 {\bm u}^T{\bm I}{\bm u}$.
Tenendo poi conto che $\bm R$ \`e una matrice unitaria,
abbiamo $|\lambda|^2=1$. Un'equazione secolare di terzo grado ammette
sempre almeno una soluzione reale e, affinch\'e sia verificato quanto 
appena asserito, unitamente al fatto che il determinante della
matrice ${\bm R}$ sia positivo, i tre valori di lambda possono essere:
$\lambda_1=\lambda_2=\lambda_3=1$, oppure
$\lambda_1=1,\, \lambda_2=\lambda_3=-1$ o,
ancora, $\lambda_1=1,\, \lambda_2=a+ib,\,\lambda_3=a-ib$, dove $|a+ib|=1$,
e i pedici 1, 2 e 3 possono naturalmente essere scambiati a piacimento.
},
che la matrice di rotazione $\bm R$
possiede sempre
un {\em autovettore} ${\bm u}_3$ reale corrispondente all'{\em autovalore}
 $\lambda_3=  1$. \`E tradizione indicizzare 
tramite il pedice ``tre'' sia l'autovalore reale
sia il corrispondente autovettore reale\footnote
{
Operando in questo modo l'asse dello spostamento elicoidale coincide
con il terzo asse (cio\`e $z'$) di una terna levogira e, in questo
nuovo riferimento, gli spostamenti dovuti alla rotazione avvengono
in piani paralleli e ortogonali a tale asse.
}. Parallela a questo autovettore esiste quindi
una retta i cui punti, a causa dello spostamento tridimensionale finito
dello spazio mobile, traslano, rimanendo per\`o sulla stessa direttrice.
\`E intuitivo pensare che, nei piani ortogonali a tale retta,
lo spostamento dello spazio mobile
sia di rotazione, fornendo  allo spostamento complessivo una possibile
connotazione {\em elicoidale}. Possibile nel senso che uno spostamento
equivalente ``pu\`o'' essere ottenuto tramite un moto elicoidale:
di come l'effettivo
spostamento tra posizione iniziale e finale sia realmente avvenuto,
nulla abbiamo affermato, quindi nulla possiamo dire.

\noindent Rimangono da chiarire due fondamentali circostanze.
La prima riguarda la giacitura della retta parallela all'autovettore
${\bm u}_3$ che, nel possibile spostamento elicoidale,
funge da asse dell'elica.
La seconda circostanza riguarda l'entit\`a della rotazione
dello spostamento elicoidale. L'angolo di rotazione $\xi$
emerge immediatamente
dalla soluzione della \ref{e117_app11}. Infatti, accanto all'autovettore
reale ${\bm u}_3$ associato a $\lambda_3=1$ troviamo altri due autovettori
complessi corrispondenti agli autovalori complessi coniugati
$\lambda_1$ e $\lambda_2$. Affermiamo, ancora una volta senza dimostrarlo,
che 
\begin{equation}
\xi = \atantwo (\mathfrak{Im}( \lambda_1),\mathfrak{Re}( \lambda_1))
\label{eq:comp_rot}
\end{equation}
\noindent rappresenta la componente di rotazione dello spostamento 
elicoidale. 
L'individuazione dell'asse del moto elicoidale risulta
leggermente pi\`u laboriosa. Essa pu\`o essere condotta 
in un piano passante per l'origine della terna fissa e
ortogonale alla direzione dell'asse incognito,
direzione che conosciamo, in quanto coincidente con quella dell'autovettore
reale ${\bm u}_3$ della matrice di rotazione.
Su questo piano, che chiameremo con la lettera greca $\sigma$, 
esister\`a (se la componente
rotatoria dello spostamento non \`e nulla) un punto fisso: tale punto
apparterr\`a all'asse di rotazione, che risulter\`a in questo modo
identificato. Per lavorare nel piano summenzionato conviene adottare
come ``nuovi'' vettori di base i tre autovettori dell'operatore
di rotazione $\bm R$.
Procediamo con ordine.
Gli autovettori dell'operatore di rotazione
$\bm R$ risulteranno, come gi\`a osservato,
${\bm u}_1$ e ${\bm u}_2$ complessi coniugati, mentre ${\bm u}_3$
reale: tutti con norma unitaria.
Associati ai primi due autovettori ci sono due autovalori complessi
coniugati, l'ultimo autovalore vale uno.
La matrice $\bm H$, che contiene, per colonna, i tre autovettori, sar\`a
la matrice che permette di passare dalle rappresentazioni nella
vecchia base dei vettori
e degli operatori, alle loro rappresentazioni nella nuova.
La nuova rappresentazione dell'operatore $\bm R$ 
sar\`a data da
\begin{equation}
{\bm R'}={\bm H}^{-1}{\bm R}{\bm H}={\bm H}^{*}{\bm R}{\bm H}=\left( \begin{array}{ccc}
a+ib & 0 & 0\\
0 & a-ib & 0\\
0 & 0 & 1\\
\end{array}
\right)
\,,
\label{eq:mat_autoval}
\end{equation}
\noindent dove col simbolo $\null^*$ indichiamo la trasposta coniugata.
I valori di $a$ e $b$ compongono i due autovalori complessi di
$\bm R$ e sono legati, come abbiamo gi\`a 
accennato, all'angolo $\xi$ di rotazione elicoidale: $b=\sin(\xi)$.
L'altro oggetto al quale dobbiamo cambiare base \`e la rappresentazione
del vettore ${{\bm s}}$, il quale
separa le origini delle due terne. Nella nuova base risulter\`a
\begin{equation}
{{\bm s}}'={\bm A}^{-1} {{\bm s}}=\
{\bm A}^* {{\bm s}}\,.
\label{eq:sjp}
\end{equation}
\noindent Abbiamo gi\`a fatto cenno alla volont\`a di operare
limitandoci al piano $\sigma$, ortogonale alla direzione di ${\bm u}_3$.
Nella base fornita dagli autovettori di $\bm R$,
il proiettore sul sottospazio generato da ${\bm u}'_3$
sar\`a dato da
\begin{equation}
{{\bm P}_u}_3={\bm u}'_3 ({{\bm u}'_3}^{*} {\bm u}'_3)^{-1}{{\bm u}'_3}^{*}=\
\left( \begin{array}{ccc}
0 & 0 & 0\\
0 & 0 & 0\\
0 & 0 & 1\\
\end{array}
\right)
\,,
\label{eq:op_proiez}
\end{equation}
\noindent mentre il proiettore sul piano $\sigma$, ortogonale 
all'autovettore reale, vale
\begin{equation}
{\bm P}_{\sigma}=({\bm I}-{\bm P})=\left( \begin{array}{ccc}
1 & 0 & 0\\
0 & 1 & 0\\
0 & 0 & 0\\
\end{array}
\right)
\,.
\label{eq:proiettore}
\end{equation}
\noindent Chiamiamo con ${{{\bm s}}}'_{\sigma}={\bm P}_{\sigma}{{\bm s}}'$ 
la proiezione, sul piano $\sigma$, del vettore che distanzia le due terne.
Si osserva senza difficolt\`a che l'applicazione dell'operatore
di proiezione ${\bm P}_{\sigma}$, a qualsiasi vettore (rappresentazione),
oscura l'ultima componente di tale vettore, autorizzandoci pertanto a lavorare
con le sole prime due componenti. Chiamiamo
$\tilde{\bm s}'=({{\bm s}'_{\sigma}}[1],{{\bm s}'_{\sigma}}[2])^T$
la proiezione su $\sigma$ del vettore $\bm s'$ privata dell'ultima
componente nulla.
Allo stesso modo, proiettiamo l'operatore di rotazione
${\bm R}'_{\sigma}={\bm P}_{\sigma}{\bm R}'$.
Anche in questo frangente rileviamo che la proiezione su  $\sigma$ annulla
la terza riga e la terza colonna di ${\bm R}'_{\sigma}$.
Riduciamo quindi la rappresentazione di tale operatore
alla sola rotazione nel piano ortogonale
all'asse dell'elica
\begin{equation}
\tilde{{\bm R}'_{\sigma}}=\left( \begin{array}{cc}
{\bm R}'_{\sigma}[1][1] & \
{\bm R}'_{\sigma}[1][2]\\
{\bm R}'_{\sigma}[2][1] & \
{\bm R}'_{\sigma}[2][2]\\
\end{array}
\right)
\,,
\label{eq:op_rot_2d}
\end{equation}
\noindent ottenendo una matrice $2\times2$.
Indicata quindi 
la posizione incognita dell'asse dell'elica con $\tilde{\bm x}'$, vettore 
bidimensionale che giace su tale piano,
possiamo scrivere
\begin{equation}
\tilde{\bm x}' - \tilde{{\bm R}'_{\sigma}}(-\tilde{\bm x}')={\tilde{{\bm s}}'}_{\sigma}\,,
\label{eq:mozzi}
\end{equation}
\noindent la quale afferma semplicemente che, se ci portassimo nella posizione
$\tilde{\bm x}'$, e applicassimo una rotazione $\tilde{{\bm R}'_{\sigma}}$ all'opposto
di tale vettore, ci ritroveremmo nel punto dove si proietta il vettore
${\bm s}'$. L'incognita $\tilde{{\bm x}}'$ si determina
risolvendo il sistema lineare di due equazioni \ref{eq:mozzi}. Abbiamo
\begin{equation}
\tilde{{\bm x}}'=({{\bm I} -\tilde{\bm R}'_{\sigma}})^{-1}{\tilde{{\bm s}}'}_{\sigma}\,.
\label{eq:sol_mozzi}
\end{equation}
\noindent Torniamo nello spazio tridimensionale, aggiungendo al 
vettore $\tilde{\bm x}'$ la terza componente,
${\bm x}'=((\tilde{\bm x}')^T,0)^T$, e riportiamoci agli originali vettori
di base mediante la trasformazione ${\bm x}={\bm H}{\bm x}'$.
La retta
\begin{equation}
\mu={\bm x} + \alpha {\bm u}_3\,,
\label{eq:asse_m}
\end{equation}
\noindent dove $\alpha$ \`e un numero reale qualsiasi,
coincide con l'asse del moto elicoidale, stabilito dal teorema
di {\em Chasles},
oppure con l'{\em asse del Mozzi}\index{asse!del Mozzi}\footnote
{
Astronomo italiano, autore dell'opera (non visionata dall'autore) {\em
Discorso Matematico Sopra il Rotamento Momentaneo dei Corpi}, 1763.
}.
Riportiamo di nuovo, data la sua importanza,
il contenuto di tale teorema,\index{teorema!di Chasles}
 gi\`a enunciato in ristretto
per i moti piani: ``Ogni spostamento rigido, non traslatorio, si riduce in
infiniti modi ad uno spostamento roto-traslatorio, sempre per\`o con lo stesso
vettore rotazione del componente rotatorio. Si riduce poi in un unico modo 
ad uno {\em spostamento elicoidale}\index{spostamento!elicoidale}'', \cite{finzi}, pag. 119.

\begin{figure}[hbt]
\centering
\begin{minipage}[b]{0.47\textwidth}
\centering
\includegraphics[width=0.95\textwidth]{part3/3d/FIG/3d/mozzi90.pdf}
\begin{picture}(0,0)(130,0)
\scriptsize{
\put(24,16){${\bm e}_2$}
\put(60,14){${\bm e}_1$}
\put(38,35){${\bm e}_3$}
\put(22,77){${{{\bm e}_2}_{R}}$}
\put(28,96){${{{\bm e}_1}_{R}}$}
\put(48,104){${{{\bm e}_3}_{R}}$}
\put(43,60){${\bm s}$}
\put(18,28){${\bm x}$}
\put(6,54){\rotatebox{90}{$\scriptscriptstyle{{\rm Mozzi}}$}}
\put(9,94){${\mu}$}
}
\end{picture}
\vskip .3cm
      \caption{\em
Asse del Mozzi con ${{\bm s}}=(2,2,4)^T$, $\theta=\psi=0$, $\phi=90^{\circ}$.
Posizione asse ${\bm x}=(0,2,0)^T$, angolo di rotazione elicoidale $\xi=90^{\circ}$.
}
 \label{fig:mozzi90}
\end{minipage}\hfill
\begin{minipage}[b]{0.45\textwidth}
\includegraphics[width=0.95\textwidth]{part3/3d/FIG/3d/mozzi303030.pdf}
\begin{picture}(0,0)(130,0)
\scriptsize{
\put(32,36){${\bm e}_2$}
\put(69,27){${\bm e}_1$}
\put(45,43){${\bm e}_3$}
\put(29,102){${{{\bm e}_2}_{R}}$}
\put(52,110){${{{\bm e}_1}_{R}}$}
\put(72,106){${{{\bm e}_3}_{R}}$}
\put(52,75){${\bm s}$}
\put(12,30){${\bm x}$}
\put(13,53){\rotatebox{60}{$\scriptscriptstyle{{\rm Mozzi}}$}}
\put(35,92){${\mu}$}
}
\end{picture}
      \caption{\em
Asse del Mozzi con ${{\bm s}}=(2,2,4)^T$, $\theta=\psi=\phi=30^{\circ}$.
Posizione asse ${\bm x}=(-1.19, 1.26, 0.64)^T$, angolo di rotazione
elicoidale $\xi=66.45^{\circ}$.
      }
 \label{fig:mozzi303030}
\end{minipage}
\end{figure}

\noindent Le figure \ref{fig:mozzi90} e \ref{fig:mozzi303030} dovrebbero
dare un'idea dell'asse dello spostamento elicoidale $\mu$ per due casi di
spostamento rigido tridimensionale. Il primo caso, quello rappresentato
in figura \ref{fig:mozzi90}, riporta una traslazione della 
terna e una rotazione molto semplice, la quale avviene nel piano
$x-y$, attorno all'asse $z$, ed \`e pari a $90^{\circ}$ in senso
antiorario. In questo caso, la retta $\mu$ giace in una posizione 
ragionevolmente intuitiva, e questo ci serve da conforto dopo tutte
le avventure matematiche affrontate. La situazione rappresentata 
invece in figura \ref{fig:mozzi303030}, a fronte di una traslazione
 dalla terna di riferimento identica a quella del caso precedente,
prevede di applicare a tale terna le rotazioni, 
esprimibili in termini di angoli di Eulero, $\theta=\psi=\phi=30^{\circ}$.
In questo frangente, l'intuizione, almeno di chi scrive, si perde.
\`E vero che, in questa figura, l'asse del {\em Mozzi} somiglia a quello
dei ``mappamondi'' sferici, e il vettore ${\bm x}$ al supporto di tale asse,
ma siamo comunque obbligati ad accordare molta fiducia alla
formulazione matematica che abbiamo utilizzato per 
concedere, alla retta  $\mu$,
le caratteristiche di asse dello spostamento elicoidale.
\vskip .4cm
\footnotesize
\noindent I calcoli eseguiti per disegnare le figure \ref{fig:mozzi90} e
\ref{fig:mozzi303030} ci mettono a disposizione, in versione numerica,
gli oggetti che compaiono nelle relazioni matematiche test\'e esposte, e
ci sembrava un peccato non riportane alcuni, che potrebbero essere di aiuto
a qualche studente volenteroso di percorrere lo stesso cammino.
Proponiamo esclusivamente le quantit\`a legate alla figura
\ref{fig:mozzi303030}, che riteniamo maggiormente significative.
Per gli angoli di Eulero, precessione, nutazione e rotazione, come gi\`a
accennato, abbiamo $\theta =  \psi=\phi=0.52360$ rad. Quindi
\begin{equation}
{\bm R}=\left( \begin{array}{ccc}
0.53349 & -0.80801 &  \phantom{-}0.25000\\
0.80801 &  \phantom{-}0.39952 & -0.43301\\
0.25000 &  \phantom{-}0.43301 &  \phantom{-}0.86603\\
\end{array}
\right)
, 
{\bm H}=\left( \begin{array}{ccc}
\phantom{-}0.62325i & - 0.62325i &  0.47235\\
\phantom{-}0.70711  &   \phantom{-}0.70711  &  0\\
-0.3340i &   \phantom{-}0.33400i &  0.88141\\
\end{array}
\right).
\nonumber
\end{equation}
\noindent Gli autovalori sono  $\lambda_1=0.39952 + 0.91672i$, 
$\lambda_2=0.39952 - 0.91672i$ e $\lambda_3= 1$. Il vettore di traslazione,
che vale ${{\bm s}}=(2,2,4)^T$, nella base costituita dagli autovettori
di ${\bm R}$ vale ${{\bm s}'}=(1.4142 - 0.2445i, 1.4142 + 0.2445i, 3.5889)^T$.
\noindent Dopo le proiezioni su $\sigma$ abbiamo
\begin{equation}
\tilde{\bm s}'_{\sigma}= \left(\begin{array}{c}
1.4142 - 0.2445i\\
1.4142 + 0.2445i\\
\end{array}
\right),
\tilde{\bm R}'_{\sigma}=\left( \begin{array}{cc}
0.39952 + 0.91672i &  0\\
0 & 0.39952 - 0.91672i\\
\end{array}
\right).
\nonumber
\end{equation}
\noindent Infine abbiamo, per la posizione di $\mu$,
$\tilde{\bm x}'=(0.89374 + 0.95725i, 0.89374 - 0.95725i)^T$ e 
${\bm x}=( -1.19322, 1.26394, 0.63945)^T$.

\vskip .4cm

\normalsize
\noindent Le cose non vanno diversamente quando gli spostamenti sono
infinitesimali, e in tal caso l'asse del {\em Mozzi} si chiama anche
{\em asse d'istantanea  rotazione}\index{asse!di istantanea rotazione}.
Le superfici rigate disegnate dalla sequenza
degli assi d'istantanea rotazione, sia nel sistema fisso sia in quello
mobile, godono di interessanti propriet\`a, analoghe a quelle delle polari nel
moto piano, \cite{sesini1}, pag. 86.

\noindent Proponiamo un'ultima considerazione, di importanza marginale in questa gi\`a
troppo lunga digressione.
Anche nel caso del moto in tre dimensioni, come gi\`a visto per
il caso piano, le velocit\`a si possono esprimere surrogando
il prodotto 
${\bm \Omega}{\bm z}_{\scriptscriptstyle{P(\Gamma)}}$ con un prodotto
vettoriale e si pu\`o riscrivere la \ref{e115_app11} come segue
\begin{equation}
{\bm v}_{\scriptscriptstyle{P}}=
{\bm v}_{\scriptscriptstyle\Gamma}+
{\bm \omega}\times {\bm z}_{\scriptscriptstyle{P}}\,,
\label{e120_app11}
\end{equation}
\noindent dove il vettore $\bm \omega$ ha le seguenti componenti
\begin{equation}
{\bm\omega}=\left( \begin{array}{c}
\omega{{\scriptscriptstyle{2,3}}}\\
\omega{{\scriptscriptstyle{3,1}}}\\
\omega{{\scriptscriptstyle{1,2}}}
\end{array}
\right)\,.
\label{e119_app11}
\end{equation}
\noindent L'equivalenza del prodotto matriciale della \ref{e115_app11} col prodotto
vettoriale della \ref{e120_app11} dipende strettamente dalla forma particolare
di $\bm \Omega$, e la verifica di questa affermazione \`e lasciata al paziente lettore.
Rimarchiamo semplicemente che la comodit\`a insita nell'introduzione del vettore
velocit\`a angolare $\bm \omega$ rende il suo uso frequentissimo. 
La possibilit\`a di rappresentare le rotazioni infinitesime, quindi 
le velocit\`a angolari, mediante un vettore \`e ristretta per\`o allo spazio tridimensionale.
Infatti con solo due dimensioni, l'angolo di rotazione \`e unico e occuperebbe
le caselle libere della matrice antisimmetrica $2\times2$
di rotazione.
Ci sarebbe perci\`o troppo spazio in un vettore a due dimensioni. \`E vero che 
abbiamo utilizzato la rappresentazione vettoriale della velocit\`a angolare
anche nel caso di moto piano, ma il vettore $\bm \omega$ risulta non appartenere al piano stesso.
Del resto
risulta cos\`i difficile concepire uno spazio che non sia quello tridimensionale!
Viene sempre la tentazione di trattare gli spazi con numero di dimensioni inferiore
a tre come se fossero immersi nel tridimensionale stesso.
Qualora poi il numero delle dimensioni $n$ dello spazio fosse maggiore
di tre,
si otterrebbe, per la matrice di 
rotazione, un numero di componenti ``libere'' pari a $(n^2 - n)/2$.
Quindi, per uno spazio quadridimensionale si dovrebbero considerare sei componenti di rotazione,
numero decisamente eccessivo per potere
entrare in un vettore quadridimensionale.
\endinput

%\newpage
\chapter{Dentature con Profilo Cicloidale
}\label{ruotecy}
\section{Una Propriet\`a Notevole della Cicloide}

\begin{figure}[hbt]
\centering
\includegraphics[width=.8\textwidth]{part3/ruote/FIG/ruote/epiipo.pdf}
\begin{picture}(0,0)(200,0)
\scriptsize{
\put(16,265){$y$}
\put(165,175){$r_r$}
\put(148,173){$\alpha$}
\put(174,161){$p_e$}
\put(-65,156){$r_1$}
\put(-33,173){$r_2$}
\put(87,93){$\alpha$}
\put(74,107){$r_r$}
\put(52,125){$t$}
\put(51,111){$\tau$}
\put(90,123){$p_i$}
\put(17,112){$O$}
\put(170,110){$x$}
}
\end{picture}
      \caption{\em Epicicloide (rosso) e ipocicloide (indaco); le circonferenze
in verde, epiciclo e ipociclo,
svolgono, rotolando, un angolo di $2\pi$.}
 \label{fig:epiipo}
\end{figure}

\noindent La curve rappresentate in figura \ref{fig:epiipo} 
di colore rosso e indaco si chiamano
rispettivamente {\em epicicloide}\index{epicicloide}
e {\em ipocicloide}\index{ipocicloide} delle circonferenze di raggio $r_1$
e $r_2$, di colori blu e arancione.
La {\em cicloide}\index{cicloide}
pu\`o essere considerata come il caso limite di queste due curve qualora la
circonferenza, sulla quale rotola l'{\em epiciclo}\index{epiciclo},
sia di raggio infinito.
L'epicicloide, rappresentata in rosso,
si ottiene facendo rotolare senza strisciamento la circonferenza verde, che
si chiama epiciclo, oppure rolletta, o anche rotella,
di raggio $r_r$ all'esterno della circonferenza di
raggio $r_1$ (colore blu) e tenendo traccia delle posizioni via via raggiunte
da un punto di tale rolletta, individuato col nome $p_e$.
Lo stesso vale per l'ipocicloide (in indaco), con la differenza che la
rolletta, ancora
di raggio $r_r$, rotola all'interno della circonferenza di raggio $r_2$
e di colore arancione: in questo caso
il punto di cui si tiene traccia si chiama $p_i$.
Abbiamo scelto di rappresentare in figura \ref{fig:epiipo} la porzione delle
due curve che si ottiene tramite  il rotolamento delle due rollette lungo
un tratto di lunghezza pari a quella della loro circonferenza.
Abbiamo altres\`i assunto arbitrariamente che ipociclo ed epiciclo abbiano
lo stesso raggio $r_r$ perch\'e, soltanto in questo caso, epicicloide e
ipocicloide saranno dotate di una interessante propriet\`a, utile nella
Meccanica delle Macchine.
\noindent Al fine di descrivere l'epicicloide tramite equazioni
parametriche, scegliamo come parametro l'angolo $t$. Le coordinate $x(t)$ e
$y(t)$ del generico punto $p_e$ si ricavano quindi tramite semplici proiezioni
\begin{equation}
\begin{cases}
{x_e}\left(t\right)=\left(r_1+r_r\right)\cos\left(t\right) +
r_r\cos\left(\pi - t - \alpha\right) \\
{y_e}\left(t\right)=\left(r_1+r_r\right)\sin\left(t\right)-
r_r\sin\left(\pi -t -\alpha\right) \,.
\end{cases}
\end{equation}
\noindent Tenendo poi conto che l'angolo $\alpha$ si pu\`o esprimere
come $\alpha = \dfrac{r_r}{r_1}t$ si ottiene
\begin{equation}
\begin{cases}
{x_e}\left(t\right)=\left(r_1+r_r\right)\cos\left(t\right) -
r_r\cos\left(\dfrac{r_1+r_r}{r_r}t\right) \\
{y_e}\left(t\right)=\left(r_1+r_r\right)\sin\left(t\right)-
r_r\sin\left(\dfrac{r_1+r_r}{r_r}t\right)\,.
\end{cases}
\label{eq:epicicloide}
\end{equation}
\noindent Nel caso dell'ipocicloide scegliamo come parametro $\tau$,
che \`e un angolo crescente in direzione opposta a $t$\footnote{
\`E usuale trovare le formule per l'ipocicloide scritte in modo analogo
a quanto si fa per l'epicicloide, cio\`e col parametro che cresce nella
direzione positiva degli angoli. In tale modo, il primo tratto della curva
occuper\`a, come avviene per l'epicicloide, il primo quadrante. L'insolita
direzione
crescente, che abbiamo scelto per $\tau$, meglio si presta alla trasformazione
di coordinate che ci
permetter\`a di mettere in relazione le due curve come profili coniugati.
}. Avremo
\begin{equation}
\begin{cases}
{x_i}\left(\tau\right)=\left(r_2-r_r\right)\cos\left(\tau\right) +
r_r\cos\left(\alpha- \tau\right) \\
{y_i}\left(\tau\right)=-\left(r_2-r_r\right)\sin\left(\tau\right)+
r_r\sin\left(\alpha -\tau\right) \,,
\end{cases}
\end{equation}
\noindent quindi
\begin{equation}
\begin{cases}
{x_i}\left(\tau\right)=\left(r_2-r_r\right)\cos\left(\tau\right) +
r_r\cos\left(\dfrac{r_2-r_r}{r_r}\tau\right) \\
{y_i}\left(\tau\right)=-\left(r_2-r_r\right)\sin\left(\tau\right)+
r_r\sin\left(\dfrac{r_2-r_r}{r_r}\tau\right)\,.
\end{cases}
\label{eq:ipocicloide}
\end{equation}
\noindent Desideriamo ora mostrare che queste due curve, epicicloide
e ipocicloide, costituiscono una coppia di profili coniugati quando
le circonferenze di raggi $r_1$ e $r_2$, alle quali debbono pensarsi solidali,
rotolano l'una sull'altra realizzando in tal modo 
le primitive di una coppia cinematica. Per la verit\`a, questa dimostrazione 
\`e gi\`a stata svolta, in modo rigoroso, nel capitolo che tratta le ruote
dentate e, in particolare, nel paragrafo che descrive l'ottenimento
di profili coniugati di assortimento. Tale dimostrazione
\`e riassunta dalle figure
\ref{fig:profili_coniugati_punto},
\ref{fig:profili_coniugati_punto_lagrange} e
\ref{fig:profili_coniugati_punto_euler},
che mostrano la coniugazione dei profili tracciati da un punto
trascinato da una primitiva che rotola su un qualsivoglia numero di altre
primitive.
Ma le dimostrazioni tanto generali, come quella riportata alla pagina
citata, nelle quali le curve in gioco 
sono generiche, dove
i procedimenti matematici tendono piuttosto all'astrazione,
dicevamo, tali dimostrazioni hanno sempre lasciato a chi scrive
un certo grado di insoddisfazione.
\begin{figure}[hbt]
\centering
\includegraphics[width=.8\textwidth]{part3/ruote/FIG/ruote/epiipoconiugate.pdf}
\begin{picture}(0,0)(200,0)
\scriptsize{
\put(31,47){$r_2$}
\put(68,167){$t$}
\put(69,170){\rotatebox{129}{$\longrightarrow$}}
\put(144,114){$\tau$}
\put(140,106){\rotatebox{25}{$\longrightarrow$}}
\put(27,92){$t_{\scriptscriptstyle B}$}
\put(35,71){$\tau$}
\put(30,83){\rotatebox{194}{$\longrightarrow$}}
\put(-85,47){$r_1$}
\put(-18,203){$r_r$}
\put(-41,178){$p_e$}
\put(-38,49){$p_i$}
\put(-5,49){$r_r$}
\put(-18,89){$O$}
\put(120,76){$O_{\scriptscriptstyle A}$}
\put(117,155){$O_{\scriptscriptstyle B}$}
\put(57,213){$O_{\scriptscriptstyle C}$}
\put(127,130){$\phi_{\scriptscriptstyle B}$}
\put(61,83){$A$}
\put(96,112){$B$}
\put(94,104){${p_e}_{\scriptscriptstyle B}\equiv {p_i}_{\scriptscriptstyle B}$}
\put(48,199){$C$}
}
\end{picture}
      \caption{\em La ruota di raggio $r_2$ rotola su quella di raggio $r_1$ inviluppando tramite l'ipocicloide a essa solidale (indaco) la epicicloide (rosso).}
 \label{fig:epiipoconiugate}
\end{figure}
Qui riportiamo una prova
diretta, sviluppata tramite il calcolo, della coniugazione tra le due curve,
l'epicicloide e l'ipocicloide, nel caso in cui il raggio
dell'ipociclo sia uguale a quello dell'epiciclo,
suggerendo, a chi \`e gi\`a pago e convinto della dimostrazione svolta
altrove e poco interessato a questo approfondimento, di saltare
direttamente al prossimo paragrafo.
Mostreremo infatti che, scegliendo di tenere ferma
la circonferenza di raggio $r_1$, dove si svolge la epicicloide
(ma tale scelta \`e ovviamente arbitraria), e facendo rotolare 
su di essa la circonferenza di raggio $r_2$, l'ipocicloide solidale 
alla $r_2$ inviluppa proprio la epicicloide solidale alla $r_1$.
Per chiarire quanto appena affermato introduciamo la figura
\ref{fig:epiipoconiugate}, della quale la descrizione non potr\`a essere
brevissima.
Cominciamo col presentare gli attori che metteremo in campo,
dando qualche indicazione circa il ruolo che essi svolgeranno.
La circonferenza di raggio $r_1$ e centro in $O$ (di colore blu) ospita la porzione di
epicicloide (di colore rosso) svolta dalla rolletta di raggio $r_r$ quando,
rotolando all'esterno della prima, percorre una lunghezza pari alla
propria circonferenza.
Il parametro $t$, dal quale dipendono le formule
\ref{eq:epicicloide}, spazier\`a quindi nell'intervallo
 $t \in[0, 2\pi r_r/r_1]$. Circonferenza blu ed epicicloide rossa saranno
ritenute fisse, aggrappate cio\`e al foglio del libro.
\noindent Una seconda circonferenza, di raggio $r_2$, reca
la porzione di ipocicloide generata da un punto dell'ipociclo
che percorre anch'esso un cammino pari alla lunghezza della propria circonferenza, quindi
il parametro $\tau$ dell'ipocicloide
assumer\`a valori compresi nell'intervallo $\tau\ \in [0, -2\pi r_r/r_2]$.
Tale circonferenza col relativo ipociclo \`e rappresentata in quattro
diverse posizioni
con centri $O$, $O_{\scriptscriptstyle A}$, $O_{\scriptscriptstyle B}$ e $O_{\scriptscriptstyle C}$.
La circonferenza di raggio $r_2$ con centro in ${O}$,
non entrando direttamente nel processo di inviluppo, viene
riportata con tratto meno marcato
insieme alla propria ipocicloide, e posizionata
nel luogo dove vale la formula \ref{eq:ipocicloide}.
Le altre circonferenze di raggio $r_2$ sono invece traslate e rotate
in modo opportuno e come vedremo invilupperanno, tramite l'ipocicloide,
l'epicicloide solidale alla $r_1$.
\noindent Chiamiamo punti
{\em omologhi} sulla epicicloide e sulla corrispondente
ipocicloide le posizioni dei punti $p_e$ e $p_i$ corrispondenti a percorsi 
della rolletta di uguale lunghezza.
Nel nostro caso avremo $r_1 t = r_r\alpha$
e $r_2 \tau= r_r\alpha$ da cui segue immediatamente il legame
tra i due parametri $\tau = \dfrac{r_1}{r_2}t$.
\noindent La posizione della ruota di raggio $r_2$ con centro in
$O_{\scriptscriptstyle A}$ si ottiene, dalla stessa di centro $O$, mediante una
rotazione di un angolo pari a $\pi$, composta con uno spostamento verso
destra pari all'interasse $i=r_1+r_2$.
Pi\`u in generale, affinch\'e la circonferenza mobile e la relativa ipocicloide
si trovino in una posizione precisa come quelle indicate 
in successione con le lettere
$B$ e $C$, sono necessari una rotazione  e uno spostamento
degli assi a cui la \ref{eq:ipocicloide} fa riferimento.
Coerentemente coi simboli gi\`a impiegati nel paragrafo dedicato
ai profili coniugati, pag. \pageref{prof_con}, scriviamo la seguente
trasformazione per i punti della ipocicloide che la conduce nella 
posizione $B$, a partire dalla posizione centrata
\begin{equation}
\begin{pmatrix} {x_i}^\star\left(\tau\right)\\{y_i}^\star\left(\tau\right) \end{pmatrix}
={\bm T}(x_{\scriptscriptstyle B}, y_{\scriptscriptstyle B}){\bm R}(\phi_{\scriptscriptstyle B} +\pi)
\begin{pmatrix} {x_i}\left(\tau\right)\\{y_i}\left(\tau\right) \end{pmatrix}\,,
 \label{eq:rot_trasl_i}
\end{equation}
\noindent avendo indicato con ${x_i}\left(\tau\right)$ e ${y_i}\left(\tau\right)$
 i punti della ipocicloide centrata in $O$, mentre con
${x_i}^\star\left(\tau\right)$ e con ${y_i}^\star\left(\tau\right)$
indichiamo le loro trasformazioni.
Una delle condizioni che deve essere soddisfatta,
affinch\'e le due curve risultino coniugate, risiede 
nella coincidenza, nella posizione $B$, dei punti delle due curve.
In altri termini, deve essere verificato che
\begin{equation}
\begin{pmatrix} {x_e}\left(t_{\scriptscriptstyle B}\right)\\{y_e}\left(t_{\scriptscriptstyle B}\right) \end{pmatrix}=
\begin{pmatrix} {x_i}^\star\left(\tau_{\scriptscriptstyle B}\right)\\{y_i}^\star\left(\tau_{\scriptscriptstyle B}\right)
\end{pmatrix}
={\bm T}(x_{\scriptscriptstyle B}, y_{\scriptscriptstyle B}){\bm R}(\phi_{\scriptscriptstyle B} +\pi)
\begin{pmatrix} {x_i}\left(\tau_{\scriptscriptstyle B}\right)\\{y_i}\left(\tau_{\scriptscriptstyle B}\right)
\end{pmatrix}\,,
 \label{eq:rot_trasl_ei}
\end{equation}
\noindent dove ${x_e}\left(t_{\scriptscriptstyle B}\right)$ e
${y_e}\left(t_{\scriptscriptstyle B}\right)$ sono le coordinate 
dell'epicicloide nel punto di interesse.
\noindent Scriviamo la
\ref{eq:rot_trasl_ei} in forma esplicita
\begin{multline}
\begin{pmatrix} {x_i}^\star\left(\tau_{\scriptscriptstyle B}\right)\\{y_i}^\star\left(\tau_{\scriptscriptstyle B}\right) \end{pmatrix} =
\begin{pmatrix} \cos(\phi_{\scriptscriptstyle B}+\pi) &-\sin(\phi_{\scriptscriptstyle B}+\pi)\\
		 \sin(\phi_{\scriptscriptstyle B}+\pi) &\cos(\phi_{\scriptscriptstyle B}+\pi) \end{pmatrix}
\begin{pmatrix} {x_i}\left(\tau_{\scriptscriptstyle B}\right)\\{y_i}\left(\tau_{\scriptscriptstyle B}\right) \end{pmatrix} + \\
+ \begin{pmatrix} (r_1+r_2)\cos(t_{\scriptscriptstyle B})\\(r_1+r_2)\sin(t_{\scriptscriptstyle B}) \end{pmatrix}\,,
\label{eq:ipocicloide_ruototraslata}
\end{multline}
\noindent nella quale la matrice, che compare
appena a destra del segno di uguaglianza, opera la rotazione piana di un angolo 
pari a $\phi_{\scriptscriptstyle B}+\pi$.
\noindent Mostriamo ora che i punti
${p_e}_{\scriptscriptstyle B}$ e ${p_i}_{\scriptscriptstyle B}$ coincidono
a valle della roto-traslazione,
fatto peraltro evidente dalla figura composta per via numerica.
Utilizziamo a questo scopo la \ref{eq:ipocicloide_ruototraslata}, inserendo per
${x_i}\left(\tau{\scriptscriptstyle B}\right)$ e ${y_i}\left(\tau{\scriptscriptstyle B}\right)$
quelli forniti dalla \ref{eq:ipocicloide}, e avendo cura di sostituire, al
parametro $\tau$, il suo omologo $t$, che origina la posizione $B$.
Ponendo quindi $\tau{\scriptscriptstyle B} = {r_1\over{r_2}}t_{\scriptscriptstyle B}$ e
$\phi_{\scriptscriptstyle B}={{r_1+r_2}\over{r_2}}t_{\scriptscriptstyle B}$,
per la prima riga delle \ref{eq:ipocicloide_ruototraslata} avremo
\begin{multline}
{x_i}^\star\left(\tau_{\scriptscriptstyle B}\right)=
[-\left(r_2-r_r\right)\cos\left({r_1\over{r_2}}t_{\scriptscriptstyle B}\right)\cos\left({{r_1+r_2}\over{r_2}}t_{\scriptscriptstyle B}\right) + \\ -
r_r\cos\left({{r_2-r_r}\over{r_r}}{r_1\over{r_2}}t_{\scriptscriptstyle B}\right)\cos\left({{r_1+r_2}\over{r_2}}t_{\scriptscriptstyle B}\right) + \\ -
\left(r_2-r_r\right)\sin\left({r_1\over{r_2}}t_{\scriptscriptstyle B}\right)\sin\left({{r_1+r_2}\over{r_2}}t_{\scriptscriptstyle B}\right) + \\ +
r_r\sin\left({{r_2-r_r}\over{r_r}}{r_1\over{r_2}}t_{\scriptscriptstyle B}\right)\sin\left({{r_1+r_2}\over{r_2}}t_{\scriptscriptstyle B}\right)] + \\ +
\left(r_1+r_2\right)\cos\left(t_{\scriptscriptstyle B}\right)\,,
\label{eq:inviluppo_x}
\end{multline}
\noindent e facendo buon uso delle formule di somma per gli archi
otteniamo infine
\begin{equation}
{x_i}^\star\left(\tau_{\scriptscriptstyle B}\right)=\left(r_1+r_r\right)\cos\left(t_{\scriptscriptstyle B}\right) -
r_r\cos\left(\dfrac{r_1+r_r}{r_r}t_{\scriptscriptstyle B}\right)=
{x_e}\left(t_{\scriptscriptstyle B}\right)\,.
\label{eq:invdimostrato}
\end{equation}
\noindent Crediamo inutile riportare anche l'espressione
di ${y_i}^\star\left(\tau_{\scriptscriptstyle B}\right)$ che coincide naturalmente
con la
${y_e}\left(t_{\scriptscriptstyle B}\right)$
riportata nelle \ref{eq:epicicloide}, per $t=t_{\scriptscriptstyle B}$.
Tali espressioni per le coordinate del punto della ipocicloide
ci assicurano che
${p_i}_{\scriptscriptstyle B}\equiv{p_e}_{\scriptscriptstyle B}$,
e la coincidenza test\'e dimostrata vale naturalmente per tutti i punti
omologhi dell'epicicloide e dell'ipocicloide.
Ci\`o per\`o non \`e sufficiente per mostrare che la ipocicloide inviluppa,
durante il suo moto, la curva epicicloide:
a tale scopo occorre infatti provare che, nei punti omologhi,
anche le tangenti alle due curve coincidono, circostanza anche
questa evidente dalla figura. Le  derivate delle coordinate dell'epicicloide
risultano essere
\begin{equation}
\begin{cases}
\dfrac{{\rm d}{x_e}\left(t\right)}{{\rm d}t} = 
\left(r_1+r_r\right)
[-\sin\left(t\right) +
\sin\left(\dfrac{r_1+r_r}{r_r}t\right)] \\
\dfrac{{\rm d}{y_e}\left(t\right)}{{\rm d}t} = 
\left(r_1+r_r\right)
[\cos\left(t\right) -
\cos\left(\dfrac{r_1+r_r}{r_r}t\right)]
\end{cases}\,,
\label{eq:depicicloide}
\end{equation}
\noindent mentre quelle delle derivate dell'ipocicloide sono
\begin{equation}
\begin{cases}
\dfrac{{\rm d}{x_i}\left(\tau\right)}{{\rm d}\tau} = 
\left(r_2-r_r\right)
[-\sin\left(\tau\right) -
\sin\left(\dfrac{r_2-r_r}{r_r}\tau\right)] \\
\dfrac{{\rm d}{y_i}\left(\tau\right)}{{\rm d}\tau} = 
\left(r_2-r_r\right)
[-\cos\left(\tau\right) +
\cos\left(\dfrac{r_2-r_r}{r_r}\tau\right)]
\end{cases}\,.
\label{eq:dipocicloide}
\end{equation}

\noindent Ruotando queste ultime derivate, calcolate in $B$, dell'angolo
$\phi_{\scriptscriptstyle B}+\pi$  otteniamo
\begin{equation}
\begin{pmatrix}
\dfrac{{\rm d}{x_i}\left(\tau\right)}{{\rm d}\tau}
\\
\dfrac{{\rm d}{y_i}\left(\tau\right)}{{\rm d}\tau}
\end{pmatrix}^\star_{\scriptscriptstyle B}
 =
\begin{pmatrix} \cos(\phi + \pi) &-\sin(\phi + \pi)\\
		 \sin(\phi + \pi) &\cos(\phi + \pi) \end{pmatrix}
\begin{pmatrix}
\dfrac{{\rm d}{x_i}\left(\tau\right)}{{\rm d}\tau}
\\
\dfrac{{\rm d}{y_i}\left(\tau\right)}{{\rm d}\tau}
\end{pmatrix}_{\scriptscriptstyle B}\,,
\label{eq:derivata_ruotata}
\end{equation}
\noindent dove si pu\`o notare l'assenza del termine relativo
 alla traslazione in quanto trattandosi di trasformare delle 
differenze tra coordinate essa avrebbe effetto nullo.
Sostituendo, come si \`e fatto in precedenza,
$\tau_{\scriptscriptstyle B} = {\dfrac{r_1}{r_2}}t_{\scriptscriptstyle B}$ e
$\phi_{\scriptscriptstyle B}={\dfrac{r_1+r_2}{r_2}}t_{\scriptscriptstyle B}$
otteniamo
\begin{multline}
\dfrac{{\rm d}{x_i}\left(\tau\right)}{{\rm d}\tau}\bigr|^\star_{\scriptscriptstyle B} = 
[\left(r_2-r_r\right)\sin\left(\dfrac{r_1}{r_2}t_{\scriptscriptstyle B}\right)
\cos\left(\dfrac{r_1+r_2}{r_2}t_{\scriptscriptstyle B}\right) + \\ +
\left(r_2-r_r\right)\sin\left(\dfrac{r_1}{r_2}\dfrac{r_2-r_r}{r_r}t_{\scriptscriptstyle B}\right)
\cos\left(\dfrac{r_1+r_2}{r_2}t_{\scriptscriptstyle B}\right) + \\ -
\left(r_2-r_r\right)\cos\left(\dfrac{r_1}{r_2}t_{\scriptscriptstyle B}\right)
\sin\left(\dfrac{r_1+r_2}{r_2}t_{\scriptscriptstyle B}\right) + \\ +
\left(r_2-r_r\right)\cos\left(\dfrac{r_1}{r_2}\dfrac{r_2-r_r}{r_r}t_{\scriptscriptstyle B}\right)
\sin\left(\dfrac{r_1+r_2}{r_2}t_{\scriptscriptstyle B}\right)]\,,
\label{eq:dx_in_t}
\end{multline}
\noindent cio\`e
\begin{equation}
\dfrac{{\rm d}{x_i}\left(\tau\right)}{{\rm d}\tau}\bigr|^\star_{\scriptscriptstyle B} = 
\left(r_2+r_r\right)[-\sin\left(t_{\scriptscriptstyle B}\right) +
\sin\left(\dfrac{r_1+r_r}{r_r}t_{\scriptscriptstyle B}\right)]\,,
\label{eq:dx}
\end{equation}
\noindent mentre, ripetendo il precedente calcolo, che preferiamo non riportare
esplicitamente, per la coordinata $y$ otterremmo
\begin{equation}
\dfrac{{\rm d}{y_i}\left(\tau\right)}{{\rm d}\tau}\bigr|\star_{\scriptscriptstyle B} = 
\left(r_2+r_r\right)[\cos\left(t_{\scriptscriptstyle B}\right) -
\cos\left(\dfrac{r_1+r_r}{r_r}t_{\scriptscriptstyle B}\right)]\,.
\label{eq:dy}
\end{equation}
\noindent Le espressioni \ref{eq:dx} e \ref{eq:dy} coincidono, a meno di
un fattore moltiplicativo, con le espressioni
delle derivate dell'epicicloide \ref{eq:depicicloide} calcolate in $B$.
La differenza tra i fattori che moltiplicano le derivate
dipende dal fatto che $t$ e $\tau$ sono parametri in generale diversi tra
loro e generano quindi differenti rapporti incrementali.
Questa circostanza per\`o nulla toglie alla nostra dimostrazione, in quanto si
avr\`a
\begin{equation}
\dfrac{\dfrac{{\rm d}{x_e}\left(t\right)}{{\rm d}t}\bigr|_{\scriptscriptstyle B}}{\dfrac{{\rm d}{y_e}\left(t\right)}{{\rm d}t}\bigr|_{\scriptscriptstyle B}} =
\dfrac{\dfrac{{\rm d}{x_i}\left(\tau\right)}{{\rm d}\tau}\bigr|^\star_{\scriptscriptstyle B}}
{\dfrac{{\rm d}{y_i}\left(\tau\right)}{{\rm d}\tau}\bigr|^\star_{\scriptscriptstyle B}} \,,
\label{eq:dtot}
\end{equation}
\noindent che prova la coincidenza delle tangenti alla curva inviluppante e
al relativo inviluppo nel generico punto $B$.
Abbiamo cos\`i mostrato che, durante il rotolamento di una ruota di
raggio $r_2$ su una di raggio $r_1$, che possiamo immaginare ferma,
l'ipocicloide, trascinata dalla ruota mobile, inviluppa un tratto di epicicloide
della ruota fissa corrispondente al medesimo epiciclo.
In definitiva, pensando alle circonferenze in gioco come alle polari del moto,
i due tratti di epicicloide e ipocicloide si comporteranno come profili
coniugati essendo costantemente a contatto e possedendo la stessa
tangente\footnote{
Nessuna ipotesi riguardante i raggi delle circonferenze \`e stata messa in campo
quindi il raggio della rolletta $r_r$ potrebbe superare quello delle
circonferenze sulle quali sviluppa le cicloidi. Nel caso di
 $r_r$ di lunghezza infinita otteniamo per le due curve archi di evolvente.
}.

\section{Ruote con Denti a Profilo Cicloidale}

\noindent Dalle conclusioni del precedente paragrafo, come anche dalle
conclusioni del paragrafo \ref{prof_con}, emerge con chiarezza la possibilit\`a
di usare i profili cicloidali per modellare i fianchi dei denti di ruote,
aventi primitive circolari,
che garantiscano la trasmissione omocinetica del moto. Contrariamente 
per\`o a quanto visto per le ruote i cui denti presentano un
profilo a evolvente di cerchio,
trattando i profili cicloidali siamo costretti a dividere i profili coniugati
in due porzioni: considerando il punto di contatto tra le
primitive, una di queste si trover\`a 
dalla parte della prima ruota mentre l'altra porzione si trover\`a
dalla parte della seconda ruota. Contenendo, poi, ciascuna delle due porzioni
dei profili ora citate, due tratti distinti, un primo di epicicloide
e un secondo di ipocicloide,
il numero di spezzoni (i quali costituiranno la sagoma dei denti) da considerare sar\`a in definitiva quattro.
Ad esempio, volendo
usare una porzione della epicicloide e della ipocicloide rappresentate in
 figura \ref{fig:epiipoconiugate}, il tratto appartenente alla epicicloide
potr\`a rappresentare la sporgenza del dente sulla ruota $r_1$, 
 mentre un tratto sulla ipocicloide
costituir\`a la rientranza sulla ruota $r_2$.
Non rimane quindi altra scelta che
ripetere quanto fatto per le due primitive a ruoli invertiti: generare cio\`e
una epicicloide sulla ruota $r_2$
e una ipocicloide sulla $r_1$. Ci preme sin d'ora sottolineare
tre aspetti importanti delle ruote a profilo cicloidale,
che ci indirizzano sulla strada della loro costruzione.
La prima considerazione riguarda proprio la coniugazione
a tratti dei profili dei denti:
teste e fianchi sono tra loro coniugati, ma non lo sono i profili
nella loro interezza.
Ci\`o ci convince facilmente che l'interasse di progetto
dovr\`a essere rigidamente rispettato e che le due circonferenze primitive
potranno solo coincidere con quelle di raggi
$r_1$ e $r_2$: le relative coppie di profili coniugati si scambieranno
il testimone proprio quando il loro punto di contatto cadr\`a su queste
circonferenze.
La seconda considerazione \`e una diretta conseguenza della
prima: gli epicicli usati per ottenere le due coppie di profili possono essere
tra loro diversi, circostanza questa
che, come vedremo pi\`u avanti, semplifica la
 costruzione di queste ruote in un determinato perimetro di applicazione che
\`e quello dell'orologeria.
\`E chiaro che risulteranno ``di assortimento'' soltanto quelle ruote
aventi il medesimo passo o, ci\`o che \`e equivalente, lo stesso modulo, 
costruite mediante lo stesso epiciclo o mediante coppie di epicicli uguali
tra loro.
Da ultimo notiamo che l'inconveniente del numero minimo di
denti -- il quale
nel caso di ruote con denti a evolvente deriva dall'impossibilit\`a
di
estendere il fianco dei denti stessi oltre le circonferenze di base,
in questo caso, scegliendo opportunamente il raggio delle rollette --
non si manifesta, essendo entrambe le coppie di profili
prolungabili (entro certi limiti) a volont\`a.

\section{Proporzionamento delle Ruote Cicloidali}

\noindent Rimane valido anche per queste ruote il proporzionamento
dei denti legato al modulo della ruota stessa? Ancora, \`e necessario,
oppure utile, introdurre anche qui tale concetto?
Come sappiamo, nell'ambito della progettazione delle ruote ad evolvente, una
stretta disciplina che lega la forma dei denti, in particolare
{\em addendum} e {\em dedendum}, al modulo $m$ rende la loro progettazione
pi\`u semplice.
Si pu\`o sicuramente definire, anche nel caso di profilo cicloidale,
il modulo $m=2 r/z$, pari cio\`e, come al solito, al diametro primitivo
diviso il numero di denti.
Una volta definito tale modulo, saremmo in grado di riferire ad esso
l'{\em addendum} e il {\em dedendum} dei denti utilizzando proporzioni che si ispirino a
 quelle delle ruote a evolvente.  Ci\`o che si ottiene
in tal caso, per un ingranaggio $z_1=11$ e $z_2=60$, \`e riportato in figura
\ref{fig:1160cy}.

\begin{figure}[hbt]
     \begin{center}
     \includegraphics[width=0.9\textwidth]{part3/ruote/FIG/ruote/1160cy.pdf}
     \end{center}
\begin{picture}(0,0)(0,0)
        \scriptsize{
        \put(112,243){$r_2$}
        \put(290,207){$r_1$}
        \put(232,142){$O_1$}
}
\end{picture}
\vskip -5mm
        \caption{
\em Generico ingranaggio cicloidale con $r_1=40$,  $z_1=11$, $z_2=60$ e
raggio degli epicicli $r_i=r_e=10$.
}
     \label{fig:1160cy}
\end{figure}

\noindent La scelta dei numeri di denti, per l'ingranaggio di questa figura,
non \`e casuale: la figura
\ref{fig:1160}, che riporta un ingranaggio tra ruote a evolvente
recanti gli stessi numeri di denti, mette infatti in evidenza, sul pignone,
la problematica che si chiama sottotaglio dei denti, oppure
interferenza. Dalla figura \ref{fig:1160cy} risulta invece evidente che tale
problema non si manifesta per le ruote cicloidali, per le quali \`e
persino possibile un numero di denti pari a due.
A dire il vero per\`o, l'ingranaggio cicloidale, sopra riportato, ci convince
fino a un certo segno. Intanto, come \`e esplicitato nella didascalia,
siamo costretti a evidenziare i raggi degli epicicli utilizzati che, nel nostro
caso, a fronte di un raggio primitivo del pignone pari a $r_1=40$,
valgono $r_e=r_i=10$ (le unit\`a di misura sono qui volutamente omesse).
Il raggio unico dell'epiciclo e dell'ipociclo sulla
prima ruota equivale naturalmente a quello dell'ipociclo e dell'epiciclo
sulla seconda. Il valore di tale raggio giunge da una raccomandazione
contenuta in \cite{punzi}, la quale segnala che la proporzione ottimale tra
il raggio della rolletta e il modulo della ruota \`e $r_e=r_i=1.375 m$,
risultando cos\`i possibile
tagliare pignoni con pochi denti senza incorrere nell'indebolimento della
base dei suoi denti. In teoria, potremmo quindi affermare che tutte
 le generiche ruote disegnate in questo modo e
 aventi lo stesso modulo risulterebbero essere
``di assortimento''. Ma \`e subito chiaro che, nell'ambito di queste
dentature, la scelta
tutto sommato arbitraria dei raggi di epiciclo e ipociclo, i raggi
dei quali si adattano forzatamente ad essere rapportati al modulo della ruota,
indebolisce molto il concetto di assortimento, e con esso quello stesso
di modulo che invece, nel caso delle ruote ad evolvente, lega bens\`i il
loro numero di denti al loro diametro primitivo, ma soprattutto 
seleziona
un utensile, usando il quale si possono tagliare
ruote che risultano tutte compatibili tra loro. Quindi, modulo s\`i, ma segnatamente
in quanto passo diametrale, il quale individua, in sostanza, le dimensioni
dell'oggetto che stiamo considerando. Vedremo a breve uno {\em standard}
che riguarda la fabbricazione di ruote per orologeria, dove il concetto
di modulo non solo \`e usato ma porta con s\'e un folto numero di
relazioni tra le quantit\`a geometriche che modellano, in quell'ambito,
i denti. Ciononostante, anche in quel mondo il concetto di assortimento
\`e molto pi\`u sfumato rispetto a quello che si trova studiando le ruote ad
evolvente. Tornando all'ingranaggio di figura \ref{fig:1160cy}, chi scrive 
non \`e sicuro che il  proporzionamento di generiche ruote cicloidali si esegua
seguendo quella via. Non siamo neppure sicuri che ruote come quelle
di figura \ref{fig:1160cy} siano mai state tagliate (nel caso, ovviamente,
con fresatura di forma), e ammettiamo senza difficolt\`a  di non
averle mai viste dal vero.


\section{Interessanti Casi di Ruote Cicloidali}

L'epilogo, relativamente privo di entusiasmo, del precedente paragrafo non
deve disincentivare la nostra indagine circa le eventuali applicazioni di questi
ingranaggi, meno comuni dei loro ``cugini'' a evolvente.
Toccheremo infatti brevemente gli aspetti geometrici
di tre tipologie particolari di
ruote con denti a profilo cicloidale, ciascuno dei quali presenta
interessanti applicazioni
concrete nella tecnica. La terza, in ordine di presentazione,
di queste famiglie di ruote gode anzi
di diffusione copiosa e capillare
perch\'e riguarda gli ingranaggi presenti negli orologi.

\noindent Un primo caso peculiare di dentatura cicloidale
emerge dalla possibilit\`a che gli ipocicli siano
le stesse circonferenze primitive. Le ipocicloidi si riducono, in tal caso,
a un punto, mentre le epicicloidi saranno tracciate dai punti di una
delle due primitive che rotola sull'altra.
Generalmente, si sceglie di far seguire a un certo numero di punti della
primitiva del pignone, punti che corrisponderanno ai suoi ``denti'',
le traiettorie di epicicloide disegnate dalla circonferenza
primitiva del  pignone stesso mentre rotola
sulla circonferenza primitiva dell'altra ruota.
I punti individuati sul pignone dovrebbero, in teoria,
 seguire fisicamente i fianchi dei 
denti con profilo epicicloidale, ma essi non potrebbero
offrire che contatti singolari,
poco idonei a trasmettere forze e, in generale, inutilizzabili nelle
applicazioni tecniche. Per questo motivo, si sostituiscono
tali punti con idonee rotelle di raggio  opportuno
e, di conseguenza,  si modifica la forma del fianco cicloidale dei
denti della ruota
condotta in modo che tali fianchi siano l'inviluppo delle
rotelle stesse quando il loro centro segue l'epicicloide.
Si ottengono in questo modo
le cosiddette 
{\em ruote a lanterna}\index{ruote!a lanterna}, il cui nome deriva 
dalla particolare forma del pignone dell'ingranaggio che pu\`o ricordare
una lanterna. Un esempio di queste ruote \`e riportato in figura \ref{fig:lanterna}.
Gli ingranaggi a lanterna vengono impiegati per movimentare ruote
 condotte
di grandi diametri e possono sopportare condizioni di lavoro gravose.
Quando il raggio
della ruota condotta diventa infinito il pignone a lanterna muove
una cremagliera. In questo caso i centri delle rotelle percorrono
dei tratti di cicloide. Quest'ultima soluzione si trova in alcuni
nastri trasportatori soggetti a carichi ingenti.
\begin{figure}[t]
     \begin{center}
     \includegraphics[width=0.9\textwidth]{part3/ruote/FIG/ruote/lanterna.pdf}
     \end{center}
\begin{picture}(0,0)(0,0)
        \scriptsize{
        \put(220,165){$r_2$}
        \put(294,165){$r_1$}
        \put(290,136){$O_1$}
        \put(127,136){$O_2$}
	\put(215,270){tratto di epicicloide, $r_e=r_1$}
}
\end{picture}
\vskip -5mm
        \caption{
\em Ruota cicloidale con pignone a lanterna, $z_1=8$, $z_2=20$.
}
     \label{fig:lanterna}
\end{figure}

\noindent Abbiamo accennato alla possibilit\`a di ottenere
ruote cicloidali con numero di denti ridottissimo.
Un particolare dispositivo atto a comprimere o pompare fluidi si
ottiene dall'ingranaggio di ruote con denti cicloidali a due o a tre 
denti.
Descriviamo il caso di ruote con due denti, che riportiamo in figura
\ref{fig:kompressor}, specificando che le proporzioni da rispettare
per ottenere tale meccanismo sono le seguenti:
$z=2$, $r_1=r_2$ e $r_e=r_i= r_1/4$.
Le ruote, disegnate mediante il proporzionamento proposto,
non sono certo in grado di trasmettere tra loro
il moto con continuit\`a. Tuttavia, qualora la loro rotazione venga garantita
da una adeguata coppia di ruote dentate esterne,
esse ingranano, da un punto di vista
puramente cinematico, correttamente.
Tali {\em ruote a lobi}\index{ruote!a lobi}, alterando periodicamente
i volumi a loro disposizione nel corpo della pompa, assicurano
fino ad un certo grado la tenuta stagna, quindi la possibilit\`a di comprimere
o di pompare i fluidi. Esistono diverse versioni commerciali di questi
dispositivi che prendono il nome di {\em pompe a lobi}\index{pompe a lobi},
quando elaborano
fluidi incomprimibile, oppure di {\em compressori volumetrici  a lobi}
quando trattano gas comprimibili.
Una nota casa automobilistica installava, fino a 
non molto tempo fa, questo tipo di compressore per sovralimentare
diversi dei suoi modelli di motore a combustione interna.
\begin{figure}[t]
     \begin{center}
     \includegraphics[width=0.9\textwidth]{part3/ruote/FIG/ruote/kompressor.pdf}
     \end{center}
\begin{picture}(0,0)(0,0)
        \scriptsize{
        \put(245,143){$O_1$}
        \put(130,143){$O_2$}
        \put(245,195){$\omega$}
        \put(130,195){$\omega$}
	\put(171,230){\rotatebox{90}{Uscita}}
	\put(171,50){\rotatebox{90}{Ingresso}}
}
\end{picture}
\vskip -5mm
        \caption{
\em Compressore a lobi ottenuto dall'ingranaggio di due ruote cicloidali a due denti.
}
     \label{fig:kompressor}
\end{figure}

\noindent Vi \`e un ambito della meccanica, molto particolare,
 dove le ruote cicloidali
la fanno da padrone, ed \`e quello dei {\em meccanismi per orologeria}\index{ruote!per orologeria}.
Si tratta, come \`e facile intuire,
 di un mondo piuttosto distante da quello della meccanica delle
macchine. L'ingombro delle ruote dentate per orologeria deve essere
infatti compatibile con le dimensioni degli orologi che le ospitano,
 quindi il loro diametro
spesso non supera i 10--15 millimetri. Altre
caratteristiche
salienti di questo speciale settore degli ingranaggi sono l'assenza sistematica
di lubrificazione e i valori minuscoli dei flussi di potenza, flussi che, non
presentando mai inversioni di direzione, rendono tollerabili (quando
non prescritti) giochi rilevanti tra le dentature.
Una volta scritta la procedura numerica
dalla quale otteniamo le figure \ref{fig:1160cy} e \ref{fig:kompressor},
la curiosit\`a di vedere tale codice all'opera
con le ruote per orologeria era elevata. La 
fortuna e la casualit\`a ci hanno fatto imbattere in un testo che potrebbe
aver rappresentato per un bel po' di tempo il riferimento degli orologiai: 
\cite{grossmann}. Per la verit\`a il libro ora citato, nell'opportuno
capitolo dedicato agli ingranaggi, considera anche le ruote con profilo
ad evolvente. Ma \`e subito abbastanza chiaro che lo fa per completezza, e che
le ruote a profilo cicloidale saranno al centro di tutte le successive
considerazioni. Siamo certi che la norma \cite{bs978}, la quale propone
una standardizzazione delle ruote per orologeria, sia stata, in larga parte,
elaborata seguendo \cite{grossmann}.
Nel paragrafo \ref{percorso-ruote}
tra i precetti che abbiamo messo in campo, nel nostro percorso intuitivo,
ve ne era uno, debole, il quale prescriveva che ci fosse una sola famiglia di 
ruote di assortimento, e non due. Ma tale precetto veniva appunto 
etichettato 
come debole, anzi, in quello stesso luogo indicavamo questo capitolo
di approfondimento come la sede
dove si descrive una serie di ingranaggi formata da due famiglie 
diverse di ruote.
In orologeria, di fatto, pignoni e ruote non possono essere confusi tra loro
e l'ingranamento pignone-pignone o ruota-ruota non \`e mai 
possibile. Addirittura,
la nomenclatura delle loro sporgenze \`e eterogenea: i pignoni ospitano
le foglie, le ruote i denti. Queste due famiglie di ruote dentate nascono
dalla volont\`a di ottenere, sia per i pignoni, sia per le ruote,
denti il cui
profilo risulti il pi\`u 
semplice possibile da tagliare. Per questo motivo, si fa in modo di utilizzare
sia sul pignone sia sulla ruota un ipociclo di raggio esattamente pari
alla met\`a
del loro raggio primitivo. L'ipocicloide cos\`i ottenuta coincide con un
diametro della primitiva, pertanto si presenta rettilineo e radiale.
\begin{figure}[b]
     \begin{center}
     \includegraphics[width=0.9\textwidth]{part3/ruote/FIG/ruote/660cy_hor.pdf}
     \end{center}
\begin{picture}(0,0)(0,0)
        \scriptsize{
        \put(150,222){$r_2$}
        \put(288,189){$r_1$}
        \put(180,162){$\alpha_a$}
        \put(193,136){$\alpha_r$}
        \put(245,145){$O_1$}
        \put(205,50){$r_{e_2}=1.949 m$}
        \put(207,40){BS 978-2}
        \put(290,59){$r_{e_1}=1.25 m$}
        \put(290,49){BS 978-2}
}
\end{picture}
\vskip -5mm
        \caption{
\em Ingranaggio cicloidale da orologeria, pignone con $z_1=6$ foglie, ruota
da $z_2=60$ denti. Accesso $\alpha_a=17^{\circ}44'$, recesso $\alpha_r=42^{\circ}16'$.
}
     \label{fig:660cy_hor}
\end{figure}
Ovviamente le epicicloidi, che costituiranno il profilo delle
teste dei denti, saranno ottenute dalle stesse rollette, gi\`a utilizzate
per gli ipocicli, a primitiva invertita.
La figura \ref{fig:660cy_hor} riporta un ingranaggio con $z_1=6$ e $z_2=60$
in cui risultano evidenti i fianchi rettilinei delle foglie del
pignone  e dei denti della ruota.
Non \`e difficile comprendere che i fianchi rettilinei dei denti
risultano di facile lavorazione, potendo essere tagliati da sottili 
fese a disco con taglienti piani. In altre parole non \`e necessario che la
sezione della fresa abbia la forma del vuoto tra i denti.
Sembra che gli appassionati di questo settore riescano a tagliare queste
minuscole ruote persino con utensili pi\`u o meno hobbistici,
 certamente aiutati
da opportuni apparecchi divisori. Le teste ({\em addendum}) delle foglie
del  pignone e dei denti della ruota meritano invece considerazioni
distinte. Cominciamo con l'osservare che i rapporti di trasmissione
utili in orologeria sono una manciata. In \cite{grossmann} ne sono
specificati sedici e l'autore dice esplicitamente di averne riportati
pi\`u di quelli realmente in uso. Le moltipliche che si eseguono
sono ancor meno delle coppie trattate, ad esempio la moltiplica per
dieci si pu\`o ottenere con gli ingranaggi
 60--6, 70--7 e 80--8. Parliamo di moltiplica
perch\'e, in generale, in quest'ambito, sono le ruote a muovere i pignoni. 
Da quest'ultima considerazione, dal fatto cio\`e che siano
le ruote ad essere motrici, si identificano naturalmente il
tratto di accesso e quello di recesso, i quali si toccano sulla
linea immaginaria che congiunge il centro della ruota a quello del
pignone.
Nel paragrafo \ref{ruote-corr1} abbiamo accennato alla dissimmetria
di cui soffrono accesso e recesso rispetto al rendimento dell'ingranaggio:
durante il recesso, il rendimento \`e sempre pi\`u elevato.
In un ambito dove la ridottissima quantit\`a di energia disponibile proviene da
una molla, o da un peso che scende, oppure da una piccola batteria, e
tale energia deve muovere il meccanismo dell'orologio il pi\`u a lungo
possibile, risulta quindi naturale 
la tendenza ad avere il minimo arco di accesso e il massimo
arco di recesso, compatibili con la geometria delle ruote coinvolte.
L'angolo di recesso maggiormente esteso
si ottiene con la massima altezza realizzabile per i denti della ruota.
Il profilo dei denti viene pertanto esteso fino ad appuntire il dente
stesso. L'ultima posizione utile del dente nel recesso si ha quando la
tangente sulla punta dell'{\em addendum} \`e parallela al fianco rettilineo della
foglia, come mostrato in figura.
Nella pratica costruttiva di queste ruote, tali fianchi che,
teoricamente, dovrebbero essere epicicloidali,
 vengono approssimati da archi di circonferenza,
facili da improntare direttamente sulla fresa a disco che taglia la
ruota.  Lo {\em standard}
 B.S. 978, \cite{bs978}, che raccoglie in compendio le raccomandazioni
di una vasta parte dell'industria svizzera di orologi, raccomanda di sostituire
la porzione epicicloidale del fianco del dente con un arco di circonferenza di
raggio medio rispetto all'epicicloide stessa. Mentre in \cite{grossmann}
il proporzionamento delle ruote e dei pignoni avviene con riferimento
al loro diametro e al  loro numero di denti, nelle raccomandazioni raccolte
in \cite{bs978}, immaginiamo
con l'intento di dare maggiore sistematicit\`a ai parametri che regolano
la geometria di queste ruote, si parte da una tabella di moduli, in mm, 
dove su circa sessanta valori standardizzati, spiccano in grassetto
0.08, 0.1, 0.15, 0.2, 0.25, 0.3, 0.4, 0.5, 0.6, 0.7, 0.8, 1.0.
Giusto per rendere l'idea del perimetro entro il quale ci muoviamo,
una ruota con 60 denti e modulo 0.1 mm avr\`a un diametro primitivo pari
a 6 mm. Parlando dei denti delle ruote, sempre in \cite{bs978}, si trova
la tabella 2 che lega i loro dettagli geometrici,
 come {\em addendum}, {\em dedendum},
raggio e posizione del centro dell'arco di circonferenza che sostituisce
lo spezzone di epicicloide e spessore del dente, al modulo
della ruota stessa.
Alcune di queste dimensioni, in particolare il raggio dell'arco
di circonferenza che costituir\`a l'{\em addendum}, dipendono per\`o, oltre che dal
modulo, anche dal rapporto di trasmissione e dal numero di foglie del
pignone col quale la ruota \`e destinata a ingranare.
Fortunatamente, come abbiamo gi\`a accennato,
sia i rapporti di trasmissione, sia i numeri delle foglie che si usano
sono limitati. Per questo motivo risulta agevole
costruire la tabella 3 di \cite{bs978},
che riporta per nove numeri delle foglie del pignone, da 6 a 16 saltando 11 e 13,
e per 14 valori del rapporto di trasmissione, i valori che identificano
geometricamente il dente della ruota in questione.
Nella nostra figura \ref{fig:660cy_hor}, la ruota ha sessanta denti mentre
il pignone presenta sei foglie, e abbiamo lasciato, sia per l'{\em addendum}
della ruota, sia per quello del pignone, i tratti di epicicloide calcolati.
\begin{figure}[b]
     \begin{center}
     \includegraphics[width=0.9\textwidth]{part3/ruote/FIG/ruote/1290cy_hor.pdf}
     \end{center}
\begin{picture}(0,0)(0,0)
        \scriptsize{
        \put(156,215){$r_2$}
        \put(295,184){$r_1$}
        \put(187,134){$\alpha_r$}
        \put(245,145){$O_1$}
}
\end{picture}
\vskip -5mm
        \caption{
\em Ingranaggio cicloidale da orologeria, pignone con $z_1=12$ foglie, ruota
da $z_2=90$ denti. Ingranaggio a tutto recesso $\alpha_r=32.5^{\circ}$.
}
     \label{fig:1290cy_hor}
\end{figure}
Riportiamo comunque in figura anche le circonferenze
approssimanti indicate in \cite{bs978}, in modo da avere un riscontro visivo
circa la loro effettiva bont\`a. 
Una volta calcolato, tramite
la nostra procedura numerica, oppure consultando \cite{grossmann}, pag. 175,
il massimo angolo di recesso, $\alpha_r$,
dobbiamo garantire un angolo di accesso che sia $\alpha_a> 2 \pi/z_1 -\alpha_r$,
in modo da poter assicurare la continuit\`a della trasmissione.
Nel nostro caso abbiamo
un recesso molto pi\`u esteso dell'accesso. L'epicicloide dell'{\em addendum}
delle foglie del pignone sar\`a pertanto impegnata in un suo breve tratto
vicino alla circonferenza primitiva. Per questo motivo l'approssimazione
del tratto di epicicloide sulle foglie dei pignoni \`e ancora pi\`u facile
di quella che si esegue sui denti delle ruote.
Tanto \`e vero che in molti casi, come quello che riportiamo in figura
\ref{fig:1290cy_hor}, l'{\em addendum} sulle foglie non lavora. Per tale ingranaggio
l'arco di accesso \`e nullo e l'arco di recesso vale $\alpha=32.5$, 
gi\`a superiore ai necessari $30^{\circ}$. La sommit\`a delle foglie presenta quindi
arrotondamenti senza alcuna funzione cinematica, ma che rendono i pignoni
maggiormente compatibili con l'ingranamento: non sarebbero apprezzabili
gli spigoli vivi sulla cima delle foglie, qualora fossero terminate
alla fine dell'ipociclo rettilineo. Tali arrotondamenti oziosi
vengono messi in evidenza, in colore pi\`u scuro,
nella stessa figura.
In entrambe le figure, \ref{fig:660cy_hor} e \ref{fig:1290cy_hor}, si pu\`o
notare un gioco notevole tra le dentature, funzionale alla
robustezza dell'ingranaggio nei confronti della possibile
intrusione di corpi estranei (polvere), il cui valore \`e
stato determinato seguendo le raccomandazioni di \cite{bs978}.
Come gi\`a accennato altrove, tale gioco non suscita inconvenienti, in quanto
la potenza trasmessa \`e minuscola e oltretutto, in generale, il
suo flusso non presenta inversioni.
\endinput

%\newpage
\chapter{Ulteriori Considerazioni sulla Profilatura delle Camme
}\label{cam2}

\section{Una Propriet\`a del Diagramma delle Accelerazioni}

\noindent Vista l'importanza che il diagramma delle accelerazioni riveste nel progetto
dei meccanismi, ci preme mettere in evidenza una sua propriet\`a,
derivante 
da una condizione sussistente nella stagrande maggioranza dei casi
nell'ambito delle camme:
la velocit\`a del cedente deve annullarsi in corrispondenza
dell'inizio ($t=0$) e della fine ($t=t_s$) dell'intervallo di avanzamento.
Si avr\`a quindi
\begin{equation}
\int_0^{t_s} \ddot y\, {\rm d} t = \big|\,\dot y\, \big|_0^{t_s} = 0\,.
\label{e9_risult}
\end{equation}
\noindent L'integrale \ref{e9_risult} evidenzia il fatto che la funzione $\ddot y(t)$ deve avere valor medio
nullo. Svolgiamo ora il seguente integrale, per parti
\begin{equation}
\int_0^{t_s} \ddot y\,t\,{\rm d} t = \big|\,\dot y t\,\big|_0^{t_s} -
\int_0^{t_s} \dot y\,{\rm d} t = 0 - \big|\,y\, \big|_0^{t_s} = -h\,.
\label{e9_moment}
\end{equation}

\begin{figure}[b]
\setlength{\unitlength}{1mm}
\centerline{
\epsfig{file=part3/camme/FIG/trave.eps,height=3cm}
}
\begin{picture}(0,0)(-30,0)
\scriptsize{
\put( 1,20){$M_i$}
\put(33, 7){d}
\put(48,29){$\dot y_{max}$} 
\put(23,13){$\dot y_{max}$} 
}
\end{picture}
\vskip -5mm
\caption
{
{\em Analogia fra diagramma delle accelerazioni e diagramma dei carichi.}
\label{f9_trave}
}
\end{figure}
\noindent Le relazioni \ref{e9_moment} assumono particolare evidenza alla luce della seguente
analogia: immaginando il diagramma delle accelerazioni $\ddot y$ come quello
di un carico distribuito lungo una trave incastrata a un estremo e di
lunghezza $t_s$, le precedenti relazioni equivalgono, in tal caso,
a imporre che tale carico
debba avere risultante nulla e momento all'incastro $M_i=h$.

\section{Un Metodo per Ridurre l'Ingombro delle Camme}\index{camme!offset}\index{camme}

\noindent Come \`e stato esposto a pagina \pageref{fig:f_ang_press}, l'angolo di pressione
massimo costituisce spesso un limite gravoso nel progetto  di una camma,
costringendo il progettista ad aumentare le dimensioni dell'eccentrico,
e spesso anche quelle della macchina, onde evitare gli
inconvenienti dovuti alle elevate forze laterali che nascono e i conseguenti
pericoli di impuntamento.

\begin{figure}[hbt]
\begin{center}
\hbox{\vspace{.5cm}\hspace{-1cm}\includegraphics[width=.9\textwidth]{part3/camme/FIG/camma/camma_offset.pdf}}
\vspace{-4cm}\hspace{7.5cm}\includegraphics[width=0.3\textwidth]{part3/camme/FIG/camma/camma_offset_ruotata.pdf}
\end{center}
\begin{picture}(0,0)(-20,-20)
\scriptsize{
\put(40,223){${\rho}_{\rm min}=24.4\; {\rm mm,\;}\;\alpha=119^{\circ}$}
\put(57,220){\rotatebox{-61}{$\rightarrow$}}
\put(93,217){\rotatebox{-77}{$\rightarrow$}}
\put(98,213){$\theta_{\rm max}=27^{\circ},\; \alpha=103^{\circ}$}
\put(226,98){$R_b=80.1\,{\rm mm}$}
\put(95,130){$\epsilon= 13\, {\rm mm}$}

\put(175,0){meccanismo ruotato di$\;\;\xi=-\arcsin({\frac{\epsilon}{R_b}})$}
}
\end{picture}
      \caption{\em Camma con punteria deviata.}
 \label{fig:camma_offset}
\end{figure}

\noindent Va detto per\`o che, mentre il contenimento di $\theta_{\rm max}$, nelle
fasi di lavoro (di solito chiamate salite),
deriva dalle cautele atte a  evitare una trasmissione
inefficiente delle forze, pericoli di impuntamenti e reazioni vincolari eccessive,
tali precauzioni si possono lasciare cadere 
nei tratti di riposo della camma, cio\`e nelle fasi  di discesa.
Pertanto, se inclinassimo il cedente in modo da diminuire l'angolo di pressione 
nelle fasi di salita, potremmo tollerare l'aumento di $\theta$
nelle fasi in cui tale cedente \`e scarico.
In figura \ref{fig:camma_offset} \`e riportata una camma tastata da un cedente a punteria con
rotella. Per inciso, tale camma \`e la stessa di figura
\ref{fig:f_camma_rotella}, dove per\`o il cedente
\`e disassato di una quantit\`a $\epsilon= 13\,{\rm mm}$ rispetto al centro
della camma stessa\footnote{Si \`e preferito lasciare inalterata la distanza del centro della rotella rispetto al centro della camma
in direzione orizzontale, che risultava essere
$80\,{\rm mm}$; per tale motivo il raggio di base attuale risulta leggermente pi\`u grande.}.
In basso a destra, nella stessa figura, mostriamo il medesimo meccanismo
ruotato di $\xi=9.34^{\circ}$, che dovrebbe mettere in luce la sostanziale
equivalenza della punteria deviata
con la scelta, per il suo movimento, di una direzione inclinata rispetto alle
 alzate fornite dalla sagoma di traslazione, cio\`e dalla legge delle alzate.
Per la precisione,
 l'angolo $\xi$ \`e variabile dovendo essere
$\sin\xi(\alpha)=\epsilon / (R_b + h(\alpha))$, dove $h(\alpha)$
\`e, come sempre, l'alzata all'angolo $\alpha$.
Utilizzando quindi il {\em cedente deviato}\index{cedente!deviato},
i valori di $\theta_{\rm max}$
 differiscono di poco meno di $7^{\circ}$,
 come si pu\`o notare osservando le figure \ref{fig:f_camma_rotella} e
 \ref{fig:camma_offset};
tale differenza
ci permetter\`a, come vedremo a breve, di diminuire l'ingombro della camma.
Per maggiore chiarezza, in figura \ref{fig:ap_offset} sono
riportati i grafici dell'angolo di pressione della camma con
cedente centrato unitamente a quello della
camma con punteria deviata, e la differenza tra i due varia
tra valori di poco inferiori a
$7^{\circ}$ e valori di poco superiori a $9^{\circ}$.

\begin{figure}[hbt]
\begin{center}
\includegraphics[width=0.7\textwidth]{part3/camme/FIG/camma/ap_offset.pdf}
\end{center}
\begin{picture}(0,0)(-150,0)
\scriptsize{
\put(15,108){rotella in asse}
\put(15,80){rotella {\em offset}}
}
\end{picture}
\vskip -5mm
      \caption{\em Confronto tra gli angoli di pressione.}
 \label{fig:ap_offset}
\end{figure}
\noindent Ancora dalla figura \ref{fig:ap_offset}
apprendiamo che l'{\em escamotage} della soluzione {\em offset} si
ripercuote in modo negativo sull'angolo di pressione non solo 
nelle fasi di discesa, ma anche nelle soste, cio\`e quando il profilo della
camma \`e circolare. Mentre il valore assunto da $\theta$ durante le 
discese \`e in generale indifferente, dobbiamo considerare che spesso,
durante le soste, la camma sta esercitando il valore massimo della forza
(situazione tipica nelle presse), quindi un valore di $\theta$ diverso
da zero in queste fasi va attentamente valutato.
Ciononostante, i benefici tecnici che si ottengono con la riduzione
dell'ingombro della
camma, quindi dell'intero corpo della macchina, mediante disassamento del cedente vengono sovente ritenuti irrinunciabili.
Tale riduzione delle dimensioni della camma si opera naturalmente
ridimensionando il raggio di 
base nella soluzione {\em offset} fino a che uno dei due limiti
costituiti dall'angolo di pressione massimo e dal minimo raggio di curvatura 
convessa non viene raggiunto.
La figura \ref{fig:camma_offset_ottimizzata} riporta tale compromesso:
$\theta_{\rm max}=31^{\circ}$ e $\rho_{\rm min}=13\, {\rm mm}$.
Questi due valori offrono sufficienti garanzie per il buon
funzionamento del meccanismo e la camma che li esprime ha un
raggio di base ridotto del 25\% rispetto a quella non ottimizzata.
Quest'ultima \`e presente nella figura in colore tenue.

\begin{figure}[hbt]
\centering
\begin{minipage}[b]{0.6\textwidth}
\centering
\begin{center}
\includegraphics[width=0.9\textwidth]{part3/camme/FIG/camma/camma_offset_ottimizzata.pdf}
\end{center}
\begin{picture}(0,0)(65,0)
\scriptsize{
\put(5,170){${\rho}_{\rm min}=13\; {\rm mm,\;}\;\alpha=121^{\circ}$}
\put(13,166){\rotatebox{-59}{$\longrightarrow$}}
\put(38,157){\rotatebox{-77}{$\rightarrow$}}
\put(44,154){$\theta_{\rm max}=31^{\circ},\; \alpha=103^{\circ}$}
}
\end{picture}
      \caption{\em Camma {\em offset} con ingombro ridotto.}
 \label{fig:camma_offset_ottimizzata}
\end{minipage}\hfill
\begin{minipage}[b]{0.38\textwidth}
\centering
\hbox{\vspace{1cm}\includegraphics[width=0.9\textwidth]{part3/camme/FIG/camma/ap_offset_ottimizzata.pdf}}
\begin{picture}(0,0)(100,0)
	\scriptsize{
\put(77,93){\tiny non ottimizzata (rosa)}
\put(45,66){\tiny ottimizzata (verde)}
\put(152,43){$\alpha$}
}
\end{picture}
	\caption{\em Confronto tra gli angoli di pressione prima e dopo l'ottimizzazione.}
     \label{fig:ap_ottimizzata}
\end{minipage}
\end{figure}

\noindent Stiamo via via approcciando la soluzione tecnica ottimale
della profilatura delle camme che entrano nel progetto di una pressa.
Occorre ora pensare alla camma coniugata o
desmodromica che dovr\`a garantire il ritorno della mazza.
La figura \ref{fig:camma_ood} riporta il gruppo camma primaria, camma coniugata
e il telaio dei cedenti con le relative rotelle. Si pu\`o notare che la soluzione {\em offset} interessa anche la camma desmodromica che presenta, laddove
si trova impegnata nel ritorno della punteria, angoli di pressione modesti,
come mostrato in figura \ref{fig:ap_ood}.
La figura \ref{fig:camma_ood} evidenzia anche
due caratteristiche particolari della camma desmodromica: un raggio di curvatura
eccessivamente ridotto e un angolo di pressione troppo elevato.
Tale circostanza per\`o non
desta alcun timore, presentandosi in zone dove l'eccentrico non \`e impegnato
da alcuna forza\footnote{Il gruppo desmodromico di figura \ref{fig:camma_ood}
\`e stato prodotto in quattro esemplari. Le otto camme di spessore pari a
$25\, {\rm mm}$ sono state ritagliate, tramite una macchina a {\em
elettroerosione a filo} e controllo numerico, direttamente da una lastra di
acciaio 35 NiCrMo 15, gi\`a bonificato. Esse azionano due piccole presse la cui
forza \`e pari a circa una tonnellata.}.
\begin{figure}[b]
\centering
\begin{minipage}[b]{0.6\textwidth}
\centering
\begin{center}
\includegraphics[width=0.9\textwidth]{part3/camme/FIG/camma/camma_ood.pdf}
\end{center}
\begin{picture}(0,0)(65,0)
\scriptsize{
\put(47,20){${\rho_c}_{\rm min}=1.4\; {\rm mm,\;}\;\alpha=274^{\circ}$}
\put(52,27.3){\rotatebox{94}{$\longrightarrow$}}
\put(63,50){\rotatebox{112}{$\rightarrow$}}
\put(62,38){$\theta_{\rm max}=57^{\circ},\;\; \alpha=291^{\circ}$}
}
\end{picture}
      \caption{\em Gruppo camma-desmodromica ottimizzato.}
 \label{fig:camma_ood}
\end{minipage}\hfill
\begin{minipage}[b]{0.38\textwidth}
\centering
\hbox{\vspace{1cm}\includegraphics[width=0.90\textwidth]{part3/camme/FIG/camma/ap_ood.pdf}}
\begin{picture}(0,0)(100,0)
	\scriptsize{
\put(82,89){\tiny coniugata}
\put(62,64){\tiny primaria}
\put(155,43){$\alpha$}
}
\end{picture}
	\caption{\em Confronto tra gli angoli di pressione.}
     \label{fig:ap_ood}
\end{minipage}
\end{figure}

\section{Cenni alla Profilatura di Camme con Cedente a Bilanciere}

\noindent Come accennato a pagina \pageref{fig:f_cambil}, il cedente
a punteria traslante viene talvolta sostituito
da un bilanciere che, anzich\'e traslare, ruota.
Questa soluzione tecnica offre il vantaggio costruttivo di eliminare 
le guide di scorrimento della punteria e di offrire in modo nativo
l'opportunit\`a di posizionare, in cascata all'azionamento a camma,
un sistema articolato, spesso
un manovellismo, la cui manovella pu\`o essere  costituita direttamente
dal bilanciere che tasta la camma.
\`E evidente che, a valle del manovellismo o del quadrilatero azionato
tramite un cedente a bilanciere, otterremo delle leggi di moto diverse
rispetto a quelle implementate sulla camma stessa.
La qual cosa non deve apparire in una luce negativa. Spesso infatti il
sistema articolato posto in serie all'azionamento a camma favorisce
l'ottenimento di una legge delle alzate che soddisfa ancora meglio le
esigenze del processo tecnologico che d\`a origine al progetto.
Entrare nel merito della progettazione di un sistema complesso, come pu\`o
essere un sistema articolato azionato da una camma, esula completamente
dallo scopo di questo lavoro. Segnaliamo solo che la maggior parte
dei programmi di calcolo commerciali che assistono il progettista di 
meccanismi permette di affrontare questo tipo di analisi.

\noindent Torniamo quindi alla nostra camma a bilanciere\index{camme!a bilanciere},
consapevoli di esporre un argomento ``zoppo'', in quanto molto difficilmente
il tema assegnato consiste solo nella rotazione di un bilanciere sia pure 
con leggi ben precise. 
In questo caso, la legge di moto $\beta(\alpha)$ riguarder\`a la rotazione del
bilanciere a partire dalla sua configurazione di riposo.
Giusto per fornire un'idea del percorso di progetto per questo tipo di camma
abbiamo confezionato una legge di moto tale che, applicata alle rotazioni
del bilanciere, non produce movimenti della rotella troppo dissimili rispetto
alla legge delle alzate di figura \ref{fig:f_legge_completa_7_h}
se non
per la natura diversa dei suoi spostamenti, che in questo
caso seguono traiettorie circolari
anzich\'e rettilinee. Il risultato della sintesi di tale camma (e della sua
coniugata) \`e riportato in figura \ref{fig:camma_bilanciere}.

\begin{figure}[hbt]
\centering
\includegraphics[width=0.9\textwidth]{part3/camme/FIG/camma/camma_bilanciere_desmo.pdf}
\begin{picture}(0,0)(150,-20)
\scriptsize{
}
\put(51,260){\rotatebox{-155}{$\curvearrowleft$}}
\put(62,256){$\beta(\alpha)$}
\put(84.5,180){$\beta_{\rm max}= 24^{\circ}$}
\put(88,220){$l= 85 {\rm mm}$}
%\put(-25,150){$\curvearrowright$}
%\put(-25,155){$\omega$}
\put(76,85){$R_b=50 {\rm mm}$}
\end{picture}
        \caption{\em Camma desmodromica con cedente a bilanciere.}
     \label{fig:camma_bilanciere}
\end{figure}

\noindent La figura \ref{fig:legge_camma_bilanciere} mostra la legge di
$\beta(\alpha)$ e quella di accelerazione $\beta''(\alpha)$, espresse rispettivamente
in gradi sessagesimali e gradi sessagesimali  su gradi radianti al quadrato.
L'azionamento della camma comporta dunque un movimento oscillatorio del
bilanciere che spazia in un arco di
$24^{\circ}$. Tale oscillazione consente alla rotella di tracciare un arco di circa $40 \,{\rm mm}$
tastando una camma avente raggio di base $R_b=50\,{\rm mm}$ tramite un bilanciere
di lunghezza $l=85\, {\rm mm}$, il che equivale all'incirca all'azionamento
operato dalla camma di figura \ref{fig:f_camma_rotella} con ingombro leggermente
pi\`u contenuto.
Naturalmente \`e quasi sempre necessario, anche per le camme a bilanciere,
progettare la camma coniugata che permette il cosiddetto {\em positive
drive}, cio\`e l'azionamento sempre condotto da forze di pressione positive.
Come abbiamo gi\`a accennato, la figura \ref{fig:camma_bilanciere} riporta
anche la camma coniugata e il relativo bilanciere.
I due bilancieri, che devono essere considerati come un unico corpo rigido,
sono rappresentati con tinte pi\`u tenui anche nelle due posizioni di sosta,
appena a valle delle due alzate prodotte dalla legge di moto.

\begin{figure}[hbt]
\centering
\includegraphics[width=0.7\textwidth]{part3/camme/FIG/camma/legge_camma_bilanciere.pdf}
\begin{picture}(0,0)(150,-20)
\put(15,91){$\beta(\alpha)$}
\put(15,67){$\beta''(\alpha)$}
\put(133,-13.5){$\alpha$}
\end{picture}
        \caption{\em Leggi di moto del bilanciere.}
     \label{fig:legge_camma_bilanciere}
\end{figure}

\noindent Mentre le cautele che riguardano il raggio di curvatura
minimo del profilo nella parte convessa rimangono fondamentali anche per
queste camme, le restrizioni che riguardano l'angolo di pressione sono
in generale pi\`u lasche. Il motivo di questa maggiore tolleranza,
soprattutto quando la camma, girando come di consueto in senso
orario, \`e disposta rispetto al cedente a bilanciere nella configurazione
riportata in figura, cio\`e quando la camma viene ``accarezzata dal cedente'', 
risiede nella
notevole diminuzione del pericolo di impuntamento e di scarso rendimento che si
verifica
nelle guide di scorrimento di punterie impegnate da notevoli azioni
laterali. Rimane tuttavia da considerare l'inefficienza di trasmettere
forze fuori asse rispetto alla direzione del moto. 

\section{Cenni alla Profilatura di Camme Tastate da Piattello}\index{camme!cedente a piattello}

\noindent Come gi\`a accennato a pagina \pageref{fig:f_campiat}, si trova
nella tecnica una vasta famiglia di camme tastate da un cedente
che striscia mediante una porzione di superficie, generalmente
piana, sul profilo della camma
stessa: le camme tastate da piattello. Queste camme popolano il mondo
dell'{\em automotive} essendo normalmente responsabili dell'apertura e della
chiusura dei condotti di aspirazione e di scarico dei gas nei motori a quattro
tempi. Il loro ingombro raramente eccede i $40\, {\rm mm}$; il cedente
 pu\`o essere
una punteria traslante oppure un bilanciere e il contatto tra camma e cedente
\`e quasi sempre garantito da una molla pre-caricata, anche se esistono esempi
notevoli di azionamenti a camme, per valvole di motori endotermici, di
tipo desmodromico.
Il percorso che abbiamo seguito per la sintesi delle camme tastate da rotella
rimane valido anche in quest'ambito bench\'e non sempre i progettisti di camme
per motori partono dalla legge dell'accelerazione. Quest'ultima rimane comunque
una caratteristica importantissima dell'azionamento in quanto anche da essa dipendono
i malfunzionamenti del motore ad alti regimi di rotazione; quando cio\`e
le valvole, trascinate dalle loro forze d'inerzia, si staccano dalla camma
nonostante la presenza delle molle, causando il cosiddetto {\em sfarfallamento}\index{sfarfallamento}.
Ciononostante, in ambito motoristico, non \`e infrequente privilegiare
la legge delle alzate partendo nel processo di sintesi da
questa, e verificare a posteriori che essa non produca
accelerazioni inaccettabili.

\noindent Se tentassimo di disegnare una camma per motore a quattro tempi
provvista delle proporzioni geometriche  caratteristiche in questo campo,
ed eseguissimo questo tentativo
adagiando ingenuamente sul tamburo della camma
una legge di moto ottenuta disponendo
in serie un paio delle solite leggi trapezoidali modificate, come quelle di figura
\ref{fig:leggi_camma_piattello_ingenua}\footnote{In
questa e nelle figure analoghe che seguiranno \`e sottintesa la misura delle
accelerazioni in millimetri su secondo quadro, delle velocit\`a
in millimetri al secondo e infine dell'alzata
in millimetri.},
otterremmo il pessimo risultato gi\`a illustrato in figura
\ref{fig:f_sottotaglio} e di nuovo proposto in figura
\ref{fig:camma_piattello_ingenua}.

\begin{figure}[hbt]
\centering
\begin{minipage}[b]{0.48\textwidth}
\centering
\includegraphics[width=0.8\textwidth]{part3/camme/FIG/camma/camma_piattello_ingenua.pdf}
\begin{picture}(0,0)(0,0)
\scriptsize{
}
\end{picture}
      \caption{\em Camma a piattello inservibile.}
 \label{fig:camma_piattello_ingenua}
\end{minipage}\hfill
\begin{minipage}[b]{0.48\textwidth}
\centering
\includegraphics[width=0.8\textwidth]{part3/camme/FIG/camma/leggi_camma_piattello_ingenua.pdf}
\begin{picture}(0,0)(80,0)
	\scriptsize{
\put(-5,70){\color{red}$y''(\alpha)$}
\put(-10,30){\color{dark-green}$y'(\alpha)$}
\put(20,47){\color{blue}$y(\alpha)$}
\put(67,2){$\alpha$}
}
\end{picture}
	\caption{\em Approccio ingenuo alla camma con piattello.}
     \label{fig:leggi_camma_piattello_ingenua}
\end{minipage}
\end{figure}
 
\noindent L'esperimento numerico
\`e stato condotto rispettando i seguenti parametri: raggio della porzione circolare della camma
$R_c=14\,{\rm mm}$\footnote{Ricordiamoci la scelta, operata a suo tempo, di
chiamare raggio di base $R_b$ quello del percorso del centro della rotella. Come
abbiamo detto in altro luogo moltissimi autori chiamano raggio di base quello
della porzione circolare ad alzata zero del profilo fisico. Qui viene indicato
con $R_c$.}, alzata massima $h=12\, {\rm mm}$, angolo di salita
$\alpha_s=75^{\circ}$ e angolo di discesa $\alpha_d=75^{\circ}$.
Questi dati corrispondono all'incirca a quelli delle camme di aspirazione
di un motore Alfa Romeo che equipaggiava la Giulia 1600 cc GTA.
Il marcato sottotaglio, che la figura
\ref{fig:camma_piattello_ingenua} evidenzia,
si spiega pensando il piattello come una rotella di raggio $r_r$ infinito 
che tasta una camma di raggio di base $R_b$ altrettanto infinito dove la 
differenza $R_b - r_r = R_c$. 
\`E opportuno a questo punto ricordare la formula \ref{r_cur} che fornisce
il raggio di curvatura del profilo della camma, una volta sottratto quello
della rotella. Avremo
\begin{equation}
\rho ={{{{[{{y'}^2\over{(R_b + y)^2}}+1]}^{3/2}(R_b+y)}\over{{1-{y''\over{(R_b + y)}}} +{2{y'}^2\over{(R_b + y)^2}}}}}-(R_b-R_c)\,,
\label{r_cur_piat0}
\end{equation}
\noindent quindi per elevati valori di $R_b$
\begin{equation}
\lim_{R_b\to\infty}\rho =
\lim_{R_b\to\infty}{(R_b+y)\over{1 -{y''\over{(R_b + y)}}}} -(R_b-Rc)\,.
\label{r_cur_piat2}
\end{equation}
\noindent Mediante pochi passaggi, che lasciamo al lettore, otteniamo
\begin{equation}
\lim_{R_b\to\infty}\rho =
R_c + y + y''\,,
\label{r_cur_piat1}
\end{equation}
\noindent riportata con altri simboli anche in \cite{ruggieri}, pag. 67.
\noindent Il pericolo di profilo intrecciato si manifester\`a dunque
per valori negativi di $y''$ e, come sempre, tale pericolo
\`e scongiurato dal crescere delle dimensioni della camma, cio\`e di $R_c$.
Occorre quindi ricercare leggi di moto con  bassi valori assoluti di
accelerazione negativa, compatibilmente
con l'esigenza di avere valori non eccessivi di accelerazione positiva.
In figura \ref{fig:acc_piat} \`e rappresentata una legge idonea a ottenere
la camma che stiamo cercando e che rappresentiamo in figura
\ref{fig:camma_piattello}.
Si noti che le due leggi di accelerazione sono profondamente asimmetriche e che
i bassi valori assoluti di accelerazione negativa risultano indispensabili
laddove anche il valore dell'alzata \`e modesto, come si evince facilmente dalla
formula \ref{r_cur_piat1}.

\begin{figure}[hbt]
\centering
\begin{minipage}[b]{0.48\textwidth}
\centering
\includegraphics[width=0.8\textwidth]{part3/camme/FIG/camma/camma_piattello.pdf}
\begin{picture}(0,0)(0,0)
\scriptsize{
}
\end{picture}
      \caption{\em Camma tastata da piattello ottenuta da leggi trapezoidali modificate.}
 \label{fig:camma_piattello}
\end{minipage}\hfill
\begin{minipage}[b]{0.48\textwidth}
\centering
\includegraphics[width=0.8\textwidth]{part3/camme/FIG/camma/leggi_camma_piattello.pdf}
\begin{picture}(0,0)(80,0)
	\scriptsize{
\put(-50,83.6){$y''_{\rm max}=84\,{\rm mm}/{\rm rad}^2$}
\put(-5,70){\color{red}$y''(\alpha)$}
\put(-6,27){\color{dark-green}$y'(\alpha)$}
\put(20,43){\color{blue}$y(\alpha)$}
\put(67,2){$\alpha$}
}
\end{picture}-
	\caption{\em Legge di moto composta da due leggi trapezoidali modificate adatta per camma-piattello.}
     \label{fig:acc_piat}
\end{minipage}
\end{figure}

\noindent 

\noindent Il cedente con piattello non
\`e azionato, in generale, dal punto del profilo che si trova sull'asse del
cedente, perch\'e il contatto tra il profilo stesso e il piattello pu\`o
trovarsi altrove. \`E il caso di avvertire il lettore che per le camme
tastate da piattello l'angolo di pressione \`e sempre costante e, nel solo
caso che abbiamo trattato, vale a dire quello che presenta
l'asse della punteria ortogonale
alla superficie del piattello stesso, risulta essere $\theta=0$.

\noindent Da alcuni rilievi, eseguiti sulla camma della nota casa
automobilistica, possiamo affermare che essa non coincide con quella di 
figura \ref{fig:camma_piattello}, anzi a un occhio appena allenato le
due camme appaiono alquanto difformi.
Ci \`e sembrato cos\`i di qualche interesse eseguire un tentativo numerico 
volto alla interpolazione dei pochi valori rilevati dal profilo della camma
originale e tratti da un
vecchio disegno a nostra disposizione.
Tale interpolazione, svolta tramite {\em spline cubiche}\index{spline},
fornisce il profilo della camma\footnote{Si tratta del profilo effettivo della
camma, cio\`e di una sagoma di traslazione che aggiunta al tamburo della
camma stessa produrr\`a lo stesso contorno dell'eccentrico rilevato dal  disegno
che abbiamo a disposizione.
Questa sagoma di traslazione non coincide pertanto con la legge delle 
alzate del piattello che in seguito riporteremo.} riportato in figura
\ref{fig:spline_profilo_completa}, dove le crocette indicano i pochi punti misurati.

\begin{figure}[hbt]
\begin{center}
\includegraphics[width=0.6\textwidth]{part3/camme/FIG/camma/spline/spline_profilo_completa.pdf}
\end{center}
\begin{picture}(0,0)(-150,0)
\scriptsize{
}
\end{picture}
      \caption{\em {\em Spline} del profilo della camma.}
 \label{fig:spline_profilo_completa}
\end{figure}

\noindent Aggiungendo questa ``finta'' sagoma di traslazione ad un tamburo con
$R_b=14\,{\rm mm}$, comportandoci cio\`e come se la camma dovesse essere in seguito
tastata da una punta, otteniamo la figura \ref{fig:cammaspline}
che mostra un eccentrico praticamente identico a quello
realmente impiegato nella citata automobile.
 
\begin{figure}[hbt]
\centering
\begin{minipage}[b]{0.48\textwidth}
\centering
\includegraphics[width=0.8\textwidth]{part3/camme/FIG/camma/spline/camma_piattello_profilo_spline.pdf}
\begin{picture}(0,0)(0,0)
\scriptsize{
}
\end{picture}
      \caption{\em Camma con profilo ottenuto dalla {\em spline} di fig. \ref{fig:spline_profilo_completa}.}
 \label{fig:cammaspline}
\end{minipage}\hfill
\begin{minipage}[b]{0.48\textwidth}
\centering
\includegraphics[width=0.8\textwidth]{part3/camme/FIG/camma/spline/leggi_piattello_profilo_spline.pdf}
\begin{picture}(0,0)(80,0)
	\scriptsize{
\put(-45,84){$y''_{\rm max}=110\,{\rm mm}/{\rm rad}^2$}
\put(4,70){\color{red}$y''(\alpha)$}
\put(3,28){\color{dark-green}$y'(\alpha)$}
\put(20,45){\color{blue}$y(\alpha)$}
\put(67,2){$\alpha$}
}
\end{picture}
	\caption{\em Effettive leggi di moto.}
     \label{fig:leggi_spline}
\end{minipage}
\end{figure}

\noindent La camma in questione \`e per\`o tastata da un piattello anzich\'e da una
punta. \`E interessante perci\`o ricavare le leggi di moto effettivamente
ottenute, che riportiamo in figura \ref{fig:leggi_spline}. 
Pur riconoscendo in tali leggi molte somiglianze con quelle canoniche di
figura \ref{fig:acc_piat}, prendiamo atto che la loro forma contiene variazioni
impossibili da ottenere
sfruttando solamente le possibilit\`a offerte dalla legge di moto trapezoidale
modificata.
\`E il caso di aggiungere che la corretta lubrificazione gioca un ruolo
fondamentale nelle camme tastate da piattello. Mentre si possono trovare
meccanismi camma-rotella sia immersi in bagno d'olio, sia lubrificati
con maggior superficialit\`a, \`e praticamente impossibile pensare a una
coppia camma-piattello che non sia fornita di lubrificazione forzata
e costante, come effettivamente avviene nei motori a combustione interna
a quattro tempi.
\newpage
\thispagestyle{empty}
\null
\endinput



\addcontentsline{toc}{chapter}{Bibliografia}
\begin{thebibliography}{99}
\thispagestyle{empty}
\bibitem{finzi} {\bfseries Bruno Finzi}, 
	{\em Meccanica Razionale},
	1 Vol.,
	terza edizione, Zanichelli, Bologna, 1959. 

\bibitem{goldstein} {\bfseries Herbert Goldstein}, 
		{\em Classical Mechanics}, 
	 second edition, Addison-Wesley Publishing Company, Reading,
	Massachusetts, 1980.

\bibitem{levicivita} {\bfseries Tullio Levi-Civita}, {\bfseries Ugo Amaldi}, 
	{\em Compendio di Meccanica Razionale}, 2 Voll.,
	seconda edizione riveduta, Zanichelli, Bologna, 1938.

\bibitem{whittaker} {\bfseries Edmund T. Whittaker}, 
	{\em A Treatise on the Analytical Dynamics of Particles and Rigid Bodies},
	fourth edition, Cambridge University Press, Cambridge, 1947.

\bibitem{aristotele} {\bfseries [Aristotele]},
		{\em Meccanica},
		 a cura di Maria Fernanda Ferrini,
		Bompiani, Milano, 2010.

\bibitem{mach} {\bfseries Ernst Mach}, 
	{\em La Meccanica nel suo Sviluppo Storico-Critico},
	Bollati Boringhieri, Torino, 1977. 

\bibitem{giordana} {\bfseries Franco Giordana}, 
	{\em Lezioni di Meccanica delle Macchine},
	Edizioni Spiegel, Milano, 2001. 

\bibitem{denavit} {\bfseries Richard S. Hartenberg}, {\bfseries Jacques
 Denavit},
        {\em Kinematic Synthesis of Linkages},
        McGraw-Hill, New York, 1964.

\bibitem{sesini1} {\bfseries Ottorino Sesini},
	{\em Meccanica Applicata alle Macchine},
	{\em Parte Prima}, {\em Cinematica},
	Casa Editrice Ambrosiana, Milano, 1964.	

\bibitem{pellicano} {\bfseries Massimo Callegari},
	{\bfseries Pietro Fanghella}, {\bfseries Francesco Pellicano},
	{\em Meccanica Applicata alle Macchine},
	seconda edizione, Citt\`a Studi Edizioni, Milano, 2017.

\bibitem{punzi} {\bfseries Giovanni Battista Punzi},
	{\em Corso Superiore Completo di Meccanica Generale ed Applicata},
	{\em Volume Secondo},
	seconda edizione riveduta, Hoepli, Milano, 1978.

\bibitem{henriot} {\bfseries Georges Henriot},
	{\em Ingranaggi}. {\em Trattato teorico e pratico}, Vol. 1 e Vol. 2,
	Tecniche Nuove, Milano, 1977.

\bibitem{guiggiani} {\bfseries Massimo Guiggiani},
	{\em Generazione per Inviluppo di Ruote Dentate ad Evolvente},
 in $<$http://docenti.ing.unipi.it/guiggiani-m/guiggiani-ruote-dentate.pdf$>$.

\bibitem{righettini} {\bfseries Guido Ruggieri},
	{\bfseries Paolo Righettini},
        {\em Ruote Dentate con Spostamento del Profilo},
      	McGraw-Hill, Milano, 2003.

\bibitem{martin} {\bfseries George H. Martin},
        {\em Kinematics and Dynamics of Machines}, second edition,
        McGraw-Hill, New York, 1982.

\bibitem{sesini3} {\bfseries Ottorino Sesini},
        {\em Meccanica Applicata alle Macchine},
        {\em Parte Terza}, {\em Le Coppie Elementari nelle Macchine},
        Casa Editrice Ambrosiana, Milano, 1964.

\bibitem{ruggieri} {\bfseries Pier Luigi Magnani},
	{\bfseries Guido Ruggieri},
        {\em Meccanismi per Macchine Automatiche},
        UTET, Torino, 1986.

\bibitem{hrones} {\bfseries John A. Hrones},
	{\bfseries George L. Nelson},
        {\em Analysis of the Four-Bar Linkage}. {\em Its Application to
        the Synthesis of Mechanisms},
        The MIT Press, Cambridge (MA), 1951.

\bibitem{hartog} {\bfseries  Jacob Pieter Den Hartog},
        {\em Mechanics},
        Dover Publications, New York, 1948.
 
\bibitem{molian} {\bfseries Samuel Molian},
        {\em The Design of Cam Mechanisms and Linkages},
        Constable, London, 1968.

\bibitem{timoshenko} {\bfseries Stephen P. Timoshenko}, {\bfseries James N. Goodier},
        {\em Theory of Elasticity},
        McGraw-Hill, New York, 1951.

\bibitem{grossmann} {\bfseries Jules Grossmann},
	{\bfseries Hermann Grossmann},
        {\em Lessons in Horology}, 1 Vol.,
        The Keystone, Philadelphia, 1905.

\bibitem{bs978} {\bfseries BS 978}, {\bfseries Part 2},
	{\em British Standard for Gears for Instruments and Clockwork
	Mechanisms}, {\em Part 2}, {\em Cycloidal type Gears},
	London, 1952.

\end{thebibliography}

\endinput

\printindex


\end{document}
